\documentclass[a4paper,11pt]{ltjsarticle}

% ========== package ========== %
% --- 数式環境 --- %
\usepackage{amsmath}
\usepackage{amssymb}
\usepackage{mathtools}
% --- 数式環境 END --- %
% ブラケット
\usepackage{braket}
% ベクトルの太文字
\usepackage{bm}
% 添え字左右上下 \fourIdx{左上}{左下}{右上}{右下}{ベース} で使う．Iは大文字
\usepackage{fouridx}
% 色
\usepackage{color}
% 余白指定
\usepackage{geometry}
% --- ページ内リンク --- %
\usepackage{hyperref}
\hypersetup{colorlinks=true,linkcolor=blue}
% --- ページ内リンク END --- %
% 花文字
\usepackage{mathrsfs}
% 画像
\usepackage{graphicx}
% tcolorbox
\usepackage{tcolorbox}
% appendix
\usepackage{appendix}
% --- 参考文献 --- %
\usepackage{etoolbox}
\usepackage[no-math]{fontspec} % fontspecは読み込んでOKですが、設定順序に注意が必要です
%% --- 参考文献フォント --- %%
  \IfFontExistsTF{Times New Roman}{
    \newfontfamily\bibfontface[Ligatures=TeX]{Times New Roman}
  }{
    \newfontfamily\bibfontface[Ligatures=TeX]{FreeSerif} % 代替フォント
  }
%% --- 参考文献フォント END --- %%
%% --- 参考文献リンク位置 --- %%
  \BeforeBeginEnvironment{thebibliography}{%
    \phantomsection
    \addcontentsline{toc}{section}{参考文献}
  }
  \AtBeginEnvironment{thebibliography}{%
    \normalfont\bibfontface
  }
%% --- 参考文献のリンク位置 --- %%
% --- 参考文献 END --- %
% ========== package END ========== %

% ========== newcommand ========== %
% --- セクションから目次へのリンク --- %
\newcommand{\sectionlink}[1]{\section[#1]{\hyperlink{toc}{#1}}}
\newcommand{\subsectionlink}[1]{\subsection[#1]{\hyperlink{toc}{#1}}}
% セクション番号リンク
\makeatletter
\renewcommand{\@seccntformat}[1]{%
  \protect\hyperlink{toc}{\csname the#1\endcsname}\quad
}
\makeatother

% 原文との照合のため，数式番号は原文と同じ通し番号を用いる．
% ==========- newcommand END ========== %
\allowdisplaybreaks
%===========================================================
% --- タイトル --- %
%===========================================================
\geometry{left=25mm, right=25mm, top=30mm, bottom=30mm}

% 原文で使用される記号と行列用の追加設定
\usepackage{euscript}
\usepackage{dcolumn}
% 長い数式は内容を保ち，式番号を含めて版面内に収める．
\usepackage{adjustbox}
\newcommand{\fitdisplay}[1]{\adjustbox{max width=\dimexpr\linewidth-3em\relax}{$\displaystyle #1$}}
\def\Tr{\mathrm{Tr}}
\def\tr{\mathrm{tr}}
\def\R{\mathrm{Re}}
\def\I{\mathrm{Im}}
\def\d{\mathrm{d}}
\def\D{\mathbf{D}}
\def\rt{(\mathbf{r},t)}
\def\rtp{(\mathbf{r},t')}
\def\qo{(\mathbf{q},\omega)}
\def\qom{(-\mathbf{q},-\omega)}
\def\Qk{\hat{Q}_{\EuScript{K}}}
\def\tauel{\tau_{\mathrm{el}}}
\hypersetup{pdftitle={Keldysh法と非線形σ模型：基本原理と応用（日本語訳）},pdfauthor={Alex Kamenev and Alex Levchenko}}
\setlength{\emergencystretch}{3em}
\tolerance=3000
\title{\bfseries Keldysh法と非線形$\sigma$模型：\\基本原理と応用}
\author{Alex Kamenev and Alex Levchenko\\
{\small ミネソタ大学物理学科，Minneapolis, MN 55455, USA}}
\date{2009年4月27日（原文第3版）}
\begin{document}
\maketitle
\begin{center}
{\small 日本語訳：\href{https://arxiv.org/abs/0901.3586v3}{arXiv:0901.3586v3}\\
原題：Keldysh technique and non-linear $\sigma$-model: basic principles and applications}
\end{center}
% 原文の参考文献番号を維持するため，先に全キーを原典順で登録する．
\nocite{Keldysh-1964,KonstantinovPerel,Schwinger,KadanoffBaym,Feynman,Wyld,MSR,DeDominicis,Lifshitz,RammerSmith,
SmithJensen-book,Mahan,Rammer-book,SchwabRaimondi,Spicka,Kamenev,Matsubara,AGD,Gardiner,Risken,
HorbachSchon,KamenevAndreev,ChamonLudwigNayak,Feigel'manLarkinSkvortsov,Edwards,Wegner,EfetovLarkinKhmelnitskii,Finkel'stein,BelitzKirkpatrick,Lerner,
Efetov,Efetov-book,Mirlin,NegeleOrland,Leggett,CaldeiraLeggett,Elgart,FeynmanHibbs,LarkinOvchinnikov,Kubo,
Landauer,Buttiker,Reviews-on-Landauer-1,Reviews-on-Landauer-2,Lesovik,Kogan,Blanter-Buttiker,Martin,Pogrebinskii,Price,
Solomon,Gramila,Sivan,Lilly,Pillarisetty,Smith,MacDonald,Kamenev-Oreg,Flensberg,Debray-1,
Debray-2,Morimoto,Yamamoto,Aguado,Kouwenhoven,Khrapai,Reznikov,LevchenkoKamenev-QPC,Glazman-QPC,Nazarov-TunnelingAction,
BelzigNazarov,FKLS,ABG,Nazarov-CBWTJ,SnymanNazarov,Usadel,Pierre,Hikami,Gor'kovLarkinKhmelnitskii,AAKL,
LeeRamakrishnan,ALW-book,AltshulerShklovskii,Mehta,AltlandKamenev,Altshuler-UCF,LeeStone,LeeStoneFukuyama,AKL-UCF,Beenakker,
LevitovLesovik-1,LevitovLesovik-2,LevitovLee,Nazarov-FCS,Levitov-FCS,Zala,Altshuler,AltshulerAronov-KE,CatelaniAleiner,KoganShulman,
Nagaev,Gutman,Bagrets,DimitrovaKtavtsov,BurmiChelk,AltshulerAronov,AAL,AltshulerAronov-review,Finkel'stein94,Nazarov-ZBA,
LevitovShytov,Kopietz,GutmanGefen,Schmid,Blanter,Drag-Third-order,SkvortsovLarkinFeigel'man,Gor'kov,Nambu,NarozhnyAleiner,
BelzigWilhelm,SchmidtSchon,LarkinOvchinnikov-F,Tinkham,Kopnin,LarkinOvchinnikov-KinEq,Zaitsev,KuprianovLukichev,Lambert,BelzigBruder,
Gueron,Zhou,Hammer,Heikkila,Wilhelm,ZhouSpivak,StoofNazarov,Virtanen,Stenberg,Crosser,
Gor'kov-GL,GinzburgLandau,Gor'kovEliashberg,HoughtonMaki,HuThompson,Kramer,SchonAmbegaokar,Hu,Krempasky,Otterlo,
Larkin-Varlamov,AltlandSimons,YurkevichLerner,LevchenkoKamenev,AslamazovLarkin,Maki,Thompson,GolubovKupriyanov,Levchenko-DOS,KulikOmelyanchuk,
Levchenko-Js,Hurault,Abrahams,VarlamovDorin,DiCastro,Reizer,Kulik,Scalapino,Martin-Balatsky,Levchenko-J-Noise}
\begin{abstract}
本総説の目的は，相互作用する非平衡フェルミオン系およびボソン系に対するKeldysh法を，教育的かつ包括的に導入することである．基礎となる微視的模型の汎関数積分表示に重点を置く．総説の大部分を，乱れた金属および超伝導体に対する非線形$\sigma$模型の導出と応用に充てる．輸送特性，メゾスコピック効果，計数統計，相互作用補正，運動論方程式などを論じる．乱れた超伝導体を扱う章では，Usadel方程式，揺らぎ補正，時間依存Ginzburg--Landau理論，近接効果，Josephson効果などを取り上げる．本総説はarXiv:cond-mat/0412296を大幅に拡張したものである．
\bigskip
\noindent\textbf{キーワード：} Keldysh法；Green関数；運動論方程式；非線形$\sigma$模型；メゾスコピック系；揺らぐ超伝導体．
\end{abstract}
\hypertarget{toc}{}
\tableofcontents
\clearpage
\sectionlink{序論}\label{sec_intro}
\subsectionlink{動機と構成}\label{sec_motivation}


本総説は，非平衡の相互作用する多体系を扱うためのKeldysh形式を解説するものである．この手法の名称は，L.~V.~Keldyshの1964年の論文~\cite{Keldysh-1964}に由来する．Keldysh法と密接に関連する先行研究として，KonstantinovとPerel~\cite{KonstantinovPerel}，Schwinger~\cite{Schwinger}，KadanoffとBaym~\cite{KadanoffBaym}，FeynmanとVernon~\cite{Feynman}の方法を挙げておくべきであろう．Keldysh法の古典的対応物も，それ自体としてきわめて有用かつ興味深い．その例には，確率系に対するWyldのダイアグラム法~\cite{Wyld}とMartin--Siggia--Rose（MSR）法~\cite{MSR}がある（関連するDeDominicisの研究~\cite{DeDominicis}も参照）．


この手法については，既存の文献に数多くの解説がある~\cite{Lifshitz,RammerSmith,SmithJensen-book,Mahan,Rammer-book,SchwabRaimondi,Spicka}．本総説はLes Houches第\textit{LXXXI}回講義~\cite{Kamenev}を大幅に拡充したものであり，汎関数積分によるアプローチに重点を置く．このアプローチによって，理論の構造と内的論理はいっそう明瞭になり，見通しもよくなる．われわれはこの手法のさまざまな応用に焦点を当て，平衡系のMatsubara法~\cite{Matsubara,AGD}や古典的なLangevin方程式，Fokker--Planck（FP）方程式~\cite{Gardiner,Risken}など，他の手法との関係を明らかにする．総説の大部分は，非線形$\sigma$模型（NLSM）~\cite{HorbachSchon,KamenevAndreev,ChamonLudwigNayak,
Feigel'manLarkinSkvortsov}の詳細な導出に充てる．これはおそらく，不規則な金属および超伝導体の理論における最も強力な計算手法である．この部分は，$\sigma$模型のレプリカ版~\cite{Wegner,EfetovLarkinKhmelnitskii,Finkel'stein,BelitzKirkpatrick,Lerner}および超対称版~\cite{Efetov,Efetov-book,Mirlin}についての従来の解説を補うものと考えられる．


本総説の目的は，多体理論のKeldysh定式化における次の応用と利点を示すことである．
\begin{itemize}
\item 外場の存在，または過渡的な状況により，熱平衡から外れた系を扱うこと．


\item 凍結した乱れをもつ系の理論において，レプリカ法および超対称法に代わる手法を与えること．


\item 量子観測量の平均値や相関関数だけでなく，その完全計数統計を計算すること．


\item Matsubara解析接続が煩雑になりうる平衡問題を扱うこと．
\end{itemize}
われわれの意図は，これまで文献に現れたこの手法の応用をすべて網羅することではない．むしろ，初学者に役立つ，系統的で自己完結した解説を目指す．したがって，引用文献の選択はきわめて限定的かつ主観的である．主な目的は，包括的な文献案内を示すことよりも，ここで選んだ例について，より深い詳細を提供することにある．


本総説の構成は次のとおりである．まず，相互作用のないボソン系（第\ref{sec_bosons}節）およびフェルミオン系（第\ref{sec_fermion}節）の簡単な例から出発して，閉じた時間経路の概念（第\ref{sec_contour}節）やGreen関数など，Keldysh法の基本要素を導入する．ボソンの相互作用，ダイアグラム法，量子運動論方程式については第\ref{sec_int_bos}節で論じる．第\ref{sec_environment}節では，散逸環境（浴）に接した粒子を扱う．この例を用いて，古典的手法（Langevin，Fokker--Planck，Martin--Siggia--Rose）および平衡系のMatsubara法とのつながりを示す．凍結した乱れの存在下にある非相互作用フェルミオンは，第\ref{sec_NLSM}節においてKeldysh非線形$\sigma$模型を用いて扱う．これを第\ref{sec_int_ferm}節ではCoulomb相互作用を含む場合へ，第\ref{sec_SuperCond}節では超伝導相関を含む場合へ一般化する．技術的な説明にはすべて応用例を添え，手法のさまざまな側面を示す．扱う題材は，メゾスコピック試料のスペクトル統計，普遍的コンダクタンス揺らぎ（UCF），電子輸送のショット雑音と完全計数統計，不規則な金属および超伝導体の輸送係数に対する相互作用補正，Coulombドラッグなどである．また，微視的なKeldysh形式から，Caldeira--Leggett，時間依存Ginzburg--Landau（TDGL），Usadelなどの有効な現象論的模型を導くことにも，十分な注意を払う．


%$$$$$$$$$$$$$$$$$$$$$$$$$$$$$$$$$$$$$$$$$$$$$$$$$$$$$$$$$$$$$$$$$$$$$$$$$$$$$$$$$$$$$$$$$$$$$$$$$$$$$$$$$$$
%$$$$$$$$$$$$$$$$$$$$$$$$$$$$$$$$$$$$$$$$$$$$$$$$$$$$$$$$$$$$$$$$$$$$$$$$$$$$$$$$$$$$$$$$$$$$$$$$$$$$$$$$$$$
\subsectionlink{閉じた時間経路}\label{sec_contour}


時間に依存していてもよいHamiltonian $\hat{H}(t)$に支配される量子多体系を考える．遠い過去$t=-\infty$において，系は多体密度行列$\hat{\rho}(-\infty)$によって指定される状態にあったと仮定しよう．この密度行列の厳密な形は重要ではない．たとえば，Hamiltonian $\hat{H}(-\infty)$に対応する平衡密度行列であってもよい．密度行列はvon Neumann方程式$\partial_{t}\hat{\rho}(t)=-i\big[\hat{H}(t),\hat{\rho}(t)\big]$に従って時間発展する．ここでは$\hbar=1$とおいた．その形式解は$\hat{\rho}(t)=\hat{\mathcal{U}}_{t,-\infty}\hat{\rho}(-\infty)
\big[\hat{\mathcal{U}}_{t,-\infty}\big]^{\dag}=
\hat{\mathcal{U}}_{t,-\infty}\hat{\rho}(-\infty)
\hat{\mathcal{U}}_{-\infty,t}$で与えられ，時間発展演算子は次の時間順序付き指数関数で表される．
\begin{eqnarray}
\hat{\mathcal{U}}_{t,t'}=\mathbb{T}\exp\left(-i\int^{\,t}_{t\,'}
\hat{H}(\tau)\mathrm{d}\tau\right)=\lim_{N\rightarrow\infty}
e^{-i\hat{H}(t)\delta_{t}}e^{-i\hat{H}(t-\delta_{t})\delta_{t}}\ldots
e^{-i\hat{H}(t\,'+\delta_{t})\delta_{t}},
\end{eqnarray}ここで，無限小の時間刻みは$\delta_{t}=(t-t')/N$である．


通常われわれが計算したいのは，ある観測量$\hat{\mathcal{O}}$（たとえば密度や電流）の時刻$t$における期待値であり，これは次のように定義される．
\begin{equation}\label{contour-O-average}
\big\langle\hat{\mathcal{O}}(t)\big\rangle\equiv
\frac{\Tr\{\hat{\mathcal{O}}\hat{\rho}(t)\}}
{\Tr\{\hat{\rho}(t)\}}=\frac{1}{\Tr\{\hat{\rho}(t)\}}
\Tr\big\{\hat{\mathcal{U}}_{-\infty,t}\hat{\mathcal{O}}
\hat{\mathcal{U}}_{t,-\infty}\hat{\rho}(-\infty)\big\}\,,
\end{equation}ここでトレースは多体Hilbert空間全体にわたってとる．最後のトレース内の式は，右から左へ読むと，初期密度行列を指定する$t=-\infty$から，観測量を計算する$t$まで前進し，その後$t=-\infty$まで後退する時間発展を表している．平衡系では，特別に工夫された方法により，この前進・後退の時間発展を避けることができる．


たとえば零温度の量子場理論~\cite{AGD}において，これがどのように働くかを思い出そう．この理論では$\langle\mathrm{GS}|\hat{\mathcal{O}}|\mathrm{GS}\rangle=
\langle0|\hat{\mathcal{U}}_{-\infty,t}\hat{\mathcal{O}}
\hat{\mathcal{U}}_{t,-\infty}|0\rangle$という型の期待値を扱う．ここで$|\mathrm{GS}\rangle=\hat{\mathcal{U}}_{t,-\infty}|0\rangle$は，相互作用をすべて含めた系の基底状態である．Hamiltonianに許される時間依存性は，遠い過去に相互作用を断熱的に入れ，遠い未来に断熱的に切ることだけである．したがって時間発展演算子は，相互作用を断熱的に導入することによって，単純な非相互作用基底状態$|0\rangle$が$|\mathrm{GS}\rangle$へ発展することを記述する．ここで工夫を施す．最も左に演算子$\hat{\mathcal{U}}_{+\infty,-\infty}$を挿入し，時間軸全体に沿った時間発展を完成させるのである．このとき$\langle0|\hat{\mathcal{U}}_{+\infty,-\infty}=\langle0|e^{iL}$が成り立つと考える．この議論は，遠い過去と遠い未来で相互作用をそれぞれゆっくり「入れる」「切る」と，系が断熱的に基底状態をたどるという仮定に基づく．したがって，非相互作用基底状態を時間軸全体に沿って発展させた結果は，位相因子$e^{iL}$を獲得することだけである．そこでこの因子で割れば，追加した時間発展区間を補償できる．その結果$\langle\mathrm{GS}|\hat{\mathcal{O}}|\mathrm{GS}\rangle=
\langle0|\hat{\mathcal{U}}_{+\infty,t}\hat{\mathcal{O}}
\hat{\mathcal{U}}_{t,-\infty}|0\rangle/e^{iL}$となり，後退区間のない，時間軸に沿った前進の時間発展を記述すればよくなる．ただし代償として，分母を処理しなければならない（これはいわゆる非連結ダイアグラムを差し引くことに相当する）．


Hamiltonianが真に時間に依存する非平衡状態では，この方法は働かない．系を平衡から駆動して離した場合，その時間発展の終状態は初期状態と一致するとは限らない．一般に，このような終状態は，相互作用の切り替え方の詳細にも，系の全履歴にも依存する．したがって，式\eqref{contour-O-average}に含まれる時間発展の履歴から，後退部分を取り除くことはできない．これが克服不可能な障害ではないことに最初に気づいたのはSchwinger~\cite{Schwinger}であった．非平衡量子場理論における時間発展は，閉じた時間経路に沿って起こることを受け入れればよい．この経路は，通常の前進経路に加えて後退経路も含む．こうすれば，$t=+\infty$における系の状態を知る必要がなくなる．


それでも，式\eqref{contour-O-average}の時間発展を$t=+\infty$まで延長し，そこから$t$まで戻すと便利である．この操作は恒等的なもので，追加の仮定をまったく要しないことを強調しておく．式\eqref{contour-O-average}の$\hat{\mathcal{O}}$の左に$\hat{\mathcal{U}}_{t,+\infty}\hat{\mathcal{U}}_{+\infty,t}=\hat{1}$を挿入すると，次を得る．
\begin{equation}
\big\langle\hat{\mathcal{O}}(t)\big\rangle=\frac{1}{\Tr\{\hat{\rho}(-\infty)\}}
\Tr\big\{\hat{\mathcal{U}}_{-\infty,+\infty}\hat{\mathcal{U}}_{+\infty,t}
\hat{\mathcal{O}}\hat{\mathcal{U}}_{t,-\infty}\hat{\rho}(-\infty)\big\}\,.
\end{equation}ここではさらに，von Neumann方程式によれば，密度行列のトレースがユニタリ時間発展のもとで変わらないことを用いた．その結果，図\ref{Fig-Contour}に示す閉じた時間経路$\mathcal{C}$に沿った時間発展が得られた．


\begin{figure}
\begin{center}\includegraphics[width=10cm]{figures/fig01.pdf}\end{center}
\caption{閉じた時間経路$\mathcal{C}$．経路の前進枝と後退枝にある点は，離散的な時刻を表す．\label{Fig-Contour}}
\end{figure}


観測量$\hat{\mathcal{O}}$は，経路の前進枝上のある時刻$t$に挿入される．$\hat{\mathcal{O}}$の右に単位演算子$\hat{\mathcal{U}}_{t,+\infty}\hat{\mathcal{U}}_{+\infty,t}=\hat{1}$を挿入すれば，同じように観測量を経路の後退枝上に置くこともできる点に注意しよう．以下で示すように，最も便利なのは，これら二つの等価な表現の和の半分をとることである．また，Hamiltonianにソース項$\hat{H}_{\mathcal{O}}(t)\equiv\hat{H}(t)\pm\hat{\mathcal{O}}\eta(t)/2$を加えて観測量を生成することもできる．ここで正号（負号）は，経路の前進（後退）部分に対応する．次に，Hamiltonian $\hat{H}_{\mathcal{O}}(t)$による経路$\mathcal{C}$に沿った時間発展演算子のトレースとして定義される生成汎関数$Z[\eta]$を計算する．このHamiltonianは二つの枝の間で非対称なので，こうして得られる閉経路の時間発展演算子は恒等的に単位演算子とはならない．観測量の期待値は，汎関数微分$\big\langle\hat{\mathcal{O}}(t)\big\rangle=\delta
Z[\eta]/\delta\eta(t)|_{\eta=0}$によって生成される．まずソース項を省き，分配関数の便利な表現を導こう．
\begin{equation}\label{contour-Z-0}
Z[0]\equiv\frac{\Tr\{\hat{\mathcal{U}}_{\mathcal{C}}\hat{\rho}(-\infty)\}}
{\Tr\{\hat{\rho}(-\infty)\}}=1\,,
\end{equation}ここで$\hat{\mathcal{U}}_{\mathcal{C}}=\hat{\mathcal{U}}_{-\infty,+\infty}
\hat{\mathcal{U}}_{+\infty,-\infty}=\hat{1}$である．前進・後退の対称性を破るソース項については，後で論じる．$Z[0]=1$なので，観測量は平衡系でよりなじみのある形$\big\langle\hat{\mathcal{O}}(t)\big\rangle=\delta\ln
Z[\eta]/\delta\eta(t)|_{\eta=0}$にも同様に書けることに注意されたい．閉じた時間経路を用いる理論では，対数をとるかどうかは任意である．


二つの枝からなる経路に沿って時間発展をたどる必要があるため，非平衡理論は平衡理論より複雑になる．しかし理論の前進・後退対称性に基づいて変数を適切に選べば，この困難を大幅に軽減できる．有利な点もある．一つの枝だけの経路を用いる理論で不可避だった分母$e^{iL}$が存在しないのである．（式(\ref{contour-Z-0})の$\Tr\{\hat\rho(-\infty)\} $を心配する必要はない．実際，この量は相互作用を断熱的に「入れる」前の$t=-\infty$だけに関係する．したがって容易に計算でき，問題を生じることはない．）分母がないことによって，凍結した乱れをもつ系の記述は劇的に簡単になる．観測量の期待値について乱れの平均をとる際，主な障害となるのは分母$e^{iL}$なのである．この障害を克服するために，レプリカ法~\cite{Edwards,Wegner,EfetovLarkinKhmelnitskii}と超対称法~\cite{Efetov,Efetov-book}が考案された．閉じた時間経路の理論には分母が存在しないので，いずれの工夫も必要ない．


%$$$$$$$$$$$$$$$$$$$$$$$$$$$$$$$$$$$$$$$$$$$$$$$$$$$$$$$$$$$$$$$$$$$$$$$$$$$$$$$$$$$$$$$$$$$$$$$$$$$$$$$$$$$
%$$$$$$$$$$$$$$$$$$$$$$$$$$$$$$$$$$$$$$$$$$$$$$$$$$$$$$$$$$$$$$$$$$$$$$$$$$$$$$$$$$$$$$$$$$$$$$$$$$$$$$$$$$$
\sectionlink{ボソン}\label{sec_bosons}
\subsectionlink{分配関数}\label{sec_bosons-1}


最も簡単な多体系として，エネルギー$\omega_0$の単一の量子状態を占有するボソン粒子を考えよう．その第二量子化されたHamiltonianは，次の形をとる．
\begin{equation}\label{boson-H}
\hat{H} = \omega_0\, \hat{b}^{\dagger} \hat{b}\, ,
\end{equation}ここで$\hat{b}^{\dagger}$と$\hat{b}$は，交換関係$[\hat{b},\hat{b}^{\dagger}]=1$を満たすボソンの生成・消滅演算子である．分配関数を次のように定義する．
\begin{equation}\label{boson-Z}
Z=\frac{\Tr\big\{\hat{\mathcal{U}}_{\mathcal{C}}\hat{\rho}\big\}}{\Tr\{\hat{\rho}\}}\,.
\end{equation}すべての外場が経路の前進枝と後退枝でまったく同じであると仮定すれば，$\hat{\mathcal{U}}_{\ \mathcal{C}}=1$であり，したがって$Z=1$である．初期密度行列$\hat{\rho}=\hat{\rho}(\hat{H})$はHamiltonianのある演算子値関数である．導出を簡単にするため，これを平衡密度行列$\hat{\rho}_0 = \exp\{-\beta(\hat{H}-\mu\hat
{N})\}=\exp\{-\beta(\omega_0-\mu)\hat{b}^\dagger\hat{b}\}$に選んでもよい．後の時刻には任意の外部摂動を入れたり切ったりできるので，初期密度行列を平衡のものに選んでも，非平衡ダイナミクスを扱う妨げにはならない．平衡初期密度行列について，次を得る．
\begin{equation}\label{boson-density-trace}
\Tr\{\hat{\rho}_0\}=\sum\limits_{n=0}^\infty
e^{-\beta(\omega_0-\mu)n}= [1-\rho(\omega_0)]^{-1}\, ,
\end{equation}ここで$\rho(\omega_0)=e^{-\beta(\omega_0-\mu)}$である．重要なのは，一般に$\Tr\{\hat{\rho}\}$が相互作用にも乱れにも依存しない定数であることである．実際，相互作用と乱れはいずれも，経路の前進（後退）部分において，$t=-\infty$より後（前）のある時刻に入れられる（切られる）と仮定される．したがって，この定数は混乱を招くことなく，しばしば省略される．


次に，図\ref{Fig-Contour}に示すように$t_1=t_{2N}=-\infty$および$t_{N}=t_{N+1}=+\infty$となるよう，経路$\mathcal{C}$を長さ$\delta_t$の$(2N-2)$個の時間刻みに分割する．そして，過完備なコヒーレント状態基底における単位演算子の分解\footnote{複素数$\phi$でパラメータ化されるボソンのコヒーレント状態$|\phi\rangle$（$\langle\phi|\,$）は，消滅（生成）演算子の右（左）固有状態として，$\hat{b}|\phi\rangle =\phi|\phi\rangle$（$\langle\phi|\hat{b}^\dagger= \langle\phi|\bar\phi\, $）により定義される．Hamiltonianのような\emph{正規順序化された}演算子の行列要素は，$\langle\phi|\hat
H(\hat{b}^\dagger,\hat{b})|\phi'\rangle = H(\bar\phi,\phi')
\langle\phi|\phi'\rangle$の形をとる．二つのコヒーレント状態の重なりは$\langle\phi|\phi'\rangle=\exp\{\bar\phi\phi'\}$である．コヒーレント状態基底は過完備なので，演算子$\hat A$のトレースは重みを付けて$\Tr\{\hat A\}=\pi^{-1}
\int\!\!\!\int d(\mbox{Re} \phi)\, d(\mbox{Im} \phi)\,
e^{-|\phi|^2}\, \langle\phi|\hat A|\phi\rangle
$と計算する．}~\cite{NegeleOrland}
\begin{equation}\label{boson-unity}
\hat{1}=\iint\frac{\mathrm{d}(\mathrm{Re}\phi_{j})\mathrm{d}(\mathrm{Im}\phi_{j})}
{\pi}\ e^{-|\phi_{j}|^{2}}|\phi_{j}\rangle\langle\phi_{j}|
\end{equation}を，経路上の各点$j=1,2,\ldots, 2N$に挿入する．たとえば$N=3$では，$\Tr\{\hat{\mathcal{U}}_{\ \mathcal{C}}\hat{\rho}_0\}$の式に次の列が現れる（右から左へ読む）．
\begin{equation}
  \langle\phi_6|\hat{\mathcal{U}}_{-\delta_t}|\phi_5\rangle
  \langle\phi_5|\hat{\mathcal{U}}_{-\delta_t}|\phi_4\rangle
  \langle\phi_4|\hat 1|\phi_3\rangle
  \langle\phi_3|\hat{\mathcal{U}}_{+\delta_t}|\phi_2\rangle
  \langle\phi_2|\hat{\mathcal{U}}_{+\delta_t}|\phi_1\rangle
  \langle\phi_1|\hat{\rho}_0|\phi_6\rangle\, ,
\end{equation}ここで$\hat{\mathcal{U}}_{\pm \delta_t}$は，時間の正（負）の向きに時間間隔$\delta_t$だけ発展させる演算子である．その行列要素は次で与えられる．
\begin{eqnarray}\label{boson-matrix-element}
\left\langle\phi_{j+1}\left|\hat{\mathcal{U}}_{\pm\delta_t}\right|\phi_j\right\rangle\equiv
\left\langle \phi_{j+1}\left| e^{\mp i\hat H(\hat{b}^\dagger,
\hat{b})\delta_t}\right|\phi_{j}\right\rangle\approx \left\langle
\phi_{j+1}\left|\big( 1\mp i\hat H(\hat{b}^\dagger,
\hat{b}\big)\delta_t\right |\phi_{j}\right\rangle\nonumber\\=
\big\langle \phi_{j+1}|\phi_{j}\big\rangle\big(1\mp  i
H(\bar\phi_{j+1},\phi_j)\delta_t\big) \approx \big\langle
\phi_{j+1}| \phi_{j}\big\rangle\, e^{\mp i
H(\bar\phi_{j+1},\phi_j)\delta_t}\, ,
\end{eqnarray}ここで近似等号は$\delta_t$の一次まで成り立つ．明らかに，この結果は単純な例\eqref{boson-H}に限られず，任意の\emph{正規順序化された}Hamiltonianに対して成り立つ．$t_N$と$t_{N+1}$の間には時間発展演算子が挿入されていないことに注意しよう．実際，これら二点は物理的に区別できず，この時間間隔では系は発展しない．コヒーレント状態の性質$\langle\phi|\phi'\rangle=\exp\{\bar\phi\phi'\}$および$\langle\phi| e^{-\beta(\omega_0-\mu) \hat{b}^\dagger
\hat{b}}|\phi'\rangle = \exp\left\{\bar\phi\phi' \rho(\omega_0)
\right\}$を用い，経路に沿ったすべての行列要素をまとめると，分配関数\eqref{boson-Z}について次を得る．
\begin{equation}\label{boson-Z-discrete}
Z=\frac{1}{\Tr\{\hat{\rho}_{0}\}}\iint\prod\limits^{2N}_{j=1}
\left[\frac{\mathrm{d}(\R\phi_{j})\mathrm{d}(\I\phi_{j})}{\pi}\right]
\exp\left(i\sum\limits^{2N}_{j,j\,'=1}\bar{\phi}_{j}\,
G^{-1}_{jj\,'}\,\phi_{j\,'}\right)\, ,
\end{equation}ここで$2N \times 2N$行列$iG^{-1}_{jj\,'}$は，次の行列を表す．
\begin{equation}\label{boson-G-matrix}
  iG^{-1}_{jj\,'}\equiv\
\left[\begin{array}{ccc|ccc}
 \,\,-1\,\,   &\,\,\,\,      &\,\,\,\,    &\,\,\,\,   &\,\,\,\,   &\,\,\,\,   \rho(\omega_0) \\
 \,\, 1\!-\!h\,\, & -1   &    &   &   &                 \\
      &  \,\,1\!-\!h \,\,& \,\,\,-1\,\,\, &   &   &                 \\ \hline
%---------------------------------------------------------
     &     &  1 & -1 &    &                 \\
     &     &    &\,\,1\!+\!h\,\, & -1 &                 \\
     &     &    &    & \,\,1\!+\!h\,\,&  -1
 \end{array} \right]\, ,
\end{equation}また$h\equiv i\omega_0\delta_t$である．このような行列の行列式は容易に評価できる．
\begin{equation}\label{boson-determinant}
\mathrm{Det}\big[i\hat{G}^{-1}\big]=(-1)^{2N} - \rho(\omega_0)(1-h^2)^{N-1}\approx
1-\rho(\omega_0)\, e^{(\omega_0\delta_t)^2(N-1)}\to 1-
\rho(\omega_0) \, ,
\end{equation}ここでは，$N\to \infty$なら$\delta_t^2N\to 0$となることを用いた（実際，仮定は$\delta_tN \to \mathrm{const}$であった）．式\eqref{boson-Z-discrete}のGaussian積分が行列$i\hat{G}^{-1}$の行列式の逆数に等しいこと（\ref{app_Gaussian}参照）と，式\eqref{boson-density-trace}を用いると，次を得る．
\begin{equation}\label{boson-Z-norm}
Z=\frac{\mathrm{Det}^{-1}\big[i\hat{G}^{-1}\big]}{\Tr\{\hat{\rho}_0\}}
= 1\, ,
\end{equation}もちろん，これは期待どおりの結果である．この規格化の恒等式を保つには，離散行列\eqref{boson-G-matrix}の右上要素を残すことが不可欠である点に注意されたい．


ここで極限$N\to \infty$をとり，連続表記$\phi_j\to \phi(t)$を用いて分配関数を形式的に次のように書くことができる．
\begin{equation}\label{boson-Z-1}
Z=\int\D[\bar{\phi}\phi]\,
\exp\left(iS[\bar{\phi},\phi]\right)=\int\D[\bar{\phi}\phi]
\exp\left(i\int_{\mathcal{C}}\d t \,
\big[\bar{\phi}(t)\,\hat{G}^{-1}\phi(t)\big] \right)\,,
\end{equation}ここで，式\eqref{boson-Z-discrete}--\eqref{boson-G-matrix}によれば作用は次で与えられる．
\begin{equation}\label{boson-S}
S[\bar{\phi},\phi] =
\sum\limits_{j=2}^{2N}\left[i\bar\phi_j\,\frac{\phi_j-\phi_{j-1}}{\delta
t_j} -\omega_0\bar\phi_j\phi_{j-1}\right]\delta
t_j+i\,\bar\phi_1\Big[\phi_1-\rho(\omega_0)\phi_{2N}\Big] \, ,
\end{equation}ただし$\delta t_j\equiv t_j-t_{j-1}=\pm \delta_t$である．したがって，演算子$\hat G^{-1}$の連続表示は次のようになる．
\begin{equation}\label{boson-G-continious}
\hat G^{-1}= i\partial_t - \omega_0 \,.
\end{equation}この連続表記は，大きな離散行列\eqref{boson-G-matrix}を表す略記にすぎないことを忘れてはならない．とりわけ，分布関数についての情報を含む行列の右上要素（式\eqref{boson-S}の最後の項）は，連続表記\eqref{boson-G-continious}では一見失われている．


閉じた時間経路に沿う積分を避けるには，ボソン場$\phi(t)$を，時間経路の前進部分と後退部分にそれぞれ存在する二つの成分$\phi_{+}(t)$と$\phi_{-}(t)$に分けると便利である．すると，連続表示の作用は次のように書き直せる．
\begin{equation} \label{boson-S-plus-minus}
S[\bar{\phi},\phi]=\int_{-\infty}^{+\infty} \d t\,
\big[\bar\phi_{+}(t)(i\partial_t - \omega_0) \phi_{+}(t)-
\bar\phi_{-}(t)(i\partial_t - \omega_0) \phi_{-}(t)\big]\, ,
\end{equation}ここで相対的な負符号は，経路の後退部分における時間積分の向きが逆であることに由来する．ここでもまた，連続表記はやや誤解を招きやすい．実際，この表記は場$\phi_{+}(t)$と$\phi_{-}(t)$がまったく相関していないという誤った印象を与える．実際には，離散行列\eqref{boson-G-matrix}に非零の非対角ブロックが存在するため，両者は相関している．したがって，適切な正則化を自動的に取り込む連続表示を構成することが望ましい．これを以下の節で実現する．まず，Green関数について論じよう．


%$$$$$$$$$$$$$$$$$$$$$$$$$$$$$$$$$$$$$$$$$$$$$$$$$$$$$$$$$$$$$$$$$$$$$$$$$$$$$$$$$$$$$$$$$$$$$$$$$$$$$$$$$$$
\subsectionlink{Green関数}\label{sec_bosons-2}


Gaussian積分の基本的な性質（\ref{app_Gaussian}参照）によれば，二つのボソン場の相関関数は次で与えられる．
\begin{equation}\label{boson-G-fun}
\big\langle\phi_{j}\,\bar\phi_{j\,'}\big\rangle \equiv \int
\D[\bar{\phi}\phi] \,\,\phi_{j}\,\bar\phi_{j\,'} \exp\left(i
\sum\limits_{j,j\,'=1}^{2N}
\bar\phi_{j}\,G^{-1}_{jj\,'}\,\phi_{j\,'}\right)= iG_{jj\,'} \, .
\end{equation}平衡理論~\cite{NegeleOrland}における対応する定義と比べて，因子$Z^{-1}$がないことに注意しよう．実際，ここでの構成では$Z=1$である．一見小さく思えるこの違いが，不規則系の理論では主要な点となる（凍結した乱れをもつフェルミオンを扱う第\ref{sec_NLSM}節の詳しい議論を参照）．式\eqref{boson-G-matrix}の離散行列の逆行列をとると，次を得る．
\begin{equation}\label{boson-G-matrix-inverted}
  iG_{jj\,'}=\frac{1}{1-\rho}
 \left[\begin{array}{lcc|ccc}
 1   &  \,\,\rho\,e^{h}\,\,    & \,\, \rho\,e^{2h}\,\,  &\,\, \rho\,e^{2h}\,\,  &\,\, \rho\,e^{h}\,\,  &\,\, \rho\, \\
 e^{-h}\,\, & 1  &  \rho\,e^{h}   & \rho\,e^{h}   & \rho  &  \,\,   \rho\,e^{-h}  \,           \\
 e^{-2h}\,\,    &  \,e^{-h} & 1 & \rho  & \rho\,e^{-h}   &   \rho\,e^{-2h}               \\ \hline
%---------------------------------------------------------
  e^{-2h}   &   \,e^{-h}   &  1 & 1 &\,\, \rho\,e^{-h}\,\,    &      \rho\,e^{-2h}            \\
   e^{-h}   &   1  &  \,e^{h}   &\,e^{h} & 1 &  \rho\,e^{-h}                \\
   1  &  \,e^{h}    &  \,e^{2h}   & \,e^{2h}    & \,e^{h}&  1
 \end{array} \right]\,,
\end{equation}ここで$\rho\equiv \rho(\omega_0)$であり，式\eqref{boson-determinant}の後の議論に従い，$(1\pm h)^j\, \approx\,
e^{\pm jh}\,$および$(1-h^2)^j\approx 1$とおいた．場$\phi_{j\pm}$を用いると（以下では$j=1,\ldots, N$とするため，上の$2N\times
2N$行列には$1,\ldots, N-1,N,N,N-1,\ldots, 1$という添字を付ける），対応する相関関数は次のように書ける．
\begin{subequations}\label{boson-corr-fun}
\begin{equation}
\hskip-3cm \langle \phi_{j+}\bar{\phi}_{j\,'-}\rangle\equiv
iG^{<}_{jj\,'}= n_B\, \exp\{-(j-j\,')h\}\,,
\end{equation}
\begin{equation}
\hskip-2.2cm \langle \phi_{j-}\bar\phi_{j\,'+}\rangle\equiv
iG^{>}_{jj\,'}=(n_B+1)\, \exp\{-(j-j\,')h\}\,,
\end{equation}
\begin{equation}
\langle\phi_{j+}\bar\phi_{j\,'+}\rangle\equiv
iG^{\mathbb{T}}_{jj\,'}= {1\over 2}\,
\delta_{jj\,'}+\theta(j-j\,')iG^{>}_{jj\,'}+
\theta(j\,'-j)iG^{<}_{jj\,'}\,,
\end{equation}
\begin{equation}
\langle\phi_{j-}\bar\phi_{j\,'-}\rangle \equiv
iG^{\widetilde{\mathbb{T}}}_{jj\,'}={1\over 2}\,\delta_{jj\,'}+
\theta(j\,'-j)iG^{>}_{jj\,'} +\theta(j-j\,')iG^{<}_{jj\,'}\,,
\end{equation}\end{subequations}
ここでボソン占有数$n_B$は$n_B(\omega_0)
\equiv \rho/(1-\rho)$を表し，記号$\mathbb{T}$（$\widetilde{\mathbb{T}})$）は時間順序（反時間順序）を表す．階段関数$\theta(j)$は$\theta(0)=1/2$となるように定義するので，$\theta(j)+\theta(-j)\equiv 1$である．


明らかに，上で定義した四つのGreen関数はすべて独立ではない．実際，直接調べれば次がわかる．
\begin{subequations}
\begin{equation}\label{boson-relation}
G^{\mathbb{T}}
+G^{\widetilde{\mathbb{T}}}-G^{>}-G^{<}=-i\delta_{jj\,'}\,,\qquad\,\,
\end{equation}
\begin{equation}\label{boson-relation-1}
G^{\mathbb{T}}-G^{\widetilde{\mathbb{T}}}=\mbox{sign}(j-j\,')\left(G^{>}-
G^{<}\right)\,,
\end{equation}\end{subequations}
ここで$\mbox{sign}(j)=\theta(j)-\theta(-j)$である．これらの関係を明示的に利用するため，場の線形変換を行いたい．これは次のKeldysh回転によって実現される．
\begin{equation}\label{boson-rotation}
\phi^{cl}_{j}={1\over \sqrt{2}}\big(\phi_{j+} +\phi_{j-}\big) \, ,
\,\,\,\,\,\,\,\,\, \phi^{q}_{j}={1\over\sqrt{2}}\big(\phi_{j+}-
\phi_{j-}\big)\,,
\end{equation}共役な場にも同様の変換を施す．上付き添字「cl」と「q」は，それぞれ場の古典成分と量子成分を表す．この名称の根拠はまもなく明らかになる．まず，式(\ref{boson-corr-fun}a)--(\ref{boson-corr-fun}d)を簡単に代数操作すると，次を得る．
\begin{equation}\label{boson-G-fun-1}
-i\big\langle\phi^{\alpha}_{j}\,\bar\phi^{\,\beta}_{j\,'}\big\rangle
= \left(\begin{array}{cc} G^K_{jj\,'} & G^{R}_{jj\,'} \\ \frac{}{} &
\frac{}{}
\\G^{A}_{jj\,'} & -{i\over 2}\,\delta_{jj\,'}
\end{array}\right)\, ,
\end{equation}ここで，以下では$\alpha, \beta = (cl,q)$とする．この行列の$(q,q)$要素の具体的な形は，恒等式\eqref{boson-relation}の表れである．上付き添字$R,A$および$K$は，それぞれGreen関数の遅延成分，先進成分，Keldysh成分を表す．これら三つのGreen関数は，Keldysh法の基本的な対象である．それぞれ次のように定義される．
\begin{subequations}\label{boson-G-fun-RAK}
\begin{equation}
\hskip-1.2cm
G^{R}_{jj\,'}=-i\big\langle\phi^{cl}_{j}\,\bar\phi^{q}_{j\,'}\big\rangle
=\theta(j-j\,')\left(G^{>}_{jj'}-G^{<}_{jj'}\right)
=-i\theta(j-j\,')\, e^{-(j-j\,')h}\, ,
\end{equation}
\begin{equation}
\hskip-1.5cm
G^{A}_{jj\,'}=-i\big\langle\phi^{q}_{j}\,\bar\phi^{cl}_{j\,'}\big\rangle
= \theta(j\,'-j)\left(G^{<}_{jj'}-G^{>}_{jj'}\right)=i\theta(j\,'-j)
e^{-(j-j\,')h}\,,
\end{equation}
\begin{equation}
G^{K}_{jj\,'}=-i\big\langle\phi^{cl}_{j}\,\bar\phi^{cl}_{j\,'}\big\rangle =
-{i\over 2}\,\delta_{jj\,'}+G^{>}_{jj'}+G^{<}_{jj'}=-{i\over
2}\,\delta_{jj\,'}-i\left(2n_B+1\right) e^{-(j-j\,')h}\,.
\end{equation}\end{subequations}
定義により$\left[G^<\right]^\dagger =-G^>$であるから［式\eqref{boson-corr-fun}参照］，次がわかる．
\begin{equation}\label{boson-G-conjugation}
G^{A} =\big[G^R\big]^{\dagger}\,, \hskip 2cm
G^K=-\big[G^K\big]^\dagger\,.
\end{equation}遅延（先進）Green関数は，時間領域で下（上）三角行列になる．三角行列をいくつ掛け合わせても再び三角行列になるため，次の簡単な規則を得る．
\begin{subequations}\label{bosons-G-traces}
\begin{equation}
G_1^{R}\circ G_2^{R}\circ \ldots \circ G_l^{R}  = G^R \,,
\end{equation}
\begin{equation}
G_1^{A}\circ G_2^{A}\circ \ldots \circ G_l^{A}  = G^A \,,
\end{equation}\end{subequations}
ここで丸印の積は時間領域での畳み込みを意味する（すなわち中間時刻に関する積分を含む）．


ここでGreen関数の連続極限（$N\to \infty$かつ$N\delta_t\to \mbox{const}$）をとることができる．そのために$t_j=j\delta_t$と定義し，$\exp\{-(j-j')h\}\to
\exp\{-i\omega_0(t-t')\}$に注意する．より非自明なのは，式\eqref{boson-G-fun-1}および\eqref{boson-G-fun-RAK}に現れる因子$\delta_{jj\,'}$を連続極限で省略できることである．その理由は二つある．（i）すべての観測量はGreen関数の\emph{非対角}要素で与えられる．たとえば時刻$t_j$における平均占有数は$\langle
n_B(t_j)\rangle =iG^{\mathbb{T}}_{jj+1}=iG^<_{jj+1}$である．（ii）計算途中の式には，$\delta_t^2\sum_{j,j\,'} \delta_{jj\,'}G_{j\,'j}\to
\delta_t^2N\to 0$という形の多重和（積分）が含まれる．その結果，上で導いた関係式の正しい連続極限は次のようになる．
\begin{equation}\label{boson-G-fun-2}
-i\big\langle\phi^{\alpha}(t)\,\bar\phi^{\,\beta}(t')\big\rangle =
G^{\alpha\beta}(t,t') =\left(\begin{array}{cc}
G^K(t,t') & G^{R}(t,t') \\
G^{A}(t,t') & 0
\end{array}\right)\, ,
\end{equation}ここで
\begin{subequations}\label{boson-G-fun-RAK-cont}
\begin{equation}
\hskip-2.5cm G^{R}=-i\theta(t-t')\,e^{-i\omega_0(t-t')}\to
(\epsilon-\omega_0+i0)^{-1}\,,
\end{equation}
\begin{equation}
\hskip-2.5cm G^{A}=\phantom{-}i\theta(t'-t)\,e^{-i\omega_0(t-t')}\to
(\epsilon-\omega_0- i0)^{-1}\,,
\end{equation}
\begin{equation}
G^{K}=-i\left[2n_B(\omega_0)+1\right] e^{-i\omega_0(t-t')}\to-2\pi i
[2n_B(\epsilon)+1]\delta(\epsilon-\omega_0)\,.
\end{equation}\end{subequations}
である．三つのGreen関数それぞれについて，$t-t'$に関するFourier変換も示した．これらのGreen関数の重要な性質は，次の関係である［式\eqref{boson-G-fun-RAK}参照］．
\begin{equation}\label{boson-GR-plus-GA}
G^R(t,t)+G^A(t,t) = 0\, .
\end{equation}三つのGreen関数に図式的な表現を導入すると便利である．そのため，場の古典成分を実線で，量子成分を破線で表そう．すると，遅延Green関数は実線から矢印を介して破線につながる線，先進Green関数は破線から矢印を介して実線につながる線，Keldysh Green関数は実線から矢印を介して実線につながる線で表される（図\ref{Fig-boson-G}参照）．$\langle \phi^q \bar\phi^q\rangle$のGreen関数を表すはずの，破線から矢印を介して破線につながる線は，連続極限では存在しないことに注意しよう．矢印は$\phi^{\alpha}$から$\bar\phi^{\,\beta}$への向きを示す．
\begin{figure}
\begin{center}\includegraphics[width=10cm]{figures/fig02.pdf}\end{center}
\caption{$G^R$，$G^A$，$G^K$の図式表現．実線は場の古典成分$\phi^{cl}$を，破線は量子成分$\phi^{q}$を表す．\label{Fig-boson-G}}
\end{figure}


遅延成分と先進成分は，スペクトルに関する情報だけを含み，占有数には依存しない一方，Keldysh成分は占有数に依存することに注意しよう．熱平衡では$\rho=e^{-\beta\epsilon}$であり，また$n_B=(e^{\beta\epsilon}-1)^{-1}$なので，次を得る．
\begin{equation}\label{boson-G-FDT}
G^{K}(\epsilon)=\left[G^{R}(\epsilon)-G^{A}(\epsilon)\right]
\mathrm{coth}\frac{\epsilon}{2T}\, .
\end{equation}最後の式は，\emph{揺動散逸定理}（FDT）の内容を表す．もちろんFDTは熱平衡の一般的な性質であり，ここで考えた単純な例に限られるものではない．これは，平衡における応答関数と相関関数の間に厳密な関係があることを意味する．


一般に，反HermiteなKeldysh Green関数［式\eqref{boson-G-conjugation}参照］を，Hermite行列$F=F^\dagger$を用いて次のようにパラメータ化すると便利である．
\begin{equation}\label{boson-G-FDT-general}
G^{K}=G^{R}\circ F-F\circ G^{A}  \, ,
\end{equation}ここで$F=F(t,t')$であり，丸印の積は畳み込みを意味する．行列$F$のWigner変換（後述）$\mathbf{f}(\tau,\epsilon)$を，\emph{分布関数}と呼ぶ．熱平衡では$\mathbf{f}(\epsilon) = \mbox{coth}(\epsilon/2T)$である［式\eqref{boson-G-FDT}参照］．


%$$$$$$$$$$$$$$$$$$$$$$$$$$$$$$$$$$$$$$$$$$$$$$$$$$$$$$$$$$$$$$$$$$$$$$$$$$$$$$$$$$$$$$$$$$$$$$$$$$$$$$$$$$$
\subsectionlink{Keldysh作用と因果性}\label{sec_bosons-3}


相関関数\eqref{boson-G-fun-2}および\eqref{boson-G-fun-RAK-cont}を正しく再現する，$\phi^{cl},\phi^{q}$で書かれた連続表示の作用がほしい．そのため，相関関数の行列\eqref{boson-G-fun-2}を形式的に逆行列にし，Gaussian作用に用いる．
\begin{equation}\label{boson-S-1}
S[\phi^{cl},\phi^{q}]=\iint_{-\infty}^{+\infty}\d t\, \d t'\,
\left(\bar{\phi}^{cl}_t,\bar{\phi}^{q}_t\right)
\left(\begin{array}{cc}
0   & \big[G^{-1}_{t,t\,'}\big]^{A}  \\
  \big[G^{-1}_{t,t\,'}\big]^{R}  & \big[G^{-1}_{t,t\,'}\big]^K
\end{array}\right)
\left(\begin{array}{c} \phi^{cl}_{t\,'} \\ \phi^{q}_{t\,'}
\end{array}\right) ,
\end{equation}ここで
\begin{subequations}\label{boson-G-fun-2-inverted}
\begin{equation}
\big[G^{-1}\big]^{R(A)} = \big[G^{R(A)}\big]^{-1} = \epsilon-\omega_0\pm i0\to
\delta_{t,t'}\left(i\partial_t-\omega_0 \pm i0\right)\,,
\end{equation}
\begin{equation}
\hskip-3.3cm \big[G^{-1}\big]^K = \big[G^R\big]^{-1}\circ F -F \circ
\big[G^A\big]^{-1}\, ,
\end{equation}\end{subequations}
である．ここでは$\epsilon$のFourier変換が$
\delta_{t,t'}i\partial_t$であることを用い，最後の行ではパラメータ表示\eqref{boson-G-FDT-general}を用いた．実際の離散行列の作用\eqref{boson-Z-discrete}--\eqref{boson-G-matrix}を，式\eqref{boson-rotation}に従って$\phi^{cl},\phi^{q}$に変換しても，式\eqref{boson-S-1}の構造には\emph{ならない}ことを指摘しておく必要がある．作用\eqref{boson-S-1}は，Gaussian積分の規則に従って相関関数の連続極限を再現するために考案された，形式的な構成とみなすべきである．しかし，次の意味で完全に自己無撞着である．（i）正則化のために離散表示に立ち戻る必要がない．（ii）摂動的な繰り込みのあらゆる次数において，その一般的構造が保たれる．


ここで作用\eqref{boson-S-1}の主要な特徴をまとめよう．よりよい用語がないので，これを\emph{因果性の構造}と呼ぶ．
\begin{itemize}
\item $cl-cl$成分は零である．これは，純粋に古典的な場の配置（$\phi^q=0$）に対して作用が零になることを反映する．実際，この場合$\phi_+=\phi_-$であり，経路の前進部分の作用は後退部分の作用によって打ち消される（連続極限で省略できる境界項は除く）．したがって，きわめて一般的に次が成り立つ．
\begin{equation}\label{boson-causality}
S\big[\phi^{cl},0\big] = 0\, .
\end{equation}明らかに，この主張は式\eqref{boson-S-1}の形のGaussian作用だけに限定されるものではない．


\item $cl-q$成分と$q-cl$成分は互いにHermite共役であり，時間表示では上三角および下三角（先進および遅延）行列である．この性質によって，応答関数の因果性が保証されるとともに，$cl-cl$成分が摂動的な繰り込みから保護される（後述）．理論の整合性には，関係式\eqref{boson-GR-plus-GA}が必要である．


\item $q-q$成分は反Hermite行列である［式\eqref{boson-G-conjugation}参照］．ここでの例では$\big[G^K\big]^{-1}=i0F$であり，$F$は正定値のスペクトルをもつHermite行列である．これが汎関数積分の収束を担う．また，分布関数についての情報も保持している．
\end{itemize}


%$$$$$$$$$$$$$$$$$$$$$$$$$$$$$$$$$$$$$$$$$$$$$$$$$$$$$$$$$$$$$$$$$$$$$$$$$$$$$$$$$$$$$$$$$$$$$$$$$$$$$$$$$$$
\subsectionlink{自由ボソン場}\label{sec_bosons-4}


この構成全体を，複数の自由度をもつボソン系へ一般化することは容易である．状態が添字$\mathbf{k}$で指定されるとしよう．この添字は，たとえば運動量ベクトルであってよい．各状態のエネルギーは関数$\omega_\mathbf{k}$で与えられる．たとえば$\omega_{\mathbf{k}}=\mathbf{k}^2/(2m)$であり，ここで$m$はボソン原子の質量である．次に，各状態$\mathbf{k}\,$に対して複素場の二重項（古典成分と量子成分）$(\phi^{cl}(\mathbf{k},t), \phi^{q}(\mathbf{k},t))$を導入し，添字$\bf{k}$に関する和を含めて作用を式\eqref{boson-S-1}の形に書く．非平衡では，Keldysh成分は添字$\mathbf{k}$について非対角になりうる．すなわち$F=F(\mathbf{k},\mathbf{k}';t,t')$である．一方，遅延（先進）成分は$[G^{R(A)}]^{-1}=i\partial_t -\omega_\mathbf{k}$という簡単な形をとる．


$\mathbf{k}$が運動量である場合，実空間へFourier変換して$(\phi^{cl}(\mathbf{r},t), \phi^q(\mathbf{r},t))$を扱うと見通しがよい．時空をまとめた添字$x=(\mathbf{r},t)$を導入すると，自由な複素ボソン場（原子）の作用は次のように書ける．
\begin{equation}\label{boson-S-2}
S_{0}[\phi^{cl},\phi^{q}] = \iint\d x\,\d
x'\,\big(\bar\phi^{cl}_x,\bar\phi^{q}_x\big) \left(\begin{array}{cc}
0   & \big[G^{A}_{x,x'}\big]^{-1}  \\
  \big[G^{R}_{x,x'}\big]^{-1}  & \big[G^{-1}_{x,x'}\big]^K
\end{array}\right)
\left(\begin{array}{c} \phi^{cl}_{x\,'} \\ \phi^{q}_{x\,'}
\end{array}\right),
\end{equation}ここで連続表記では
\begin{equation}\label{boson-G-Real-Space}
\big[G^{R(A)}\big]^{-1}(x,x\,')=\delta(x-x\,')\left( i\partial_{t}
+{1\over 2m}
\partial_{\mathbf{r}}^2 +\mu\right)\, ,
\end{equation}であり，離散表示では，時間について（空間についてではなく）下（上）三角行列である．自由場の$\big[G^{-1}\big]^K$成分は，時間の境界項に由来する正則化因子にすぎない．これは一般には$x$および$x'$について非局所的だが，純粋な境界項なので，しばしば省略される．ここでは，相関関数行列の逆$\hat{G}$が因果性の構造\eqref{boson-G-fun-2}をもたなければならないことを思い出すために残しておく．式\eqref{boson-G-Real-Space}には化学ポテンシャル$\mu$を導入した．これは，有効Hamiltonian $\hat{H}-\mu\hat{N}$を考える場合も想定したものであり，$\hat{N}$は全粒子数演算子である．この新しい項は，Lagrange乗数$\mu$を用いて特定の粒子数を課すための項とみなせる．実ボソン場については\ref{app_SPQM}を参照されたい．


%$$$$$$$$$$$$$$$$$$$$$$$$$$$$$$$$$$$$$$$$$$$$$$$$$$$$$$$$$$$$$$$$$$$$$$$$$$$$$$$$$$$$$$$$$$$$$$$$$$$$$$$$$$$
%$$$$$$$$$$$$$$$$$$$$$$$$$$$$$$$$$$$$$$$$$$$$$$$$$$$$$$$$$$$$$$$$$$$$$$$$$$$$$$$$$$$$$$$$$$$$$$$$$$$$$$$$$$$
\sectionlink{ボソンの衝突と運動論方程式}\label{sec_int_bos}
\subsectionlink{相互作用}\label{sec_int_bos-1}


ボソン原子の短距離二体衝突は，局所的な\emph{四ボソン}Hamiltonian $\hat{H}_{\mathrm{int}}=\lambda\sum_{\mathbf{r}}
\hat{b}^{\dagger}_{\mathbf{r}} \hat{b}^{\dagger}_{\mathbf{r}}
\hat{b}_{\mathbf{r}} \hat{b}_{\mathbf{r}}$によって記述される．ここで添字$\mathbf{r}$は空間位置を「数え上げる」．相互作用定数$\lambda$は，よく用いられる$s$波散乱長$a_s$と，$\lambda=4\pi a_s/m$という関係にある（文献\cite{Leggett}参照）．連続表示のKeldysh作用における対応する項は，次の形をとる．
\begin{equation}\label{int-bos-S}
S_{\mathrm{int}}[\phi_{+},\phi_{-}]=-\lambda\int\d\mathbf{r}\int_{\cal
C}\d t\,(\bar\phi
\phi)^2=-\lambda\int\d\mathbf{r}\int_{-\infty}^{+\infty}\d t\,
\big[(\bar{\phi}_{+}\phi_{+})^2-(\bar{\phi}_{-}\phi_{-})^2\big]\,.
\end{equation}遠い過去$t=-\infty$には相互作用が存在しないことを忘れてはならない（一方，未来$t=+\infty$には存在する）．相互作用は，前進枝では断熱的に入り，後退枝では断熱的に切れるものとする．これにより，行列\eqref{boson-G-matrix}の非対角ブロックはそのまま保たれる．相互作用が変えるのは，$t=-\infty$から離れたところにある時間発展演算子\eqref{boson-matrix-element}の行列要素だけである．また，離散時間表示では，Keldysh経路${\cal C}$に沿って，場$\bar\phi$が場$\phi$より一時間刻み$\delta_t$だけ\emph{後}にとられることも覚えておくとよい．Keldysh回転\eqref{boson-rotation}を行うと，次を得る．
\begin{equation}\label{int-bos-S-1}
S_{\mathrm{int}}[\phi^{cl},\phi^{q}]=-\lambda\int\d\mathbf{r}
\int^{+\infty}_{-\infty}\d t\,\Big[ \bar{\phi}^{cl}\bar{\phi^{q}}
\big[\big(\phi^{cl}\big)^2+\big(\phi^{q}\big)^{2}\big]+c.c.\Big] \,,
\end{equation}ここで$c.c.$は第一項の複素共役を表す．衝突作用\eqref{int-bos-S-1}は，明らかに因果性条件\eqref{boson-causality}を満たす．ダイアグラムで表すと，作用\eqref{int-bos-S-1}は図\ref{Fig-Phi4}に示す二種類の頂点を生成する（さらに，矢印の向きを逆にした二つの複素共役頂点もある）．一つは三つの古典場（実線）と一つの量子場（破線）をもち，もう一つは一つの古典場と三つの量子場をもつ．
\begin{figure}
\begin{center}\includegraphics[width=10cm]{figures/fig03.pdf}\end{center}
\caption{$|\phi|^4$理論の二つの相互作用頂点の図式表現．すべての矢印の向きを逆にした，二つの複素共役頂点も存在する．\label{Fig-Phi4}}
\end{figure}


作用に衝突項を加えても，基本的な規格化$Z=1$が破られないことを示そう．そのため，$\exp(iS_{\mathrm{int}})$を$\lambda$のべきで展開し，各項をGaussian作用\eqref{boson-S-2}で平均する．規格化$Z=1$が衝突によって変わらないことを示すには，$\langle
S_{\mathrm{int}}\rangle = \langle S_{\mathrm{int}}^{\,2}\rangle =
\ldots=0$を示せばよい．Wickの定理を適用すると，$\lambda$の一次の項として$\big\langle\bar\phi^q
\bar\phi^{cl}\big(\phi^{cl}\big)^2 +c.c.\big\rangle \sim \big[
G^R(t,t)+G^A(t,t)\big] G^K(t,t)=0$および$\big\langle \bar\phi^{q}
\bar\phi^{cl} \big(\phi^{q}\big)^{2}+c.c\big\rangle=0$を得る．第一項は恒等式\eqref{boson-GR-plus-GA}によって消え，第二項は$\big\langle
\phi^{q}\bar\phi^{q}\big\rangle=0$のために消える（$\big\langle\phi^{q}_{j}\bar\phi^{q}_{j\,'}\big\rangle=-i\delta_{jj\,'}/2\neq
0$である離散表示\eqref{boson-G-fun-1}に立ち戻っても，この項は$\sum_{jj\,'}\delta_{jj\,'}(G^A_{j\,'j}+G^R_{j\,'j})=0$で与えられるため，やはり恒等的に零である．式\eqref{boson-GR-plus-GA}参照）．$\lambda$の二次の項は二つの群に分かれる．第一は
\[\fitdisplay{\big\langle\bar\phi^{q}_{1}
\bar\phi^{cl}_{1}\big(\phi^{cl}_{1}\big)^2
\phi^{q}_{2}\phi^{cl}_{2}\big(\bar\phi^{cl}_{2}\big)^2\big\rangle\sim G^R(t_2,t_1)
G^A(t_2,t_1)[G^K(t_1,t_2)]^2}\]
であり，第二は
\[\fitdisplay{\big\langle\bar\phi^{q}_{1}
\bar\phi^{cl}_{1}\big(\phi^{cl}_{1}\big)^2
\phi^{q}_{2}\phi^{cl}_{2}\big(\bar\phi^{q}_{2}\big)^2\big\rangle
\sim [G^R(t_{1},t_{2})]^2G^R(t_{2},t_{1})G^A(t_{2},t_{1})}\]
である．ここで$\phi^{\alpha}_{1,2}\equiv\phi^\alpha_{j_{1,2}}$である．これらの項はいずれも零である．なぜなら$G^R(t_2,t_1)\sim\theta(t_2-t_1)$である一方$G^A(t_2,t_1)\sim G^R(t_1,t_2)^*\sim\theta(t_1-t_2)$であり，その積の台が存在しないからである\footnote{厳密には，$G^R(t_{2},t_{1})$と$G^A(t_{2},t_{1})$は対角部分$t_{1}=t_{2}$で同時に非零となる．しかし，対角部分から積分への寄与は，$N\to \infty$のとき$\sim
 \delta_t^2N\to 0$である．}．まったく同じ理由で，すべての高次の項も消えることが容易にわかる．したがって，少なくとも摂動展開の範囲では，規格化は変わらない．


もう一つの例として，次の三次の非線形性をもつ実ボソン場（\ref{app_SPQM}参照）を考えよう．
\begin{equation}\label{int-bos-S-Phi-3}
S_{\mathrm{int}}=
\frac{\kappa}{6}\int\d\mathbf{r}\int_{\mathcal{C}}\d t\,
\phi^{3}=\frac{\kappa}{6}\int\d\mathbf{r} \int^{+\infty}_{-\infty}\d
t\, \big[\phi^{3}_{+}-\phi^{3}_{-}\big]=
\kappa\int\d\mathbf{r}\int^{+\infty}_{-\infty}\d t\,
\left[\big(\phi^{cl}\big)^{2}\phi^{q}+\frac{1}{3}
\big(\phi^{q}\big)^{3}\right].
\end{equation}この場合も，因果性条件\eqref{boson-causality}が満たされる．ダイアグラムでは，三次の非線形性は図\ref{Fig-Phi3}に示す二種類の頂点を生成する．一方は二つの古典場（実線）と一つの量子場（破線）をもち，もう一方は三つの量子場をもつ．前者の頂点には因子$\kappa$が，後者には重み$\kappa/3$が付く．実場の場合，線の向きを矢印で指定しないことに注意しよう．
\begin{figure}
\begin{center}\includegraphics[width=10cm]{figures/fig04.pdf}\end{center}
\caption{$\phi^3$理論の二つの相互作用頂点の図式表現．両者の間に3分の1の相対係数があることに注意．\label{Fig-Phi3}}
\end{figure}


%$$$$$$$$$$$$$$$$$$$$$$$$$$$$$$$$$$$$$$$$$$$$$$$$$$$$$$$$$$$$$$$$$$$$$$$$$$$$$$$$$$$$$$$$$$$$$$$$$$$$$$$$$$$
\subsectionlink{鞍点方程式}\label{sec_int_bos-2}


摂動論をさらに進める前に，作用の鞍点について論じる必要がある．式\eqref{boson-causality}によれば，作用には$\bar\phi^{q}$と$\phi^{q}$の両方について零次となる項が存在しない．明らかに$\delta S/ \delta \bar\phi^{cl}$についても同じことがいえるので，鞍点方程式の一つ
\begin{equation}
{\delta S\over \delta \bar\phi^{cl} } = 0\,
\end{equation}は，常に
\begin{equation}\label{int-bos-Phi-q-0}
\Phi^{q}=0 \, ,
\end{equation}によって解かれる．これは古典成分$\Phi^{cl}$が何であるかに依存しない．大文字$\Phi^{cl(q)}$は鞍点方程式の解を表す．たとえば式\eqref{boson-S-2}と\eqref{int-bos-S-1}の和で与えられる作用について，実際にそうなることを確認できる．条件\eqref{int-bos-Phi-q-0}のもとで，第二の鞍点方程式は次の形になる．
\begin{equation}\label{int-bos-GP}
 {\delta S\over \delta \bar\phi^{q} }=
\left(\big[G^R\big]^{-1} - \lambda\,|\Phi^{cl}|^2\right)
\Phi^{cl}=\left(i\partial_t +
\frac{1}{2m}\partial^{2}_{\mathbf{r}}+\mu - \lambda\,
|\Phi^{cl}|^2\right)\Phi^{cl}=0 \, .
\end{equation}これは非線形な時間依存Gross--Pitaevskii方程式であり，初期条件と境界条件を指定すれば，古典場の配置を決定する．


要点は，Keldysh作用の鞍点方程式の可能な解の中には，量子成分が零で，古典成分が古典的な（非線形）運動方程式に従うものが必ず存在する，ということである．このような鞍点を「古典的」鞍点と呼ぶ．式\eqref{boson-causality}および\eqref{int-bos-Phi-q-0}により，古典的鞍点の場の配置における作用は恒等的に零である．上で論じたように，古典的鞍点のまわりの微小揺らぎによる摂動展開は，正しく規格化された分配関数$Z=1$を与える．一見すると，これは他の鞍点が存在する可能性を排除しているように思える．しかし，この結論は早計である．系には，$\Phi^{q}\neq 0$を満たす「非古典的」鞍点が存在しうる．このような鞍点は分配関数に寄与せず（したがって基本的な規格化$Z=1$も変えない），観測量や相関関数には寄与しうる．一般に，\emph{非古典的}鞍点における作用は非零である．したがって，その寄与は指数関数的に小さい項，あるいは振動する項と結び付く．その例には，トンネル現象，熱活性化（次節で扱う），準位統計への振動的な寄与などがある．


ここで，\emph{古典的}鞍点からのずれ$\phi^{cl}=\Phi^{cl}+\delta\phi^{cl}$および$\phi^q=0+\delta\phi^q$について，系統的な摂動展開を組み立てよう．上で論じたように，これによって分配関数について新たな情報が得られるわけではない．しかし，Green関数，したがってさまざまな観測量についての情報は得られる．特に重要なのは，分布関数に対する運動論方程式を導けることである．以下の考察を簡単にするため，Bose凝縮が存在しない場合，すなわち$\Phi^{cl}=0$が古典的鞍点方程式\eqref{int-bos-GP}の適切な解となる場合に限定する．この場合$\phi^\alpha =\delta\phi^\alpha$であるので，記号$\delta$は省略できる．


%$$$$$$$$$$$$$$$$$$$$$$$$$$$$$$$$$$$$$$$$$$$$$$$$$$$$$$$$$$$$$$$$$$$$$$$$$$$$$$$$$$$$$$$$$$$$$$$$$$$$$$$$$$$
\subsectionlink{Dyson方程式}\label{sec_int_bos-3}


次の目標は，次式で定義される\emph{着衣した}Green関数を計算することである．
\begin{equation}\label{int-bos-G-dressed}
\mathbf{G}^{\alpha\beta}(t,t')=-i\int\D[\bar{\phi}\phi]
\,\phi^\alpha(t)\,\bar{\phi}^{\,\beta}(t')\,
\exp\big(iS_0+iS_{\mathrm{int}}\big)\, ,
\end{equation}ここで$\alpha,\beta=(cl,q)$であり，作用は式\eqref{boson-S-2}と\eqref{int-bos-S-1}で与えられる．そのため，指数関数を$S_{\mathrm{int}}$のべきで展開する．残ったGaussian作用による汎関数積分は，Wickの定理を用いて行う（\ref{app_Gaussian}参照）．これによって通常のダイアグラム級数を得る．すべての一粒子既約ダイアグラムを自己エネルギー行列$\hat{\Sigma}$にまとめると，次を得る．
\begin{equation}\label{int-bos-G-expansion}
\hat{{\bf G}}=\hat{G}+\hat{G}\circ\hat{\Sigma}\circ
\hat{G}+\hat{G}\circ\hat{\Sigma}\circ \hat{G}\circ\hat{\Sigma}\circ
\hat{G} +\ldots= \hat{G}\circ\left(\hat 1+\hat{\Sigma}\circ\hat{{\bf
G}}\right) \, ,
\end{equation}ここで$\hat G$は式\eqref{boson-G-fun-2}で与えられ，丸印の積は時間・空間領域での畳み込みと$2\times 2$行列の積を意味する．教科書にあるダイアグラム展開~\cite{Mahan,AGD,NegeleOrland}との唯一の違いは，$2\times 2$のKeldysh行列構造が存在する点である．級数が行列積の列としてまとまることは不思議ではない．実際，Keldysh添字$\alpha=(cl,q)$は，時間，空間，スピンなどに追加された，もう一つの添字にすぎない．したがって，他の添字と同じように，中間のすべての値について和をとり，それが行列の積となる．自己エネルギー行列$\hat{\Sigma}$の具体的な形はKeldysh法に固有であり，以下でやや詳しく論じる．


式\eqref{int-bos-G-expansion}の両辺に左から$\hat{G}^{-1}$を掛けると，厳密な着衣Green関数$\hat{{\bf G}}$に対するDyson方程式を次の形で得る．
\begin{equation}\label{int-bos-Dyson}
\left(\hat{G}^{-1} -\hat{\Sigma}\right)\circ \hat{{\bf G}}=\hat 1\, ,
\end{equation}ここで$\hat 1$は単位行列である．Keldysh法のきわめて非自明な特徴は，自己エネルギー行列$\hat{\Sigma}$が$\hat{G}^{-1}$と同じ因果性の構造をもつことである［式\eqref{boson-S-1}参照］．すなわち，
\begin{equation}\label{int-bos-Sigma}
\hat{\Sigma}=\left(\begin{array}{cc}
0   & \Sigma^{A}  \\
  \Sigma^{R}  & \Sigma^K
\end{array}\right)\, ,
\end{equation}ここで$\Sigma^{R(A)}$は時間領域の下（上）三角行列であり，$\Sigma^K$は反Hermite行列である．この事実は以下で示す．$\hat{G}^{-1}$と$\hat{\Sigma}$が同じ構造をもつため，着衣Green関数$\hat{{\bf G}}$も式\eqref{boson-G-fun-2}のような因果性の構造をもつことがわかる．したがって，Dyson方程式は次の形になる．
\begin{equation}\label{int-bos-Dyson-matrix}
\left(\begin{array}{cc}
0   & \big[G^{A}\big]^{-1}- \Sigma^{A}  \\
  \big[G^{R}\big]^{-1}- \Sigma^{R}  & -\Sigma^K
\end{array}\right)\circ
\left(\begin{array}{cc}
{\bf G}^K   & {\bf G}^R  \\
  {\bf G}^A  & 0
\end{array}\right) =\hat 1 ,
\end{equation}ここでは，$\big[G^{-1}\big]^K$が純粋な正則化項（$\sim i0F$）であり，非零の$\Sigma^K$があれば省略できることを考慮した．$\big[G^{R(A)}\big]^{-1}$の具体形［式\eqref{boson-G-Real-Space}参照］を用いると，遅延（先進）成分について次を得る．
\begin{equation}\label{int-bos-Dyson-RA}
\left( i\partial_t + {1\over 2m}\partial_\mathbf{r}^2 +\mu -
\Sigma^{R(A)}\right)\circ {\bf G}^{R(A)} =
\delta(t-t')\delta(\mathbf{r}-\mathbf{r}') \, .
\end{equation}自己エネルギー成分$\Sigma^{R(A)}$が何らかの近似で既知であれば，式\eqref{int-bos-Dyson-RA}は着衣Green関数の遅延（先進）成分に対する閉じた方程式となる．このGreen関数には，相互作用する系のスペクトルに関する情報が含まれる．


Keldysh成分に対する方程式を書くため，式\eqref{boson-G-FDT-general}にならって${\bf G}^K={\bf G}^R\circ {\bf F}- {\bf F}\circ {\bf G}^A$とパラメータ化する．ここで${\bf F}$は時間領域におけるHermite行列である．すると，Keldysh成分の方程式は$\big(\big[G^R\big]^{-1} -\Sigma^R\big)\circ
\big({\bf G}^R\circ {\bf F}-{\bf F} \circ {\bf
G}^A\big)=\Sigma^K\circ{\bf G}^A$となる．これに右から$\big(\big[G^A\big]^{-1}-\Sigma^A\big)$を掛け，式\eqref{int-bos-Dyson-RA}を用いると，最終的に次を得る．
\begin{equation}\label{int-bos-Dyson-F}
\left[{\bf F}\,,\, \left( i\partial_t + {1\over
2m}\partial^{2}_\mathbf{r} \right) \right]=\,
\Sigma^K-\left(\Sigma^{R}\circ {\bf F} - {\bf F}
\circ\Sigma^A\right)\, ,
\end{equation}ここで$[\,\,,\,]\,$は交換子を表す．この式は分布行列${\bf F}$の量子運動論方程式である．その左辺を\emph{運動項}と呼び，右辺を，ある因子を除いて\emph{衝突積分}と呼ぶ．以下で示すように，$\Sigma^K$は「流入」項，$\Sigma^{R}\circ {\bf F} - {\bf F} \circ\Sigma^A$は「流出」項の意味をもつ．平衡では，これら二つの過程が互いに打ち消し合い（運動項は消える），自己エネルギーはGreen関数と同じ構造$\Sigma^K=\Sigma^{R}\circ {\bf F}
- {\bf F} \circ\Sigma^A$をもつ．しかし，平衡から離れるとそうではない．


\begin{figure}
\begin{center}\includegraphics[width=10cm]{figures/fig05.pdf}\end{center}
\caption{$\phi^3$理論の自己エネルギーダイアグラム．\label{Fig-Sigma}}
\end{figure}


%$$$$$$$$$$$$$$$$$$$$$$$$$$$$$$$$$$$$$$$$$$$$$$$$$$$$$$$$$$$$$$$$$$$$$$$$$$$$$$$$$$$$$$$$$$$$$$$$$$$$$$$$$$$
\subsectionlink{自己エネルギー}\label{sec_int_bos-4}


自己エネルギー行列$\hat{\Sigma}$が，実際に因果性の構造\eqref{int-bos-Sigma}をもつことを示そう．そのため，$\kappa\phi^3$の非線形性\eqref{int-bos-S-Phi-3}をもつ実ボソン場を考え，パラメータ$\kappa$の二次まで計算する．図\ref{Fig-Phi3}の二つの頂点を用いると，次の結果を得る．


（i）\textit{$cl-cl$}成分は，図\ref{Fig-Sigma}aに描いた一つのダイアグラムで与えられる．対応する解析的な式は$\Sigma^{cl-cl}\!(t,t')=4i \kappa^2
G^R\!(t,t')G^A\!(t,t')\!=\!0$である．実際，積$G^R(t,t')G^A(t,t')$には台が存在しない（第\ref{sec_int_bos-1}節の脚注参照）．


（ii）\textit{$cl-q$}（先進）成分は，図\ref{Fig-Sigma}bの一つのダイアグラムで与えられる．対応する式は
\begin{equation}\label{int-bos-Sigma-A}
\Sigma^A(t,t') = 4i \kappa^2 G^A(t,t')G^K(t,t')\,.
\end{equation}である．$\Sigma^A(t,t')\sim G^A(t,t')\sim\theta(t'-t)$であるから，これは確かに先進（上三角）行列である．このダイアグラムには組合せ因子4が付く（外線の選び方が4通り$\times$内部の置換が2通り$\times$同じ頂点が二つあることによる1/(2!)）．


（iii）\textit{$q-cl$}（遅延）成分は，図\ref{Fig-Sigma}cのダイアグラムによって与えられる．
\begin{equation}\label{int-bos-Sigma-R}
\Sigma^R(t,t') = 4i\kappa^2 G^R(t,t')G^K(t,t')\,,
\end{equation}もちろん，これは式\eqref{boson-G-conjugation}を用いて，式\eqref{int-bos-Sigma-A}のHermite共役$\Sigma^R=\big[\Sigma^A\big]^\dagger$をとることでも得られる．$\Sigma^R(t,t')\sim
G^R(t,t')\sim\theta(t-t')$なので，確かに遅延（下三角）行列である．


（iv）\textit{$q-q$}（Keldysh）成分は，図\ref{Fig-Sigma}d--\ref{Fig-Sigma}fの三つのダイアグラムで与えられる．対応する式（これらのダイアグラムの和）は
\begin{eqnarray}\label{int-bos-Sigma-K}
\Sigma^K(t,t') = 2i\kappa^2 \big[G^K(t,t')\big]^2 +
6i\left({\kappa\over 3}\right)\kappa \big[G^A(t,t')\big]^2
+ 6i\left({\kappa\over 3}\right)\kappa \big[G^R(t,t')\big]^2\nonumber
\\
= 2i\kappa^2\left(\big[G^K(t,t')\big]^2+\big[G^R(t,t') -
G^A(t,t')\big]^2\right)\,.
\end{eqnarray}である．組合せ因子は図\ref{Fig-Sigma}dでは2，図\ref{Fig-Sigma}eおよびfでは6である．最後の等号では，時間領域で台をもたないことから$G^R(t,t')G^A(t,t')=0$となることを再び用いた．式\eqref{boson-G-conjugation}を用いると$\Sigma^K=-\big[\Sigma^K\big]^\dagger$を得る．


これで，自己エネルギー$\hat{\Sigma}$が$\hat{G}^{-1}$と同じ構造をもつことが示された．この主張は高次でも成り立つことを確認できる．式\eqref{int-bos-Sigma-A}--\eqref{int-bos-Sigma-K}では空間座標を省略したが，それらは明らかな仕方で復元できる．


%$$$$$$$$$$$$$$$$$$$$$$$$$$$$$$$$$$$$$$$$$$$$$$$$$$$$$$$$$$$$$$$$$$$$$$$$$$$$$$$$$$$$$$$$$$$$$$$$$$$$$$$$$$$
\subsectionlink{運動論方程式}\label{sec_int_bos-5}


運動論方程式の議論をさらに進めるには，Wigner変換（WT）\footnote{行列$A(\mathbf{r},\mathbf{r}')$のWigner変換は，$
a(\mathbf{R},\mathbf{k})\equiv\int\d \mathbf{r}_1  \,
A\left(\mathbf{R}+{\mathbf{r}_1\over 2},
\mathbf{R}-{\mathbf{r}_1\over 2}\right)\,\exp\{i\mathbf{kr}_1\}.$と定義される．行列$C=A\circ B$，すなわち$C(\mathbf{r},\mathbf{r}')=\int\d\mathbf{r}''
A(\mathbf{r},\mathbf{r}'')B(\mathbf{r}'',\mathbf{r}')$のWigner変換は，$$
c(\mathbf{R},\mathbf{k})=\iint\d\mathbf{r}_1\d\mathbf{r}_2 \iint{\d
\mathbf{k}_1\d\mathbf{k}_2\over (2\pi)^{2d}}\,\,
a\left(\mathbf{R}+{\mathbf{r}_1\over
2},\mathbf{k}+\mathbf{k}_1\right)
b\left(\mathbf{R}+{\mathbf{r}_2\over
2},\mathbf{k}+\mathbf{k}_2\right)
\exp\{i(\mathbf{k}_1\mathbf{r}_2-\mathbf{k}_2\mathbf{r}_1)\}\,.
$$に等しいことを示せる．積分内の関数を$\mathbf{k}_i$と$\mathbf{r}_i$について展開すると，$c(\mathbf{R},\mathbf{k})=
a(\mathbf{R},\mathbf{k})\,b(\mathbf{R},\mathbf{k})+(2i)^{-1}\big(\nabla_{\bf{R}}
a\nabla_\mathbf{k} b - \nabla_\mathbf{k} a\nabla_{\bf{R}} b\big)
+\ldots.$を得る．}を行うと便利である．分布関数行列$\mathbf
{F}(\mathbf{r},\mathbf{r}';t,t')$のWTは関数$\mathbf{f}(\mathbf{R},\mathbf{k};\tau,\epsilon)$であり，$\tau$と$\bf{R}$はそれぞれ中心時刻と中心座標である．定義\eqref{boson-G-FDT-general}によれば，関数${\bf f}$は$G^R-G^A$との積として現れる．後者は，自由粒子では$\epsilon=\omega_\mathbf{k}$に鋭いピークをもつ関数であり，相互作用系でも準粒子が明確に定義される限りこの性質が保たれる．したがって，$\epsilon=\omega_\mathbf{k}$であることを了解したうえで，しばしば${\bf f}(\bf{R},\mathbf{k},\tau)$と書く．


運動項［式\eqref{int-bos-Dyson-F}の左辺］をWigner表示で書き直すため，$i\partial_t$のWTが$\epsilon$であり，$\partial^2_\mathbf{r}$のWTが$-\mathbf{k}^2$であることに注意する．すると，たとえば$[{\bf
F},\partial_\mathbf{r}^2]_-\to [\mathbf{k}^2,\mathbf{f}]_-
+i\nabla_\mathbf{k} \mathbf{k}^2\nabla_{\mathbf{R}}
\mathbf{f}=2i\mathbf{k}\nabla_{\mathbf{R}}\mathbf{f}$となる．ここでWTどうしは可換なので交換子は消える．同様に$[{\bf
F},i\partial_t]_-\to -i\partial_\tau {\bf f}$である．Hamiltonianにスカラーポテンシャル$V(\mathbf{r})\hat{b}^\dagger_\mathbf{r}\hat{b}_\mathbf{r}$がある場合，作用には項$-V(\bar\phi^{cl}\phi^{q}+\bar\phi^{q}\phi^{cl})$が生じ，したがって$\big[G^{R(A)}\big]^{-1}$に$-V(\mathbf{r})$が加わる．これはDyson方程式\eqref{int-bos-Dyson-F}の左辺に項$-[{\bf F},V]_-$をもたらし，WT後には$i\mathbf{E}\nabla_\mathbf{k}{\mathbf{f}}$となる．ここで$\mathbf{E}\equiv-\nabla_{\mathbf{R}} V$は電場である．結果として，Dyson方程式\eqref{int-bos-Dyson-F}のWTは次の形をとる．
\begin{equation}\label{int-bos-kinetic-eq}
\Big(\partial_\tau-v_\mathbf{k}\nabla_{\mathbf{R}} -
\mathbf{E}\nabla_\mathbf{k}\Big)\,\mathbf{f}(\mathbf{R},\mathbf{k},\tau)
= I_{\mathrm{coll}}[\mathbf{f}]\, ,
\end{equation}ここで$v_\mathbf{k}\equiv \mathbf{k}/m$であり，$I_{\mathrm{coll}}[{\bf
f}]$は式\eqref{int-bos-Dyson-F}の右辺のWTに$i$を掛けたものである．これが分布関数に対する運動論方程式である．


分散関係$\epsilon=\omega_\mathbf{k}$をもつ実ボソン（\ref{app_SPQM}参照）では，運動項は$[\epsilon^2 -
\omega_\mathbf{k}^2,{\bf F}]_-\to2i\big(\epsilon\,
\partial_\tau - \omega_\mathbf{k}(\nabla_\mathbf{k}\omega_\mathbf{k})
\nabla_{\mathbf{R}}\big)\,{\bf f}= 2i\epsilon\big(\partial_\tau -
v_\mathbf{k}\nabla_{\mathbf{R}}\big)\,{\bf f}$の形をとる．ここで$v_\mathbf{k}\equiv \nabla_\mathbf{k}\omega_\mathbf{k}$は群速度である．したがって運動論方程式は$\big(
\partial_\tau - v_\mathbf{k}\,\nabla_{\mathbf{R}}\big)\,
\mathbf{f}(\mathbf{R},\mathbf{k},\tau)= I_{\mathrm{coll}}[{\bf f}]$となる．衝突積分$I_{\mathrm{coll}}[{\bf f}]$は，式\eqref{int-bos-Dyson-F}の右辺のWTを$-2i\epsilon$で割ったものである．


ここで，第\ref{sec_int_bos-4}節における$\phi^{3}$理論の計算を例に，衝突積分を論じよう．計算を短くするため，空間的に一様で，運動量空間では等方的な系を考える．すなわち，エネルギー緩和だけに注目する．この場合，分布関数は$\mathbf{f}(\mathbf{R},\mathbf{k},\tau)=\mathbf{f}(\tau,\omega_\mathbf{k})
=\mathbf{f}(\tau,\epsilon)$であり，運動量の大きさへの依存性は，引数$\omega_\mathbf{k}=\epsilon$で置き換えられる．式\eqref{int-bos-Sigma-A}--\eqref{int-bos-Sigma-K}を用いると，式\eqref{int-bos-Dyson-F}の右辺のWTについて，次を得る\footnote{WTの積だけを残し，すべての勾配項は無視する．特に${\bf G}^K\to {\bf f}\,({\bf g}^R-{\bf g}^A)$とする．式\eqref{int-bos-Sigma-A}--\eqref{int-bos-Sigma-K}の時空表示の代わりにエネルギー・運動量表示を用い，$\Sigma^{R}\circ {\bf F} -{\bf F} \circ\Sigma^A$の式では，引数$\omega$と$\epsilon-\omega$の間で対称化を行う．}．
\begin{subequations}\label{int-bos-Sigma-WT}
\begin{equation}
\Sigma^{R}\circ {\bf F} -{\bf F} \circ\Sigma^A \to-2i\,{\bf
f}(\tau,\epsilon)\int\d\omega\, M(\tau,\epsilon,\omega)\, \Big[{\bf
f}(\tau,\epsilon-\omega)+{\bf f}(\tau,\omega)\Big]\,,
\end{equation}
\begin{equation}
\hskip-2.5cm \Sigma^{K}\to  -2i\int\d\omega\,
M(\tau,\epsilon,\omega)\, \Big[{\bf f}(\tau,\epsilon-\omega){\bf
f}(\tau,\omega)+1\Big]\,,
\end{equation}\end{subequations}
ここで遷移率は，次で与えられる．
\begin{equation}\label{int-bos-M}
M(\tau,\epsilon,\omega) = 2\pi\kappa^2\sum_{q}\, {\bf
\Delta_{g}}(\tau,\epsilon-\omega;\mathbf{k}-\mathbf{q})\,{\bf
\Delta_{g}}(\tau,\omega;\mathbf{q})\, .
\end{equation}ここで${\bf \Delta_{g} }\equiv i({\bf g}^R - {\bf g}^A)/(2\pi)$であり，${\bf g}^{R(A)}(\tau,\epsilon;\mathbf{k})$は遅延（先進）Green関数$\mathbf{G}^{R(A)}$のWTである．ダイアグラム級数の部分的な再総和を行うため，式\eqref{int-bos-Sigma-A}--\eqref{int-bos-Sigma-K}には，裸のGreen関数の代わりに着衣Green関数を代入した．（この方法は\emph{自己無撞着Born近似}と呼ばれることもある．この近似でも頂点補正は無視している．）明確に定義された準粒子が常に存在すると仮定すれば，${\bf \Delta_{g}}(\tau,\epsilon,\mathbf{k})$は$\epsilon=\omega_\mathbf{k}$に鋭いピークをもつ関数とみなせる．この場合，式\eqref{int-bos-M}は，$\epsilon=\omega_\mathbf{k}$をもつ初期粒子が，エネルギー$\omega=\omega_\mathbf{q}$および$\epsilon-\omega=\omega_{\mathbf{k}-\mathbf{q}}$をもつ二つの実粒子（質量殻上の粒子）へ崩壊するという事実を表すにすぎない．したがって，運動論方程式として最終的に次を得る．
\begin{equation}\label{int-bos-kinetic-eq-fin}
{\partial {\bf f}(\epsilon)\over \partial\tau}= \int\!\! {\d\omega}\,\,
{M(\epsilon,\omega)\over \epsilon}\,\, \Big\{ {\bf
f}(\epsilon-\omega){\bf f}(\omega)+1 -{\bf f}(\epsilon)\big[{\bf
f}(\epsilon-\omega)+{\bf f}(\omega)\big]\Big\},
\end{equation}ここでは簡潔さのため時刻の引数を省略した．恒等式$\coth(a-b)\,\coth( b)+1=\coth(a)\big(\coth
(a-b)+\coth(b)\big)$により，衝突積分は${\bf f}(\epsilon)=\coth(\epsilon/2T)$によって恒等的に零になる．ここで$T$は温度である．これは熱平衡の分布関数である．運動論方程式\eqref{int-bos-kinetic-eq-fin}によれば，この分布は任意の温度で安定である（温度は外部リザーバーによって決まるか，閉じた系では全エネルギー保存から決まる）．平衡分布は明らかに運動項を零にするため，式\eqref{int-bos-Dyson-F}によれば，\emph{厳密な}自己エネルギーは$\Sigma^K=\coth(\epsilon/2T)\big[\Sigma^R-\Sigma^A\big]$を満たす．裸のGreen関数も同じ関係\eqref{boson-G-FDT}に従うので，熱平衡における\emph{厳密な}着衣Green関数について，次が成り立つと結論できる．
\begin{equation}\label{int-bos-FDT}
{\bf G}^K=\big({\bf G}^R - {\bf G}^A\big)\,\coth\, {\epsilon \over
2\,T}\, .
\end{equation}これが\emph{揺動散逸定理}（FDT）の内容である．その帰結として，平衡ではKeldysh成分は遅延成分に加えて新たな情報を含まない．したがって，原理的にはKeldysh法を，より簡潔な構成であるMatsubara法に置き換えられる．もちろん，後者は非平衡では機能しない．


運動論方程式\eqref{int-bos-kinetic-eq-fin}に戻ると，衝突積分の「流入」項と「流出」項を同定できる．そのためには，${\bf
f}_{\epsilon}=1+2\,{\bf n}_{\epsilon}$で定義される占有数$\mathbf{n}_{\epsilon}$を用いて衝突積分を書けばよい．式\eqref{int-bos-kinetic-eq-fin}の右辺の波括弧内は，$4\left[{\bf n}_{\epsilon-\omega}{\bf n}_\omega- {\bf
n}_\epsilon({\bf n}_{\epsilon-\omega}+{\bf n}_\omega+1)\right]$の形になる．第一項${\bf n}_{\epsilon-\omega}{\bf n}_\omega$は，エネルギー$\epsilon-\omega$の粒子がエネルギー$\omega$の粒子を吸収し，エネルギー$\epsilon$の状態を占有する確率を与える．これが衝突積分の「流入」項である．この項は自己エネルギーの$\Sigma^K$部分に由来する．第二項$ -{\bf n}_\epsilon({\bf n}_{\epsilon-\omega}+{\bf
n}_\omega+1)$は，エネルギー$\epsilon$の状態から，エネルギー$\epsilon-\omega$および$\omega$の粒子の誘導放出，または自発放出によって粒子が失われることを示す．これが「流出」項であり，$\Sigma^{R(A)}$の寄与に由来する．


最後に，Wigner変換に含まれる近似について論じよう．式\eqref{int-bos-Dyson-F}は形式的には厳密だが，そこから有用な情報を引き出すことは非常に難しい．したがって，式\eqref{int-bos-kinetic-eq}や\eqref{int-bos-kinetic-eq-fin}のような，近似的ではあるがはるかに扱いやすい形へ移ることが望ましい．その際には，WTの近似形を用いなければならない．実際，演算子$\nabla_\mathbf{k}\nabla_{\mathbf{R}}$について形式的には無限である級数を，通常は最初の非零項で打ち切る．この手続きは$\delta \mathbf{k}\,\delta \mathbf{R}\gg 1$が成り立つ限り正当化される．ここで$\delta \mathbf{k}$は${\bf f}$の運動量依存性の特徴的な微視的スケールであり，$\delta\mathbf{R}$はその空間変化の特徴的なスケールである．時間領域にも同様の条件$\delta\epsilon\,
\delta \tau\gg 1$があるのか，という疑問が生じるかもしれない．ここで$\delta\epsilon$と$\delta \tau$は，それぞれ${\bf f}$の特徴的なエネルギースケールと時間スケールである．通常$\delta\epsilon \approx T$なので，この条件はきわめて厳しく，低温では$\delta\tau \gg
1/T$を満たす非常に遅い過程しか扱えなくなる．幸い，実際にはそうならない．${\bf \Delta_g}(\epsilon,\mathbf{k})$がピーク構造をもつため，そのピークが鋭い限り，エネルギー引数$\epsilon$は$\omega_\mathbf{k}$に固定され，独自のダイナミクスをもたない．したがって実際の基準は，$\delta\epsilon$が${\bf \Delta_g}(\epsilon,\mathbf{k})$のピーク幅より十分大きいことである．後者は定義により準粒子寿命$\tau_{qp}\,$に対応するので，条件は$\tau_{qp}\gg 1/T$となる．相互作用があまり強くない多くの系では，この条件は実際に満たされる．


%$$$$$$$$$$$$$$$$$$$$$$$$$$$$$$$$$$$$$$$$$$$$$$$$$$$$$$$$$$$$$$$$$$$$$$$$$$$$$$$$$$$$$$$$$$$$$$$$$$$$$$$$$$$
%$$$$$$$$$$$$$$$$$$$$$$$$$$$$$$$$$$$$$$$$$$$$$$$$$$$$$$$$$$$$$$$$$$$$$$$$$$$$$$$$$$$$$$$$$$$$$$$$$$$$$$$$$$$
\sectionlink{環境に接した粒子}\label{sec_environment}
\subsectionlink{量子散逸作用}\label{sec_environment-1}


ポテンシャル$U(\Phi)$中に置かれ，調和的な弦$\varphi(x,t)$に結合した，座標$\Phi(t)$の粒子を考える．この粒子は，Josephson接合の位相や量子ドットの電荷など，集団的な自由度を表していてもよい．一方，弦は散逸環境を模型化する役割を果たす．一次元の弦の利点は，低エネルギーで一定の状態密度をもつ，最も簡単な連続体であることにある．この性質により，たとえばFermi海との相互作用を模擬できる．低エネルギーで一定の状態密度をもつ連続体のリザーバーは，「Ohmic」環境（または浴）と呼ばれることがある．環境は熱平衡にあるものとする．


この系のKeldysh作用は，三つの項$S=S_{\mathrm{p}}+S_{\mathrm{str}}+S_{\mathrm{int}}$で与えられる．ここで（\ref{app_SPQM}参照），
\begin{subequations}\label{environment-S-Str-Particle}
\begin{equation}
S_{\mathrm{p}}[\Phi]=\int_{-\infty}^{+\infty}\d
t\left[-2\,\Phi^{q}{\d^{\,2} \Phi^{cl}\over \d t^2}
-U\left(\Phi^{cl}+\Phi^{q}\right)+ U(\Phi^{cl}-\Phi^{q}) \right]\,,
\end{equation}
\begin{equation}
\hskip-2.9cm S_{\mathrm{str}}[\varphi]=\int_{-\infty}^{+\infty} \d
t\int\d x\ \vec{\varphi}^{T}(x,t) \hat{D}^{-1}\vec{\varphi}(x,t)\,,
\end{equation}
\begin{equation}
\hskip-2cm S_{\mathrm{int}}[\Phi,\varphi]=
2\sqrt{\gamma}\int_{-\infty}^{+\infty}\d t
\left.\vec{\Phi}^{T}(t)\,\hat{\sigma}_x\,\partial_x\vec{\varphi}(x,t)\right|_{x=0}\,.
\end{equation}\end{subequations}
である．ここでは，たとえば$\vec{\Phi}^T\equiv (\Phi^{cl},\Phi^{q})$のように古典成分と量子成分からなるベクトル，および弦の相関関数$\hat{D}^{-1}$を導入した．後者はボソンに典型的な形\eqref{boson-S-2}をもち，$\big[D^{R(A)}\big]^{-1}=-\partial^{2}_{t}+v^{2}_{s}\partial^{2}_{x}$である．これは式\eqref{app-SPQM-S-Phonons-1}から従う．$S_{\mathrm{p}}$は粒子を表す（\ref{app_SPQM}の式\eqref{app-SPQM-S-3}に関する議論を参照）．$S_{\mathrm{str}}$は弦の作用\eqref{app-SPQM-S-Phonons-1}である．粒子と弦の相互作用項は，粒子座標と$x=0$における弦の応力との局所的な積にとる（したがって粒子に働く力は，弦の局所的な応力に比例する）．時間領域では相互作用は瞬間的であり，Keldysh経路上では$\Phi(t)\partial_x\varphi(x,t)|_{x=0}\to
\Phi_+\partial_{x}\varphi_+-\Phi_-\partial_{x}\varphi_-$となる．古典・量子成分の表記へ変換すると$2(\Phi^{cl}\partial_{x}\varphi^{q}+\Phi^{q}\partial_{x}\varphi^{cl})$を得て，これは因果性条件\eqref{boson-causality}を満たす．行列表記では，式(\ref{environment-S-Str-Particle}c)の形になる．相互作用定数は$\sqrt{\gamma}$と表す．


ここで調和的な弦の自由度を積分消去して，問題を粒子座標だけに帰着させることができる．Gaussian積分の標準的な規則（\ref{app_Gaussian}参照）に従うと，粒子に対するいわゆる散逸作用を得る．
\begin{subequations}\label{environment-S-diss}
\begin{equation}
\hskip-.7cm S_{\mathrm{diss}}=\gamma\iint_{-\infty}^{+\infty}\d
t\,\d t' \, \vec{\Phi}^T(t) \hat{\mathfrak{D}}^{-1}(t-t')
\hat\Phi(t') \, ,
\end{equation}
\begin{equation}
\hat{\mathfrak{D}}^{-1}(t-t\,')=-\left.\hat{\sigma}_{x}\,\partial_x\,\partial_{x\,'}
\hat{D}(x-x\,';t-t\,')\right|_{x=x'=0}\, \hat{\sigma}_{x}\,.
\end{equation}\end{subequations}
行列の積を直接計算すると，散逸相関関数$\hat{\mathfrak{D}}^{-1}$が標準的な因果性の構造をもつことがわかる．その遅延（先進）成分をFourier変換すると，次を得る．
\begin{equation}\label{environment-L-RA}
\big[\mathfrak{D}^{R(A)}(\epsilon)\big]^{-1}=
-\sum\limits_k\frac{k^2}{(\epsilon\pm i0)^2-k^2} = \pm\, {i\over
2}\,\epsilon + \mbox{const}\, ,
\end{equation}ここでは簡潔さのため$v_s=1$とおいた．$\epsilon$に依存しない定数（$R$成分と$A$成分で共通）は，ポテンシャル$U(\Phi)=\mbox{const}\, \Phi^2+\ldots$の調和部分の再定義に吸収できるので，省略してよい．平衡における相関関数のKeldysh成分は，FDTによって定まる．
\begin{equation}\label{environment-L-K}
\big[\mathfrak{D}^{-1}\big]^{K}(\epsilon)=\left(
\big[\mathfrak{D}^{R}\big]^{-1}-\big[\mathfrak{D}^{A}\big]^{-1}
\right)\coth{\epsilon\over 2\,T} = i\epsilon\,\coth{\epsilon\over
2\,T}\, .
\end{equation}これは正定値の虚部をもつ反Hermite演算子であり，$\Phi$についての汎関数積分の収束を保証する．


時間表示では，相関関数の遅延（先進）成分は，時間について局所的な形$\big[\mathfrak{D}^{R(A)}\big]^{-1}= \mp {1\over
2}\,\delta(t-t')\,\partial_{t}$をとる．一方，低温ではKeldysh成分は非局所的な関数であり，式\eqref{environment-L-K}の逆Fourier変換から求められる．
\begin{equation}\label{eqenvironment-L-K-time}
\big[\mathfrak{D}^{-1}\big]^{K}(t-t') =i\left[(2T+C)\delta(t-t')-
\frac{\pi T^2}{\sinh^2[\pi T(t-t')]}\right]\,,
\end{equation}ここで無限大の定数$C$は，条件$\int\d t\big[\mathfrak{D}^{-1}(t)\big]^{K}=
\big[\mathfrak{D}^{-1}(\epsilon=0)\big]^{K}=2iT$を満たすためのものである．最後に，弦に結合した粒子のKeldysh作用として，次を得る．
\begin{eqnarray}\label{environment-S-diss-fin}
S[\Phi]=\int_{-\infty}^{+\infty}\d t\left[-2\,\Phi^q\left({\d^{\,2}
\Phi^{cl}\over \d t^2} +{\gamma\over 2}{\d \Phi_{cl}\over \d
t}\right)- U\left(\Phi^{cl}+\Phi^{q}\right)+ U(\Phi^{cl}-\Phi^{q})
\right]\nonumber
\\
 +2i\gamma\int^{+\infty}_{-\infty}\d t\left[T\big(\Phi^{q}(t)\big)^2+
  \frac{\pi T^2}{2} \int_{-\infty}^{+\infty}\d t'\,
\frac{\big(\Phi^q(t)-\Phi^q(t')\big)^2}{\sinh^2[\pi
T(t-t')]}\right]\,.
\end{eqnarray}この作用は，第\ref{sec_bosons-3}節に列挙した因果性の条件をすべて満たす．この場合のKeldysh（$q-q$）成分は，単なる正則化因子ではなく，弦との結合に由来する量子揺らぎの減衰項であることに注意しよう．弦のもう一つの効果は，$R$成分と$A$成分に摩擦項$\sim \gamma\partial_t$が存在することである．平衡では，摩擦係数と揺らぎの振幅はFDTによって厳密に関係付けられる．量子散逸作用\eqref{environment-S-diss-fin}は，さまざまな近似や他のアプローチとのつながりを示すのに便利な題材である．


%$$$$$$$$$$$$$$$$$$$$$$$$$$$$$$$$$$$$$$$$$$$$$$$$$$$$$$$$$$$$$$$$$$$$$$$$$$$$$$$$$$$$$$$$$$$$$$$$$$$$$$$$$$$
\subsectionlink{古典極限}\label{sec_environment-2}


\emph{古典的}鞍点方程式（$\Phi^{q}(t)=0$とするもの）は，次の形をとる．
\begin{equation}\label{environment-Newton-eq}
\left. -{1\over 2}\frac{\delta S[\Phi]}{\delta
\Phi^{q}}\right|_{\Phi^{q}=0} = {\d^{\,2} \Phi^{cl}\over \d t^2}
+{\gamma\over 2}{\d \Phi^{cl}\over \d t} + {\partial
U(\Phi^{cl})\over
\partial \Phi^{cl} } =0 \, .
\end{equation}これは決定論的な古典運動方程式である．ここでは，粘性力$-(\gamma/2) \dot \Phi^{cl}$を伴うNewton方程式になっている．この近似では，\emph{量子}揺らぎと\emph{熱}揺らぎの両方を無視している．


量子揺らぎを完全に無視しながら，熱揺らぎを残すこともできる．そのためには，作用\eqref{environment-S-diss-fin}にPlanck定数を復元してから，極限$\hbar\to 0$をとると便利である．次元の関係から，作用の前には因子$\hbar^{-1}$が付くはずである．古典運動方程式\eqref{environment-Newton-eq}を与える作用の部分にPlanck定数が現れないようにするには，$\Phi^{q} \to \hbar
\Phi^{q}$と変数の尺度を変更するとよい．最後に，単位を正しく保つため，式\eqref{environment-S-diss-fin}の最終項で$T\to T/\hbar$と置き換える必要がある．これで極限$\hbar\to 0$は容易にとれる．（i）$U(\Phi^{cl}\pm \hbar\Phi^{q})$を$\hbar\Phi^{q}$の一次まで展開し，それより高次の項はすべて無視する．（ii）$T\to\infty$と等価な$\hbar\to0$極限では，作用\eqref{environment-S-diss-fin}の非局所部分は消え，局所項$\propto \big(\Phi^{q}(t)\big)^2$は残る．その結果，散逸作用の古典極限は次のようになる．
\begin{equation}\label{environment-S-diss-class}
S[\Phi]=2\int_{-\infty}^{+\infty}\d t\left[ -\Phi^{q}\left(
 {\d^{\,2} \Phi^{cl}\over \d t^2}+{\gamma\over 2}{\d
\Phi^{cl}\over \d t}+{\partial U(\Phi^{cl})\over
\partial \Phi^{cl} } \right)+i\gamma\, T\, \big(\Phi^{q}\big)^2 \right].
\end{equation}物理的には，極限$\hbar\to 0$は$\hbar\Omega\ll T$を意味する．ここで$\Omega$は，粒子の特徴的な古典振動数である．この条件は，式\eqref{environment-S-diss-fin}の最後の項が時間について局所的な形をとるために必要である．$U$の高階微分を無視するための条件は$\hbar\ll \gamma  \big(\tilde{\Phi}^{cl}\big)^2$であり，$\tilde\Phi^{cl}$は粒子運動の特徴的な古典振幅である．


%$$$$$$$$$$$$$$$$$$$$$$$$$$$$$$$$$$$$$$$$$$$$$$$$$$$$$$$$$$$$$$$$$$$$$$$$$$$$$$$$$$$$$$$$$$$$$$$$$$$$$$$$$$$
\subsectionlink{Langevin方程式}\label{sec_environment-3}


古典作用\eqref{environment-S-diss-class}からさらに進む方法の一つは，その最後の項に$i$を掛けたものの指数関数が，次のように恒等的に書き直せることに注目することである．
\begin{equation}
\exp\left(-2\gamma T\int^{+\infty}_{-\infty}\d
t\,\big[\Phi^{q}(t)\big]^{2}\right)=\int\D[\xi]\
\exp\left(-\int^{+\infty}_{-\infty}\d t\left[{\xi^2(t)\over 2\gamma
T} - 2i \xi(t)\Phi^{q}(t)\right]\right)\, .
\end{equation}この恒等式はHubbard--Stratonovich変換と呼ばれ，$\xi(t)$は補助的なHubbard--Stratonovich場である．この恒等式は，右辺の指数の中を平方完成し，各時刻でGaussian積分を行えば証明できる．積分測度$\D[\xi]$には，定数の乗法因子が含まれている．


$\xi$と$\Phi$についての汎関数積分の順序を交換すると，分配関数について次を得る．
\begingroup\interdisplaylinepenalty=10000
\begin{eqnarray}\label{environment-Langevin-Z}
 Z\!&=&\!\int\D[\xi]\,
 \exp\left(-{1\over 2\gamma T}\int^{+\infty}_{-\infty}\d t\, \xi^2(t)\right)\nonumber\\
 &\times&\! \int\D\big[\Phi^{cl}\big]\int\D\big[\Phi^{q}\big]\,
 \exp\left(-2i\int^{+\infty}_{-\infty}\d t\, \Phi^{q}(t)
 \left[ {\d^{\,2}\Phi^{cl}\over \d t^2} +{\gamma\over 2}{\d
\Phi^{cl}\over \d t}+ {\partial U(\Phi^{cl})\over
\partial \Phi^{cl} }-\xi(t) \right]\right).
\end{eqnarray}
\endgroup{}指数は$\Phi^{q}(t)$に線形に依存するので，$\D\big[\Phi^{q}\big]$の積分を行うと，角括弧内の式を引数とする$\delta$関数が生じる．この汎関数的な$\delta$関数は，各時刻でその引数が零となることを課す．したがって，すべての可能な軌道$\Phi^{cl}(t)$のうち，次の方程式を満たすものだけが分配関数に寄与する．
\begin{equation}\label{environment-Langevin-eq}
 {\d^{\,2} \Phi^{cl}\over \d t^2} +{\gamma\over 2}{\d \Phi^{cl}\over
\d t} + {\partial U(\Phi^{cl})\over \partial \Phi^{cl} } = \xi(t) \,.
\end{equation}これは，時間依存する外力$\xi(t)$を伴うNewton方程式である．同じ議論は，たとえば$\big\langle
\Phi^{cl}(t)\Phi^{cl}(t')\big\rangle$のような古典場の任意の相関関数にも適用できるので，解法の手順は次のようになる．（i）$\xi(t)$のある実現を選ぶ．（ii）式\eqref{environment-Langevin-eq}を，たとえば数値的に解く．（iii）得られた解$\Phi^{cl}(t)$を用いて相関関数を計算する．（iv）外力$\xi(t)$の実現の集団について結果を平均する．この外力の統計は，式\eqref{environment-Langevin-Z}の汎関数積分$\D[\xi]$の重み因子で決まる．それによれば，$\xi(t)$は次の相関関数をもつGaussianの短距離相関雑音，すなわち白色雑音である．
\begin{equation}\label{environment-Langevin-noise}
\langle \xi(t) \rangle = 0\,, \hskip 1cm \langle \xi(t)
\xi(t')\rangle = \gamma T \delta(t-t')\,.
\end{equation}右辺に白色雑音をもつ式\eqref{environment-Langevin-eq}は，Langevin方程式と呼ばれる．これは確率的な熱揺らぎの存在下での古典的なNewton力学を記述する．雑音振幅が摩擦係数$\gamma$および温度に比例するという事実は，FDTの表れである．FDTは，環境（弦）が熱平衡にある限り成り立つ．


%$$$$$$$$$$$$$$$$$$$$$$$$$$$$$$$$$$$$$$$$$$$$$$$$$$$$$$$$$$$$$$$$$$$$$$$$$$$$$$$$$$$$$$$$$$$$$$$$$$$$$$$$$$$
\subsectionlink{Martin--Siggia--Rose法}\label{sec_environment-4}


前節では，$\Phi^{cl}$と別の場$\Phi^q$で書かれた作用から，古典座標$\Phi^{cl}$のLangevin方程式を導いた．逆に，Langevin方程式から有効作用を導く手続きは，Martin--Siggia--Rose（MSR）法~\cite{MSR}として知られている．ここではDeDominicis~\cite{DeDominicis}が提案した形でその概要を示す．


次のLangevin方程式を考える．
\begin{equation}
\hat {\cal O}\big[\Phi^{cl}\big] = \xi(t)\, ,
\end{equation}ここで$\hat {\cal O}\big[\Phi^{cl}\big]$は座標$\Phi^{cl}(t)$に作用する非線形微分演算子であり，$\xi(t)$は式\eqref{environment-Langevin-noise}で指定される白色雑音の力である．「分配関数」を次のように定義する．
\begin{equation}\label{environment-Z-MSR}
Z[\xi]=\int\D[\Phi^{cl}]\, {\cal J}[\hat {\cal O }
]\,\delta\big(\hat {\cal O}\big[\Phi^{cl}\big] - \xi(t)\big) \equiv
1\, .
\end{equation}${\cal J}[\hat {\cal O}]$が演算子$\hat {\cal O}\big[\Phi^{cl}\big]$のJacobianであれば，$\delta$関数の積分によって，これは恒等的に1に等しい．式\eqref{environment-Z-MSR}を解釈するには，時間軸を離散化して$N$次元ベクトル$\Phi^{cl}_j=\Phi^{cl}(t_j)$と$\xi_j=\xi(t_j)$を導入すればよい．演算子は${\cal O}_i=\mathcal{O}_{ij}\Phi^{cl}_j+{1\over
2}\Gamma_{ijk}\Phi^{cl}_j\Phi^{cl}_k +\ldots$という形になる．ここで繰り返された添字について和をとる．Jacobian ${\cal J}$は，次の$N\times N$行列の行列式の絶対値で与えられる．すなわち$J_{ij}\equiv
\partial {\cal O}_i/ \partial\Phi^{cl}_j=
\mathcal{O}_{ij}+\Gamma_{ijk}\Phi^{cl}_k+\ldots$である．適切な遅延型の正則化を選び，行列$J_{ij}$を，主対角要素がすべて1の下三角行列にすることができる（この主対角要素はすべて$\mathcal{O}_{ii}=1$の項から来る）．すると，この場合${\cal J}=1$となる．実際，例として$\hat {\cal O}\big[\Phi^{cl}\big]=\partial_t\Phi^{cl}-U(\Phi^{cl})$を考えよう．Langevin方程式の遅延型に正則化した形は$\Phi^{cl}_{i}
=\Phi^{cl}_{i-1}+\delta_t(U(\Phi^{cl}_{i-1})+\xi_{i-1})$である．この場合，明らかに$J_{ii}=1$かつ$J_{i,i-1}=-1-\delta_t
U'(\Phi^{cl}_{i-1})$であり，他のすべての成分は零である．結果として${\cal J}=1$を得る．


分配関数\eqref{environment-Z-MSR}自体は自明であるが，意味のある観測量と相関関数はすべて，一連の因子$\Phi^{cl}(t)\Phi^{cl}(t')\ldots$を汎関数積分\eqref{environment-Z-MSR}に挿入して得られることは明らかである．このことを念頭に置き，分配関数の計算を進めよう．補助場$\Phi^{q}(t)$を用いた$\delta$関数の積分表示を使うと，次を得る．
\begin{equation}
Z[\xi]=\int\D[\Phi^{cl},\Phi^{q}]\,\exp\left(-2i\int\d
t\,\Phi^{q}(t) \big[\hat {\cal O}^R[\Phi^{cl}(t)] -
\xi(t)\big]\right)\, ,
\end{equation}ここで$\hat {\cal O}^R$は演算子$\hat {\cal O}$の遅延型正則化を表すので，${\cal J}=1$とおく．ここで$\xi$についてGaussian積分を行うことにより，白色雑音\eqref{environment-Langevin-noise}について平均できる．
\begin{eqnarray}\label{environment-Z-MSR-average}
Z&=&\int\D[\xi] \exp\left(-{1\over 2\gamma T}\int\d
t\,\xi^2(t)\right)Z[\xi]\nonumber
\\
&=&\int\D[\Phi^{cl},\Phi^{q}]\,
 \exp\left(-\int\d t \,\left[ 2i\,\Phi^{q}(t)\hat{\cal O}^R\big[\Phi^{cl}(t)\big]+2\gamma
T\big[\Phi^q(t)\big]^{2}\right] \right)\, .
\end{eqnarray}式\eqref{environment-Z-MSR-average}の指数は，微分演算子の遅延型正則化も含め，Keldysh作用の古典極限そのものである［式\eqref{environment-S-diss-class}参照］．すなわち，MSR作用はKeldysh作用の古典（高温）極限にほかならない．MSR法は，古典的な確率問題を，対応する適切な汎関数表示へ変換する簡単な方法を与える．この表示は解析計算に有用である．以下にその一例を示す．


%$$$$$$$$$$$$$$$$$$$$$$$$$$$$$$$$$$$$$$$$$$$$$$$$$$$$$$$$$$$$$$$$$$$$$$$$$$$$$$$$$$$$$$$$$$$$$$$$$$$$$$$$$$$
\subsectionlink{熱活性化}\label{sec_environment-5}


図\ref{Fig-Activation}aに示す準安定なポテンシャル井戸中の粒子を考える．ポテンシャルは$\Phi=0$に準安定極小をもち，$\Phi=1$に，極小からの相対的な高さ$U_0$の極大をもつ．また，粒子の運動が過減衰，すなわち$\gamma\gg \sqrt{U''}$であると仮定する．この場合は慣性項を無視でき，粘性的な緩和ダイナミクスだけが残る．古典散逸作用\eqref{environment-S-diss-class}は，次の形になる．
\begin{equation}\label{environment-S-diss-viscous}
S[\Phi] = 2\int_{-\infty}^{+\infty}\d t\left[ -\Phi^{q}(t)
\left({\gamma\over 2}{\d \Phi_{cl}\over \d t}+ {\partial
U(\Phi^{cl})\over
\partial \Phi^{cl} } \right) + i\gamma\, T\, \big[\Phi^{q}(t)\big]^2\, \right]\, .
\end{equation}対応する鞍点方程式は，次である．
\begin{subequations}\label{environment-escape-eq}
\begin{equation}
{\gamma\over 2}\,\dot\Phi^{cl}= - {\partial U(\Phi^{cl})\over
\partial \Phi^{cl} } + 2i\gamma T\, \Phi^q\,,
\end{equation}
\begin{equation}
\hskip-1.2cm
{\gamma\over 2}\, \dot\Phi^{q}=  \Phi^q \,{\partial^2
U(\Phi^{cl})\over
\partial (\Phi^{cl})^2 } \,.
\end{equation}\end{subequations}
これらの方程式には\emph{古典的}な解がある．すなわち$\Phi^q(t)\equiv 0$であり，$\Phi^{cl}(t)$は古典運動方程式${\gamma\over 2}\,\dot\Phi^{cl}= -\partial
U(\Phi^{cl})/\partial \Phi^{cl}$を満たす．初期条件$\Phi^{cl}(0)<1$に対して，この方程式は$\Phi^{cl}=0$の極小へ向かう粘性緩和を予言する．この方程式によれば，この極小から脱出する可能性はない．したがって，式\eqref{environment-escape-eq}の古典解は熱活性化を記述\emph{しない}．そこで，式\eqref{environment-escape-eq}の別の可能な解，すなわち$\Phi^q\neq 0$をもつ解を探さなければならない．
\begin{figure}
\begin{center}\includegraphics[width=10cm]{figures/fig06.pdf}\end{center}
\caption{a）準安定極小をもつポテンシャル．b）Hamilton系\eqref{environment-activation-H}の位相図．太線は零エネルギーに対応し，矢印は時間発展の向きを示す．\label{Fig-Activation}}
\end{figure}


そのため，線形な変数変換$\Phi^{cl}(t)=q(t)$および$\Phi^q(t)=p(t)/(i\gamma)$を行おう．すると，散逸作用\eqref{environment-S-diss-viscous}はHamilton作用$iS=-\int\d t\,\big(p\,\dot{q} - H(p,q)
\big)$の形になる．ここで，有効Hamiltonian
\begin{equation}\label{environment-activation-H}
H(p,q)\equiv {2\over \gamma}\left[ -p\, {\partial U(q)\over
\partial q} + Tp^{\,2} \right] ,
\end{equation}を導入した．新しい変数を用いると，運動方程式\eqref{environment-escape-eq}がHamilton方程式$\dot q=\partial H/\partial
p$および$\dot p=-\partial H/\partial q$の形になることは容易にわかる．したがって，Hamiltonian \eqref{environment-activation-H}をもつHamilton系を調べればよい．これを可視化するには，$(p,q)$平面上に等エネルギー線$E=H(p(t),q(t))$からなる位相図を描けばよい（図\ref{Fig-Activation}b参照）．そのトポロジーは，二本の零エネルギー線$p=0$と$Tp=\partial U(q)/\partial q$によって決まる．これらはポテンシャルの二つの停留点$q=0$と$q=1$で交わる．$p=0$の線は，Langevin雑音のない古典ダイナミクスに対応し（この線に沿う運動では作用が恒等的に零であることに注意），したがって$q=0$は安定点，$q=1$は不安定点である．Liouvilleの定理により，どの固定点も一つの安定方向と一つの不安定方向をもたなければならない．したがって，「非古典的」な線$p=T^{-1}\partial U(q)/\partial q$に沿っては状況が逆転し，$q=0$が不安定，$q=1$が安定となる．ポテンシャル井戸の底$q=0$から脱出するには，系は零エネルギーの非古典的な線に沿って障壁の頂上$q=1$まで発展し，その後は古典運動方程式に従って運動を続けなければならない（すなわち古典的な線$p=0$に沿って動く）ことが，これで明らかになった．非古典的な線に沿う運動には，非零の作用
$$ iS=
-\int\d t\, p\dot q=-\int_0^1 p(q)\d q=-{1\over T}\int_0^1 {\partial
U(q)\over \partial q}\d q = -\,{U_0\over T}\,,$$
が伴う．ここでは軌道に沿って$H=0$であることを用いた．その結果，熱的な脱出確率は$e^{iS}=e^{-U_0/T}$に比例する．これはまさに熱活性化の指数因子である．


驚くべきことに，粘性的（あるいは拡散的）ダイナミクスをHamilton形式へ書き直すこの方法は，広い範囲の問題で機能する（たとえば文献\cite{Elgart}参照）．その代償は，自由度の数を二倍にすることである．Hamilton形式では$q$と$p$，Keldysh形式では「古典」成分と「量子」成分が必要になる．


%$$$$$$$$$$$$$$$$$$$$$$$$$$$$$$$$$$$$$$$$$$$$$$$$$$$$$$$$$$$$$$$$$$$$$$$$$$$$$$$$$$$$$$$$$$$$$$$$$$$$$$$$$$$
\subsectionlink{Fokker--Planck方程式}\label{sec_environment-6}


作用\eqref{environment-S-diss-viscous}に対する別のアプローチは，それが$\Phi^q$について二次であり，したがって$\D[\Phi^q]$の積分を明示的に実行できることに注目するものである．表記を短くし，古典座標との関係を強調するため，前節にならって$\Phi^{cl}(t)\equiv q(t)$と書く．式\eqref{environment-S-diss-viscous}で与えられる$S\big[\Phi^{cl},\Phi^q\big]$を用いて，$\exp\big(iS[\Phi]\big)$を$\Phi^q$についてGaussian積分すると，$\Phi^{cl}\equiv q$だけに依存する作用を得る．
\begin{equation}\label{environment-S-FP}
iS[q] = -{1\over 2\gamma T}\int_{-\infty}^{+\infty}\d t\,
\left({\gamma\over 2}\,\dot q+U'_q\right)^2 \, .
\end{equation}ここで，Feynman経路積分からSchr\"odinger方程式へ移る際と同じ方法~\cite{FeynmanHibbs}を用いることができる．すなわち，時刻$t+\delta_t$に点$q_N\equiv q$を通るすべての軌道について$\exp(iS[q])$の汎関数積分を行った結果として，「波動関数」$\mathcal{P}(q,t)$を導入する．最後の時間刻み$\delta_t$の積分を明示的に取り出すと，${\cal P}(q_N,t+\delta_t)$を${\cal
P}(q_{N-1},t)={\cal P}(q-\delta_q,t)$の$\delta_q\equiv
q-q_{N-1}$についての積分として，次のように書ける．
\begin{equation}
\fitdisplay{\begin{aligned}
 {\cal
P}(q,t+\delta_t)&=C\int_{-\infty}^{\infty}
\d[\delta_q] \exp\left(-{\delta_t\over 2\gamma T}\left[{\gamma\over
2}\,{\delta_q\over\delta_t}+ U_q'(q-\delta_q)\right]^2\right){\cal
P}(q-\delta_q,t)
\\
&=C\int_{-\infty}^{\infty} \d[\delta_q]
\exp\left(-{\gamma\over 8 T}
{\delta_q^2\over\delta_t}\right)\left[\exp\left(\,-{\delta_q\over
2T}\, U'_q(q-\delta_q) - {\delta_t\over 2\gamma T} \left( U'_q
\right)^2\right){\cal P}(q-\delta_q,t)\right]\,,
\end{aligned}}
\end{equation}ここで積分測度$C$は，条件$C\int\d[\delta_q]\,
\exp\big(-\gamma\delta_q^2/(8T\delta_t)\big)=1$によって定まる．最後の式の右辺の角括弧内を，$\delta_q$の二次および$\delta_t$の一次まで展開すると，次を得る．
\begin{eqnarray}
{\cal P}(t+\delta_t)
&=&
\left( 1+{\langle \delta_q^2\rangle\over
2T}U''_{qq} + {1\over 2} {\langle \delta_q^2\rangle\over 4T^2}
\left( U'_q \right)^2 - {\delta_t\over 2\gamma T} \left( U'_q
\right)^2\right){\cal P}+{\langle \delta_q^2\rangle\over
2T}U'_q{\cal P}'_q \nonumber
+{\langle \delta_q^2\rangle\over 2}{\cal P}''_{qq}  \\
&=&
{\cal P}(t)+\frac{2\delta_t}{\gamma}\left(U''_{qq}\,{\cal P}+U'_q {\cal
P}'_q+ T{\cal P}''_{qq}\right)\,,
\end{eqnarray}ここで
$$\langle \delta_q^2\rangle\equiv C\int_{-\infty}^{\infty}
\d[\delta_q] \, \delta^2_q \,
\exp\left(-\frac{\gamma\delta_q^2}{8T\delta_t}\right) =
4T\delta_t/\gamma\,.$$
である．最後に，上式を微分形に書き直すと，次を得る．
\begin{equation}
{\partial {\cal P}\over\partial t}={2\over \gamma}\left[
 {\partial \over \partial
q}{\partial U\over \partial q}  + T{\partial^2\over \partial
q^2}\right] {\cal P}
 ={2\over \gamma}\,
 {\partial \over \partial
q}\left[ {\partial U\over \partial q}\,{\cal P}  + T\,{\partial
{\cal P}  \over
\partial q}\right]\,.
\end{equation}これは確率分布関数${\cal P}(q,t)$の時間発展に対するFokker--Planck（FP）方程式である．この分布関数は，時刻$t$に点$q=\Phi^{cl}$で粒子を見いだす確率を記述する．初期に鋭い（決定論的な）分布${\cal P}(q,0)= \delta(q-q(0))$から出発すると，FP方程式の右辺第一項は，ポテンシャル$U(q)$中での粒子の粘性的なドリフトを表す．実際，第二項がない場合（$T=0$），この方程式の解は${\cal P}(q,t)= \delta(q-q(t))$であり，$q(t)$は決定論的な運動方程式$(\gamma/2)\dot q(t)=-\partial
U(q(t))/\partial q$を満たす\footnote{この主張を確かめるには，${\cal P}(q,t)= \delta(q-q(t))$を$T=0$のFP方程式に代入すればよい．すると$\delta'_q(q-q(t))(-\dot q(t)) =(2/\gamma)\left[U''_{\!
qq}\delta(q-q(t)) + U'_q\delta'_q(q-q(t))\right]$となる．次にこの式の両辺に$q$を掛け，$\d q$について部分積分すると，$\dot q(t) =
-(2/\gamma)U'_q(q(t))$を得る．}．右辺第二項は，熱的な確率雑音$\xi(t)$による確率分布の拡散的な広がりを表す．閉じ込めポテンシャル$U(q)$（$U(\pm\infty)\to \infty$を満たすもの）では，FP方程式の定常解は平衡Boltzmann分布${\cal P}(q)\sim \exp\{-U(q)/T\}$となる．


FP方程式は，（虚時間の）Schr\"odinger方程式$\dot {\cal P}=\hat H {\cal P}$とみなすことができる．ここでHamiltonian $\hat H$は，前節で導入した古典Hamiltonian \eqref{environment-activation-H}を「量子化」したものにほかならない．「量子化」の規則は$p\to \hat p\equiv -\partial/\partial q$であり，正準交換関係$[q,\hat p]=1$が成り立つ．この量子化規則を適用する前に，対応する古典Hamiltonianを\emph{正規順序化}しなければならないことに注意しよう．すなわち，運動量$\hat p$を座標$q$の左に置く必要がある［式\eqref{environment-activation-H}参照］．交換関係を用いると，量子化されたHamiltonianは
$$\hat H
=T\hat p^2-\hat pU'_q= T\left(\hat p - U'_q/(2T)\right)\left(\hat p
- U'_q/(2T)\right) - (U'_q)^2/(4T) + U''_{qq}/2$$
と書き直せる（$\gamma/2=1$とおいた）．そして，正準変換$Q=q$および$\hat P =\hat p -U'_q/(2T)$を行うことができる．新しい変数では，Hamiltonianはなじみ深い形$\hat H = T\hat P^2 +V(Q)$をとる．ここで$V(Q)= - (U'_Q)^2/(4T) + U''_{QQ}/2$であり，「波動関数」は$\tilde {\cal P}(Q,t) = e^{U(Q)/(2T)} {\cal
P}$と変換される．


%$$$$$$$$$$$$$$$$$$$$$$$$$$$$$$$$$$$$$$$$$$$$$$$$$$$$$$$$$$$$$$$$$$$$$$$$$$$$$$$$$$$$$$$$$$$$$$$$$$$$$$$$$$$
\subsectionlink{MatsubaraからKeldyshへ}\label{sec_environment-7}


応用によっては，平衡系の\emph{Matsubara}法~\cite{Matsubara,AGD}で作用を導き，後の段階でKeldysh表示へ移って非平衡問題を扱うと便利なことがある．本節では，この変換をどのように行えるかを示す．そのため，次のボソンのMatsubara作用を考えよう．
\begin{equation}\label{environment-S-Matsubara}
S[\Phi_m]=\gamma\, T\!\! \sum\limits_{m=-\infty}^\infty\!{1\over
2}\,  |\epsilon_m||\Phi_m|^2\, ,
\end{equation}ここで$\epsilon_m=2\pi T m$であり，$\Phi_m=\Phi_{-m}^*=\int_0^{\beta}
d\tau \Phi(\tau) e^{i\epsilon_m \tau}$は周期境界条件$\Phi(0)=\Phi(\beta)$を満たす実ボソン場$\Phi(\tau)$のMatsubara成分である．絶対値があるため，$|\epsilon_m|\neq i\partial_\tau$であることに注意されたい．実際，虚時間表示では，作用\eqref{environment-S-Matsubara}の核$\mathbb{K}_{m}=|\epsilon_{m}|$は$\mathbb{K}(\tau)=\sum_{m}|\epsilon_m|e^{-i\epsilon_{m}\tau}=C\delta(\tau)-\pi
T^2/\sin^{2}(\pi T\tau)$の形になる．無限大の定数$C$は，$\int^{\beta}_{0}\d\tau\mathbb{K}(\tau)=\mathbb{K}_0=0$を満たすように選ばれる．したがって，虚時間表示における作用\eqref{environment-S-Matsubara}は，次の非局所的な形をとる．
\begin{eqnarray}\label{environment-S-CL}
S[\Phi]&=&\frac{\gamma
T}{2}\iint^{\beta=1/T}_{0}\d\tau\d\tau'\,\Phi(\tau)\mathbb{K}(\tau-\tau')\Phi(\tau')
\nonumber\\
&=&{\gamma\over 4\pi}\iint_{0}^{\beta} \d\tau\, \d\tau' \,
\frac{\pi^2 T^2}{\sin^2[\pi T(\tau-\tau')]}
\,\Big(\Phi(\tau)-\Phi(\tau')\Big)^2\, .
\end{eqnarray}この作用は，散逸が量子トンネル現象に与える影響を調べるために用いたCaldeiraとLeggett~\cite{CaldeiraLeggett}にちなんで呼ばれることが多い．


Keldysh表示へ変換するには，自由度の数を$\Phi\to
\vec{\Phi}=(\Phi^{cl},\Phi^q)^T$と二倍にする必要がある．すると，第\ref{sec_bosons-4}節の因果性の構造に従って，時間並進不変なKeldysh作用の一般形は次のようになる．
\begin{equation}
S\big[\Phi^{cl},\Phi^{q}\big] = \gamma\int{\d\epsilon\over 2\pi}\,
\big(\Phi^{cl}_{\epsilon},\Phi^q_{\epsilon}\big)
\left(\begin{array}{cc}
0   & \big[\mathfrak{D}^{A}(\epsilon)\big]^{-1}  \\
 \big[\mathfrak{D}^{R}(\epsilon)\big]^{-1}  & \big[\mathfrak{D}^{-1}(\epsilon)\big]^{K}
\end{array}\right)
\left(\begin{array}{c} \Phi^{cl}_{\epsilon} \\ \Phi^q_{\epsilon}
\end{array}\right) \, ,
\end{equation}ここで$[\mathfrak{D}^{R(A)}(\epsilon)]^{-1} $は，Matsubara相関関数$|\epsilon_m|/2$を，複素変数$\epsilon_m$の\emph{上（下）半平面}から実軸へ，$-i\epsilon_m\to \epsilon$により解析接続したものである（文献\cite{AGD}参照）．その結果$\big[\mathfrak{D}^{R(A)}(\epsilon)\big]^{-1}=\pm
i\epsilon/2$となる．平衡では，Keldysh成分はFDTから$\big[\mathfrak{D}^{-1}(\epsilon)\big]^{K}=
([\mathfrak{D}^{R}]^{-1}-[\mathfrak{D}^{A}]^{-1}) \coth\,
(\epsilon/2\,T)=i\epsilon\,\coth\, (\epsilon/2\,T)$と定まる［式\eqref{environment-L-RA}および\eqref{environment-L-K}参照］．したがって，Matsubara作用\eqref{environment-S-Matsubara}または\eqref{environment-S-CL}のKeldysh版は，すでに見た散逸作用\eqref{environment-S-diss-fin}である（もちろん，ポテンシャル項と慣性項は除く）．ここに外場を含めれば，系を平衡から離れさせることができる．


%$$$$$$$$$$$$$$$$$$$$$$$$$$$$$$$$$$$$$$$$$$$$$$$$$$$$$$$$$$$$$$$$$$$$$$$$$$$$$$$$$$$$$$$$$$$$$$$$$$$$$$$$$$$
\subsectionlink{散逸する鎖と膜}\label{sec_environment-8}


ここでは浴に結合した単一粒子の代わりに，それぞれが浴に結合し，互いにも結合した粒子からなる鎖または格子を考えよう．そのため，（i）場に空間添字${\bf r}$を加えて$\Phi(t)\to \Phi(\mathbf{r},t)$とし，（ii）隣接粒子間の調和相互作用ポテンシャルを加える．これは連続極限で$\sim D(\Phi(\mathbf{r},t)-\Phi(\mathbf{r}
+\mathbf{1},t))^2\to D(\partial_\mathbf{r}\Phi)^2$となり，$D$は鎖または膜の剛性である．古典・量子成分に変換し，空間について部分積分すると［式\eqref{app-SPQM-S-Phonons-1}参照］，勾配項は$D\left(\Phi^q\partial^2_\mathbf{r}\Phi^{cl} +
\Phi^{cl}\partial^2_\mathbf{r}\Phi^{q}\right)$になる．したがって，これは相関関数の遅延成分と先進成分を変えるが，$(q-q)$のKeldysh成分には影響\emph{しない}．
\begin{equation}
\big[\mathfrak{D}^{R(A)}\big]^{-1}={1\over 2}\,
\delta(t-t')\,\delta(\mathbf{r}-\mathbf{r}')\big( \mp
\partial_{t} + D\,\partial^2_{\mathbf{r}} \big)\, .
\end{equation}Fourier表示では$\big[\mathfrak{D}^{R(A)}(\mathbf{k},\epsilon)\big]^{-1}={1\over 2}
\big(\pm i\epsilon -D\mathbf{k}^2\big)$である．平衡ではKeldysh成分は勾配項の影響を受けず，式\eqref{environment-L-K}で与えられる（実空間表示では因子$\delta(\mathbf{r}-\mathbf{r}')$が付く）．特に，その古典極限は$\big[\mathfrak{D}^{-1}\big]^K=i2T\delta(t-t')\delta(\mathbf{r}-\mathbf{r}')$である［式\eqref{eqenvironment-L-K-time}参照］．その結果，浴に接した古典的な弾性膜の作用は次のようになる．
\begin{equation}\label{environment-S-chain}
S[\Phi^{cl},\Phi^{q}]=2\iint\d\mathbf{r}\d t\left[-\Phi^q
\left(\partial_t \Phi^{cl}-D\partial^2_\mathbf{r}\Phi^{cl}+{\partial
U(\Phi^{cl})\over
\partial \Phi^{cl} }\right)+i2T\big[\Phi^{q}\big]^2\right]\,,
\end{equation}ここでは慣性項を無視し，簡潔さのため$\gamma/2=1$とおいた．次に，補助的なHubbard--Stratonovich場$\xi(\mathbf{r},t)$を導入し，第\ref{sec_environment-4}節に従ってLangevin方程式を書くことができる．
\begin{equation}
\partial_{t}\Phi^{cl} -D\partial^2_\mathbf{r}\Phi^{cl}+{\partial U(\Phi^{cl})\over
\partial \Phi^{cl} } =\xi(\mathbf{r},t)\, ,
\end{equation}ここで$\xi$は，短距離相関$\langle
\xi(\mathbf{r},t)\xi(\mathbf{r}',t')\rangle=2T\delta(t-t')
\delta(\mathbf{r}-\mathbf{r}')$をもつGaussian雑音である．


図\ref{Fig-Activation}aに描いた，$\mathbf{r}$に依存しない準安定なポテンシャル井戸の底に置かれた弾性鎖を考えよう．鎖の十分大きな部分が熱的に井戸から脱出すると，ポテンシャルを下っていくことが有利になり，鎖全体を井戸から引き出す可能性がある．このような最適な大きさの臨界領域の形状とその作用を求めるため，第\ref{sec_environment-7}節のHamilton変数$q(\mathbf{r},t)\equiv
\Phi^{cl}(\mathbf{r},t) $および$p(\mathbf{r},t)\equiv
2i\Phi^q(\mathbf{r},t)$に移ろう．作用\eqref{environment-S-chain}はHamilton形式$iS=-\iint\d\mathbf{r}\d t\,\big(p\,\dot
q - H(p,q) \big)$となり，
\begin{equation}\label{environment-chain-H}
H\equiv p\, D\partial_{\mathbf{r}}^2 q - p\, {\partial U(q)\over
\partial q} +  Tp^{\,2}  ,
\end{equation}である．対応する運動方程式は次のようになる．
\begin{subequations}
\begin{equation}
\dot q = \frac{\delta H}{\delta p}\, =\, D\partial^2_\mathbf{r} q -
U'_q(q)
+2Tp\,,
\end{equation}
\begin{equation}
\hskip-.4cm \dot p =-\frac{\delta H}{\delta q} =
-D\partial^2_\mathbf{r} p + p\,U''_{\!qq}(q)\,.
\end{equation}\end{subequations}
これらは複雑な偏微分方程式であり，一般には解けない．幸い，最適な臨界領域の形状は求めることができる．第\ref{sec_environment-7}節で論じたように，最小作用の軌道は零エネルギー$H=0$の運動に対応する．式\eqref{environment-chain-H}によれば，これは$p=0$（作用が零の古典軌道），または$Tp= U'_q(q)-D
\partial_\mathbf{r}^2 q $（有限の作用をもつ脱出軌道）のいずれかの場合に成り立つ．後者の場合，$q(\mathbf{r},t)$の運動方程式は\emph{時間を反転した}古典方程式$\dot q =-
D\partial^2_\mathbf{r} q+U'_q(q)=Tp\,$の形をとる．最後の等号により，$p(\mathbf{r},t)$の運動方程式も自動的に満たされる\footnote{実際，$T\dot p =\partial_t \dot
q=-D\partial^2_\mathbf{r} \dot q+\dot q
U''_{\!qq}=T(-D\partial^2_\mathbf{r} p + pU''_{\!qq} )$である．この非自明な事実は，偶発的な保存則$H\big(p\rt,q\rt\big)=0$が\emph{局所的}に存在することを反映している．一般原理から要請されるのは，大域的な全エネルギーの保存だけである．}．時間を反転したダイナミクスでは，配置$q(\mathbf{r},t)=0$は不安定であるため，鎖には「舌状」の突出部が生じ，次の定常形状に達するまで成長する．
\begin{equation}\label{environment-domain-eq}
- D\partial^2_\mathbf{r} q + U'_q(q)=0\, .
\end{equation}この方程式を境界条件$q(\pm\infty)=0$のもとで解けば，臨界領域の形状が得られる．いったん形成されると，それは古典方程式$\dot q = D\partial^2_\mathbf{r} q - U'_q(q)$および$p=0$に従い，作用零でさらに成長できる．この「舌状」領域を形成するために必要な，非古典的脱出軌道に沿う作用は，$H(p,q)=0$であることから，次のようになる．
\begin{equation}\label{environment-S-chain-static}
iS=-\iint\d\mathbf{r}\d t\,p\, \dot q
=-{1\over T} \iint\d\mathbf{r}\d t\,
\left(-D \partial^2_\mathbf{r} q + U'_q(q)\right) \dot q
=-{1\over T}\int\d\mathbf{r}\,\Big({D\over 2} (\partial_\mathbf{r}
q)^2 + U(q)\Big)\, ,
\end{equation}最後の等号では，時間についての積分を明示的に行った．したがって脱出作用は，弾性エネルギーとポテンシャルエネルギーの両方を含む静的な活性化の式で与えられる．最適な領域\eqref{environment-domain-eq}は，この静的作用\eqref{environment-S-chain-static}を最小化することで求められる．こうして，一次相転移の熱力学的なLandau型の記述に到達する．この有効な熱力学的記述は，$H(p,q)=0$を仮定し，すべての過程に無限に長い時間がかかることから生じる点に注意されたい．



%$$$$$$$$$$$$$$$$$$$$$$$$$$$$$$$$$$$$$$$$$$$$$$$$$$$$$$$$$$$$$$$$$$$$$$$$$$$$$$$$$$$$$$$$$$$$$$$$$$$$$$$$$$$
%$$$$$$$$$$$$$$$$$$$$$$$$$$$$$$$$$$$$$$$$$$$$$$$$$$$$$$$$$$$$$$$$$$$$$$$$$$$$$$$$$$$$$$$$$$$$$$$$$$$$$$$$$$$


\sectionlink{フェルミ粒子}\label{sec_fermion}
\subsectionlink{分配関数}\label{sec_fermion-1}



エネルギー$\epsilon_0$をもつ単一の量子状態を考えよう．この状態にはスピンのないフェルミ粒子（Pauliの排他原理に従う粒子）が入る．実際，この状態に存在できる粒子数はゼロまたは1である．このような系の第二量子化されたHamiltonianは次の形をとる：
\begin{equation}\label{fermion-H}
\hat{H}=\epsilon_0\,\hat{c}^\dagger \hat{c}\, ,
\end{equation}
ここで，$\hat{c}^\dagger$と$\hat{c}$は，状態$\epsilon_0$におけるフェルミ粒子の生成演算子と消滅演算子である．これらは通常の反交換関係$\{\hat{c}\,,\hat{c}^\dagger\}=1$および$\{\hat{c}\,,\hat{c}\}=\{\hat{c}^\dagger\,,\hat{c}^\dagger\}=0$に従う．ただし，$\{\,\, ,\,\}$は反交換子を表す．



ここで，Keldysh経路$\mathcal{C}$に沿った発展演算子と，それに対応する分配関数$Z=1$を考えることができる．これらはボソン系の場合の式~\eqref{boson-Z}とまったく同じように定義される．平衡密度行列のトレースは$\Tr\{\hat{\rho}_0\}=1+\rho(\epsilon_0)$であり，2つの項は空の状態と1粒子が占有する状態を表す．Keldysh経路を，長さ$\delta_t\sim 1/N\to 0$の$(2N-2)$個の時間区間に分割し，Keldysh経路$\mathcal{C}$上の$2N$個の点で単位演算子の分解を挿入する（図~\ref{Fig-Contour}参照）．第~\ref{sec_bosons-1}節のボソンの場合との唯一の違いは，ここではフェルミ粒子のコヒーレント状態基底による単位演算子の分解を用いることである\footnote{Grassmann数$\psi$（$\{\psi,\psi'\}=\{\psi,\hat{c}\}=0$を満たす）でパラメータ表示されるフェルミ粒子のコヒーレント状態$|\psi\rangle\equiv (1-\psi
\hat{c}^\dagger)|0\rangle$は，消滅演算子の固有状態である：$\hat{c}|\psi\rangle=\psi|\psi\rangle$．同様に，$\langle\psi|\hat{c}^\dagger=\langle\psi|\bar{\psi}$である．ここで，$\bar{\psi}$は$\psi$とは{\em 無関係な}別のGrassmann数である．Hamiltonianなどの{\em 正規順序化された}演算子の行列要素は，$\langle\psi|\hat{H}(\hat{c}^\dagger,\hat{c})|\psi'\rangle =
H(\bar{\psi},\psi') \langle\psi|\psi'\rangle$の形をとる．任意の2つのコヒーレント状態の重なりは$\langle\psi|\psi'\rangle=1+\bar\psi\psi'=\exp\{\bar\psi\psi'\}$である．演算子$\hat{\mathcal{O}}$のトレースは$\Tr\big\{\hat{\mathcal{O}}\big\}= \langle 0| \hat{\mathcal{O}}|0
\rangle + \langle 1| \hat{\mathcal{O}}|1 \rangle = \langle 0|
\hat{\mathcal{O}}|0 \rangle + \langle 0|\hat{c}\,
\hat{\mathcal{O}}\, \hat{c}^\dagger |0 \rangle=  \iint\d\bar{\psi}\,
\d\psi\,
e^{-\bar{\psi}\psi}\langle-\psi|\hat{\mathcal{O}}|\psi\rangle $と計算される．ここで，Grassmann積分は$\int\! \d\psi\,
1=0$および$\int\! \d\psi\, \psi =1$により{\em 定義}される．}：
\begin{equation}\label{fermion-unity}
\hat{1}=\iint\d\bar{\psi}_j\, \d\psi_j\,\,
e^{-\bar{\psi}_{j}\,\psi_{j}}\, |\psi_j\rangle\langle\psi_j|\, ,
\end{equation}
ここで，$\bar{\psi}_j$と$\psi_j$は{\em 互いに独立な}Grassmann変数である．残りの代数計算はボソンの場合とまったく同様に進む（第~\ref{sec_bosons-1}節参照）．その結果，次式を得る：
\begin{equation}\label{fermion-Z}
Z={1\over \Tr\{\hat{\rho}_{0}\}}\iint
\prod\limits_{j=1}^{2N}\left[\d\bar{\psi}_j\, \d\psi_j \right] \,\,
\exp\left(\,i \sum\limits_{j,j\,'=1}^{2N} \bar{\psi}_{j}\,
G^{-1}_{jj\,'}\,\psi_{j\,'}\right)\, ,
\end{equation}
ここで，$2N \times 2N$行列$G^{-1}_{jj\,'}$は
\begin{equation}\label{fermion-G-matrix}
  iG^{-1}_{jj\,'}\equiv
\left[\begin{array}{rrr|rrr}
 -1   &      &    &   &   &   -\rho(\epsilon_0) \\
  1\!-\!h & -1   &    &   &   &                 \\
      &  1\!-\!h & -1 &   &   &                 \\ \hline
%---------------------------------------------------------
     &     &  1 & -1 &    &                 \\
     &     &    &1\!+\!h & -1 &                 \\
     &     &    &    & 1\!+\!h&  -1
 \end{array} \right]\, ,
\end{equation}
であり，$h\equiv i\epsilon_0\delta_t$である．ボソンの場合との唯一の違いは，行列要素$\rho(\epsilon_0)$の前の負号である．これはフェルミ粒子のトレースの式に現れるコヒーレント状態$\langle-\psi_{2N}|$の負号に由来する．規格化を確認するため，この行列の行列式を計算しよう：
\begin{equation}\label{fermion-determinant}
\mathrm{Det}\big[i\hat{G}^{-1}\big]=1+\rho(\epsilon_0)(1-h^2)^{N-1}
\approx1+ \rho(\epsilon_0)\,e^{(\epsilon_0\delta_t)^2(N-1)}\to
1+\rho(\epsilon_0)\,.
\end{equation}
フェルミ粒子のGauss積分は，相関行列の行列式で与えられる（ボソンでは行列式の逆数となる）のを用いると（詳細は\ref{app_Gaussian}参照），次式を得る：
\begin{equation}\label{fermion-Z-unity}
Z=\frac{\mathrm{Det}\big[i\hat{G}^{-1}\big]}{\Tr\{\hat{\rho}_0 \} }
= 1\, ,
\end{equation}
これは期待どおりの結果である．ここでも，正しい規格化を保つには離散行列~\eqref{fermion-G-matrix}の右上の要素が不可欠である．$N\to \infty$の極限をとり，連続表記$\psi_j\to \psi(t)$を導入すると，次式を得る：
\begin{equation}\label{fermion-Z-1}
Z=\int\D[\bar{\psi} \psi]\, \exp\left(iS[\bar{\psi},\psi]\right) =
\int\D[\bar{\psi} \psi] \, \exp\left(i \int_{\mathcal{C}}\d t\,\big[
\bar\psi(t)\,\hat{G}^{-1}\psi(t)\big] \right),
\end{equation}
ここで，式~\eqref{fermion-Z}および\eqref{fermion-G-matrix}より，作用は次式で与えられる：
\begin{equation}\label{fermion-S}
S[\bar{\psi},\psi]=
\sum\limits_{j=2}^{2N}\left[i\bar{\psi}_j\,\frac{\psi_j-\psi_{j-1}}{\delta
t_j} -\epsilon_0\bar{\psi}_{j}\,\psi_{j-1}\right]\delta
t_j\,+i\,\bar{\psi}_1\Big[\psi_1+\rho(\epsilon_0)\psi_{2N}\Big]\, ,
\end{equation}
ただし，$\delta t_j\equiv t_j-t_{j-1}=\pm \delta_t$である．したがって，演算子$\hat G^{-1}$の連続表現はボソンの場合の式~\eqref{boson-G-continious}と同じであり，$\hat G^{-1}=
i\partial_t-\epsilon_0$である．ここでも，分布関数の情報を含む離散行列の右上の要素（式~\eqref{fermion-S}の最後の項）は，連続表記では見かけ上失われている．



Grassmann場$\psi(t)$を，時間経路の順方向部分と逆方向部分にそれぞれ属する2つの成分$\psi_{+}(t)$と$\psi_{-}(t)$に分けると，作用は次のように書き直せる：
\begin{equation}\label{fermion-S-1}
S[\bar{\psi},\psi]=\int_{-\infty}^{+\infty}\d t\,\left[
\bar{\psi}_{+}(t)(i\partial_t-\epsilon_0)
\psi_{+}(t)-\bar{\psi}_{-}(t)(i\partial_t-\epsilon_0)
\psi_{-}(t)\right] ,
\end{equation}
ここで，離散行列~\eqref{fermion-G-matrix}に非ゼロの非対角ブロックが存在するため，$\psi_{+}$と$\psi_{-}$のダイナミクスは実際には互いに独立では{\em ない}．


%$$$$$$$$$$$$$$$$$$$$$$$$$$$$$$$$$$$$$$$$$$$$$$$$$$$$$$$$$$$$$$$$$$$$$$$$$$$$$$$$$$$$$$$$$$$$$$$$$$$$$$$$$$$

\subsectionlink{Green関数とKeldysh回転}\label{sec_fermion-2}



4つのフェルミ粒子Green関数$G^{\mathbb{T}(\widetilde{\mathbb{T}})}$および$G^{<(>)}$は，対応するボソンのGreen関数と同様に定義される（式~\eqref{boson-corr-fun}参照）：
\begin{subequations}\label{fermion-corr-fun}
\begin{equation}
\hskip-2cm \langle\psi_{+}(t)\bar{\psi}_{-}(t\,')\rangle\equiv
iG^{<}(t,t')=-n_{F}\exp\{-i\epsilon_{0}(t-t')\}\,,
\end{equation}
\begin{equation}
\hskip-1.7cm \langle\psi_{-}(t)\bar{\psi}_{+}(t\,')\rangle\!\equiv
iG^{>}(t,t')=(1\!-\!n_{F})\exp\{-i\epsilon_{0}(t-t')\}\,,
\end{equation}
\begin{equation}
\langle\psi_{+}(t)\bar{\psi}_{+}(t\,')\rangle\equiv
iG^{\mathbb{T}}(t,t')
=\theta(t-t')iG^{>}(t,t')+\theta(t'-t)iG^{<}(t,t')\,,
\end{equation}
\begin{equation}
\langle\psi_{-}(t)\bar{\psi}_{-}(t\,')\rangle\equiv
iG^{\widetilde{\mathbb{T}}}(t,t')=
\theta(t'-t)iG^{>}(t,t')+\theta(t-t')iG^{<}(t,t')\,.
\end{equation}
\end{subequations}
ただし，反交換関係のために$G^<$の式に負号が現れる点と，Bose占有数がFermi占有数$n_B\to
n_F\equiv \rho(\epsilon_0)/(1+\rho(\epsilon_0))$に置き換わる点が異なる．式~\eqref{boson-relation}および\eqref{boson-relation-1}は，フェルミ粒子のGreen関数に対しても成り立つ．



フェルミ粒子の場合には，ボソンの場合とは異なる方法でKeldysh回転を行うのが慣例である．新しい場を次のように定義する：
\begin{equation}\label{fermion-rotation}
\psi_{1}(t)={1\over \sqrt{2}}\big(\psi_{+}(t)+\psi_{-}(t)\big)\,,
\qquad
\psi_{2}(t)={1\over\sqrt{2}}\big(\psi_{+}(t)-\psi_{-}(t)\big)\,.
\end{equation}
LarkinとOvchinnikov~\cite{LarkinOvchinnikov}に従い，バーの付いた場は別の仕方で変換することにする：
\begin{equation}\label{fermion-rotation-1}
\bar{\psi}_{1}(t)={1\over\sqrt{2}}\big(\bar{\psi}_{+}(t)-
\bar{\psi}_{-}(t)\big)\,,\qquad \bar{\psi}_{2}(t)=
{1\over\sqrt{2}}\big(\bar{\psi}_{+}(t)+\bar\psi_{-}(t)\big)\,.
\end{equation}
要点は，Grassmann場$\bar\psi$が$\psi$の共役では{\em なく}，完全に独立な場だということである．したがって，（変換行列の行列式がゼロでない限り）任意の方法で変換できる．Grassmann積分は常に収束するため，積分の収束に関する問題はないことに注意しよう．また，Grassmann変数は決して古典的な意味をもたないため，添字$cl$と$q$も用いない．実際，$\bar\psi,\psi$を用いて鞍点方程式やその他の方程式を書くことはできず，これらは計算のいずれかの段階で必ず積分消去しなければならない．



式~\eqref{fermion-rotation}および\eqref{fermion-rotation-1}を式~\eqref{fermion-corr-fun}とともに用いると，次式を得る：
\begin{equation}\label{fermion-G-fun}
-i\big\langle\psi_{a}(t)\bar{\psi}_{b}(t\,')\big\rangle=
G_{ab}(t,t') = \left(\begin{array}{cc}
G^R(t,t') & G^{K}(t,t') \\
0                & G^{A}(t,t')
\end{array}\right)\, ,
\end{equation}
ここで，以後$a,b  = (1,2)$とする．この行列の$(2,1)$要素がゼロであることは，恒等式~\eqref{boson-relation}の現れである．Green関数~\eqref{fermion-G-fun}の{\em 遅延}，{\em 先進}，{\em Keldysh}成分は，ボソンに対する式~\eqref{boson-G-fun-RAK}とまったく同じように，$G^{\mathbb{T}(\widetilde{\mathbb{T}})}$と$G^{<(>)}$で表される．したがって，式~\eqref{boson-G-conjugation}--\eqref{boson-GR-plus-GA}と同じ対称性をもつ．式~\eqref{bosons-G-traces}と\eqref{boson-GR-plus-GA}から得られる重要な帰結は
\begin{equation}\label{fermion-traces}
\Tr\left\{G_{ab}^{(1)}\circ G_{bc}^{(2)}\circ\ldots\circ G_{za}^{(l)}
\right\}(t,t) =0\, ,
\end{equation}
である．ここで，丸付きの積の記号は，$2\times 2$行列の積に加えて時間領域の畳み込みも含む．引数$(t,t)$は，$G^{(1)}$の最初の時間引数と$G^{(l)}$の最後の時間引数が等しいことを示す．



フェルミ粒子のGreen関数は，対応するボソンのGreen関数~\eqref{boson-G-fun-2}とは異なる構造をもつことに注意しよう．行列中の$R,A$成分と$K$成分の位置が入れ替わっている．もちろん，これはバーの付いた場の変換規約が異なるためである．フェルミ粒子の規約をボソンの規約と同じに選び（ただし，その逆は{\em できない}），フェルミ粒子に対してもボソンと同じ構造~\eqref{boson-G-fun-2}を得ることは可能である．Larkin--Ovchinnikovの選択~\eqref{fermion-G-fun}の利点は，逆Green関数$\hat{G}^{-1}$とフェルミ粒子の自己エネルギー$\hat{\Sigma}_F$が，$\hat{G}$と同じ形をもつことである．すなわち，
\begin{equation}\label{fermion-G-inverse}
\hat{G}^{-1}=\left(\begin{array}{cc}
\big[G^{R}\big]^{-1} & \big[G^{-1}\big]^{K} \\
0                & \big[G^{A}\big]^{-1}
\end{array}\right),\qquad
\hat{\Sigma}_F = \left(\begin{array}{cc}
\Sigma_F^R       &  \Sigma_F^{K} \\
0                &  \Sigma_F^{A}
\end{array}\right)\, ,
\end{equation}
となる．一方，ボソンの場合には，$\hat{G}^{-1}$（式~\eqref{boson-S-1}参照）と$\hat{\Sigma}$（式~\eqref{int-bos-Sigma}参照）は，$\hat{G}$（式~\eqref{boson-G-fun-2}参照）とは異なる形をとる．この事実により，式~\eqref{fermion-G-fun}と\eqref{fermion-G-inverse}の形には計算上の利点がある．



単一のフェルミ粒子状態については，Keldysh回転の後，相関関数~\eqref{fermion-corr-fun}から行列~\eqref{fermion-G-fun}の各成分が求まる：
\begin{subequations}\label{fermion-G-fun-RAK}
\begin{equation}
\hskip-1.45cm
G^{R}(t,t\,')=-i\theta(t-t\,')e^{-i\epsilon_{0}(t-t')}\to
(\epsilon-\epsilon_{0}+i0)^{-1}\,,
\end{equation}
\begin{equation}
\hskip-1.7cm G^{A}(t,t\,')=i\theta(t'-t)e^{-i\epsilon_{0}(t-t')}\to
(\epsilon-\epsilon_{0}-i0)^{-1}\,,
\end{equation}
\begin{equation}
G^{K}(t,t\,')=-i(1-2n_{F})e^{-i\epsilon_{0}(t-t')}\to -2\pi
i(1-2n_{F})\delta(\epsilon-\epsilon_{0})\,,
\end{equation}
\end{subequations}
ここで，右辺にはFourier変換も併記した．熱平衡では，次式を得る：
\begin{equation}\label{fermion-FDT}
G^{K}(\epsilon)
=\left[G^{R}(\epsilon)-G^{A}(\epsilon)\right]\tanh{\epsilon \over
2\,T} \, .
\end{equation}
これがフェルミ粒子に対する揺動散逸定理（FDT）である．ボソンの場合と同じく，FDTはこの単純な模型に限らない平衡系の一般的性質である．一般に，反HermiteなKeldysh Green関数を，Hermite行列$F=F^{\dagger}$を用いて次のようにパラメータ表示すると便利である：
\begin{equation}\label{fermion-FDT-general}
G^{K}=G^{R}\circ F-F\circ G^{A} \, .
\end{equation}
$F(t,t')$のWigner変換はフェルミ粒子の分布関数の役割を果たす．


%$$$$$$$$$$$$$$$$$$$$$$$$$$$$$$$$$$$$$$$$$$$$$$$$$$$$$$$$$$$$$$$$$$$$$$$$$$$$$$$$$$$$$$$$$$$$$$$$$$$$$$$$$$$

\subsectionlink{自由フェルミ場とその作用}\label{sec_fermion-3}



次に，添字$\mathbf{k}$で区別される多数の自由度をもつ系に進もう．そのためには，$\epsilon_0\to \epsilon_\mathbf{k}$と置き換え，$\mathbf{k}$について和をとればよい．$\mathbf{k}$が運動量であり，$\epsilon_\mathbf{k}=\mathbf{k}^2/(2m)$である場合には，座標空間表示$\psi(\mathbf{k},t)\to \psi(\mathbf{r},t)$に移ると見通しがよい．このとき，$\epsilon_\mathbf{k}=\mathbf{k}^2/(2m)\to
-\partial^2_\mathbf{r}/(2m)$となる．最終的に，相互作用のないフェルミ気体のKeldysh作用は次の形をとる：
\begin{equation}\label{fermion-S-2}
S_{0}[\bar{\psi},\psi]=\iint \d x\, \d x'\sum\limits_{a,b=1}^2
\bar{\psi}_a(x) \big[\hat G^{-1}(x,x')\big]_{ab} \,\psi_{b}(x')\, ,
\end{equation}
ここで，$x=(\mathbf{r},t)$であり，相関行列$[\hat
G^{-1}]_{ab}$は式~\eqref{fermion-G-inverse}の構造をもち，
\begin{equation}\label{fermion-G-inverse-1}
\big[G^{R(A)}(x,x')\big]^{-1} =\delta(x-x')\left(i\partial_{t} +
{1\over 2m}
\partial_{\mathbf{r}}^2 + \mu \right)\, .
\end{equation}
である．連続表記では$R$成分と$A$成分は見かけ上同じに見えるが，離散時間表示ではそれぞれ主対角線より下側と上側に構造をもつ行列であることを忘れてはならない．Keldysh成分は連続極限をもたないという意味で，純粋な正則化項である（自己エネルギーのKeldysh成分は非ゼロの連続表現をもつ）．しかし，これらの情報はすべてGreen関数~\eqref{fermion-G-fun}の構造にすでに正しく取り込まれている．


%$$$$$$$$$$$$$$$$$$$$$$$$$$$$$$$$$$$$$$$$$$$$$$$$$$$$$$$$$$$$$$$$$$$$$$$$$$$$$$$$$$$$$$$$$$$$$$$$$$$$$$$$$$$

\subsectionlink{外場とソース}\label{sec_fermion-4}



Keldysh法の基本的な考え方によれば，分配関数$Z=1$は構成上規格化されている（式~\eqref{fermion-Z-unity}参照）．理論全体を意味あるものにするには，粒子密度や電流などのさまざまな観測量を計算できる補助的なソース場を導入しなければならない．例えば，経路$\mathcal{C}$に沿って定義された，時間依存する外部スカラーポテンシャル$V\rt$を導入できる．これは$S_{V}=\int\d\mathbf{r}\int_{\mathcal{C}}\d t\,
V\rt\bar{\psi}\rt\psi\rt$の形でフェルミ粒子と相互作用する．経路の順方向および逆方向の枝に属する場の成分でこれを表すと，次式を得る：
\begin{eqnarray}\label{fermion-SV}
S_{V}&=&\int\d\mathbf{r}\int^{+\infty}_{-\infty}\d
t\,\big[V_{+}\bar{\psi}_{+}\psi_{+}-V_{-}\bar{\psi}_{-}\psi_{-}\big]\nonumber
\\
&=&\int\d\mathbf{r} \int^{+\infty}_{-\infty}\d
t\,\big[V^{cl}(\bar{\psi}_{+}\psi_{+} - \bar{\psi}_{-}\psi_{-})+
V^{q}(\bar{\psi}_{+}\psi_{+} + \bar{\psi}_{-}\psi_{-})\big]\nonumber
\\
&=&\int\d\mathbf{r}\int^{+\infty}_{-\infty}\d t\
[V^{cl}(\bar{\psi}_{1}\psi_{1}+\bar{\psi}_{2}\psi_{2})+
V^{q}(\bar{\psi}_{1}\psi_{2}+\bar{\psi}_{2}\psi_{1})]\,,
\end{eqnarray}
ここで，$V^{cl(q)}$成分は，実ボソン場に対する通常の方法で$V^{cl(q)} = (V_{+}\pm V_{-})/2$と定義される．また，式~\eqref{fermion-rotation}と\eqref{fermion-rotation-1}に従い，$\psi_{\pm}$から$\psi_{1(2)}$への回転を行った．物理的なフェルミ粒子密度（Keldysh経路の2つの枝について対称化したもの）$\varrho = {1\over 2} \big(
\bar\psi_{+}\psi_{+}+ \bar\psi_{-}\psi_{-} \big)$は，ソース場の量子成分$V^q$に結合することに注意しよう．一方，ソースの古典成分$V^{cl}$は，2つの枝で同一の，物理的な外部スカラーポテンシャルにほかならない．



2つの頂点$\hat \gamma$行列を導入すると，表記を大幅に簡潔化できる：
\begin{equation}\label{fermion-gammas}
\hat{\gamma}^{cl} \equiv \left( \begin{array}{cc}
 1 & 0   \\
 0 & 1
 \end{array} \right)\, , \hskip 1cm
\hat{\gamma}^{q} \equiv \left( \begin{array}{cc}
 0 & 1   \\
 1 & 0
 \end{array} \right)\, .
\end{equation}
これらの定義を用いると，ソース作用~\eqref{fermion-SV}は次のように書ける：
\begin{eqnarray}\label{fermion-SV-1}
S_{V}=\int\d\mathbf{r}\int^{+\infty}_{-\infty}\d
t\sum^{2}_{a,b=1}\left[V^{cl}\bar{\psi}_{a}\gamma^{cl}_{ab}\psi_{b}+
V^{q}\bar{\psi}_{a}\gamma^{q}_{ab}\psi_{b}\right]
=\Tr\big\{\vec{\bar{\Psi}}\hat{V}\vec{\Psi}\big\},
\end{eqnarray}
ここでは，次のように定義されるKeldysh二重項$\vec{\Psi}$と行列$\hat{V}$を導入した：
\begin{equation}
\vec{\Psi}=\left(\begin{array}{c}\psi_{1} \\
\psi_{2}\end{array}\right)\,,\hskip 1cm
\hat{V}=V^{\alpha}\hat\gamma^{\alpha}= \left(\begin{array}{cc}V^{cl} & V^{q} \\
V^{q} & V^{cl}\end{array}\right)\,,
\end{equation}
ただし，$\alpha=(cl,q)$である．



同様に，外部ベクトルポテンシャルもこの形式に導入できる．対応する作用の部分\footnote{ここで用いるベクトルソース$\mathbf{A}\rt$は，実際のベクトルポテンシャルとは係数$e/c$だけ異なる．ただし，これをベクトルポテンシャルと呼び，最終的な式で電子の電荷を復元する．}$S_{A}=\int\d\mathbf{r}\int_{\mathcal{C}}\d
t\, \mathbf{A}\rt\mathbf{j}\rt$は，$\mathbf{A}\rt$とフェルミ粒子電流$\mathbf{j}\rt=\frac{1}{2mi}[\bar{\psi}\rt\partial_{\mathbf{r}}\psi\rt-
\partial_{\mathbf{r}}\bar{\psi}\rt\psi\rt]$との結合を表す．$\int_{\mathcal{C}}\d t$を順方向と逆方向の部分に分けてKeldysh回転を行うと，スカラーポテンシャルの場合の式~\eqref{fermion-SV}と同様に，次式を得る：
\begin{equation}\label{fermion-SA}
S_{A}=\Tr\big\{\vec{\bar{\Psi}} \hat{\mathbf{A}}\mathbf{v}_{F}
 \vec{\Psi}\big\}\,,\qquad
\hat{\mathbf{A}}=\mathbf{A}^{\alpha}\hat{\gamma}^{\alpha}=
\left(\begin{array}{cc}\mathbf{A}^{cl} & \mathbf{A}^{q} \\
\mathbf{A}^{q} & \mathbf{A}^{cl}\end{array}\right)\,.
\end{equation}
ここでは，Fermiエネルギーの近傍でフェルミ粒子の分散関係を線形化し，$-i\partial_{\mathbf{r}}\approx \mathbf{p}_{F}$および$\mathbf{v}_{F}=\mathbf{p}_{F}/m$を用いた．



ここで，生成関数を次のように定義しよう：
\begin{equation}\label{fermion-ZV}
Z\big[V^{cl},V^{q}\big]\equiv\left\langle\exp\big(iS_{V}\big)\right\rangle\,,
\end{equation}
ここで，山括弧は，フェルミ粒子の作用~\eqref{fermion-S-2}によって定まる重み$\exp(iS_{0})$を用いた，Grassmann場$\bar{\psi}$および$\psi$にわたる汎関数積分を表す．量子成分が存在しない場合，すなわち$V^{q}=0$の場合には，ソース場は時間経路の両方の枝で同一である．したがって，経路に沿った発展によって，系は厳密にもとの状態へ戻る．このため，古典成分だけでは基本的な規格化$Z=1$は変わらないと期待される．その結果，
\begin{equation}\label{fermion-Z-normalization}
Z[V^{cl}, 0]\equiv1\, ,
\end{equation}
となる．これは第~\ref{sec_bosons}節ですでに議論したとおりである（式~\eqref{boson-causality}参照）．実際，分配関数~\eqref{fermion-ZV}を$V^{cl}$のべきで展開し，Wickの定理を用いれば，この主張を明示的に確認できる．例えば，1次では$Z[V^{cl},0]=1+\int \d t\,
\Tr\big[\hat{\gamma}^{cl}\hat{G}(t,t)\big]=1$となる．ここで，$\hat{\gamma}^{cl}=\hat 1$と式~\eqref{fermion-traces}を用いた．まったく同じ理由で，$V^{cl}$のすべての高次項も消えることは容易にわかる．



式~\eqref{fermion-Z-normalization}からわかるのは，{\em 量子}ソース（経路の順方向と逆方向の枝で符号が変わるソース）を必ず導入しなければならないということである．このようなソース場の存在は因果律を明示的に破り，その結果，生成関数を変化させる．一方，これらの場は通常，物理的意味をもたず，補助的な役割だけを果たす．多くの場合，適切な微分によって観測量を生成するためだけに用いられる．実際，前述のように，物理的密度はソースの量子成分に結合する．最後に量子ソースをゼロとすれば，作用の因果性が回復する．古典成分$V^{cl}$をゼロとする必要は{\em ない}ことに注意しよう．



これがどのように機能するか見てみよう．ある物理的スカラーポテンシャル$V^{cl}(t)$が存在するもとで，時刻$t$における平均フェルミ粒子密度$\varrho$を求めたいとする．式~\eqref{fermion-SV}と\eqref{fermion-ZV}によれば，これは次式で与えられる：
\begin{equation}\label{fermion-density}
\varrho(x;V^{cl}) = -{i\over 2}\,{\delta \over \delta V^{q}(x)}\,
Z[V^{cl},V^q] \Big|_{V^{q}=0}\, ,
\end{equation}
ここで，$x=(\mathbf{r},t)$である．外場$V^{cl}$が何らかの意味で弱い場合には，問題は簡単になる．そのときは，感受率を次のように定義して，線形応答に限って考えることができる：
\begin{equation}\label{fermion-Pi-R}
\Pi^{R}(x,x\,') \equiv {\delta \over \delta V^{cl}(x\,')}\,
\varrho(x;V^{cl})\Big|_{V^{cl}=0} = -{i\over 2}\,\left. {\delta^2 \,
Z[V^{cl},V^q] \over \delta V^{cl}(x\,') \delta V^{q}(x)}\, \right|_{
V^q=V^{cl}=0}\, .
\end{equation}
応答関数は{\em 遅延}でなければならないという物理的な要請（因果律）を見越して，添字$R$を付けた．このことはすぐ後で示す．まず，{\em 分極}行列を次のように導入しよう：
\begin{equation} \label{fermion-Pi-matrix-def}
\hat \Pi^{\alpha\beta}(x,x\,') \equiv -{i\over 2}\, \left. {\delta^2 \ln
Z[\hat V] \over \delta V^{\beta}(x\,') \delta V^{\alpha}(x)}
\right|_{\hat V=0} = \left(\begin{array}{cc}
0   & \Pi^A(x,x\,')  \\
\Pi^R(x,x\,')  & \Pi^K(x,x\,')
\end{array}\right)\, .
\end{equation}
基本的な規格化条件~\eqref{fermion-Z-normalization}により，$R$成分と$A$成分では対数を省いてもよく，したがって2つの定義~\eqref{fermion-Pi-R}と\eqref{fermion-Pi-matrix-def}は矛盾しない．$\Pi^{cl,cl} =0$であることは式~\eqref{fermion-Z-normalization}から明らかである．分極行列$\hat{\Pi}$を計算するため，Gauss作用~\eqref{fermion-S-2}を考えよう．ソース項~\eqref{fermion-SV-1}を加えると，$S_0+S_V = \int \d x\,\,
\vec{\bar\Psi}[\hat{G}^{-1} + V^{\alpha}\hat{\gamma^{\alpha}}]
\vec{\Psi}$となる．\ref{app_Gaussian}のフェルミ粒子のGauss積分の規則に従ってフェルミ場$\bar\psi$，$\psi$を積分消去すると，次式を得る：
\begin{equation} \label{fermion-Z-trlog}
 Z[V^{cl},V^{q}]={1\over \Tr{\{\hat\rho_0}\}}\mathrm{Det}
 \left[i\hat{G}^{-1}+ iV^\alpha\hat{\gamma}^\alpha\right]=
\mathrm{Det}\left[\hat 1+\hat{G}\, V^{\alpha}\hat{\gamma}^{\alpha}
\right]=\exp\left\{\Tr\ln[\hat 1 + \hat{G}\,
V^{\alpha}\hat{\gamma}^\alpha]\right\} ,
\end{equation}
ここで，式~\eqref{fermion-Z-unity}を用いた．$Z[0]=1$なので，規格化は厳密に正しい．次に，$\ln[\hat 1+
\hat{G}\, V^{\alpha}\hat{\gamma}^\alpha]$を$V^{\alpha}$の2次まで展開すればよい．その結果，分極行列について次式を得る：
\begin{equation}\label{fermion-Pi-matrix}
\hat \Pi^{\alpha\beta}(x,x\,') = -{i\over 2}\,\, \Tr\left\{\hat{\gamma}^{\alpha}
 \hat{G}(x,x\,')\hat{\gamma}^{\beta}\hat{G}(x\,',x) \right\}\, ,
\end{equation}
これは明快な図式表現をもつ（図~\ref{Fig-Pi}参照）．
\begin{figure}
  \begin{center}\includegraphics[width=10cm]{figures/fig07.pdf}\end{center}
  \caption{分極演算子$\hat \Pi^{\alpha\beta}(x,x\,')$．各実線はフェルミ粒子の行列Green関数~\eqref{fermion-G-fun}を，波線は外部の古典または量子ポテンシャル$V^{cl(q)}$を表す．また，$x=(\mathbf{r},t)$である．ループ図は，式~\eqref{fermion-Pi-matrix}のトレースの図式表現である．}\label{Fig-Pi}
\end{figure}



ガンマ行列~\eqref{fermion-gammas}とGreen関数~\eqref{fermion-G-fun}の具体的な形を代入すると，{\em 応答}成分と{\em 相関}成分について次式を得る：
\begin{subequations}\label{fermion-Pi-RAK}
\begin{equation}
\hskip-3.1cm \Pi^{R(A)}(x,x\,')=-{i\over
2}\left[G^{R(A)}(x,x\,')G^{K}(x\,',x)+
G^{K}(x,x\,')G^{A(R)}(x\,',x)\right]\,,
\end{equation}
\begin{equation}
\Pi^{K}(x,x\,')= -{i\over 2}\left[G^{K}(x,x\,') G^{K}(x\,',x)+
G^{R}(x,x\,')G^{A}(x\,',x)+G^{A}(x,x\,')G^{R}(x\,',x)\right]\,.
\end{equation}
\end{subequations}
第1行より，$\Pi^{R(A)}(x,x\,')$が時間領域で実際に下三角（上三角）行列であることは明らかであり，上付き添字の選び方が正当化される．さらに，フェルミ粒子Green関数の対称性から，$\Pi^R=[\Pi^A]^\dagger$および$\Pi^K=-[\Pi^K]^\dagger$を得る．その結果，分極行列$\hat\Pi$はボソンの自己エネルギー$\hat\Sigma$の対称性をすべて備える（式~\eqref{int-bos-Sigma}参照）．



$\Pi^R$に対する式~\eqref{fermion-Pi-RAK}は，密度・密度応答関数に対するKubo公式~\cite{Mahan,Kubo}を与える．平衡では，これはMatsubara法を用いて導出できる．Matsubara法では，離散的な虚周波数$\omega_m=2\pi i m T $から実周波数$\omega$への解析接続が必要となる．具体的な応用では，この手順が煩雑になることがある．以上の議論の目的は，線形応答の問題をKeldysh形式でいかに簡潔に定式化できるかを示すことである．この形式では解析接続を回避し，実周波数領域で直接結果を得ることができる．


%$$$$$$$$$$$$$$$$$$$$$$$$$$$$$$$$$$$$$$$$$$$$$$$$$$$$$$$$$$$$$$$$$$$$$$$$$$$$$$$$$$$$$$$$$$$$$$$$$$$$$$$$$$$

\subsectionlink{応用I：量子輸送}\label{app_Part-I}
\subsubsection{Landauer公式}\label{app_Part-I-1}



量子ポイントコンタクト（QPC）のLandauerコンダクタンス~\cite{Landauer}の計算に，Keldysh法をどのように適用できるかを示そう．そのため，左（$L$）および右（$R$）と呼ぶ2つのリザーバーに接続された，準一次元の断熱的な狭窄部を考える．リザーバー中の電子の分布関数はFermi分布$n_{L(R)}(\epsilon_k)=\big[\exp[(\epsilon_k -
\mu_{L(R)})/T]+1\big]^{-1}$であり，電気化学ポテンシャルは電圧によって$\mu_L-\mu_R=eV$だけずれている．QPC内では電子の運動を横方向成分と縦方向成分に分離できる．閉じ込めのために横方向の運動は量子化されるので，横方向の伝導チャネルを区別する量子数$n$を導入する．対応する横方向波動関数を$\phi_{n}(\mathbf{r}_{\perp})$とする．縦方向の運動は，広がった散乱状態，すなわちメゾスコピックな散乱領域に左から入射する規格化された電子平面波
\begin{equation}\label{Part-I-u-L}
u^{L}_{n}(k,\mathbf{r})=\phi_{n}(\mathbf{r}_{\perp})\left\{
\begin{array}{ll}
e^{ikx}+\mathbf{r}_{n}(k)e^{-ikx} & x\to-\infty \\
\mathbf{t}_{n}(k)e^{ikx} & x\to+\infty
\end{array}
\right.\,,
\end{equation}
および右から入射する平面波
\begin{equation}\label{Part-I-u-R}
u^{R}_{n}(k,\mathbf{r})=\phi_{n}(\mathbf{r}_{\perp})\left\{
\begin{array}{ll}
e^{-ikx}+\mathbf{r}_{n}(k)e^{ikx} & x\to+\infty \\
\mathbf{t}_{n}(k)e^{-ikx} & x\to-\infty
\end{array}
\right.\,,
\end{equation}
によって記述される（図~\ref{Fig-QPC}）．
\begin{figure}
  \begin{center}\includegraphics[width=10cm]{figures/fig08.pdf}\end{center}
  \caption{量子ポイントコンタクトにおける2端子の散乱問題．}\label{Fig-QPC}
\end{figure}
ここで，$k(\epsilon)$は電子の波数ベクトルであり，$\mathbf{t}_{n}(k)$と$\mathbf{r}_{n}(k)$はチャネルごとの透過振幅と反射振幅である．第二量子化された電子場の演算子を，通常の方法で次のように導入する：
\begin{equation}
\hat{\Psi}(\mathbf{r},t)=\sum_{nk}\left[\hat{\psi}^{L}_{n}(k,t)u^{L}_{n}(k,\mathbf{r})+
\hat{\psi}^{R}_{n}(k,t)u^{R}_{n}(k,\mathbf{r})\right]\,,
\end{equation}
ここで，$\hat{\psi}^{L(R)}_{n}(k,t)$はそれぞれ左および右のリザーバーにおけるフェルミ粒子の消滅演算子である．後で用いるため，電流演算子も次のように定義する：
\begin{equation}\label{Part-I-Current}
\hat{I}(x,t)=\sum_{nk,n'k'}M^{ab}_{nn'}\hat{\psi}^{\dag
a}_{n}(k,t)\hat{\psi}^{b}_{n'}(k',t)\,,
\end{equation}
その行列要素は
\begin{equation}
M^{ab}_{nn'}(x;k,k')=\frac{e}{2im}\int
d\mathbf{r}_{\perp}\left[u^{*a}_{n}(k,\mathbf{r})\partial_{x}u^{b}_{n'}(k',\mathbf{r})-
[\partial_{x} u^{*a}_{n}(k,\mathbf{r})]u^{b}_{n'}(k',\mathbf{r})\right]\,
, \quad a=L,R\,,
\end{equation}
であり，散乱状態~\eqref{Part-I-u-L}--\eqref{Part-I-u-R}から構成される．直交条件$\int
\d\mathbf{r}_{\perp}\phi^{\phantom{*}}_{n}(\mathbf{r}_{\perp})
\phi^{*}_{n'}(\mathbf{r}_{\perp})=\delta_{nn'}$に基づいて，$x>0$に対する$\hat{M}_{nn'}(x;k,k')$を直接計算すると，次式を得る\footnote{式~\eqref{Part-I-M}は，いくつかの近似の結果として得られる．電流行列の厳密な式は，座標$x$に明示的に依存する．項には2つの種類がある．第1の種類は，$\exp(\pm i(k+k')x)\approx\exp(\pm2ik_{F}x)$の形で$x$に依存する．ここで，$k_{F}$はFermi運動量であり，これはFriedel振動を表す．これらの電流への寄与は$(k-k')/k_{F}\ll 1$と小さいため，無視する．第2の種類の項は$\exp(\pm
i(k-k')x)\approx 1$を含む．これは$|k-k'|\sim L_T^{-1}\ll x^{-1}$であり，$L_{T}=v_{F}/T$が弾道的熱長で，座標$x$が試料サイズ$L\ll L_T$の範囲に制限されるためである．対応する議論については文献~\cite{Buttiker}を参照されたい．}：
\begin{equation}\label{Part-I-M}
\hat{M}_{nn'}(k,k')=ev_{F}\delta_{nn'} \left(\begin{array}{cl}
\mathbf{t}^{*}_{n}(k)\mathbf{t}_{n}(k') &
\mathbf{t}^{*}_{n}(k)\mathbf{r}_{n}(k') \\
\mathbf{r}^{*}_{n}(k)\mathbf{t}_{n}(k') &
\mathbf{r}^{*}_{n}(k)\mathbf{r}_{n}(k')-1
\end{array}\right)\approx ev_{F}\delta_{nn'}
\left(\begin{array}{lr}|\mathbf{t}_{n}|^{2} &
\mathbf{t}^{*}_{n}\mathbf{r}_{n}\\ \mathbf{r}^{*}_{n}\mathbf{t}_{n}
& -|\mathbf{t}_n|^{2}
\end{array}\right) \,,
\end{equation}
ここで，$v_{F}=k_{F}/m$はFermi速度である．$x<0$の場合の$\hat{M}$の式も同様であり，式~\eqref{Part-I-M}とは全体の符号と複素共役だけが異なる．右辺の2つ目の近似関係は，温度または印加バイアスが定めるスケールで，透過振幅の波数$k$への依存性が弱く，運動量依存性を無視できる場合に対して書かれている．



この輸送問題の分配関数は，次のように構成できる：
\begin{equation}\label{Part-I-Z}
Z[A]=\frac{1}{\Tr\{\hat{\rho}_{0}\}}\int\D[\bar{\psi}\psi]
\exp\left\{i\,  \vec{\bar
\Psi}[\hat{\mathbf{G}}^{-1}+\hat{A}\hat{M}]\vec{\Psi}\right\}\,,
\end{equation}
ここで，$\vec{\bar \Psi}=(\bar{\psi}^{L},\bar{\psi}^{R})$であり，$\hat{\mathbf{G}}=\mathrm{diag}\{\hat{G}_{L},\hat{G}_{R}\}$は$4\times4$のGreen関数行列，$\hat{G}_{a}$はKeldysh空間の$2\times2$行列，$\hat{A}$は補助的なベクトルポテンシャルである（式~\eqref{fermion-SA}参照）．式~\eqref{Part-I-Z}のフェルミ場に関する汎関数積分は2次形式なので，Gauss積分によって次式を得る：
\begin{equation}\label{Part-I-logZ}
\ln Z[A]=\Tr\ln\big[\hat{1}+\hat{\mathbf{G}}\hat{A}\hat{M}\big]\,.
\end{equation}
式~\eqref{fermion-density}と同様に，平均電流は，ベクトルポテンシャルの量子成分について$Z[A]$を汎関数微分することで生成される：$\langle
I\rangle=-(i/2)\delta\ln Z[A]/\delta A^{q}(t)|_{A^{q}=0}$．対数のトレースを$\hat{A}$の1次まで$\Tr\ln[\hat{1}+\hat{\mathbf{G}}\hat{A}\hat{M}]\approx
\Tr[\hat{\mathbf{G}}\hat{A}\hat{M}]$のように展開すると，電流について次式を得る：
\begin{equation}
\fitdisplay{
\langle I\rangle=-\frac{iev_{F}}{2}\Tr\left\{
\left(\begin{array}{cc}\hat{G}_{L}\hat{\gamma^{q}} & 0
\\ 0 & \hat{G}_{R}\hat{\gamma}^{q}\end{array}\right)
\left(\begin{array}{lr}|\mathbf{t}_{n}|^{2}
& \mathbf{t}^{*}_{n}\mathbf{r}_{n}\\
\mathbf{r}^{*}_{n}\mathbf{t}_{n} & -|\mathbf{t}_n|^{2}
\end{array}\right)\right\}=-\frac{iev_{F}}{2}\sum_{nk}T_{n}(\epsilon_{k})
\int\frac{\d \epsilon}{2\pi}\,
[G^{K}_{L}(\epsilon,k)-G^{K}_{R}(\epsilon,k)]\,,
}
\end{equation}
ここで，Keldyshトレース$$\Tr\{\hat{G}_{a}\hat{\gamma}^{q}\}=G^{K}_a(t,t,k)=\int
\frac{\d\epsilon}{2\pi}G^{K}_a(\epsilon,k)\,,$$を用い，QPCの透過確率$T_{n}(\epsilon_k)=|\mathbf{t}_{n}(k)|^{2}$を導入した．最後に，$\epsilon_{k}=v_{F}k$としてGreen関数のKeldysh成分$G^{K}_{a}(\epsilon,k)=-2\pi i\delta(\epsilon-\epsilon_{k}+
\mu_{a})[1-2n_F(\epsilon)]$（式~\eqref{fermion-G-fun-RAK}参照）を用い，運動量積分を実行する．この積分は$G^{K}$のデルタ関数のおかげで容易に行える．結果は
\begin{equation}\label{Part-I-I}
\langle I\rangle=\frac{e}{2\pi}\sum_{n}\int\d\epsilon\,
T_{n}(\epsilon)[n_L(\epsilon)-n_{R}(\epsilon)]\,.
\end{equation}
となる．温度と印加電圧が小さい場合，式~\eqref{Part-I-I}はコンダクタンス$\langle I\rangle=\mathrm{g}V$を与える．ここで，
\begin{equation}\label{Part-I-g}
\mathrm{g}=\frac{e^{2}}{2\pi\hbar}\sum_{n}T_{n}\, ,
\end{equation}
であり，すべての透過確率はFermiエネルギーで評価する：$T_{n}=T_{n}(\epsilon_{F})$（コンダクタンスの最終的な式ではPlanck定数$\hbar$を復元したことに注意）．式~\eqref{Part-I-g}は多チャネルLandauer公式として知られている（この話題の詳細な総説については文献~\cite{Reviews-on-Landauer-1,Reviews-on-Landauer-2}参照）．


%$$$$$$$$$$$$$$$$$$$$$$$$$$$$$$$$$$$$$$$$$$$$$$$$$$$$$$$$$$$$$$$$$$$$$$$$$$$$$$$$$$$$$$$$$$$$$$$$$$$$$$$$$$$

\subsubsection{ショットノイズ}\label{app_Part-I-2}



前の例をさらに進めて，電流揺らぎの2次モーメント，いわゆるノイズパワーを計算できる．これは電流相関のFourier変換として次のように定義される：
\begin{equation}\label{Part-I-noise-def}
\mathcal{S}(\omega,V)=\int \d t\, e^{i\omega t} \langle\delta
\hat{I}(t)\delta\hat{I}(0)+\delta\hat{I}(0)\delta\hat{I}(t)\rangle,\quad
\delta\hat{I}(t)=\hat{I}(t)-\langle I\rangle\,.
\end{equation}
Keldysh法では，この相関関数を$Z[A]$から導くことができる（式~\eqref{Part-I-logZ}参照）．実際，今度は式~\eqref{Part-I-logZ}の対数のトレースを補助的ベクトルポテンシャル$\hat{A}$の2次まで展開し，$\ln
Z[A]\propto\Tr\big[\hat{\mathbf{G}}\hat{A}\hat{M}\hat{\mathbf{G}}\hat{A}\hat{M}\big]$を量子成分$A^q$について2回微分すればよい：
\begin{equation}
\mathcal{S}(\omega,V)=-\frac{1}{2}\left.\frac{\delta^{2}\ln
Z[A]}{\delta A^{q}(\omega)\delta A^{q}(-\omega)}\right|_{A^q=0}\,.
\end{equation}
この式は，適切に対称化されたノイズパワー~\eqref{Part-I-noise-def}を自動的に与える．微分の結果，次式を得る：
\begin{eqnarray}
&&\hskip-1cm \mathcal{S}(\omega,V)=\frac{1}{2}
\mathrm{Tr}\left\{\hat{\mathbf{G}}(\epsilon_{+})\hat{\gamma}^{q}
\hat{M}\hat{\mathbf{G}}(\epsilon_{-})\hat{\gamma}^{q}
\hat{M}\right\}=
\frac{e^{2}v^{2}_{F}}{2}\sum_{nkk'}\int\frac{\d\epsilon}{2\pi}\,
\left[T^{2}_{n}\Tr\{\hat{G}_{L}(\epsilon_{+})\hat{\gamma}^{q}
\hat{G}_{L}(\epsilon_{-})\hat{\gamma}^{q}\}\right.\nonumber
\\
&&\hskip-1cm\left.+
T_{n}R_{n}\Tr\{\hat{G}_{L}(\epsilon_{+})\hat{\gamma}^{q}\hat{G}_{R}(\epsilon_{-})\hat{\gamma}^{q}\}
+T_{n}R_{n}\Tr\{\hat{G}_{R}(\epsilon_{+})\hat{\gamma}^{q}\hat{G}_{L}(\epsilon_{-})\hat{\gamma}^{q}\}
+T^{2}_{n}\Tr\{\hat{G}_{R}(\epsilon_{+})\hat{\gamma}^{q}\hat{G}_{R}(\epsilon_{-})\hat{\gamma}^{q}\}
\right]\,,
\end{eqnarray}
ここでは，透過確率がエネルギーに依存しないと仮定して，左・右の部分空間における部分トレースをすでに計算し，$\epsilon_{\pm}=\epsilon\pm\omega/2$および$R_{n}=1-T_n$という表記を用いた．Keldyshトレースの計算には式~\eqref{fermion-G-fun}と\eqref{fermion-gammas}を用い，次式を得る：
\begin{equation}
\Tr\{\hat{G}_{a}\hat{\gamma}^{q}\hat{G}_{b}\hat{\gamma}^{q}\} =
G^{K}_{a}G^{K}_{b} + G^{R}_{a}G^{A}_{b} + G^{A}_{a}G^{R}_{b}\,.
\end{equation}
残るのは運動量積分である．式~\eqref{fermion-G-fun-RAK}における$G^{R(A)}_{a}(\epsilon,k)=(\epsilon-v_{F}k + \mu_{a}\pm i0)^{-1}$および$G^{K}_{a}(\epsilon,k)=-2\pi
i\delta(\epsilon-v_{F}k+\mu_{a})[1-2n_F(\epsilon)]$を用いると，$\sum_{kk'}\int\d
\epsilon
\Tr\{\hat{G}_{a}\hat{\gamma}^{q}\hat{G}_{b}\hat{\gamma}^{q}\}=v^{-2}_{F}\int\d
\epsilon\, [1-(1-2n_{a})(1-2n_{b})]$となる．その結果，Lesovik~\cite{Lesovik}によって得られたノイズパワーの最終的な式は次のように書かれる：
\begin{equation}\label{Part-I-noise}
\mathcal{S}(\omega,V)=\frac{e^{2}}{2\pi\hbar}\sum_{n}\int\d
\epsilon\left[T^{2}_{n}B_{LL}(\epsilon)+T_{n}R_{n}B_{LR}(\epsilon)+
T_{n}R_{n}B_{RL}(\epsilon)+T^{2}_{n}B_{RR}(\epsilon)\right]\,,
\end{equation}
ここで，統計因子は$B_{ab}(\epsilon)=n_{a}(\epsilon_{+})[1-n_{b}(\epsilon_{-})]+
n_{b}(\epsilon_{-})[1-n_{a}(\epsilon_{+})]$であり，最後に再び$\hbar$を復元した．見かけは複雑だが，式~\eqref{Part-I-noise}の$\epsilon$積分は閉じた形で実行できる\footnote{式~\eqref{Part-I-noise-general}を導くには，統計因子を$B_{ab}(\epsilon)=\frac{1}{2}\big[1-\tanh[(\epsilon_{+} - \mu_a)/2T]
\tanh[(\epsilon_{-} - \mu_b)/2T]\big]$と書き，積分$\int^{+\infty}_{-\infty}\d x\,
[1-\tanh(x+y)\tanh(x-y)]=4y\coth(2y)$を用いる．}：
\begin{equation}\label{Part-I-noise-general}
\mathcal{S}(\omega,V)=\frac{e^2}{2\pi\hbar}\sum_{n}
\left[T^{2}_{n}\omega\coth\left(\frac{\omega}{2T}\right)+
T_{n}(1-T_{n})(eV+\omega)
\coth\left(\frac{eV+\omega}{2T}\right)+\{\omega\to-\omega\}\right]\,.
\end{equation}


式~\eqref{Part-I-noise-general}からは，興味深い2つの極限の場合が容易に得られる．第1は，熱平衡の電流揺らぎに対応する$V\to0$の極限である．この場合，
\begin{equation}\label{Part-I-noise-FDT}
\mathcal{S}(\omega,0)=2\mathrm{g}\omega\coth\left(\frac{\omega}{2T}\right)\,,
\end{equation}
となる．ここで，コンダクタンス$\mathrm{g}$に対する式~\eqref{Part-I-g}を用いた．この結果は，よく知られた電流揺らぎの揺動散逸関係にほかならない．式~\eqref{Part-I-noise}では透過振幅への依存性が複雑であるにもかかわらず，平衡ノイズパワー~\eqref{Part-I-noise-FDT}がコンダクタンス~\eqref{Part-I-g}だけで書けることに注意しよう．もう一方の極限は，ゼロ温度$T\to0$で有限のバイアス$V$を印加した，完全に非平衡なノイズである．この場合，式~\eqref{Part-I-noise-general}からノイズの超過分について次式を得る：
\begin{equation}\label{Part-I-noise-shot}
\mathcal{S}(\omega,V) - \mathcal{S}(\omega,0) =
\frac{e^2}{2\pi\hbar} \Big(|eV+\omega|+
|eV-\omega| -2|\omega| \Big)  \sum_{n}T_{n}(1-T_n)\,,
\end{equation}
これは{\em ショットノイズ}と呼ばれる．ここで重要なのは，平衡ノイズ~\eqref{Part-I-noise-FDT}とは異なり，ショットノイズはコンダクタンス$\mathrm{g}$だけでは書けないという点である．すべての透過確率が小さい$T_n\ll 1$というトンネル接合の場合に限り，式~\eqref{Part-I-noise-shot}は，Schottky公式として知られる$\mathcal{S}(0,V)=2eV\mathrm{g}=2e\langle I\rangle$に帰着する（さまざまな系のショットノイズの総説については，例えば文献~\cite{Kogan,Blanter-Buttiker,Martin}参照）．


%$$$$$$$$$$$$$$$$$$$$$$$$$$$$$$$$$$$$$$$$$$$$$$$$$$$$$$$$$$$$$$$$$$$$$$$$$$$$$$$$$$$$$$$$$$$$$$$$$$$$$$$$$$$

\subsubsection{Coulombドラッグ}\label{app_Part-I-3}



Pogrebinskii~\cite{Pogrebinskii}とPrice~\cite{Price}によって提案されたドラッグ効果は，現在では電子間散乱を調べ，測定する標準的な方法の1つとなっている．バルクの二次元系（絶縁体によって隔てられた，平行な2つの二次元電子気体）では，ドラッグ効果は実験的に十分確立され~\cite{Solomon,Gramila,Sivan,Lilly,Pillarisetty}，理論的にも研究されている~\cite{Smith,MacDonald,Kamenev-Oreg,Flensberg}．近年，量子細線~\cite{Debray-1,Debray-2,Morimoto,Yamamoto}，量子ドット~\cite{Aguado,Kouwenhoven}，QPC~\cite{Khrapai}など，量子閉じ込めのある形状におけるCoulombドラッグを調べる実験が数多く行われた．これらの系では，ソース・ドレイン電圧$V$を印加して{\em 駆動回路}に電流を流し，{\em ドラッグ回路}に誘起される電流（または電圧）を測定する．このようなドラッグ電流は，駆動電圧$V$に加え，一方または両方の回路の透過率を制御するゲート電圧$V_{g}$の関数でもある．図~\ref{Fig-drag}aは，駆動回路とドラッグ回路をそれぞれQPCで構成した，このような配置の一例を示す．



\begin{figure}
\begin{center}\includegraphics[width=8cm]{figures/fig09.pdf}\end{center}
\caption{a）結合した2つのQPCと周囲の電気回路．Coulomb結合は相互容量$C_c$によって生じる．ゲート電圧$V_{g}$は，例えば駆動QPCの透過率を制御する．b）ゲート電圧の関数としての，駆動QPCのコンダクタンスとドラッグ電流の模式図．}\label{Fig-drag}
\end{figure}



Keldysh法は，線形応答領域と，駆動回路に比較的大きなバイアスを印加した平衡から離れた領域の両方で，ドラッグ問題を扱う効率的な方法である．各QPC内では電子は相互作用しないものとし，その運動を，量子化された横方向の運動と，広がった縦方向の運動に分離する（第~\ref{app_Part-I-1}節参照）．相互作用のないポイントコンタクトを記述する作用は
\begin{equation}\label{Part-I-Action-QPC}
iS_{\mathrm{QPC}}=i\, \Tr\big\{\vec{\bar \Psi}
\hat{\mathbf{G}}^{-1}\vec{\Psi}\big\}\,,
\end{equation}
である．ここで，$\vec{\bar \Psi}=(\bar{\psi}^{L}_{jn},\bar{\psi}^{R}_{jn})$および$\hat{\mathbf{G}}=\delta_{jj'}
\mathrm{diag}\{\hat{G}_{L},\hat{G}_{R}\}$である．添字$j=1,2$はそれぞれQPC$_{1(2)}$を表し，$n$は各QPC内の横方向チャネルの添字，$\hat{G}_{L(R)}$は式~\eqref{fermion-G-fun}で与えられる$2\times2$のKeldysh行列である．



2つのQPC間の相互作用項は
\begin{equation}\label{Part-I-Action-int}
iS_{\mathrm{int}}=\sum_{ab\alpha\beta}\iint^{+\infty}_{-\infty}dtdt'\,
I^{\alpha}_{1a}(t)\mathbf{K}^{\alpha\beta}_{ab}(t-t')I^{\beta}_{2b}(t')\,,
\end{equation}
である．ここで，$I_{jR(L)}(t)$はQPC$_j$の右側（左側）の電流演算子であり，回路の電磁環境を表す核$\hat{\mathbf{K}}_{ab}(t-t')$によって結合している．相互作用核の遅延成分と先進成分は，トランスインピーダンス行列と$\mathbf{K}^{R(A)}_{ab}(\omega)=\mathbf{Z}^{R(A)}_{ab}(\omega)/(\omega\pm
i0)$の関係にある．後者は$\mathbf{Z}^{R(A)}_{ab}(\omega)=\partial\Phi_{a}(\pm
\omega)/\partial I_{b}(\mp \omega)$によって定義される．対応する局所的な揺らぐ電流$I_a$と電圧$\Phi_a$を図~\ref{Fig-drag}aに示した．相互作用核のKeldysh成分は，揺動散逸定理$\mathbf{K}^{K}_{ab}(\omega)=[\mathbf{K}^{R}_{ab}(\omega)-
\mathbf{K}^{A}_{ab}(\omega)]\coth(\omega/2T)$によって定まる．すなわち，周囲の電気的環境は平衡に近いと仮定する．最後に，電流演算子は式~\eqref{Part-I-Current}および\eqref{Part-I-M}で与えられる．



ドラッグ電流は，フェルミ粒子の自由度について$I_{2}$を平均することで求まる：
\begin{equation}\label{Part-I-ID-def}
I_{D}=
\int\mathbf{D}[\psi\bar{\psi}]\,\mathrm{Tr}\left[\bar{\psi}_{2}M\psi_{2}\right]\,
\exp\big(iS_{\mathrm{QPC}}[\bar{\psi}\psi]+iS_{\mathrm{int}}[\bar{\psi}\psi]\big)\,.
\end{equation}
指数関数を相互作用項$S_{\mathrm{int}}$の2次まで展開すると，次式を得る：
\begin{equation}\label{Part-I-ID-def-perturb}
I_{D}=\frac{1}{2}\int\mathbf{D}[\psi\bar{\psi}]\,
\mathrm{Tr}\left[\bar{\psi}_{2}M\psi_{2}\right]
\mathrm{Tr}\left[I_{1}\mathbf{K}I_{2}\right]
\mathrm{Tr}\left[I_{1}\mathbf{K}I_{2}\right]
\exp\big(iS_{\mathrm{QPC}}[\bar{\psi}\psi]\big)\,.
\end{equation}
残ったフェルミ場のGauss積分は，Wickの定理を用いて計算する．式~\eqref{Part-I-M}で与えられる$M$行列を用いて，電流演算子の式~\eqref{Part-I-Current}を代入し，$\psi$場の間の可能なWick収縮をすべて考慮する．これらの収縮はGreen関数~\eqref{fermion-G-fun}で与えられる．このようにして，ドラッグ電流について次式を得る：
\begin{equation}\label{Part-I-ID-result}
I_D(V)=
\int\frac{d\omega}{4\pi\omega^{2}}\mathrm{Tr}\left[\hat{\mathbf{Z}}(\omega)
\hat{\mathcal{S}}_{1}(\omega,V)\hat{\mathbf{Z}}(-\omega)
\hat{\Gamma}_{2}(\omega)\right]\,.
\end{equation}
駆動回路は，電流・電流相関行列$\mathcal{S}_{ab}(\omega,V)=\int dt\, e^{i\omega
t}\big\langle\big\langle\delta\hat{I}_{a}(t)\delta\hat{I}_{b}(0)+
\delta\hat{I}_{b}(0)\delta\hat{I}_{a}(t)\big\rangle\big\rangle$の{\em 超過}部分$\mathcal{S}^{ab}_{1}(\omega,V)=\mathcal{S}_{ab}(\omega,V)-\mathcal{S}_{ab}(\omega,0)$によって特徴づけられる．相関行列は，例えば次式で与えられる：
\begin{eqnarray}\label{Part-I-noise-Srr}
\mathcal{S}_{RR}(\omega,V)=\frac{2}{R_{Q}}\sum_{n}\int\d\epsilon\,
\big[B_{LL}(\epsilon)|\mathbf{t}^{L}_{n}(\epsilon_{+})|^{2}
|\mathbf{t}^{L}_{n}(\epsilon_{-})|^{2}+B_{LR}(\epsilon)
|\mathbf{t}^{L}_{n}(\epsilon_{+})|^{2}|\mathbf{r}^{R}_{n}(\epsilon_{-})|^{2}\nonumber\\
+B_{RL}(\epsilon)|\mathbf{r}^{R}_{n}(\epsilon_{+})|^{2}
|\mathbf{t}^{L}_{n}(\epsilon_{-})|^{2}+B_{RR}(\epsilon)
[1-\mathbf{r}^{*R}_{n}(\epsilon_{+})\mathbf{r}^{R}_{n}(\epsilon_{-})]
[1-\mathbf{r}^{R}_{n}(\epsilon_{+})\mathbf{r}^{*R}_{n}(\epsilon_{-})]\big]\,,
\end{eqnarray}
ここで，$\epsilon_\pm = \epsilon \pm \omega/2\,\,$，$\mathbf{t}^{L(R)}_{n}(\epsilon_\pm)=\mathbf{t}^{L(R)}_{n}(\epsilon_\pm+eV_{L(R)})$および$\mathbf{r}^{L(R)}_{n}(\epsilon_\pm)=\mathbf{r}^{L(R)}_{n}(\epsilon_\pm+eV_{L(R)})$であり，$R_{Q}=2\pi\hbar/e^2$は量子抵抗である．また，統計的占有の形状因子$B_{ab}(\epsilon)$は式~\eqref{Part-I-noise}で与えられる．ここでの$\mathcal{S}_{RR}(\omega,V)$は，式~\eqref{Part-I-noise}をエネルギーに依存する透過確率の場合に一般化したものである~\cite{Buttiker}．ノイズ行列の他の成分$\mathcal{S}_{LL}$，$\mathcal{S}_{LR}$，$\mathcal{S}_{RL}$の式も同様である（文献~\cite{Buttiker,LevchenkoKamenev-QPC}参照）．



式~\eqref{Part-I-ID-result}のドラッグ回路は，（平衡に近い）ドラッグQPC$_2$に加わる交流電圧揺らぎの整流係数$\hat{\Gamma}_{2}(\omega)=\Gamma_{2}(\omega)\hat{\varsigma}_{z}$によって特徴づけられる．ここで，$\hat{\varsigma}_{z}$は左・右の部分空間に作用する第3のPauli行列である．整流係数は次式で与えられる\footnote{Keldysh行列を用いると，整流係数は次のトレースで与えられる：$\Gamma_{2}(\omega)=\Tr\big[\hat{\mathbf{G}}\hat{\gamma}^{q}\hat{M}
\hat{\mathbf{G}}\hat{\gamma}^{cl}\hat{M}\hat{\mathbf{G}}\hat{\gamma}^{cl}\hat{M}\big]$．$\Gamma_{2}(\omega)$を式~\eqref{Part-I-Gamma}の形で求めるには，Keldyshトレース$\mathrm{Tr}\left[\hat{G}\hat{\gamma}^{q}\hat{G}\hat{\gamma}^{cl}
\hat{G}\hat{\gamma}^{cl}\right]
=\sum_{\pm}\left[G^{R}(\epsilon)G^{R}(\epsilon\pm\omega)G^{K}(\epsilon)+
G^{R}(\epsilon)G^{K}(\epsilon\pm\omega)G^{A}(\epsilon)+
G^{K}(\epsilon)G^{A}(\epsilon\pm\omega)G^{A}(\epsilon)\right]$を用いる．この式をさらに簡単にするには，揺動散逸関係$G^{K}(\epsilon)=\big[G^{R}(\epsilon)-G^{A}(\epsilon)\big][1-2n(\epsilon)]$を用いてGreen関数の各Keldysh成分を分解し，得られた式のうち，適切な因果性をもつ項のみを残せばよい．すなわち，$G^{A}G^{A}G^{A}$や$G^{R}G^{R}G^{R}$のように同じ種類のGreen関数を3つ含む組合せは寄与しない．このようにして，Keldyshトレースについて$\mathrm{Tr}\left[\hat{G}\hat{\gamma}^{q}\hat{G}
\hat{\gamma}^{cl}\hat{G}\hat{\gamma}^{cl}\right]
\propto\big[n_F(\epsilon_{-})-n_F(\epsilon_{+})\big]$を得る．残った左・右部分空間での電流頂点行列$\hat{M}$のトレースは，シフトしたエネルギーでの透過確率に帰着する．具体的には$\mathrm{Tr}\big[\hat{M}\hat{M}\hat{M}\big]\propto
|\mathbf{t}_{n}(\epsilon_{+})|^{2}-
|\mathbf{t}_{n}(\epsilon_{-})|^{2}$となり，式~\eqref{Part-I-Gamma}を得る．}：
\begin{equation}\label{Part-I-Gamma}
\Gamma_{2}(\omega)=\frac{2e}{R_{Q}}\sum_{n}\int d\epsilon\,
\big[n_F(\epsilon_{-})-n_F(\epsilon_{+})\big]
\Big[|\mathbf{t}_{n}(\epsilon_{+})|^{2}-|\mathbf{t}_{n}(\epsilon_{-})|^{2}\Big]\,.
\end{equation}
QPC$_2$の特性は，エネルギー依存する透過確率$|\mathbf{t}_{n}(\epsilon)|^{2}$を通じて現れる．この式は明快に解釈できる．例えばQPCの左側にある周波数$\omega$のポテンシャル揺らぎは，右向きに進む粒子の枝に，エネルギー$\epsilon_\pm$の電子・正孔対を生成する．その結果，電子は確率$|\mathbf{t}_{n}(\epsilon_{+})|^{2}$でQPCを通過し，正孔は確率$|\mathbf{t}_{n}(\epsilon_{-})|^{2}$で通過する．両者の差が，QPCを横切る直流電流を与える．電子と正孔の非対称性，したがって非ゼロの整流係数$\Gamma_{2}(\omega)$を得るには，ドラッグQPCの透過確率のエネルギー依存性が不可欠であることに注意しよう．式~\eqref{Part-I-ID-result}は，図~\ref{Fig-drag-diagram}に示す明快な図式表現をもつ．



\begin{figure}
 \begin{center}\includegraphics[width=8cm]{figures/fig10.pdf}\end{center}
  \caption{回路間相互作用$\mathbf{K}=\mathbf{Z}/\omega$（波線）の2次におけるドラッグ電流$I_D$．ドラッグ回路は三角形の整流頂点$\Gamma_2(\omega)$で表され，駆動回路は非平衡の電流・電流相関関数$\mathcal{S}_1(\omega, V)$（ループ）で表される．}\label{Fig-drag-diagram}
\end{figure}



滑らかなQPCにおける，部分的に開いた単一チャネルに注目すると，それを横切るポテンシャル障壁は，ほぼ放物線形とみなせる．この場合，透過確率は次式で与えられる：
\begin{equation}\label{Part-I-t}
|\mathbf{t}(\epsilon)|^2=\left(\exp\{(eV_{g}-\epsilon)/\Delta_j\}+1\right)^{-1}\,,
\end{equation}
ここで，$\Delta_j$はQPC$_j$の放物線形障壁の曲率に対応するエネルギースケールであり，ゲート電圧$V_g$はFermiエネルギーに対する障壁頂上の位置を移動させる．この透過確率の形は，QPCのコンダクタンス量子化の説明に用いられ~\cite{Glazman-QPC}，Coulombドラッグ問題への応用にも有用であることがわかる．式~\eqref{Part-I-t}を式~\eqref{Part-I-Gamma}に代入し，エネルギー積分を実行すると，
\begin{equation}\label{Part-I-Gamma-1}
\Gamma_2(\omega)=\frac{2e\Delta_2}{R_{Q}}\,
\ln\left(1+\frac{\sinh^2(\omega/2\Delta_2)}
{\cosh^2(eV_{g}/2\Delta_2)}\right)
\end{equation}
を$T\ll\Delta_2$の場合に得る．もう一方の極限$T\gg\Delta_2$では，式~\eqref{Part-I-Gamma-1}において$\Delta_2\to T$と置き換えればよい．低周波数$\omega\ll \Delta_2$では$\Gamma_2\sim
\omega^2$となるので，式~(\ref{Part-I-ID-result})の積分は$\omega\to 0$の領域で収束することに注意しよう．



\textbf{\textit{線形ドラッグ領域．}}印加電圧$V$が小さいとき，応答電流$I_D$は$V$に比例すると期待される．$\hat{\mathcal{S}}_{1}(\omega,V)$を$V$の1次まで展開すると，電流・電流相関行列の対角成分だけが線形応答に寄与し，その結果，
\begin{equation}\label{Part-I-noise-S-linear}
\hat{\mathcal{S}}_{1}(\omega,V)=V\, \frac{\partial}{\partial\omega}
\left[\coth\frac{\omega}{2T}\right]\Gamma_{1}(\omega)\hat{\varsigma}_{z} +O(V^3)\,,
\end{equation}
となる．ここで，$\Gamma_{1}(\omega)$は，式~\eqref{Part-I-Gamma}のQPC$_2$の透過確率をQPC$_1$の透過確率に置き換えることで得られる．式~\eqref{Part-I-noise-S-linear}を式~\eqref{Part-I-ID-result}に代入すると，次式を得る：
\begin{equation}\label{Part-I-ID-linear}
I_{D}=V\, \frac{R^{2}_{Q}}{4\pi}\int
d\omega\,\frac{\alpha_{+}(\omega)}{\omega^{2}}\,
\frac{\partial}{\partial\omega}\left[\coth\frac{\omega}{2T}\right]
\Gamma_{1}(\omega)\, \Gamma_{2}(\omega)\,,
\end{equation}
ここで，無次元の相互作用核$\alpha_{+}(\omega)$は，トランスインピーダンス行列を用いて$$\alpha_{+}(\omega)=\frac{1}{2R^{2}_{Q}}\Tr\big[
\hat{\mathbf{Z}}(\omega)\hat{\varsigma}_{z}
\hat{\mathbf{Z}}(\omega)\hat{\varsigma}_{z}\big]\,.$$と表される．式~\eqref{Part-I-ID-linear}は，バルク二次元系のドラッグ電流の式~\cite{Kamenev-Oreg,Flensberg}と同じ一般的構造をもつ．$1\leftrightarrow 2$の交換に対して対称なので，線形応答係数に対するOnsager関係を満たす．式~\eqref{Part-I-ID-linear}に残る周波数積分を行う際には，相互作用核をゼロ周波数で評価すれば十分である．実際，$\alpha_{+}(\omega)$が変化する周波数スケールは，回路の$RC$時定数の逆数で決まる．ほとんどの実験のように，ドラッグ回路の負荷インピーダンスが駆動回路のものに比べて大きく$Z_{1}\ll Z_2\ll R_Q$，かつ両回路の相互容量が小さい$C_c\ll C_{R,L,s}$場合（図~\ref{Fig-drag}a参照），$\tau^{-1}_{RC}=(Z_1 C_s)^{-1}\gg
T$となる．式~\eqref{Part-I-ID-linear}の$I_{D}$は$\omega\lesssim T$によって決まるため，$\alpha_{+}(\omega)\approx\alpha_{+}(0)$という近似が正当化される\footnote{図~\ref{Fig-drag}に示す回路では，トランスインピーダンス核の低周波数極限は$$\alpha_{\pm}(0)=\frac{Z^{2}_{1}}{8R^{2}_{Q}}\frac{C^{2}_{c}}{C^{2}_{L}C^{2}_{R}}
\left\{\begin{array}{l}2C^{2}_{L}+2C_{L}C_{R}+2C^{2}_{R}\\
C^{2}_{L}-C^{2}_{R}
\end{array}\right.$$となる．}．式~\eqref{Part-I-Gamma-1}を式~\eqref{Part-I-ID-linear}に代入すると，例えば低温領域$T\ll \Delta_{1,2}$について次式を得る：
\begin{equation}\label{Part-I-ID-linear-1}
I_D= \frac{V}{R_{Q}}\,   \frac{\alpha_{+}(0)\pi^2}{6}\, {T^{2}\over
\Delta_1\Delta_2}\,  {1\over \cosh^{2}(eV_{g}/2\Delta_1)}\,,
\end{equation}
ここでは，QPC$_2$のゲート電圧を，その障壁の頂上がFermiエネルギーに一致するように調整したと仮定し，$I_D$をQPC$_1$のゲート電圧の関数として書いた．得られた式は，図~\ref{Fig-drag}bと同様に$V_g=0$でピークを示す．この式は，平衡に近い熱揺らぎの整流を記述しており，それゆえ係数$T^2$が現れる．この整流は電子・正孔非対称性によるものであり，そのため$V_g$に対して非単調な依存性を示す．



\textbf{\textit{非線形領域．}}駆動電圧が大きくなると，ドラッグ電流は$V$に比例しなくなる．さらに線形応答の場合とは異なり，$\hat{\mathcal{S}}_{1}(\omega,V)$では透過確率のエネルギー依存性は必要なく，エネルギーに依存しない$|\mathbf{t}_n|^2$に対して評価できる（これは$T,eV\ll \Delta_1$の場合に妥当な仮定である）．さらに$T\ll eV$を仮定すると，$\hat{\mathcal{S}}^{ab}_{1}(\omega,V)=
\big[\mathcal{S}_{ab}(\omega,V)-\mathcal{S}_{ab}(\omega,0)\big]\hat{\varsigma}_{0}$を得る．ここで，$\mathcal{S}_1(\omega,V)$は式~\eqref{Part-I-noise-shot}で与えられる（$T_{n}\equiv|\mathbf{t}_{n}|^{2}$であることを思い出そう）．これを式~\eqref{Part-I-ID-result}に代入し，電圧で制限された周波数積分を行うと，ドラッグ電流について次式を得る：
\begin{equation}\label{Part-I-ID-nonlinear}
I_{D}=\frac{eV^2}{\Delta_{2}R_{Q}}\,  \alpha_{-}(0) \,
\sum_{n}T_{n}(1-T_{n})\,.
\end{equation}
ここでも，検出器となるQPC$_2$をプラトー間の遷移に調整したと仮定した．また，$$\alpha_{-}(\omega)=\frac{1}{2R^{2}_{Q}}\Tr\big[
\hat{\mathbf{Z}}(\omega)\hat{\varsigma}_{0}
\hat{\mathbf{Z}}(\omega)\hat{\varsigma}_{z}\big]$$をその直流値$\alpha_{-}(0)$に置き換えるため，$eV\ll
(Z_1C_s)^{-1}$と仮定した．$\alpha_{+}>0$である一方，$\alpha_{-}\propto
C^{2}_{L}-C^{2}_{R}$なので$\alpha_{-}$の符号は任意であることに注意しよう（図~\ref{Fig-drag}aおよび注12参照）．完全に対称な回路では$\alpha_{-}=0$となり，極端に非対称な回路では$|\alpha_{-}|\approx \alpha_{+}/2$となる．式~\eqref{Part-I-ID-nonlinear}の導出は$T\ll eV$について示したが，$T\ll \mathrm{min}\{\Delta_1, (Z_1C_s)^{-1}\}$である限り，任意の温度でも有効であることを示せる．



式~\eqref{Part-I-ID-nonlinear}は，ドラッグ電流が量子ショットノイズの整流の結果であり，したがって駆動回路のFano因子~\cite{Lesovik}に比例することを示す．これは図~\ref{Fig-drag}bに示した一般的な振る舞いを示すが，その理由は，線形領域での同様の振る舞いとはかなり異なる．非線形ドラッグ電流の向きは，駆動電流の向きではなく，回路の反転非対称性によって（$\alpha_{-}$の符号を通じて）決定される．その結果，ある駆動電圧の極性に対しては，ドラッグ電流が{\em 負}となる．最後に，一般的な回路では$\alpha_+\sim \alpha_-$と仮定し，式~\eqref{Part-I-ID-linear-1}と\eqref{Part-I-ID-nonlinear}を比較すると，$T\ll \Delta_1$の場合には，線形領域から非線形領域への移行が$V\approx V^*$で起こり，$eV^*=T^2/\Delta_1\ll T$であることがわかる．反対の極限$T>\Delta_1$では，クロスオーバー電圧は温度によって$eV^*=T$と与えられる．さらなる詳細と議論は文献~\cite{LevchenkoKamenev-QPC}にある．



%$$$$$$$$$$$$$$$$$$$$$$$$$$$$$$$$$$$$$$$$$$$$$$$$$$$$$$$$$$$$$$$$$$$$$$$$$$$$$$$$$$$$$$$$$$$$$$$$$$$$$$$$$$$
%$$$$$$$$$$$$$$$$$$$$$$$$$$$$$$$$$$$$$$$$$$$$$$$$$$$$$$$$$$$$$$$$$$$$$$$$$$$$$$$$$$$$$$$$$$$$$$$$$$$$$$$$$$$

\sectionlink{乱れのあるフェルミ粒子系}\label{sec_NLSM}



静的な（凍結した），空間に依存する乱れのポテンシャル$U_{\mathrm{dis}}(\mathbf{r})$が存在するもとで，例えば密度・密度応答関数や電流・電流応答関数を計算したいことが多い．さらに，乱れのポテンシャルの正確な形は一般には未知なので，$U_{\mathrm{dis}}(\mathbf{r})$の実現のアンサンブルにわたる平均を求めたい．Keldysh形式では，応答関数は，生成関数の{\em 対数ではなく}生成関数そのものの変分として定義できる．より正確には，基本的な規格化$Z=1$のおかげで，対数を用いる定義と用いない定義が一致する．平衡の形式ではそうはならず，正しい規格化のためには対数が不可欠であり，微分後に因子$Z^{-1}$が生じる．このような乱れに依存する因子$Z^{-1}=Z^{-1}[U_{\mathrm{dis}}]$は，$U_{\mathrm{dis}}$についての平均を著しく複雑にする．この平均を実行するために，レプリカ法~\cite{Edwards,Wegner,EfetovLarkinKhmelnitskii,Finkel'stein,BelitzKirkpatrick}と超対称性法~\cite{Efetov,Efetov-book}という2つの手法が考案された．前者は$\ln Z=\lim_{n\to 0}(Z^n-1)/n$という関係を利用し，同じ系の整数個$n$のレプリカについて計算し，最後に$n\to 0$の極限をとる．後者は，相互作用のないフェルミ粒子系の$Z^{-1}$が，同じランダムポテンシャル中のボソン系の$Z$に等しいことに基づく．そこで，対象とするフェルミ粒子系に，ボソンのレプリカを追加する．Keldysh形式は，構成上$Z=1$を保証することによって，これら2つの手法に代わる方法を提供する~\cite{HorbachSchon,KamenevAndreev,ChamonLudwigNayak}．本節の目的は，非線形$\sigma$模型（NLSM）として知られる，乱れた電子気体の有効場理論を，Keldysh形式の中でどのように構築するかを示すことである．


%$$$$$$$$$$$$$$$$$$$$$$$$$$$$$$$$$$$$$$$$$$$$$$$$$$$$$$$$$$$$$$$$$$$$$$$$$$$$$$$$$$$$$$$$$$$$$$$$$$$$$$$$$$$

\subsectionlink{乱れの平均}\label{sec_NLSM-1}



フェルミ粒子の作用に，乱れに依存する項$S_{\mathrm{dis}}[\bar{\psi},\psi]=\int_{\mathcal{C}}\d
t\int\d\mathbf{r}U_{\mathrm{dis}}(\mathbf{r}) \bar{\psi}\rt\psi\rt$を加える．ここで，$U_{\mathrm{dis}}(\mathbf{r})$は不純物のランダムな配置によって生じる静的スカラーポテンシャルである．通常，不純物は短距離のポテンシャルをもち，系内に一様に分布すると仮定するのが妥当であり，相関関数は$\langle
U_{\mathrm{dis}}(\mathbf{r})U_{\mathrm{dis}}(\mathbf{r}')\rangle\sim
\delta(\mathbf{r}-\mathbf{r}')$の形をとる．さらに，不純物ポテンシャルがGauss分布に従うと仮定すると，乱れの平均は次の汎関数積分によって実行される：
\begin{equation}\label{NLSM-dis-average-def}
\langle\ldots\rangle_{\mathrm{dis}}=\int\D[U_{\mathrm{dis}}] \ldots
\exp\left\{-\pi\nu\tauel\int\d\mathbf{r}\,U^{2}_{\mathrm{dis}}(\mathbf{r})\right\},
\end{equation}
ここで，乱れの強さは弾性平均自由時間$\tauel$によって特徴づけられ，$\nu$はFermiエネルギーにおける電子の状態密度である．乱れのポテンシャルは古典成分しかもたないため，Keldysh経路の両方の枝で厳密に同じである．したがって，頂点行列$\hat{\gamma}^{cl}=\hat
1$にのみ結合する．次に，分配関数の乱れに依存する項について$U_{\mathrm{dis}}$にわたるGauss積分を実行し（この段階では規格化因子が存在しないことが本質的である），次式を得る：
\begin{eqnarray}\label{NLSM-dis-average}
\int\D[U_{\mathrm{dis}}]\exp\left(-\int\d\mathbf{r}\left[\pi\nu\tauel
U^{2}_{\mathrm{dis}}(\mathbf{r})-iU_{\mathrm{dis}}(\mathbf{r})\int^{+\infty}_{-\infty}
\d t\, \bar{\psi}^{a}\rt\hat
\gamma^{cl}_{ab}\psi^{b}\rt\right]\right)\nonumber
\\ =\exp\left(-\frac{1}{4\pi\nu\tauel}\int\d\mathbf{r}
\iint^{+\infty}_{-\infty}\d t\d t'
\big[\bar{\psi}^{a}\rt\psi^{a}\rt\big]
\big[\bar{\psi}^{b}\rtp\psi^{b}\rtp\big]\right)\,,
\end{eqnarray}
ここで，$a,b=1,2$であり，繰り返し現れるすべての添字について和をとるものとする．直前の式の右辺の指数中で$[\bar{\psi}^a\rt\psi^{a}\rt][
\bar{\psi}^{b}\rtp\psi^{b}\rtp]=-[\bar{\psi}^{a}\rt
\psi^{b}\rtp][\bar{\psi}^{b}\rtp\psi^{a}\rt]$と並べ替え（負号はGrassmann数の反交換性に由来する），行列値のHubbard--Stratonovich場$\hat{Q}=Q^{ab}_{tt'}(\mathbf{r})$を用いると，（時間について非局所的な）4フェルミオン項を次のように分離できる\footnote{ここでは，時間反転対称性，すなわちHamiltonianが実演算子であることを追跡しないため，以下の考察は，例えば外部磁場によって時間反転不変性が破れている場合（複素HermiteなHamiltonian）に限られる．これはいわゆる{\em ユニタリー}NLSMである．時間反転対称性を回復した{\em 直交}NLSMは，乱れた超伝導体を扱う第~\ref{sec_SuperCond}節で考察する．}：
\begin{eqnarray}\label{NLSM-HS-transform}
\exp\left(\frac{1}{4\pi\nu\tauel}\int\d\mathbf{r}
\iint^{+\infty}_{-\infty}\d t\d t'
[\bar{\psi}^{a}\rt\psi^{b}\rtp][\bar{\psi}^{b}\rtp\psi^{a}\rt]\right)\nonumber
\\
=\int\D[\hat{Q}]\,\exp\left(-\frac{\pi\nu}{4\tauel}\Tr\{
\hat{Q}^{2}\}-\frac{1}{2\tauel}\int\d\mathbf{r}\iint^{+\infty}_{-\infty}\d
t \d t'
Q^{ab}_{tt'}(\mathbf{r})\bar{\psi}^{b}\rtp\psi^{a}\rt\right)\,.
\end{eqnarray}
ここでは，$\hat{Q}^{2}$のトレースに行列添字の和と時間・空間積分を含めるという表記を導入した：
\begin{equation}
\Tr\big\{\hat{Q}^{2}\big\}=\int\d\mathbf{r}\iint^{+\infty}_{-\infty} \d t\d
t'\sum^{2}_{a,b=1}Q^{ab}_{tt'}(\mathbf{r})Q^{ba}_{t't}(\mathbf{r}).
\end{equation}
これで，{\em 平均された}作用はGrassmann変数の2次形式$S[\Psi,\hat{Q}]=\Tr\big\{\vec{\bar{\Psi}}\big[\hat{G}^{-1}+
\frac{i}{2\tau_{\mathrm{el}}}\hat{Q}]\vec{\Psi}\big\}$となり，これらを明示的に積分消去できる．その結果，対応する2次形式$\hat{G}^{-1}+\frac{i}{2\tau_{\mathrm{el}}}\hat{Q}$の行列式を得る．ここでの行列はすべて，$2\times 2$のKeldysh構造と，離散時間に関する$N\times N$の構造を併せもつと理解すべきである．したがって，乱れについて平均した生成関数$\mathcal{Z}=\langle Z\rangle_{\mathrm{dis}}$に対して次式を得る：
\begin{eqnarray}\label{NLSM-action}
&&\mathcal{Z}=\int\D[\hat Q]\,\exp\big(iS[\hat Q]\big)\,,\nonumber\\
&&iS[\hat Q]=-\frac{\pi\nu}{4\tauel}\Tr\big\{\hat{Q}^{2}\big\}+
\Tr\ln\left[\hat{G}^{-1}+\frac{i}{2\tauel}\hat{Q}\right]\,.
\end{eqnarray}
このようにして，当初の静的な場$U_{\mathrm{dis}}(\mathbf{r})$にわたる汎関数積分を，動的な行列場$\hat{Q}_{tt'}(\mathbf{r})$にわたる汎関数積分に置き換えた．一見，それほど優れた発想には思えないかもしれない．それでも，この手順には大きな簡単化が隠れている．要点は，乱れのポテンシャルが$\delta$相関をもつ，急激に振動する関数だということである．一方，以下に示すように，$\hat Q$行列場は（空間的にも時間的にも）ゆっくり変化する関数である．したがって，これは系の真の巨視的な（あるいは流体力学的な）自由度，すなわち拡散的に伝播するモードを表す．


%$$$$$$$$$$$$$$$$$$$$$$$$$$$$$$$$$$$$$$$$$$$$$$$$$$$$$$$$$$$$$$$$$$$$$$$$$$$$$$$$$$$$$$$$$$$$$$$$$$$$$$$$$$$

\subsectionlink{非線形$\sigma$模型}\label{sec_NLSM-2}



先に進むため，式~\eqref{NLSM-action}の作用$S[\hat{Q}]$の停留配置を求めよう．$\hat{Q}_{tt'}(\mathbf{r})$について変分すると，鞍点方程式
\begin{equation}\label{NLSM-saddle-point-eq}
\underline{\hat{Q}}_{tt'}(\mathbf{r})= \frac{i}{\pi\nu}
\left(\hat{G}^{-1}+\frac{i}{2\tauel}\underline{\hat{Q}}\right)^{-1}_{tt',\mathbf{r}\mathbf{r}}
\,,
\end{equation}
を得る．ここで，$\underline{\hat{Q}}_{tt'}(\mathbf{r})$は揺らぐ場$\hat{Q}_{tt'}(\mathbf{r})$の停留配置を表す．方針は，まず式~\eqref{NLSM-saddle-point-eq}の空間的に一様かつ時間並進不変な解$\underline{\hat{Q}}_{\,
t-t'} $を求め，次にその解からの空間・時間依存するずれを考えることである．この方針は磁性系の理論から取り入れたものであり，そこではまず一様で静的な磁化配置を求め，次いでその静的で一様な解の上の滑らかな摂動としてスピン波を扱う．式~\eqref{NLSM-saddle-point-eq}の構造から，停留配置$\underline{\hat{Q}}$はフェルミ粒子の自己エネルギー~\eqref{fermion-G-inverse}と同じ形をもつと期待される（より正確には，可能な停留配置の中に，式~\eqref{fermion-G-inverse}の因果的構造をもつ{\em 古典的な}配置があると期待される）．したがって，式~\eqref{NLSM-saddle-point-eq}の解を次の行列の形で求める：
\begin{equation}\label{NLSM-Lambda-trial}
\underline{\hat{Q}}_{\, t-t'}=\hat{\Lambda}_{t-t'}=\left(
\begin{array}{cc}\Lambda^{R}_{t-t'} & \Lambda^{K}_{t-t'}\\
0 & \Lambda^{A}_{t-t'}\end{array}\right)\,.
\end{equation}
この式を，エネルギー・運動量表示では$\hat{\Lambda}_{\epsilon}=\frac{i}{\pi\nu}\sum_{p}\big(
\epsilon-\epsilon_{p}+\frac{i}{2\tauel}\hat{\Lambda}_{\epsilon}\big)^{-1}$と書かれる式~\eqref{NLSM-saddle-point-eq}に代入する．ただし，$\epsilon_{p}\equiv p^{2}/2m-\epsilon_{F}$である．すると，次式を得る：
\begin{equation}
\Lambda^{R(A)}_{\epsilon}=\frac{i}{\pi\nu}\sum_{p}
\frac{1}{\epsilon-\epsilon_{p}+\frac{i}{2\tauel}
\Lambda^{R(A)}_{\epsilon}}=\pm1\,,
\end{equation}
ここで，$\sum_{p}\ldots\to\nu\int\d\epsilon_{p}$という規約を用いた．右辺の符号は因果性を満たすように選んである．すなわち，遅延（先進）Green関数は，複素エネルギー$\epsilon$の上（下）半平面全体で解析的である．また，状態密度$\nu$を一定としたままエネルギー積分を負の無限大まで延長するため，$1/\tauel\ll\epsilon_{F}$を仮定した．Keldysh成分はいつものように，Hermiteな分布関数を用いて$\Lambda^{K}=\Lambda^{R}\circ
F-F\circ\Lambda^{A}$とパラメータ表示できる．分布関数$F$は鞍点方程式~\eqref{NLSM-saddle-point-eq}だけでは決まらず，境界条件から決める必要がある．ただし平衡では，$F$は熱的なフェルミ粒子分布関数$F^{eq}_{\epsilon}=\tanh\frac{\epsilon}{2T}$にほかならず，したがって$\Lambda^{K}_{\epsilon}=(\Lambda^{R}_{\epsilon}-\Lambda^{A}_{\epsilon})
F^{eq}_{\epsilon}= 2F^{eq}_{\epsilon}$となる．最終的に，停留した$\underline{\hat{Q}}$行列配置について次式を得る：
\begin{equation}\label{NLSM-Lambda}
\hat{\Lambda}_{\epsilon}=\left(
\begin{array}{cc} 1^{R}_{\epsilon} & 2F_{\epsilon}\\
0 & -1^{A}_{\epsilon}\end{array}\right)\,,
\end{equation}
ここでは，因果的構造を明示するため，遅延および先進の単位行列を導入し，簡潔さのため，分布関数$F$の上付き添字「$eq$」を省略した．時間表示に戻すと$\Lambda^{R(A)}_{t-t'}=\pm\delta(t-t'\mp0)$を得る．ここで，$\mp0$は$\delta$関数が主対角線$t=t'$の下側（上側）へずれていることを示す．その結果，$\Tr\{\hat{\Lambda}\}=0$および$S[\hat
\Lambda]=0$となる．もちろん，これは任意の純粋に古典的な場の配置~\eqref{NLSM-Lambda-trial}が満たすべき性質である．ただし，鞍点解~\eqref{NLSM-Lambda}のこの特定の形は，1粒子状態密度$\nu$をエネルギーに依存しないとした近似の結果であることに注意しよう．一般には状態密度は$\epsilon$に依存するので，$\hat{\Lambda}_{\epsilon}$の遅延（先進）成分は上（下）半平面で解析的なエネルギーの関数であり，Fermiエネルギー$\epsilon_{F}$程度のスケールでエネルギーに依存する．したがって，$\Lambda^{R(A)}_{t-t'}=\pm\delta(t-t'\mp0)$に現れる無限小だけずれた$\delta$関数は，$\delta_{t\mp0}=f_{\pm}(t)\theta(\pm t)$と理解すべきである．ここで，$\theta(\pm t)$はHeaviside階段関数であり，$f_{\pm}(t)$は$|t|\lesssim\epsilon^{-1}_{F}$に鋭いピークをもち，規格化条件$\int^{\pm\infty}_{0}\d t f_{\pm}(t)=\pm 1$を満たす関数である．計算で偽の非物理的な定数を生じさせないためには，この高エネルギー正則化を念頭に置くことが重要である．特に，このために$1^{R}_{t-t'}M^{R}_{t',t}=0$および$1^{A}_{t-t'}M^{A}_{t',t}=0$となる．ここで，$M^{R(A)}_{t',t}$は時間空間における任意の遅延（先進）行列である．



これで，鞍点~\eqref{NLSM-Lambda}のまわりの揺らぎを調べる準備が整った．$\hat{Q}$の揺らぎは，大きく2種類に分類される：（i）質量$\propto
\nu/\tauel$をもつ有質量の揺らぎと，（ii）無質量の揺らぎ，すなわち作用がこれらの自由度の勾配または時間微分にのみ依存する揺らぎである．有質量モードに沿った揺らぎはGauss近似で積分消去でき，作用のパラメータに重要ではない繰り込みを与える．無質量モード，すなわちGoldstoneモードは電子の拡散運動を記述する．これらの無質量モードに沿った$\hat{Q}$行列の揺らぎは小さくないので，ある非線形な拘束条件を満たす行列によってパラメータ表示しなければならない．関連するGoldstoneモードを特定するため，式~\eqref{NLSM-action}の作用$S[\hat Q]$の第1項を考えよう．式~\eqref{NLSM-Lambda}で与えられる停留配置は
\begin{equation}\label{NLSM-Q-constraint}
\hat{Q}^{2}=\left(\begin{array}{cc}1^{R}_{\epsilon} & 0
\\ 0 & 1^{A}_{\epsilon}\end{array}\right)=\hat{1}\,.
\end{equation}
を満たす．遅延・先進の単位行列の定義により，$\Tr\big\{\hat{Q}^{2}\big\}=\Tr\,\{\hat{1}^R\}+
\Tr\,\{\hat{1}^A\}=0$であることに注意しよう．式~\eqref{NLSM-Q-constraint}を満たさない$\hat{Q}$の揺らぎは有質量である．拘束条件~\eqref{NLSM-Q-constraint}に従う$\hat{Q}$行列配置の集合は，停留行列$\hat{\Lambda}_{\epsilon}$の回転によって生成され，次のようにパラメータ表示できる：
\begin{equation}\label{NLSM-Q-Rotated}
\hat{Q}=\hat{\mathcal{R}}^{-1}\circ
\hat{\Lambda}\circ\hat{\mathcal{R}}\,.
\end{equation}
$\hat{\mathcal{R}}$の具体的な形は現段階では重要ではなく，後で選ぶことにする．$\hat{\mathcal{R}}_{tt'}(\mathbf{r})$に対応する無質量モード（磁性との類推を用いればスピン波）は，$t+t'$と$\mathbf{r}$のゆっくり変化する関数であり，その勾配は小さい．ここでの目標は，式~\eqref{NLSM-Q-constraint}および\eqref{NLSM-Q-Rotated}で与えられるソフトモードの$\hat{Q}$場の配置に対する作用を導くことである．



そのため，式~\eqref{NLSM-Q-Rotated}を式~\eqref{NLSM-action}に代入し，トレースの中で$\hat{\mathcal{R}}$行列を巡回置換する．こうして$\hat{\mathcal{R}}\circ
\hat{G}^{-1}\circ\hat{\mathcal{R}}^{-1}=\hat{G}^{-1}+\hat{\mathcal{R}}
\circ[\hat{G}^{-1}\stackrel{\circ}{,}\hat{\mathcal{R}}^{-1}]
=\hat{G}^{-1}+i\hat{\mathcal{R}}\partial_{t}\hat{\mathcal{R}}^{-1}+
i\hat{\mathcal{R}}\mathbf{v}_{F}\partial_{\mathbf{r}}\hat{\mathcal{R}}^{-1}$を得る．ここでは，Fermi面近傍の分散関係を$\epsilon_{p}=p^{2}/2m-\epsilon_{F}\approx\mathbf{v}_{F}\mathbf{p}\to
-i\mathbf{v}_{F}\partial_{\mathbf{r}}$と線形化した．その結果，求める作用は次の形をとる：
\begin{equation}\label{NLSM-action-trlog-R}
iS[\hat Q]=\Tr\ln\left[\hat{1}+i\hat{\mathcal{G}}\hat{\mathcal{R}}
\partial_{t}\hat{\mathcal{R}}^{-1}+
i\hat{\mathcal{G}}\hat{\mathcal{R}}\mathbf{v}_{F}
\partial_{\mathbf{r}}\hat{\mathcal{R}}^{-1}\right]\,,
\end{equation}
ここでは簡潔さのため，丸付きの積の記号を省略した．$\hat{\mathcal{G}}$は，Dyson方程式$\big(\hat{G}^{-1}+\frac{i}{2\tauel}\hat{\Lambda}\big)\hat{\mathcal{G}}=\hat{1}$で定義される，{\em 不純物による効果を繰り込んだ}Green関数行列である．実際の計算では，$\hat{\mathcal{G}}$を次の形で書くと便利である：
\begin{equation}\label{NLSM-G-impurity-dressed}
\hat{\mathcal{G}}=\left(\begin{array}{cc}\mathcal{G}^{R} &
\mathcal{G}^{K}\\ 0 &
\mathcal{G}^{A}\end{array}\right)=\frac{1}{2}\mathcal{G}^{R}[\hat{1}+\hat{\Lambda}]+
\frac{1}{2}\mathcal{G}^{A}[\hat{1}-\hat{\Lambda}]\,,
\end{equation}
遅延，先進，Keldysh成分は次式で与えられる：
\begin{equation}\label{NLSM-G-impurity-dressed-RAK}
\mathcal{G}^{R(A)}(\mathbf{p},\epsilon)=\big[\epsilon-\epsilon_{p}\pm
i/2\tauel\big]^{-1},\qquad
\mathcal{G}^{K}(\mathbf{p},\epsilon)=\mathcal{G}^{R}(\mathbf{p},\epsilon)
F_{\epsilon}-F_{\epsilon}\mathcal{G}^{A}(\mathbf{p},\epsilon)\,.
\end{equation}
次に，式~\eqref{NLSM-action-trlog-R}の対数を，回転行列$\hat{\mathcal{R}}$の勾配について，$\partial_{t}\hat{\mathcal{R}}^{-1}$の1次および$\partial_{\mathbf{r}}\hat{\mathcal{R}}^{-1}$の2次まで展開する（空間勾配の1次の寄与は角度積分により消える）．その結果，
\begin{equation}\label{NLSM-action-trlog-Approx}
iS[\hat Q]\approx
i\Tr\big\{\hat{\mathcal{G}}\hat{\mathcal{R}}\partial_{t}\hat{\mathcal{R}}^{-1}\big\}
+\frac{1}{2}\Tr\big\{
\hat{\mathcal{G}}\big(\hat{\mathcal{R}}\mathbf{v}_{F}
\partial_{\mathbf{r}}\hat{\mathcal{R}}^{-1}\big)
\hat{\mathcal{G}}\big(\hat{\mathcal{R}}\mathbf{v}_{F}
\partial_{\mathbf{r}}\hat{\mathcal{R}}^{-1}\big)\big\}\,.
\end{equation}
となる．鞍点方程式~\eqref{NLSM-saddle-point-eq}から直接従う$\sum_{p}\hat{\mathcal{G}}(\mathbf{p},\epsilon)=-i\pi\nu\hat{\Lambda}_{\epsilon}$という関係により，作用の$\partial_{t}$項に対して$i\Tr\{\hat{\mathcal{G}}\hat{\mathcal{R}}\partial_{t}\hat{\mathcal{R}}^{-1}\}=
\pi\nu\Tr\{\partial_{t}\hat{Q}\}$を得る．$\partial_{\mathbf{r}}$項に対しては$-\frac{1}{4}\pi\nu D\Tr\big\{(\partial_{\mathbf{r}}
Q)^{2}\big\}$を得る．ここで，$D=v^{2}_{F}\tauel/d$は拡散定数，$d$は空間次元である．実際，Green関数の積に対して$\sum_{p}\mathcal{G}^{R}(\mathbf{p},\epsilon)\mathbf{v}_{F}
\mathcal{G}^{A}(\mathbf{p},\epsilon)\mathbf{v}_{F}=2\pi\nu\tauel
v^{2}_{F}/d=2\pi\nu D$を用いる一方，対応する$R--R$および$A--A$の項は$\epsilon_{p}$積分を実行すると消える．その後，式~\eqref{NLSM-G-impurity-dressed}を用いると，$\Tr\big\{[\hat{1}+\hat{\Lambda}]
(\hat{\mathcal{R}}\partial_{\mathbf{r}} \hat{\mathcal{R}}^{-1})
[\hat{1}-\hat{\Lambda}](\hat{\mathcal{R}}
\partial_{\mathbf{r}}\hat{\mathcal{R}}^{-1})\big\}=-\frac{1}{2}
\Tr\big\{\big(\partial_{\mathbf{r}}(\hat{\mathcal{R}}^{-1}
\hat{\Lambda}\hat{\mathcal{R}})\big)^{2}\big\}=
-\frac{1}{2}\Tr\big\{(\partial_{\mathbf{r}}\hat{Q})^{2}\big\}$に到達する．最終的に，ソフトモード配置の作用として次式を得る~\cite{HorbachSchon,KamenevAndreev,ChamonLudwigNayak}：
\begin{equation}\label{NLSM-action-noninteracting}
iS[\hat Q]=-\frac{\pi\nu}{4}\,\Tr\Big\{D(\partial_{\mathbf{r}}
\hat{Q})^{2}-4\partial_{t}\hat{Q}\Big\}\,.
\end{equation}
見かけは単純だが，作用~\eqref{NLSM-action-noninteracting}は拘束条件$\hat{Q}^{2}=\hat{1}$のために強い非線形性をもつ．式~\eqref{NLSM-Q-constraint}と\eqref{NLSM-action-noninteracting}で規定される理論を，{\em 行列非線形$\sigma$模型}と呼ぶ．この名称は磁性の理論に由来する．そこでは，長さ1のベクトル$\vec \sigma(\mathbf{r})$が局所的な（古典）スピンを表し，球面$\vec \sigma^{2}=1$上を回転できる．



次に，ソース項$S_{V}$と$S_{A}$（式~\eqref{fermion-SV}および\eqref{fermion-SA}参照）を，作用のフェルミ粒子部分に取り込むことができる：
$$\Tr\left\{\vec{\bar{\Psi}}\left[\hat{G}^{-1}+\frac{i}{2\tauel}\hat{Q}+
\hat{V}+\mathbf{v}_{F}\hat{\mathbf{A}}\right]\vec{\Psi}\right\}\,.$$
$\bar{\Psi}$と$\Psi$についてGauss積分を行うと，ソース場に依存する分配関数として次式を得る（式~\eqref{NLSM-action}と比較されたい）：
\begin{eqnarray}\label{NLSM-action-SAV}
&&\mathcal{Z}[\mathbf{A},V]=\int\D[\hat Q]
\exp\big(iS[\hat Q,\mathbf{A},V]\big)\,,\nonumber\\
&& iS[\hat
Q,\mathbf{A},V]=-\frac{\pi\nu}{4\tauel}\Tr\{\hat{Q}^{2}\}+
\Tr\ln\left[\hat{G}^{-1}+\frac{i}{2\tauel}\hat{Q}+
\hat{V}+\mathbf{v}_{F}\hat{\mathbf{A}}\right]\,.
\end{eqnarray}
式~\eqref{NLSM-Q-Rotated}を用いて，対数のトレースを$\hat{Q}$の勾配で展開する．この際，ソース場$\hat{V}$と$\hat{\mathbf{A}}$は何らかの意味で小さく，停留配置~\eqref{NLSM-Lambda}を乱さないと仮定する（この点の議論については第~\ref{sec_int_ferm}節参照）．すると，式~\eqref{NLSM-action-noninteracting}と同様に，式~\eqref{NLSM-action-SAV}から次式を得る：
\begin{eqnarray}\label{NLSM-action-general}
iS[\hat Q,\mathbf{A},V]=\frac{i\nu}{2}\, \Tr\big\{\hat{V}\hat{\sigma}_{x}\hat{V}\big\}
-\frac{\pi\nu}{4}\, \Tr\Big\{D(\hat{\bm{\partial}}_{\mathbf{r}}\hat{Q})^{2}-
4\partial_{t}\hat{Q}+4i\hat{V}\hat{Q}\Big\}\,,
\end{eqnarray}
ここで，$\hat{\sigma}_{x}$はKeldysh空間に作用するPauli行列であり，共変微分
\begin{equation}\label{NLSM-covariant-deriv}
\hat{\bm{\partial}}_{\mathbf{r}}\hat{Q}=
\partial_{\mathbf{r}}\hat{Q}-i[\hat{\mathbf{A}},\hat{Q}]\,.
\end{equation}
を導入した．式~\eqref{NLSM-action-general}について，いくつか補足しておこう．まず，この式も依然として$\hat{Q}^{2}=\hat{1}$を満たす$\hat{Q}$行列の多様体に制限される．式~\eqref{NLSM-action-general}の右辺の第2のトレースは$\hat{Q}$を含み，対数の展開に現れる$\sum_{p}\mathbf{v}_{F}\mathcal{G}^{R}\mathbf{v}_{F}\mathcal{G}^{A}$および$\sum_{p}\mathcal{G}^{R(A)}$の組合せに由来する．一方，式~\eqref{NLSM-action-general}の右辺の第1項は，$\sum_{p}\mathcal{G}^{R}\mathcal{G}^{R}$および$\sum_{p}\mathcal{G}^{A}\mathcal{G}^{A}$の組合せに由来する．行列$V^{\alpha}(\epsilon-\epsilon')\hat{\gamma}^{\alpha}$はFermiエネルギー近傍の幅$1/\tauel$のシェルに制限されないため，これらの項を保持しなければならない．これは，スカラーポテンシャルがエネルギー帯域$|\epsilon|,|\epsilon'|<1/\tauel$だけでなく，電子バンド全体をシフトさせるためである．したがって，電荷中性を保つには，電子スペクトルの変化をバンドの底まで追うことが不可欠である．$\Tr\{\hat{V}\hat{\sigma}_{x}\hat{V}\}$を導くには，任意の物理的なフェルミ粒子分布関数について$F_{\epsilon\to\pm\infty}\to\pm1$であることを用いなければならない．式~\eqref{NLSM-action-general}および\eqref{NLSM-covariant-deriv}は，式~\eqref{NLSM-action-noninteracting}で与えられる有効$\sigma$模型の作用を一般化する．式~\eqref{NLSM-action-SAV}から式~\eqref{NLSM-action-general}を導くために必要な追加の計算上の詳細は，\ref{app_GradientExpansion}に示した．


%$$$$$$$$$$$$$$$$$$$$$$$$$$$$$$$$$$$$$$$$$$$$$$$$$$$$$$$$$$$$$$$$$$$$$$$$$$$$$$$$$$$$$$$$$$$$$$$$$$$$$$$$$$$

\subsectionlink{トンネル作用}\label{sec_NLSM-3}



トンネル障壁で隔てられた2つの金属リードを考えよう．外部電圧を印加すると，両者の間に電流が流れる．この場合，系のHamiltonianに対応するトンネル項を加えなければならない：
$$\hat{H}_{T}=\int_{\mathbf{r}\in
L}\d\mathbf{r}\int_{\mathbf{r}'\in
R}\d\mathbf{r}'\big[T^{\phantom{*}}_{\mathbf{rr}'}
\hat{\psi}^{\dag}_{L}(\mathbf{r})
\hat{\psi}^{\phantom{\dag}}_{R}(\mathbf{r}')+T^{*}_{\mathbf{rr}'}
\hat{\psi}^{\dag}_{R}(\mathbf{r}')\hat{\psi}^{\phantom{\dag}}_{L}(\mathbf{r})\big]\,,$$
ここで，$\hat{\psi}_{L(R)}$はトンネル障壁の左側（右側）における電子の消滅演算子であり，$\hat{\psi}^{\dag}_{L(R)}$は対応する生成演算子である．$T^{\phantom{*}}_{\mathbf{rr}'}$と$T^{*}_{\mathbf{rr}'}$はトンネル行列要素であり，電子波動関数の重なりが接合から離れるにつれて指数関数的に減衰するため，その作用範囲は接合の近傍に限られる．トンネルHamiltonianは，次のフェルミ粒子のトンネル作用に対応する：
$$iS_{T}=\int_{\mathcal{C}}\d
t\iint\d\mathbf{r}\d\mathbf{r}' \big[T^{\phantom{*}}_{\mathbf{rr}'}
\bar{\psi}_{L}(\mathbf{r},t)
\psi_{R}(\mathbf{r}',t)+T^{*}_{\mathbf{rr}'}
\bar{\psi}_{R}(\mathbf{r}',t)\psi_{L}(\mathbf{r},t)\big]\,.$$
$S_{T}$は依然としてフェルミ場の2次形式なので，そのGauss積分は容易に行え，次の形の，乱れについて平均された作用が得られる：
\begin{eqnarray}\label{NLSM-tunnel-Z}
&&\mathcal{Z}=\int\D[\hat Q_{L}, \hat Q_{R}]\,\exp\big(iS[\hat Q_{L},\hat Q_{R}]\big)\,,\nonumber
\\
&&iS[\hat Q_{L}, \hat
Q_{R}]=-\frac{\pi\nu}{4\tauel}\sum_{a=L,R}\Tr\big\{\hat{Q}^{2}_{a}\big\}
+\Tr\ln\left(\begin{array}{cc}
\hat{G}_L^{-1}+\frac{i}{2\tauel}\hat{Q}_{L} & \hat{T} \\
\hat{T}^{\dag} & \hat{G}_R^{-1}+\frac{i}{2\tauel}\hat{Q}_{R}
\end{array}\right)\,.
\end{eqnarray}
式~\eqref{NLSM-tunnel-Z}の導出では，各リードの乱れによる4フェルミオン項~\eqref{NLSM-HS-transform}を独立に分離するため，2つの$\hat Q$行列を導入する必要がある．その際，簡単のため，両方の乱れた試料は同じ平均自由時間と裸の電子状態密度をもつと仮定した．式~\eqref{NLSM-tunnel-Z}には，それぞれ$\hat{Q}_{L(R)}$で記述される左・右の電子部分系の空間に，追加の$2\times2$行列構造が含まれている．また，$S[\hat{Q}_{L},\hat{Q}_{R}]$に現れるトンネル行列要素はKeldysh部分空間では単位行列に比例し，$\hat{T}_{\mathbf{rr}'}=T_{\mathbf{rr}'}\hat{\sigma}_{0}$であることにも注意しよう．



表記$\hat{\mathbf{G}}^{-1}_{a}=\hat{G}^{-1}_a +
\frac{i}{2\tauel}\hat{Q}_{a}$を導入すると，式~\eqref{NLSM-tunnel-Z}の作用$S[\hat{Q}_{L},\hat{Q}_{R}]$の最後の項は，恒等的に次のように書き直せる：
\begin{equation}\label{NLSM-trlog-transform}
\Tr\ln\left(\begin{array}{cc}
\hat{\mathbf{G}}^{-1}_{L} & \hat{T} \\
\hat{T}^{\dag} & \hat{\mathbf{G}}^{-1}_{R}
\end{array}\right)=\Tr\ln\left(\begin{array}{cc}
\hat{\mathbf{G}}^{-1}_{L} & 0 \\
0 & \hat{\mathbf{G}}^{-1}_{R}\end{array}\right)
+\Tr\ln\left[\hat 1+\left(\begin{array}{cc}
0 & \hat{\mathbf{G}}_{L}\hat{T} \\
\hat{\mathbf{G}}_{R}\hat{T}^{\dag} & 0
\end{array}\right)\right]\,.
\end{equation}
ここで，鞍点$\hat{\Lambda}_{a}$のまわりで，$\Tr\ln\hat{\mathbf{G}}^{-1}_{a}$を$\hat{Q}_{a}$行列の勾配によって展開すると，2つのリードのそれぞれに対して独立に，σ模型の作用~\eqref{NLSM-action-noninteracting}を得る．両者の結合は，式~\eqref{NLSM-trlog-transform}の右辺の第2項で記述され，これによってトンネル作用$S_{T}[\hat{Q}_{L},\hat{Q}_{R}]$が定義される．透過率の小さいトンネル接合では，対数のトレースを$\hat{T}$の最低次（2次）まで展開して，次式を得る：
\begin{equation}\label{NLSM-action-S-T-log-expansion}
iS_{T}[\hat Q_{L},\hat Q_{R}]=\Tr\ln\left[\hat 1+\left(\begin{array}{cc}
0 & \hat{\mathbf{G}}_{L}\hat{T} \\
\hat{\mathbf{G}}_{R}\hat{T}^{\dag} & 0
\end{array}\right)\right]
\approx-\Tr\big\{\hat{\mathbf{G}}_{L}\hat{T}\hat{\mathbf{G}}_{R}\hat{T}^{\dag}\big\}+\ldots\,.
\end{equation}
行列要素$T_{{\bf r r}'}$の局所性と，ソフトモード多様体上では$\hat{Q}_{a}=\frac{i}{\pi\nu}\hat{\mathbf{G}}_{a}({\bf r},{\bf r})$であること（式~\eqref{NLSM-saddle-point-eq}参照）を用いると，作用のトンネル部分として次式を得る：
\begin{equation}\label{NLSM-action-S-T}
iS_{T}[\hat Q_{L},\hat Q_{R}]=\frac{
\mathrm{g}_{T}}{4\mathrm{g}_{Q}}\Tr\big\{\hat{Q}_{L}\hat{Q}_{R}\big\}
= - \frac{\mathrm{g}_{T}}{8\mathrm{g}_{Q}}
\Tr\big\{(\hat{Q}_{L} - \hat{Q}_{R})^2 \big\}\,.
\end{equation}
ここでは，トンネル行列を$T_{\mathbf{rr}'}=T_{0}\delta(\mathbf{r}-\mathbf{r}')$と近似し，トンネルコンダクタンス$\mathrm{g}_{T}=4\pi^2
e^{2}|T_{0}|^{2}\nu^2$と量子コンダクタンス$\mathrm{g}_Q=e^2/(2\pi \hbar )$を導入した．トンネル作用~(\ref{NLSM-action-S-T})は，NLSM作用~(\ref{NLSM-action-noninteracting})の$\Tr\big\{D(\partial_{\bf r} Q)^2\big\}$項を，トンネル構造に一般化したものである．



トンネル振幅$T_{\mathbf{rr}'}$が小さくない場合には，式~(\ref{NLSM-action-S-T-log-expansion})の対数の展開で高次項を保持する必要がある．偶数個のトンネル振幅$T_{\mathbf{rr}'}$の積を，各横方向チャネルの透過確率$T_n$で表すと便利である（例えば文献~\cite{ABG}の付録C参照）．式~\eqref{NLSM-saddle-point-eq}を用いると，式~(\ref{NLSM-action-S-T-log-expansion})の対数の展開は，各次数ごとに次の作用の展開と等価であることを示せる~\cite{Nazarov-TunnelingAction,BelzigNazarov,FKLS}：
\begin{equation}\label{NLSM-action-S-T-general}
iS_{T}[\hat Q_{L},\hat Q_{R}]=\frac{
1}{2}\sum_{n}\Tr\ln\left[\hat 1 - \frac{T_{n}}{4}
\left(\hat{Q}_{L}-\hat{Q}_{R}\right)^2 \right]\,.
\end{equation}
すべての透過確率が小さく$T_{n}\ll 1$であれば，式~\eqref{NLSM-action-S-T-general}を$T_{n}$の最低次まで展開し，トンネルコンダクタンスを$\mathrm{g}_{T}=\mathrm{g}_{Q}\sum_n T_n$と同定して，式~\eqref{NLSM-action-S-T}を再現できる（式~(\ref{Part-I-g})参照）．式~(\ref{NLSM-action-S-T-general})はこの極限を超えて，任意の2端子形状におけるメゾスコピック輸送を扱うことを可能にする．多端子の場合への一般化も，Nazarovらによって展開された~\cite{Nazarov-TunnelingAction,Nazarov-CBWTJ,SnymanNazarov}．


%$$$$$$$$$$$$$$$$$$$$$$$$$$$$$$$$$$$$$$$$$$$$$$$$$$$$$$$$$$$$$$$$$$$$$$$$$$$$$$$$$$$$$$$$$$$$$$$$$$$$$$$$$$$

\subsectionlink{Usadel方程式}\label{sec_NLSM-4}



式~\eqref{NLSM-action-noninteracting}で規定される作用に戻ろう．目標はNLSMの物理的帰結を調べることである．第一歩として，ソフトモード多様体~\eqref{NLSM-Q-constraint}上で最も確からしい（停留した）配置$\underline{\hat{Q}}_{tt'}(\mathbf{r})$を決める必要がある．そのため，$\underline{\hat{Q}}_{tt'}(\mathbf{r})$からのずれを$\hat{Q}=\hat{\mathcal{R}}^{-1}\circ\underline{\hat{Q}}\circ\hat{\mathcal{R}}$とパラメータ表示し，$\hat{\mathcal{R}}=\exp(\hat{\mathcal{W}}/2)$を選ぶ．ここで，$\hat{\mathcal{W}}_{tt'}(\mathbf{r})$は回転の生成子である．$\hat{\mathcal{W}}$の1次まで展開すると，$\hat{Q}=\underline{\hat{Q}}-[\hat{\mathcal{W}}\stackrel{\circ}{,}\underline{\hat{Q}}]/2$を得る．この$\hat Q$行列を作用~\eqref{NLSM-action-noninteracting}に代入し，$\hat{\mathcal{W}}$に比例する項が消えることを要求すればよい．これにより，$\hat{\underline{Q}}$に対する鞍点方程式が得られる．式~\eqref{NLSM-action-noninteracting}の右辺の波括弧内の第1項からは，$\frac{1}{2}\Tr\big\{\hat{\mathcal{W}}\partial_{\mathbf{r}}
D\,\big[(\partial_{\mathbf{r}}\underline{\hat{Q}})
\underline{\hat{Q}}-\underline{\hat{Q}}
\partial_{\mathbf{r}}\underline{\hat{Q}}\big]\big\}=
-\Tr\big\{\hat{\mathcal{W}}\partial_{\mathbf{r}}
D\,\big(\underline{\hat{Q}}\partial_{\mathbf{r}}
\underline{\hat{Q}}\big)\big\}$を得る．ここでは，$\underline{\hat{Q}}^{2}=\hat{1}$より$\partial_{\mathbf{r}}\underline{\hat{Q}}\circ
\underline{\hat{Q}}+\underline{\hat{Q}}\circ
\partial_{\mathbf{r}}\underline{\hat{Q}}=0$となることを用いた．第2項については$\Tr\big\{\hat{\mathcal{W}}_{tt'}\big(\partial_{t}+\partial_{t'}\big)
\underline{\hat{Q}}_{t't}\big\}=
\Tr\big\{\hat{\mathcal{W}}\{\partial_{t},\underline{\hat{Q}}\}\big\}$を得る．$\hat{\mathcal{W}}$の1次の項が消えることを要求すると，次式を得る：
\begin{equation}\label{NLSM-Usadel}
\partial_{\mathbf{r}}
\big(D\,
\underline{\hat{Q}}\circ\partial_{\mathbf{r}}\underline{\hat{Q}}\big)-
\{\partial_{t},\underline{\hat{Q}}\}=0\,.
\end{equation}
これが，停留した$\underline{\hat{Q}}$行列に対するUsadel方程式~\cite{Usadel}である．Usadel方程式の解を，因果的構造をもつ「古典的な」配置の部分空間で求めるなら，まだ指定していない分布関数$F_{tt'}(\mathbf{r})$を含む$\underline{\hat{Q}}=\hat{\Lambda}$をとる．したがって，この場合Usadel方程式は，分布関数$F_{tt'}(\mathbf{r})$に対する1つの方程式に帰着する．式~(\ref{NLSM-Lambda})の$\hat{\Lambda}$を式~\eqref{NLSM-Usadel}に代入し，Wigner変換
\begin{equation}\label{NLSM-Wigner}
F_{tt'}(\mathbf{r})=\int\frac{\d\epsilon}{2\pi}\,F_{\epsilon}
\left(\mathbf{r},\tau\right)\,e^{-i\epsilon(t-t')},\quad \quad
\tau=\frac{t+t'}{2}\, ,
\end{equation}
を行うと，次式を得る：
\begin{equation}\label{NLSM-Diffusion-Eq}
\partial_{\mathbf{r}}\big[D(\mathbf{r}) \partial_{\mathbf{r}} F_{\epsilon}(\mathbf{r},\tau)\big]-
\partial_{\tau}F_{\epsilon}(\mathbf{r},\tau)=0\,,
\end{equation}
ここでは，拡散定数の（滑らかな）空間依存性を許した．これは，相互作用のない極限での，乱れた系のフェルミ粒子分布関数に対する運動論方程式であり，拡散方程式にほかならない．どのエネルギー$\epsilon$に対しても同じ方程式となり，異なるエネルギー間の結合がないことに注意しよう．これは相互作用のない系として自然である．相互作用が存在すると，右辺に，異なるエネルギーを互いに混ぜ合わせる衝突積分が加わる．弾性散乱はこの衝突積分には現れないことを指摘しておこう．それはUsadel方程式の導出ですでに完全に考慮されており，拡散項に取り込まれている．



例として，異なる電圧に保たれた2つのリードに接続された，長さ$L$の，乱れのある準一次元細線を考えよう~\cite{Pierre}．細線に沿う座標を$x$とし，$D\,\partial^{2}_{x}F_{\epsilon}(x)=0$を満たす空間依存の定常関数$F_{\epsilon}(x)$を求める．境界条件は$F_{\epsilon}(x=0)=F_{L}(\epsilon)$および$F_{\epsilon}(x=L)=F_{R}(\epsilon)$であり，$F_{R(L)}(\epsilon)$は左右のリードの分布関数である．適切な解は
\begin{equation}
F_{\epsilon}(x)=F_{L}(\epsilon)+[F_{R}(\epsilon)-F_{L}(\epsilon)]\frac{x}{L}\,.
\end{equation}
である．細線内の分布関数は，2つの分布の間を線形に補間する．低温では2段の階段関数のような形をとり，段差間のエネルギー差は印加電圧$eV$で与えられ，相対的な高さは位置$x$に依存する．式~\eqref{NLSM-Diffusion-Eq}を連続の式と比較すると，（所定のエネルギー$\epsilon$での）電流密度が$j(\epsilon)=D\,\partial_{x}
F_{\epsilon}(x)=D[F_{R}(\epsilon)-F_{L}(\epsilon)]/L$で与えられることがわかる．したがって，全電流は
$$I=e\nu\int\d\epsilon
j(\epsilon)=\frac{e\nu
D}{L}\int\d\,\epsilon[F_{R}(\epsilon)-F_{L}(\epsilon)]=e^{2}\frac{\nu
D}{L}V=\sigma_D V/L\,,$$
となる．ここで，拡散細線のDrude伝導率は$\sigma_{D}=e^{2}\nu D$で与えられる．


%$$$$$$$$$$$$$$$$$$$$$$$$$$$$$$$$$$$$$$$$$$$$$$$$$$$$$$$$$$$$$$$$$$$$$$$$$$$$$$$$$$$$$$$$$$$$$$$$$$$$$$$$$$$

\subsectionlink{揺らぎ}\label{sec_NLSM-5}



前節までの議論に従い，停留解$\underline{\hat{Q}}_{tt'}(\mathbf{r})=\hat{\Lambda}_{t-t'}$（式~\eqref{NLSM-Lambda}参照）の近傍の揺らぎを考える．$\hat{Q}^{2}=\hat{1}$を満たすソフトモードの揺らぎに限定し，この多様体の外にある有質量モードをすべて無視する．$\hat Q$行列の無質量の揺らぎは，次のようにパラメータ表示できる：
\begin{equation}\label{NLSM-Q-U-W-parametrization}
\hat{Q}=\hat{\mathcal{U}}\circ
e^{-\hat{\mathcal{W}}/2}\circ\hat{\sigma}_{z}\circ
e^{\hat{\mathcal{W}}/2}\circ\hat{\mathcal{U}}^{-1}\,,
\end{equation}
ここで，回転の生成子は次式で与えられる：
\begin{equation}\label{NLSM-U-W}
\hat{\mathcal{W}}=
\left(\begin{array}{cc}0&d\\
\bar{d}&0\end{array}\right)\,,\; \quad\quad \quad
\hat{\mathcal{U}}=\hat{\mathcal{U}}^{-1}=
\left(\begin{array}{cc}1&F\\0&-1\end{array}\right)\, .
\end{equation}
ここで，$d_{tt'}(\mathbf{r})$と$\bar{d}_{tt'}(\mathbf{r})$は時間空間における2つの独立なHermite行列である．したがって，式~\eqref{NLSM-action-SAV}における$\hat{Q}_{tt'}(\mathbf{r})$の汎関数積分は，時間領域の互いに独立な2つのHermite行列$d_{tt'}(\mathbf{r})$と$\bar{d}_{tt'}(\mathbf{r})$についての積分と理解される．$d_{tt'}(\mathbf{r})$の物理的意味は，フェルミ粒子分布関数$F_{tt'}(\mathbf{r})$の定常値からのずれである．一方，$\bar{d}_{tt'}(\mathbf{r})$には古典的解釈はない．これはおおむね$d_{tt'}(\mathbf{r})$の量子的な対応物の役割を果たし，図式では内線としてのみ現れる．$\hat{Q}$を式~\eqref{NLSM-Q-U-W-parametrization}の形に選ぶ理由は，次のように説明できる．まず，$\underline{\hat{Q}}\equiv\hat{\Lambda}=
\hat{\mathcal{U}}\,\hat{\sigma}_{z}\,\hat{\mathcal{U}}^{-1}$に注意する．次に，$\underline{\hat{Q}}$と可換な$\hat{\mathcal{W}}$の部分はいかなる揺らぎも生成しないため，$\hat{\mathcal{W}}$を$\hat{\mathcal{W}}\,\hat{\sigma}_{z}+
\hat{\sigma}_{z}\,\hat{\mathcal{W}}=0$を満たすものに制限する．したがって，$\hat{\mathcal{W}}$は非対角でなければならず，最も一般的には，式~\eqref{NLSM-U-W}のように2つの独立な場$d$と$\bar{d}$によってパラメータ表示される．



次に，作用~\eqref{NLSM-action-noninteracting}を$\bar{d}_{tt'}(\mathbf{r})$と$d_{tt'}(\mathbf{r})$のべきで展開する．$\underline{\hat{Q}}_{tt'}$は停留点として選んだため，展開は2次から始まる．定常分布$F_{t,t'}({\bf r})$が空間的に一様であれば，次式を得る：
\begin{equation}\label{NLSM-S-W}
iS[\hat{\mathcal{W}}]=-\frac{\pi\nu}{2}\int\d\mathbf{r} \iint\d t\d t'\,
\bar{d}_{tt'}(\mathbf{r})
\left[-D\,\partial^{2}_{\mathbf{r}}+\partial_{t}+\partial_{t\,'}\right]
d_{t't}(\mathbf{r})\,.
\end{equation}
この2次形式は，エネルギー・運動量表示への変換
$$\hat{\mathcal{W}}_{\epsilon\epsilon'}(\mathbf{q})=\int\d\mathbf{r}
\iint\d t\d t'\hat{\mathcal{W}}_{tt'}(\mathbf{r})\exp(i\epsilon
t-i\epsilon't')\exp(-i\mathbf{qr})\,.$$
によって対角化できる．その結果，小さな$\hat Q$行列の揺らぎの伝播関数は
\begin{equation}\label{NLSM-d-d}
\langle d_{\epsilon_{2}\epsilon_{1}}(\mathbf{q})
\bar{d}_{\epsilon_{3}\epsilon_{4}}(-\mathbf{q})\rangle_{\mathcal{W}}=
-\frac{2}{\pi\nu}\,\frac{\delta_{\epsilon_{1}\epsilon_{3}}\delta_{\epsilon_{2}\epsilon_{4}}}
{Dq^2+i\omega}\equiv -\frac{2}{\pi\nu}\,
\delta_{\epsilon_{1}\epsilon_{3}}\delta_{\epsilon_{2}\epsilon_{4}}\,
\mathcal{D}^A(\mathbf{q},\omega)\,,
\end{equation}
となる．ここで，$\omega\equiv\epsilon_{1}-\epsilon_{2}=\epsilon_3-\epsilon_4$であり，$\mathcal{D}^{R(A)}(\mathbf{q},\omega)=\mathcal{D}^{R(A)}(\mathbf{q},\epsilon_{1}-\epsilon_{2})=
\big[Dq^{2}\mp i(\epsilon_{1}-\epsilon_{2})\big]^{-1}$を{\em ディフューゾン}と呼ぶ．作用~\eqref{NLSM-action-noninteracting}を$d_{tt'}(\mathbf{r})$と$\bar{d}_{tt'}(\mathbf{r})$について展開した高次項は，{\em Hikamiボックス}と呼ばれる頂点をもつ拡散モード間の非線形相互作用を記述する~\cite{Hikami,Gor'kovLarkinKhmelnitskii}．これらの非線形項は弱局在補正をもたらす~\cite{Gor'kovLarkinKhmelnitskii,AAKL,LeeRamakrishnan,ALW-book}．分布関数$F_{tt'}(\mathbf{r})$が空間的に非一様であれば，2次の作用に追加の項$-(\pi\nu D/2)\Tr\big\{\bar{d}(\partial_{\mathbf{r}}
F)\bar{d}(\partial_{\mathbf{r}} F)\big\}$が現れる．この項は$\langle
dd\rangle_{\mathcal{W}}$型の非ゼロの相関を生成し，いくつかの応用で重要となる．



%$$$$$$$$$$$$$$$$$$$$$$$$$$$$$$$$$$$$$$$$$$$$$$$$$$$$$$$$$$$$$$$$$$$$$$$$$$$$$$$$$$$$$$$$$$$$$$$$$$$$$$$$$$$

\subsectionlink{応用II：メゾスコピック効果}\label{app_Part-II}
\subsubsection{Kubo公式と線形応答}\label{app_Part-II-1}



第~\ref{sec_fermion-4}節では，Keldysh法において線形応答理論をどのように定式化するかを示した．ここでは，乱れた電子気体の分極演算子がNLSM作用からどのように得られるかを見てみよう．そのため，式~\eqref{fermion-Pi-matrix-def}による密度応答関数$\Pi^{R}(x,x\,')$の一般的な定義と，乱れについて平均した作用~\eqref{NLSM-action-general}を用いると，次式を得る：
\begin{equation}\label{Part-II-Pi}
\fitdisplay{
\Pi^{R}(x,x\,')=-\frac{i}{2}\left.\frac{\delta^{2}
\mathcal{Z}[V^{cl},V^{q}]}{\delta V^{cl}(x\,')\delta
V^{q}(x)}\right|_{\hat{V}=0}=\nu\delta(\mathbf{r}-\mathbf{r}')\delta(t-t')
+\frac{i}{2}(\pi\nu)^{2}
\left\langle\Tr\big\{\hat{\gamma}^{q}\hat{Q}_{tt}(\mathbf{r})\big\}
\Tr\big\{\hat{\gamma}^{cl}\hat{Q}_{t\,'t\,'}(\mathbf{r}')\big\}\right\rangle_{Q}\,,
}
\end{equation}
ここで，$x=(\mathbf{r},t)$であり，山括弧は作用~\eqref{NLSM-action-noninteracting}による平均を表す．式~\eqref{Part-II-Pi}の右辺の第1項は，作用~\eqref{NLSM-action-general}の$\Tr\big\{\hat{V}\hat{\sigma}_{x}\hat{V}\big\}$の部分を微分することで生じ，第2項は$\Tr\big\{\hat{V}\hat{Q}\big\}$の微分から生じる．式~\eqref{Part-II-Pi}は，線形密度応答のKubo公式に対応する$\sigma$模型での表現である．



Fourier表示では，直前の式は次の形をとる：
\begin{equation}\label{Part-II-Pi-Fourier}
\Pi^{R}\qo=\nu+\frac{i}{2}(\pi\nu)^{2}\iint
\frac{\d\epsilon\d\epsilon'}{4\pi^{2}}
\left\langle\Tr\big\{\hat{\gamma}^{q}\hat{Q}_{\epsilon+\omega,\epsilon}(\mathbf{q})\big\}
\Tr\{\hat{\gamma}^{cl}\hat{Q}_{\epsilon',\epsilon'+\omega}(-\mathbf{q})\}
\right\rangle_{Q}\,.
\end{equation}
式~\eqref{NLSM-Q-U-W-parametrization}と\eqref{NLSM-U-W}を用いると，拡散揺らぎの1次で次式を得る（0次の寄与は$\nu$のみであり，実際$\Tr\{\hat{\gamma}^{cl}\hat{\Lambda}\}=0$である）：
\begin{eqnarray}
&&\Tr\big\{\hat{\gamma}^{cl}\hat{Q}_{\epsilon',\epsilon'+\omega}(-\mathbf{q})\big\}=
\bar{d}_{\epsilon',\epsilon'+\omega}(-\mathbf{q})
(F_{\epsilon'+\omega}-F_{\epsilon'})\,,\nonumber\\
&&\Tr\big\{\hat{\gamma}^{q}\hat{Q}_{\epsilon+\omega,\epsilon}(\mathbf{q})\big\}=
\bar{d}_{\epsilon+\omega,\epsilon}(\mathbf{q})
(1-F_{\epsilon}F_{\epsilon+\omega})-d_{\epsilon+\omega,\epsilon}(\mathbf{q})\,.
\end{eqnarray}
$\langle\bar{d}\bar{d}\rangle_{\mathcal{W}}\equiv0$なので，直前の式の最後の項だけが式~(\ref{Part-II-Pi-Fourier})の平均に寄与する．結果は
\begin{equation}\label{Part-II-Pi-diffusive}
\fitdisplay{
\Pi^{R}\qo=\nu+\frac{i\pi\nu^{2}}{4}
\int^{+\infty}_{-\infty}\d\epsilon\,\big(F_{\epsilon}-F_{\epsilon+\omega}\big)
\left\langle d_{\epsilon+\omega,\epsilon}(\mathbf{q})
\bar{d}_{\epsilon,\epsilon+\omega}(-\mathbf{q})\right\rangle_{\mathcal{W}}
=\nu\left[1+\frac{i\omega}{Dq^{2}-i\omega}\right]=\frac{\nu
Dq^{2}}{Dq^{2}-i\omega},
}
\end{equation}
となる．ここで，ディフューゾンの伝播関数~\eqref{NLSM-d-d}と積分$\int\d\epsilon\,(F_{\epsilon}-F_{\epsilon+\omega})=-2\omega$を用いた．$\Pi^{R}(0,\omega)=0$という事実は，粒子数保存の帰結である．こうして，密度・密度応答関数の拡散形が得られた．また，この関数は要請どおり，実際に遅延関数（複素$\omega$の上半平面で解析的）であることにも注意しよう．電流・電流応答関数$K^{R}\qo$も同様に求められる．しかし，より直接的には，連続の式$\mathbf{q}\cdot\mathbf{j}+\omega\varrho=0$を用いればよい．これから密度応答関数と電流応答関数の間に$K^{R}\qo=\omega^{2}\Pi^{R}\qo/q^{2}$という関係が従う．その結果，伝導率は次式で与えられる：
\begin{equation}
\sigma\qo=\frac{e^{2}}{i\omega}\,K^{R}\qo =
e^{2}\, \frac{-i\omega}{q^2}\, \Pi^R\qo
= e^{2}\nu D\, \frac{-i\omega}{Dq^{2}-i\omega}\, ,
\end{equation}
これは一様極限$\mathbf{q}\to0$で，Drudeの結果$\sigma_{D}\equiv\sigma(0,\omega)=e^{2}\nu D$に帰着する．


%$$$$$$$$$$$$$$$$$$$$$$$$$$$$$$$$$$$$$$$$$$$$$$$$$$$$$$$$$$$$$$$$$$$$$$$$$$$$$$$$$$$$$$$$$$$$$$$$$$$$$$$$$$$

\subsubsection{スペクトル統計}\label{app_Part-II-2}



$L\gg
l$を満たす，サイズ$L$の乱れた金属片を考えよう．ここで，$l\equiv v_F\tauel$は弾性平均自由行程である．Schr\"odinger方程式のスペクトルは，離散的な準位$\epsilon_n$からなり，{\em 試料固有の}状態密度（DOS）$\nu(\epsilon)=
\sum_n\delta(\epsilon-\epsilon_n)$によって特徴づけられる．この量は大きく揺らぎ，通常は解析的に計算できないし，計算する必要もない．乱れの実現について平均すると，平均状態密度$\langle\nu(\epsilon)\rangle_{\mathrm{dis}}$を得る．後者はFermiエネルギーのスケールで滑らかなエネルギーの関数なので，定数$\langle\nu(\epsilon_F)\rangle_{\mathrm{dis}}\equiv\nu$とみなせる．これこそ，前節までで用いていた状態密度である．



試料固有の状態密度$\nu(\epsilon)$の揺らぎをどう捉えられるか，特に，あるエネルギー$\epsilon$でのスペクトルが別のエネルギー$\epsilon'$での自身とどのように相関しているかを知りたくなるだろう．この問いに答えるには，スペクトル相関関数
\begin{equation}\label{Part-I-R-def}
R(\epsilon,\epsilon')\equiv\langle\nu(\epsilon)\nu(\epsilon')\rangle_{\mathrm{dis}}
-\nu^2\, .
\end{equation}
を計算すればよい．この関数は，AltshulerとShklovskiiの先駆的な論文~\cite{AltshulerShklovskii}で計算された．ここでは，Keldysh NLSMを用いて導出する．



状態密度は
$$\nu(\epsilon) =
i\sum_{k}(G^R(\mathbf{k},\epsilon)-
G^A(\mathbf{k},\epsilon))/(2\pi)=(\langle \psi_1\bar\psi_1\rangle
-\langle \psi_2\bar\psi_2\rangle)/(2\pi) =
-\big\langle\vec{\bar{\Psi}}\hat{\sigma}_z\vec{\Psi}\big\rangle/(2\pi)\,,$$
と定義される．ここで，山括弧は（乱れの平均ではなく）量子平均を表し，添字はKeldysh空間のものである．任意のエネルギーで状態密度を生成するには，フェルミ粒子の作用~\eqref{NLSM-action-noninteracting}にソース項
$$iS_{\mathrm{DOS}}=-\int\frac{\d\epsilon}{2\pi} J_\epsilon\int
\d\mathbf{r}\,\vec{\bar{\Psi}}(\epsilon,\mathbf{r})\hat{\sigma}_z
\vec{\Psi}(\epsilon,\mathbf{r})=-\iint\d t\d t'\int\d\mathbf{r}\,
\vec{\bar{\Psi}}\rt
J_{t-t'}\hat{\sigma}_z\vec{\Psi}(\mathbf{r},t')\,,$$
を加える．乱れについて平均し，$\hat Q$行列表示に移ると，状態密度のソース項は
$$iS_{\mathrm{DOS}}=\pi\nu\int\frac{\d\epsilon}{2\pi}
J_\epsilon\int\d\mathbf{r}\,\Tr\{\hat{Q}_{\epsilon\epsilon}(\mathbf{r})\hat
\sigma_z\}\,.$$
に変換される．すると，状態密度は$\nu(\epsilon)=\delta\mathcal{Z}[J]/\delta J_\epsilon$によって生成される．このとき，$\langle\nu(\epsilon)\rangle_{\mathrm{dis}}={1\over 2}\nu
\langle \Tr\{\hat{Q}_{\epsilon\epsilon}\hat{\sigma}_z\}\rangle_Q$であることは明らかである．$\hat{Q}_{\epsilon\epsilon}=\hat{\Lambda}_\epsilon$を代入すると，当然期待されるように，$\langle\nu(\epsilon)\rangle_{\mathrm{dis}}=\nu$を得る．また，$\hat{\Lambda}$のまわりの揺らぎが結果を変えないことも容易に確かめられる（因果性の拘束によって，揺らぎの図式はすべて打ち消し合う）．これで，相関関数~\eqref{Part-I-R-def}を計算する準備が整った：
\begin{equation}\label{Part-I-R}
R(\epsilon,\epsilon')\equiv\frac{\delta^{2}\mathcal{Z}[J]}{\delta
J_\epsilon \delta J_{\epsilon'}} -\nu^2=\nu^{2}\left[\frac{1}{4}
\left\langle \Tr\{\hat{Q}_{\epsilon\epsilon}\hat{\sigma}_{z}\}\Tr\{
\hat{Q}_{\epsilon'\epsilon'}\hat{\sigma}_{z}\}\right\rangle_{Q}-1\right]\,.
\end{equation}
式~\eqref{NLSM-Q-U-W-parametrization}のパラメータ表示を用いると，拡散揺らぎ$\hat{\mathcal{W}}$の2次までで，次式を得る：
\begin{equation}
\Tr\big\{\hat{Q}\hat{\sigma}_{z}\big\}=\frac{1}{2}\,\big[4-2\,F\circ\bar{d}-2\,\bar{d}\circ
F+d\circ\bar{d}+\bar{d}\circ d\big]\,.
\end{equation}
$\langle\bar{d}\bar{d}\rangle_{\mathcal{W}}=0$なので，式~\eqref{Part-I-R}に寄与する非ゼロの項は，$d$と$\bar{d}$をまったく含まない項（これは$\nu^{2}$項を打ち消す）と，$\langle d\bar{d}d\bar{d}\rangle_{\mathcal{W}}$型の項だけである．後者を集めると，次式を得る：
\begin{equation}\label{Part-I-R-dd}
R(\epsilon,\epsilon')=\frac{\nu^{2}}{16}\int\d\mathbf{r}
\iint\frac{\d\epsilon_{1}\d\epsilon_{2}}{(2\pi)^{2}}
\,\left\langle\left(d_{\epsilon\epsilon_{1}}\bar{d}_{\epsilon_{1}\epsilon}+
\bar{d}_{\epsilon\epsilon_{1}}d_{\epsilon_{1}\epsilon}\right)
\left(d_{\epsilon'\epsilon_{2}}\bar{d}_{\epsilon_{2}\epsilon'}+
\bar{d}_{\epsilon'\epsilon_{2}}d_{\epsilon_{2}\epsilon'}\right)\right\rangle_{\mathcal{W}}\,.
\end{equation}
ここで，式~\eqref{NLSM-d-d}から従う相関関数$\langle
d_{\epsilon\epsilon'}\bar{d}_{\epsilon'\epsilon}\rangle_{\mathcal{W}}
\propto\mathcal{D}^{R}(\epsilon-\epsilon')$を用いてWick収縮を行う必要がある．また，$\epsilon_{1}$の上（下）半平面全体で解析的な関数の積分であるために，$\int\d\epsilon_{1}[\mathcal{D}^{R(A)}(\mathbf{q},\epsilon-\epsilon_{1})]^{2}=0$となることも考慮する．その結果，
\begin{equation}
R(\epsilon,\epsilon')=\frac{1}{4\pi^{2}}\sum_{q}
\left[\big(\mathcal{D}^{R}(\mathbf{q},\epsilon-\epsilon')\big)^{2}+
\big(\mathcal{D}^{A}(\mathbf{q},\epsilon-\epsilon')\big)^{2}\right]\,,
\end{equation}
を得る．ここで，運動量の和は，金属の境界で電流がゼロ（微分がゼロ）という条件を満たす拡散演算子$D\partial^{2}_{\mathbf{r}}$の離散モードについての和を表す．これはユニタリー対称性クラスに対するAltshulerとShklovskiiの結果~\cite{AltshulerShklovskii}である．相関関数$R(\epsilon,\epsilon')$はエネルギー差$\omega=\epsilon-\epsilon'$のみに依存することに注意しよう．$R(\epsilon,\epsilon')$の図式表現を図~\ref{Fig-AS-DOS}に示す．拡散伝播関数の具体的な形を用いると，スペクトル相関関数は次の形となる：
\begin{equation}
R(\epsilon-\epsilon')=-\frac{1}{2\pi^{2}}\,\mathrm{Re}\sum_{n}\frac{1}
{\big(\epsilon-\epsilon'+iDq^{2}_{n}\big)^{2}}\,,
\end{equation}
ここで，$q^{2}_{n}=\sum_{\mu}\pi^{2}n^{2}_{\mu}/L^{2}_{\mu}$，$\mu=x,y,z\,$，$n_{\mu}=0, 1, 2\ldots$であり，$L_{\mu}$はメゾスコピック試料の各方向の寸法である．
\begin{figure}[h!]
\begin{center}\includegraphics[width=8cm]{figures/fig11.pdf}\end{center}
\caption{状態密度のメゾスコピック揺らぎ$R(\epsilon,\epsilon')$を計算する図式（式~\eqref{Part-I-R}参照）．これはWick収縮$ \langle
d_{\epsilon\epsilon_{1}}\bar{d}_{\epsilon_{1}\epsilon}
\bar{d}_{\epsilon'\epsilon_{2}}d_{\epsilon_{2}\epsilon'}\rangle_{\mathcal{W}}\to
\langle d_{\epsilon\epsilon_{1}}
\bar{d}_{\epsilon'\epsilon_{2}}\rangle_{\mathcal{W}}
\langle\bar{d}_{\epsilon_{1}\epsilon}d_{\epsilon_{2}\epsilon'}
\rangle_{\mathcal{W}}
\propto[\mathcal{D}^{R}(\mathbf{q},\epsilon-\epsilon')]^{2}
\delta_{\epsilon_1\epsilon'}\delta_{\epsilon_2\epsilon}$から生成される（式~\eqref{Part-I-R-dd}参照）．先進ディフューゾンを含む同様の図式も存在する．\label{Fig-AS-DOS}}
\end{figure}



エネルギー差が小さい$\omega\ll E_{Th}=D/L^{2}$の場合，拡散演算子の最も低い一様モード$q_{n}=0$（いわゆるゼロモード）のみを残せばよく，したがって$R(\omega)=
-1/(2\pi^2\omega^2)$となる．これはランダム行列理論の普遍的な結果である．相関関数が負であることは，エネルギー準位が互いに小さな間隔$\omega$で見つかる可能性が{\em 低い}ことを意味する．これはエネルギー準位反発の現れである．相関がエネルギー間隔の2乗に反比例して，非常にゆっくり減衰することに注意しよう．一方，ランダム行列理論の結果~\cite{Mehta}
\begin{equation}
R_{RMT}(\omega)=-\frac{1-\cos(2\pi\omega/\delta)}{2\pi^2\omega^2}\,,
\end{equation}
には，エネルギー差の振動関数も含まれる．ここで，$\delta$は平均準位間隔である．これらの振動は，もとのエネルギースペクトルの離散性を反映している．$\hat{\Lambda}$という「点」の近傍の小さな揺らぎを扱う摂動論では，これらを求めることはできない．しかし，追加の停留点（因果的構造をもたないもの）を考慮すれば，再現できる~\cite{AltlandKamenev}．鞍点法と摂動論は，$\omega\gg
\delta$である限り有効である．現時点では，$\omega\lesssim \delta$でKeldysh NLSMをどのように扱うかは知られていない．



%$$$$$$$$$$$$$$$$$$$$$$$$$$$$$$$$$$$$$$$$$$$$$$$$$$$$$$$$$$$$$$$$$$$$$$$$$$$$$$$$$$$$$$$$$$$$$$$$$$$$$$$$$$$

\subsubsection{普遍的コンダクタンス揺らぎ}\label{Part-II-3}



前節の議論と同様に，サイズ$L$が電子の位相コヒーレンス長と同程度，すなわち$L\sim L_{\varphi}$である，小さな金属試料のアンサンブルを考える．各試料の乱れのポテンシャルの具体的な実現が異なるため，コンダクタンスは試料ごとに揺らぐ．これらの再現性のある揺らぎは，{\em 普遍的コンダクタンス揺らぎ}（UCF）として知られている．UCFの理論研究は，Altshuler~\cite{Altshuler-UCF}，およびLeeとStone~\cite{LeeStone}によって始められた．より詳細な研究は，その後，文献~\cite{AltshulerShklovskii,LeeStoneFukuyama}で行われた．本節で展開する，Keldysh非線形$\sigma$模型の枠組みでUCFを扱う手法は，Altshuler，Kravtsov，Lerner~\cite{AKL-UCF}がゼロ温度のレプリカ非線形σ模型の枠組みで展開した理論と，密接に対応している．



出発点は，線形応答における$\mathrm{dc}$伝導率の次の表式である：
\begin{equation}\label{Part-II-sigma-def}
\sigma_{\mu\nu}=-\frac{e^2}{2}\lim_{\Omega\to 0}\frac{1}{\Omega}
\left(\frac{\delta^{2}\mathcal{Z}[\mathbf{A}^{cl},\mathbf{A}^{q}]}
{\delta
A^{cl}_{\nu}(\Omega)A^{q}_{\mu}(-\Omega)}\right)_{\mathbf{A}^{cl}=\mathbf{A}^{q}=0}\,,
\end{equation}
ここで，添字$\mu,\nu$は空間の直交座標を表す．式~\eqref{NLSM-covariant-deriv}を用いて作用~\eqref{NLSM-action-general}をベクトルポテンシャルの2次まで展開すると，生成関数の対応する項は$\mathcal{Z}[\mathbf{A}^{cl},\mathbf{A}^{q}]=\frac{\pi\nu
D}{2}\big\langle\Tr\big\{\hat{\mathbf{A}}
\hat{Q}\hat{\mathbf{A}}\hat{Q}\big\}\big\rangle_{Q}$と書かれる．鞍点$\hat{Q}=\hat{\Lambda}$において，式~\eqref{Part-II-sigma-def}のベクトルポテンシャルに関する微分を順次行うと，平均伝導率について次式を得る：
\begin{equation}\label{NLSM-sigma-Drude}
\langle\sigma_{\mu\nu}\rangle_{\mathrm{dis}}=
\delta_{\mu\nu}\lim_{\Omega\to 0}\frac{\pi\sigma_{D}}{4\Omega}
\Tr\left\{\hat{\gamma}^{cl}\hat{\Lambda}_{\epsilon+\Omega}
\hat{\gamma}^{q}\hat{\Lambda}_{\epsilon-\Omega}\right\}=
\delta_{\mu\nu}\frac{\pi\sigma_{D}}{2}\lim_{\Omega\to0}
\frac{1}{\Omega}\int\frac{\d\epsilon}{2\pi}
\big(F_{\epsilon+\Omega}-F_{\epsilon-\Omega}\big)=\sigma_{D}\delta_{\mu\nu}
\end{equation}
ここで，当然のことながら$\sigma_D=e^2\nu D$である．この段階で，鞍点$\hat{\Lambda}$のまわりの$\hat Q$行列の揺らぎ$\hat{\mathcal{W}}$を保持すれば，平均伝導率に対する弱局在補正~\cite{Hikami,Gor'kovLarkinKhmelnitskii,AAKL,LeeRamakrishnan,ALW-book}を計算できる．以下では，次のように定義される伝導率揺らぎの既約相関関数の計算を考える：
$$\langle\delta\sigma_{\mu_{1}\nu_{1}}
\delta\sigma_{\mu_{2}\nu_{2}}\rangle_{\mathrm{dis}}=
\left\langle\big(\sigma_{\mu_{1}\nu_{1}}-\langle\sigma_{\mu_{1}\nu_{1}}\rangle\big)
\big(\sigma_{\mu_{2}\nu_{2}}-
\langle\sigma_{\mu_{2}\nu_{2}}\rangle\big)\right\rangle_{\mathrm{dis}}\,.$$
式~\eqref{Part-II-sigma-def}より，この既約相関関数は，$\hat Q$行列を用いて次のように表せる：
\begin{eqnarray}\label{Part-II-UCF-def}
\langle\delta\sigma_{\mu_{1}\nu_{1}}
\delta\sigma_{\mu_{2}\nu_{2}}\rangle_{\mathrm{dis}}=
\left(\frac{\pi\sigma_{D}}{4}\right)^{2}
\prod^{2}_{i=1}\left(\lim_{\Omega_{i}\to 0}
\frac{1}{\Omega_{i}}\frac{\delta^{2}}{\delta
A^{cl}_{\nu_{i}}(\Omega_{i})\delta
A^{q}_{\mu_{i}}(-\Omega_{i})}\right)
\left\langle\Tr\{\hat{\mathbf{A}}\hat{Q}\hat{\mathbf{A}}\hat{Q}\}
\Tr\{\hat{\mathbf{A}}\hat{Q}\hat{\mathbf{A}}\hat{Q}\}\right\rangle_{Q}\nonumber\\
-\sigma^{2}_{D} \delta_{\mu_{1}\nu_{1}}\delta_{\mu_{2}\nu_{2}}\,,
\end{eqnarray}
ここでは，式~\eqref{NLSM-action-general}を用い，$\exp(iS[\hat{Q},\mathbf{A}])$をベクトルポテンシャルの4次まで展開した．次に，$\hat Q$行列の揺らぎを，生成子$\hat{\mathcal{W}}$の2次まで考慮しなければならない．ここには2つの可能性がある．式~\eqref{Part-II-UCF-def}の右辺の各トレースの中で，各$\hat Q$行列を$\hat{\mathcal{W}}$の1次まで展開すると，$\mathcal{T}_{1}[\hat{\mathcal{W}}]=\Tr\{\hat{\mathbb{A}}\hat{\sigma}_{z}
\hat{\mathcal{W}}\hat{\mathbb{A}}\hat{\sigma}_{z}\hat{\mathcal{W}}\}$を得る．あるいは，一方の$\hat Q$行列を$\hat{\Lambda}$とし，他方を2次まで展開すると，$\mathcal{T}_{2}[\hat{\mathcal{W}}]=
\Tr\{\hat{\mathbb{A}}\hat{\sigma}_{z}\hat{\mathbb{A}}\hat{\sigma}_{z}
\hat{\mathcal{W}}^{2}\}$を得る．ここで，$\hat{\mathbb{A}}=\hat{\mathcal{U}}^{-1}
\hat{\mathbf{A}}\hat{\mathcal{U}}$である．その結果，式~\eqref{Part-II-UCF-def}は次の形をとる：
\begin{eqnarray}\label{Part-II-UCF-W}
&&\langle\delta\sigma_{\mu_{1}\nu_{1}}
\delta\sigma_{\mu_{2}\nu_{2}}\rangle_{\mathrm{dis}}=
\left(\frac{\pi\sigma_{D}}{4}\right)^{2}
\prod^{2}_{i=1}\left(\lim_{\Omega_{i}\to 0}
\frac{1}{\Omega_{i}}\frac{\delta^{2}}{\delta
A^{cl}_{\nu_{i}}(\Omega_{i})\delta
A^{q}_{\mu_{i}}(-\Omega_{i})}\right) \nonumber \\
&&\left[\big\langle\mathcal{T}_{1}[\hat{\mathcal{W}}]
\mathcal{T}_{1}[\hat{\mathcal{W}}]\big\rangle_{\mathcal{W}}+
\big\langle\mathcal{T}_{2}[\hat{\mathcal{W}}]
\mathcal{T}_{2}[\hat{\mathcal{W}}]\big\rangle_{\mathcal{W}}\right]-
\sigma^{2}_{D} \delta_{\mu_{1}\nu_{1}}\delta_{\mu_{2}\nu_{2}}\,.
\end{eqnarray}
ここでの各平均は，図式で表すと便利である（図~\ref{Fig-UCF}参照）．図~\ref{Fig-UCF}aの菱形は$\mathcal{T}_{1}[\hat{\mathcal{W}}]$を含む項に対応し，向かい合う頂点は行列$\hat{\mathbb{A}}$を表す．一方，図~\ref{Fig-UCF}bの隣接する頂点をもつ長方形は，$\mathcal{T}_{2}[\hat{\mathcal{W}}]$を含む項に対応する．頂点は，場$\hat{\mathcal{W}}$のディフューゾン伝播関数によって結ばれる．式~\eqref{Part-II-UCF-W}には交差寄与$2\langle\mathcal{T}_{1}[\hat{\mathcal{W}}]\mathcal{T}_{2}[\hat{\mathcal{W}}]
\rangle_{\mathcal{W}}$も含まれるはずだが，これは$\hat{\mathcal{W}}$について平均すると消える．
\begin{figure}
\begin{center}\includegraphics[width=10cm]{figures/fig12.pdf}\end{center}
\caption{コンダクタンス揺らぎの分散に対する図式．\label{Fig-UCF}}
\end{figure}



式~\eqref{Part-II-UCF-W}の各項を個別に微分し，行列を掛け合わせ，式~\eqref{NLSM-d-d}のディフューゾン伝播関数を用いると，式~\eqref{Part-II-UCF-W}に対して次式を得る：
\begin{eqnarray}\label{Part-II-UCF-AS}
&& \langle\delta\sigma_{\mu_{1}\nu_{1}}
\delta\sigma_{\mu_{2}\nu_{2}}\rangle_{\mathrm{dis}}=
\left(\frac{4\sigma_{D}}{\pi\nu}\right)^{2}\!\!
\iint^{+\infty}_{-\infty}\frac{\d\epsilon_{1}\d\epsilon_{2}}
{\big[2T\cosh(\epsilon_{1}/2T)\cosh(\epsilon_{2}/2T)\big]^{2}}\nonumber
\\
&&\sum_{q}\left[
|\mathcal{D}^{R}(\mathbf{q},\epsilon_{1}-\epsilon_{2})|^{2}
\big(\delta_{\mu_{1}\mu_{2}}\delta_{\nu_{1}\nu_{2}}+
\delta_{\mu_{1}\nu_{2}}\delta_{\nu_{1}\mu_{2}}\big)+\mathrm{Re}\,
\big[\mathcal{D}^{R}(\mathbf{q},\epsilon_{1}-\epsilon_{2})\big]^{2}
\delta_{\mu_{1}\nu_{1}}\delta_{\mu_{2}\nu_{2}}\right]\,.
\end{eqnarray}
式~\eqref{Part-II-UCF-AS}の角括弧内の第1項は図~\ref{Fig-UCF}aに，第2項は図~\ref{Fig-UCF}bに対応する．$\epsilon_{1}-\epsilon_{2}=\omega$と$\epsilon_{1}+\epsilon_{2}=2\epsilon$を導入し，$\epsilon$について積分すると，式~\eqref{Part-II-UCF-AS}は次の形に書き換えられる：
\begin{equation}
\langle\delta\sigma_{\mu_{1}\nu_{1}}
\delta\sigma_{\mu_{2}\nu_{2}}\rangle_{\mathrm{dis}}=
\sigma^{2}_{1}\big(\delta_{\mu_{1}\mu_{2}}\delta_{\nu_{1}\nu_{2}}+
\delta_{\mu_{1}\nu_{2}}\delta_{\nu_{1}\mu_{2}}\big)+
\sigma^{2}_{2}\,\delta_{\mu_{1}\nu_{1}}\delta_{\mu_{2}\nu_{2}},
\end{equation}
ここで，
\begin{subequations}
\begin{equation}
\sigma^{2}_{1}=\left(\frac{4\sigma_{D}}{\pi\nu}\right)^{2}\int^{+\infty}_{-\infty}\
\frac{\d\omega}{2T}\
\mathcal{F}\left(\frac{\omega}{2T}\right)\sum_{q}
\frac{1}{\big(Dq^{2}\big)^{2}+\omega^{2}}\,,
\end{equation}
\begin{equation}
\sigma^{2}_{2}=\left(\frac{4\sigma_{D}}{\pi\nu}\right)^{2}\int^{+\infty}_{-\infty}\
\frac{\d\omega}{2T}\
\mathcal{F}\left(\frac{\omega}{2T}\right)\mathrm{Re}\sum_{q}
\frac{1}{\big(Dq^{2}-i\omega\big)^{2}}\,,
\end{equation}
\end{subequations}
であり，無次元関数は$
\mathcal{F}(x)=[x\coth(x)-1]/\sinh^{2}(x)$で与えられる．ここで，$\sigma^{2}_{1}$は拡散係数のメゾスコピック揺らぎからの寄与（図~\ref{Fig-UCF}a），$\sigma^{2}_{2}$は状態密度の揺らぎからの対応する寄与（図~\ref{Fig-UCF}b）とみなせる．$\langle\mathcal{T}_{1}[\hat{\mathcal{W}}]\mathcal{T}_{2}[\hat{\mathcal{W}}]
\rangle_{\mathcal{W}} = 0$という事実は，拡散係数と状態密度のメゾスコピック揺らぎが統計的に独立であることを意味する．一般に，$\sigma^{2}_{1}$と$\sigma^{2}_{2}$の寄与は異なる．しかし，ゼロ温度$\omega\to 0$では両者は等しくなり，その結果，
\begin{equation}
\langle\delta\sigma_{\mu_{1}\nu_{1}}
\delta\sigma_{\mu_{2}\nu_{2}}\rangle=c_{d}
\left(\frac{e^{2}}{2\pi\hbar}\right)^{2}
\big(\delta_{\mu_{1}\mu_{2}}\delta_{\nu_{1}\nu_{2}}+
\delta_{\mu_{1}\nu_{2}}\delta_{\nu_{1}\mu_{2}}+
\delta_{\mu_{1}\nu_{1}}\delta_{\mu_{2}\nu_{2}}\big)\,,
\end{equation}
となる．ここで，$c_{d}=(4/\pi)^2\sum_{n_{\mu}}(\pi n_{\mu}n_{\mu})^{-2}$は次元と形状に依存する係数である（最終的な答えではPlanck定数を復元したことに注意）．この式はコンダクタンス揺らぎの普遍性を反映し，当然ながら，もともと不純物図式法によって得られた結果~\cite{AltshulerShklovskii,LeeStone}と一致する．総説については文献~\cite{ALW-book,Beenakker}を参照されたい．


%$$$$$$$$$$$$$$$$$$$$$$$$$$$$$$$$$$$$$$$$$$$$$$$$$$$$$$$$$$$$$$$$$$$$$$$$$$$$$$$$$$$$$$$$$$$$$$$$$$$$$$$$$$$

\subsubsection{完全計数統計}\label{Part-II-4}



導体に電流$I(t)$が流れるとき，通常，電流は平均値$\langle I\rangle$のまわりで揺らぐ．電流揺らぎの2次，あるいはさらに高次のモーメントを計算したいことが多い．この種の例は第~\ref{app_Part-I-2}節ですでに扱った．注目すべきことに，場合によっては，揺らぐ電流の特定のモーメントを計算するだけでなく，電流揺らぎの分布関数全体を復元できる．Keldysh法を用いた電子輸送の完全計数統計（FCS）への理論的アプローチは，Levitov，Lesovikと共同研究者によって切り開かれた~\cite{LevitovLesovik-1,LevitovLesovik-2,LevitovLee}．以下では，Nazarov~\cite{Nazarov-FCS}によって展開された，拡散的電子輸送への応用を考える．



外部から印加した電圧$V$によって化学ポテンシャルがずれた，2つのリザーバーを考える．両リザーバーは，長さ$L$の拡散的な準一次元細線で接続されているとする．細線のコンダクタンスは$\mathrm{g}_{D}=\sigma_{D}\mathcal{A}/L$であり，$\mathcal{A}$は細線の断面積である．細線を通る拡散的電子輸送を記述するには，乱れについて平均した生成関数$\mathcal{Z}[\chi]=\int\D[\hat Q]\exp(iS[\hat Q,A_{\chi}])$から出発する．作用は式~\eqref{NLSM-action-general}で与えられ，補助ベクトルポテンシャル$\hat{A}_{\chi}$は共変微分~\eqref{NLSM-covariant-deriv}を通じて現れる．$\hat{A}_{\chi}$は古典成分をもたない純粋な量子成分として，次のように選ぶ：
\begin{equation}\label{Part-II-FCS-A}
\hat{A}_{\chi}(t)=\frac{\hat{\gamma}^{q}}{2L}\left\{\begin{array}{cc}\chi\quad
& 0<t<t_{0}
\\ 0\quad & \text{それ以外}\end{array}\right.\,.
\end{equation}
ここで，量子Keldysh行列$\hat{\gamma}^{q}$は式~\eqref{fermion-gammas}で与えられ，$\chi$は{\em 計数場}と呼ばれる．作用$S[\hat{Q},A_{\chi}]$には，細線の両端での$\hat Q(x)$行列の境界条件
\begin{eqnarray}\label{Part-II-Q-boundary}
\hat Q(0)=\left(\begin{array}{cc} 1\quad & 2F_{\epsilon}
\\ 0\quad & -1
\end{array}\right)\,, \qquad
\hat Q(L)=
\left(\begin{array}{cc} 1\quad & 2F_{\epsilon-eV}
\\ 0\quad & -1
\end{array}\right)\,.
\end{eqnarray}
が付随する．$\mathcal{Z}[\chi]$がわかれば，計数場$\chi$について$\mathcal{Z}[\chi]$を微分することで，測定時間$t_0$の間にリザーバー間を移動した電子数の任意のモーメント$\langle
q^n\rangle$を求められる．既約相関関数は$\mathcal{C}_{1}=\langle q\rangle=q_0$および$\mathcal{C}_{n}=\langle
(q-q_{0})^n\rangle$（$n=2,3,\ldots$）と定義される．ここで，$q=\frac{1}{e}\int_0^{t_0}I(t)dt$および$q_0=t_{0}\mathrm{g}_{D}V/e=t_{0}\langle I\rangle/e$であり，$\mathrm{g}_D$は平均の拡散コンダクタンスである．これらは，$\mathcal{Z}[\chi]$の対数を計数場のべきで展開することで求まる：
\begin{equation}
\label{Part-II-log-Z-chi}
\ln\mathcal{Z}[\chi]=\sum_{n=0}^{\infty}\frac{(i\chi)^n}{n!}\mathcal{C}_n\,.
\end{equation}


$S[\hat Q,A_{\chi}]$を極値にする鞍点$\underline{\hat{Q}}=\hat{\Lambda}_{\chi}$で作用を評価することにより，$\mathcal{Z}[\chi]$を計算する．難点は，作用$S[\hat Q,A_{\chi}]$が計数場$\chi$に明示的に依存し，対応する鞍点方程式の解が任意の$A_{\chi}$に対しては知られていないことである．この障害は，ベクトルポテンシャル~\eqref{Part-II-FCS-A}が純粋ゲージであり，次の変換によって作用から$S[\hat
Q,A_{\chi}]\to S[\hat Q_{\chi}]$とゲージ消去できることに気づけば克服できる：
\begin{equation}
\hat{Q}(x\,;t,t')=\exp\big\{ix\hat{A}_{\chi}(t)\big\}\,
\hat{Q}_{\chi}(x\,;t,t')\,\exp\big\{-ix\hat{A}_{\chi}(t')\big\}\,.
\end{equation}
ただし，代わりに境界条件~\eqref{Part-II-Q-boundary}が，それに応じて次のように変わる：
\begin{equation}\label{Part-II-Q-boundary-twisted}
\hat{Q}_{\chi}(0)=\hat{Q}(0), \qquad
\hat{Q}_{\chi}(L)=\exp\big(-i\chi\hat{\gamma}^{q}/2\big)\hat{Q}(L)
\exp\big(i\chi\hat{\gamma}^{q}/2\big)\,.
\end{equation}
この変換の利点は，$\hat{Q}_{\chi}$に対する鞍点方程式，すなわちUsadel方程式~\eqref{NLSM-Usadel}
\begin{equation}
D\frac{\partial}{\partial
x}\left(\hat{Q}_{\chi}\circ\frac{\partial\hat{Q}_{\chi}}{\partial
x}\right)=0\,,
\end{equation}
を，今度は明示的に解けることである．そのため，$\hat{Q}_{\chi}\circ\partial_x\hat{Q}_{\chi}=-
\partial_x\hat{Q}_{\chi}\circ \hat{Q}_{\chi} =\hat{J}$が定数行列，すなわち$x$に依存しない行列であることに注意しよう．$\hat{Q}_{\chi}^2=\hat{1}$であるから，$\hat J$は$\hat{Q}_{\chi}$と反可換であり，すなわち$\hat{Q}_{\chi}\circ\hat{J}+\hat{J}\circ\hat{Q}_{\chi}=0$を満たす．その結果，$\hat{Q}_{\chi}(x)=\hat{Q}_{\chi}(0)\exp\big(\hat{J}x\big)$を得る．$x=L$とおき，左から$\hat{Q}_{\chi}(0)$を掛けると，まだ未知である行列$\hat J$を境界条件~\eqref{Part-II-Q-boundary-twisted}で表せる：$\hat{J}=\frac{1}{L}\ln\big[\hat{Q}_{\chi}(0)\hat{Q}_{\chi}(L)\big]$．



計数場$\chi$の任意の選択に対して$\hat{Q}_{\chi}$行列の鞍点配置が決まったので，これを作用$S[\hat
Q_{\chi}]$に戻して，生成関数
$$\ln\mathcal{Z}[\chi]=iS[\hat{Q}_{\chi}]=-\frac{\pi\nu D}{4}\,
\Tr\{(\partial_x\hat{Q}_{\chi})^{2}\}= \frac{\pi \nu D}{4}\,
\Tr\{\hat{J}^2\}\,,$$
を得る．ここでは，反交換関係$\{\hat{Q}_{\chi}(0),\hat{J}\}=0$を用いた．時間積分を計算する際に，Wigner変換$\iint\d t\d t'\to
t_{0}\int\frac{\d\epsilon}{2\pi}$に移る．ここで，$t_{0}$は中心時刻についての積分から生じる．すると，次式を得る：
\begin{equation}\label{Part-II-FCS-Z}
\ln\mathcal{Z}[\chi]=\frac{t_{0}\mathrm{g}_{D}}{8e^{2}}\int\d\epsilon\
\mathrm{Tr}\ln^{2}\left[\hat{Q}(0)\exp\big(-i\chi\hat{\gamma}^{q}/2\big)
\hat{Q}(L)\exp\big(i\chi\hat{\gamma}^{q}/2\big)\right]\,.
\end{equation}


以下では，式~\eqref{Part-II-FCS-Z}をゼロ温度極限$T=0$で解析する．このとき，$F_{\epsilon}=\tanh(\epsilon/2T)\to \mbox{sign}(\epsilon)$である．次の行列を用いて回転$\hat{\mathcal{Q}}=\hat{\mathcal{O}}^{-1}\hat{Q}\hat{\mathcal{O}}$を行えば，代数計算をさらに大幅に短縮できる：
\begin{equation}
\hat{\mathcal{O}}=\frac{1}{\sqrt{2}}\left(\begin{array}{cr}1 & -1
\\ 1 & 1\end{array}\right)\,.
\end{equation}
$\hat{\mathcal{O}}^{-1}\exp(\pm
i\chi\hat{\gamma^{q}}/2)\hat{\mathcal{O}}=\exp(\pm
i\chi\hat{\sigma}_{z}/2)$にも注意しよう．$T=0$では，式~\eqref{Part-II-FCS-Z}のトレースに寄与するエネルギー領域は，$0<\epsilon<eV$の範囲だけであることを容易に示せる．さらに，このエネルギー領域では，回転後の$\mathcal{Q}$行列はエネルギーに依存せず，次式で与えられる：
\begin{equation}
\hat{\mathcal{Q}}(0)=\left(\begin{array}{cc}-1 & -2 \\ \phantom{-}0
& \phantom{-}1\end{array}\right)\,,\qquad
\hat{\mathcal{Q}}(L)=\left(\begin{array}{cc}\phantom{-}1 &
\phantom{-}0
\\ -2 & -1\end{array}\right)\,.
\end{equation}
その結果，式~\eqref{Part-II-FCS-Z}の$\epsilon$積分は因子$eV$を与え，$\hat{\mathcal{Q}}$を$\ln\mathcal{Z}[\chi]$に代入すると，後者は次式に帰着する：
\begin{equation}\label{Part-II-FCS-Z-T-0}
\ln\mathcal{Z}[\chi]=\frac{t_{0}\mathrm{g}_{D}V}{8e}\,\Tr
\ln^{2}\left(\begin{array}{cc} -1+4e^{i\chi} & \phantom{-}2 \\
-2e^{i\chi} & -1
\end{array}\right)\,.
\end{equation}
トレースは基底の選択に依存しないので，式~\eqref{Part-II-FCS-Z-T-0}の対数の引数となる行列が対角となる基底で評価すると便利である．固有値問題を解き，トレースを計算すると，最終結果として次式を得る：
\begin{equation}
\ln\mathcal{Z}[\chi]=\frac{t_{0}\mathrm{g}_{D}V}{4e}
\ln^{2}\left[p_{\chi}+\sqrt{p^{2}_{\chi}-1}\right]\,,\quad
p_{\chi}=2e^{i\chi}-1\,.
\end{equation}
$\ln \mathcal{Z}[\chi]$がわかれば，$\chi$のべきで展開し，式~(\ref{Part-II-log-Z-chi})を用いることで，必要なキュムラントをすべて取り出せる．例えば，$\mathcal{C}_{1}=q_{0}$，$\mathcal{C}_{2}=q_0/3$，$\mathcal{C}_{3}=q_0/15$などである．FCSの総説については文献~\cite{Levitov-FCS}を参照されたい．



%$$$$$$$$$$$$$$$$$$$$$$$$$$$$$$$$$$$$$$$$$$$$$$$$$$$$$$$$$$$$$$$$$$$$$$$$$$$$$$$$$$$$$$$$$$$$$$$$$$$$$$$$$$$
%$$$$$$$$$$$$$$$$$$$$$$$$$$$$$$$$$$$$$$$$$$$$$$$$$$$$$$$$$$$$$$$$$$$$$$$$$$$$$$$$$$$$$$$$$$$$$$$$$$$$$$$$$$$

\sectionlink{フェルミ粒子の相互作用と運動論方程式}\label{sec_int_ferm}
\subsectionlink{相互作用}\label{sec_int_ferm-1}

瞬時の密度間相互作用
$$\hat{H}_{\mathrm{int}}=-{1\over 2} \iint\d\mathbf{r}\d\mathbf{r}'
:\hat\varrho(\mathbf{r}) U_{0}(\mathbf{r}-\mathbf{r}')
\hat{\varrho}(\mathbf{r}'):\,,$$
を通じて相互作用する電子液体を考える．ここで$\hat{\varrho}(\mathbf{r})
=\hat{\psi}^\dagger(\mathbf{r})\hat{\psi}(\mathbf{r})$は局所密度演算子，$U_{0}(\mathbf{r}-\mathbf{r}')$は裸のCoulomb相互作用ポテンシャル，$:\ldots:$は正規順序を表す．対応するKeldysh経路上の作用は次の形になる．
\begin{equation}\label{int-ferm-S-int}
S_{\mathrm{int}}[\bar{\psi},\psi]=-{1\over 2}\int_{\cal C} \d t
\iint\d\mathbf{r}\d\mathbf{r}'\,\bar{\psi}(\mathbf{r},t)
\bar{\psi}(\mathbf{r}',t)U_{0}(\mathbf{r}-\mathbf{r}')
\psi(\mathbf{r}',t)\psi(\mathbf{r},t)\,.
\end{equation}
ここで，経路上に定義された実ボソン場$\phi(\mathbf{r},t)$を用いてHubbard--Stratonovich変換を行い，相互作用項を分離できる．
\begin{eqnarray}\label{int-ferm-HS}
\exp\big(iS_{\mathrm{int}}[\bar{\psi},\psi]\big)=\int\D[\phi]\!\!\!\!\!\!\!\!&&\!\!
\exp\left(\frac{i}{2}\int_{\mathcal{C}}\d t
\iint\d\mathbf{r}\d\mathbf{r}'\phi\rt\,
U^{-1}_{0}(\mathbf{r}-\mathbf{r}')\,\phi(\mathbf{r}',t)\right)\nonumber \\
\times \!\!&&\!\!\!\!\!\!\! \exp\left(i\int_{\mathcal{C}}\d
t\int\d\mathbf{r}\ \phi\rt \bar{\psi}\rt\psi\rt \right)\, ,
\end{eqnarray}
ここで$U^{-1}_{0}$は相互作用核の逆であり，すなわち$\int\d\mathbf{r}'' U_{0}(\mathbf{r}-\mathbf{r}'')
U^{-1}_{0}(\mathbf{r}''-\mathbf{r}')=\delta(\mathbf{r}-\mathbf{r}')$である．補助ボソン場$\phi\rt$は，スカラーソース場~\eqref{fermion-SV}と全く同じ形でフェルミ粒子の作用に入ることに注意する．式~\eqref{fermion-SV-1}に従って$\phi^{cl(q)}\equiv(\phi_{+}\pm\phi_{-})/2$を導入すると，フェルミ粒子とボソンの相互作用項を$\bar\psi_{a}\phi^{\alpha}\hat \gamma^\alpha_{ab}\psi_b\,$と書き直せる．ここでは$a,b=(1,2)$および$\alpha=(cl,q)$についての和を暗黙に取るものとし，ガンマ行列$\hat{\gamma}^{\alpha}$は式~\eqref{fermion-gammas}で定義される．自由ボソンの項は${1\over 2} \phi U^{-1}_{0}\phi\to \phi^\alpha U^{-1}\hat
\sigma^{\alpha\beta}_{x}\phi^\beta$の形になる．式~\eqref{int-ferm-HS}に従ってフェルミ粒子を陽に積分すると，相互作用する乱れた電子液体の分配関数が得られる．
\begin{eqnarray}\label{int-ferm-NLSM-action}
&&\mathcal{Z}=\int\D[\Phi]\,\exp
\left(i\Tr\{\vec{\Phi}^{T}U^{-1}_{0}\hat{\sigma}_{x}\vec{\Phi}\}\right)
\int\D[\hat Q]\,\exp\left(iS[\hat Q,\Phi]\right)\,,\nonumber\\
&&iS[\hat Q,\Phi]=-\frac{\pi\nu}{4\tauel}\Tr\{\hat{Q}^{2}\}+
\Tr\ln\left[\hat{G}^{-1}+\frac{i}{2\tauel}\hat{Q}+
\hat{\Phi}+\mathbf{v}_{F}\hat{\mathbf{A}}\right]\,,
\end{eqnarray}
ここでは二重項$\vec{\Phi}^{T}=(\phi^{cl},\phi^{q})$および行列$\hat{\Phi}=\phi^{\alpha}\hat{\gamma}^{\alpha}$を導入した．これを式~\eqref{NLSM-action-general}で与えられた非相互作用系の作用と比較されたい．相互作用に由来する新たな複雑さは，式~\eqref{int-ferm-NLSM-action}に現れる動的なボソン場$\hat{\Phi}$についての汎関数積分が加わることである．

外部ベクトルポテンシャルを$\hat{\mathbf{A}}=0$として，式~\eqref{int-ferm-NLSM-action}の作用を$\hat Q$行列について変分し，$\delta S[\hat Q,\Phi]/\delta \hat{Q}=0$と置くと，鞍点行列$\underline{\hat{Q}}=\underline{\hat{Q}}[\Phi]$に対する次の方程式が得られる．
\begin{equation}\label{int-ferm-NLSM-saddle-point-eq}
\underline{\hat{Q}}_{tt'}(\mathbf{r})= \frac{i}{\pi\nu}
\left(\hat{G}^{-1}+\frac{i}{2\tauel}\underline{\hat{Q}}+
\hat{\Phi}\right)^{-1}_{tt',\mathbf{r}\mathbf{r}} \,,
\end{equation}
これは式~\eqref{NLSM-saddle-point-eq}を相互作用のある場合へ一般化したものである．以下では，揺らぐボソン場$\hat{\Phi}$のある実現に対して式~\eqref{int-ferm-NLSM-saddle-point-eq}の定常解を求め，その解からの空間および時間に依存するずれを考える方針を取る．

ここで概念上の問題は，任意の$\hat{\Phi}\rt$に対して鞍点方程式（\ref{int-ferm-NLSM-saddle-point-eq}）を厳密には解けないことである．ただし，ボソン場の実現が空間的に一様な$\hat{\Phi}=\hat{\Phi}(t)$である特別な場合には，式~\eqref{int-ferm-NLSM-saddle-point-eq}を解ける．これは，非相互作用系の鞍点に次のゲージ変換を施すことで達成される．
\begin{equation}\label{int-ferm-NLSM-Q-gauge}
\underline{\hat{Q}}_{tt'}[\Phi(t)]=\exp\left(i\int^{\,t}\d
t\,\hat{\Phi}(t)\right)
\hat{\Lambda}_{t-t'}\,\exp\left(-i\int^{\,t'}\d
t\,\hat{\Phi}(t)\right)\,.
\end{equation}
この解の正しさは，式~\eqref{int-ferm-NLSM-saddle-point-eq}の両辺に演算子$\hat{G}^{-1}+i/(2\tauel)\underline{\hat{Q}}+\hat{\Phi}$を作用させ，$\hat{\Lambda}_{t-t'}$が$\hat{\Phi}=0$とした式~\eqref{int-ferm-NLSM-saddle-point-eq}の解であることを利用すれば確認できる．解を式~\eqref{int-ferm-NLSM-Q-gauge}の形に書く際には，頂点行列の可換性$[\hat{\gamma}^{cl},\hat{\gamma}^{q}]=0$も用いている．この例は，ゲージを適切に選ぶと，鞍点を見いだし，そのまわりで摂動展開を行う作業が大幅に簡単になることを示している．以下では，相互作用効果の計算に適した，特に便利なゲージ（$\EuScript{K}$ゲージ）があることを示す．

%$$$$$$$$$$$$$$$$$$$$$$$$$$$$$$$$$$$$$$$$$$$$$$$$$$$$$$$$$$$$$$$$$$$$$$$$$$$$$$$$$$$$$$$$$$$$$$$$$$$$$$$$$$$
\subsectionlink{$\EuScript{K}$ゲージ}\label{sec_int_ferm-2}

元の$\hat Q$行列から，$\Qk$行列と呼ぶ新しい行列へのゲージ変換を行おう．その定義は
\begin{equation}\label{int-ferm-NLSM-Q-K-gauge}
\hat{Q}_{\EuScript{K}}(\mathbf{r};t,t')=\exp\left(-i\hat{\EuScript{K}}\rt\right)
\,\hat{Q}_{tt'}(\mathbf{r})\,
\exp\left(i\hat{\EuScript{K}}\rtp\right),
\end{equation}
である．ここで行列$\hat{\EuScript{K}}\rt=\EuScript{K}^{\alpha}\rt\hat{\gamma}^{\alpha}$は，以下で指定する$\alpha=(cl,q)$を持つ2つのスカラー場$\EuScript{K}^{\alpha}\rt$によって定義される．作用~\eqref{int-ferm-NLSM-action}に$\hat
Q =e^{i \hat{\EuScript{K}}} \hat{Q}_{\EuScript{K}}
e^{-i\hat{\EuScript{K}}}$を代入し，巡回置換に対するトレースの不変性を用いると，作用を次のように書き直せる~\footnote{式~\eqref{int-ferm-NLSM-action-K-gauge}の導出では，トレースの間の明らかな等式$\Tr\{\Qk^{2}\}=\Tr\{\hat{Q}^{2}\}$を用いる．対数項については$\Tr\left\{e^{-i\hat{\EuScript{K}}}
\ln\left[\hat{G}^{-1}+\hat{\Phi}+\mathbf{v}_{F}\hat{\mathbf{A}}+\frac{i}{2\tauel}
e^{i\hat{\EuScript{K}}}\Qk
e^{-i\hat{\EuScript{K}}}\right]e^{i\hat{\EuScript{K}}}\right\}= \Tr
\ln\left[e^{-i\hat{\EuScript{K}}}\hat{G}^{-1}e^{i\hat{\EuScript{K}}}+
\hat{\Phi}+\mathbf{v}_{F}\hat{\mathbf{A}}+\frac{i}{2\tauel}
\Qk\right]$と書く．ここでは，行列$\hat{A}$の任意の解析関数$f$に対して成り立つ，よく知られた代数的恒等式$\Tr\{\hat{L}f(\hat{A})\hat{L}\}=\Tr\{
f(\hat{L}\hat{A}\hat{L}^{-1})\}$を用いた．最後に$e^{-i\hat{\EuScript{K}}}\hat{G}^{-1}e^{i\hat{\EuScript{K}}}=\hat{G}^{-1}+
e^{-i\hat{\EuScript{K}}}[\hat{G}^{-1},e^{i\hat{\EuScript{K}}}]$と書き直し，交換子$[\hat{G}^{-1},e^{i\hat{\EuScript{K}}}]=
e^{i\hat{\EuScript{K}}} \left(-\partial_{t}
\hat{\EuScript{K}}-\mathbf{v}_{F}\partial_{\mathbf{r}}\hat{\EuScript{K}}
-\frac{1}{2m}
(\partial_{\mathbf{r}}\hat{\EuScript{K}})^{2}\right)$を計算する．}．
\begin{equation}\label{int-ferm-NLSM-action-K-gauge}
iS[\hat
Q_{\EuScript{K}},\Phi]=-\frac{\pi\nu}{4\tauel}\Tr\big\{\Qk^{2}\big\}+
\Tr\ln\left[\hat{G}^{-1}+\hat{C}+\frac{i}{2\tauel}\Qk-\frac{1}{2m}(\partial_{\mathbf{r}}
\hat{\EuScript{K}})^{2}\right]\,,
\end{equation}
ここでは$\hat{C}\rt = \hat{\Phi}_{\EuScript{K}}\rt+
\mathbf{v}_{F}\hat{\mathbf{A}}_{\EuScript{K}}\rt $という記法を導入し，ゲージ変換した電磁ポテンシャルを次のように定義した．
\begin{equation}\label{int-ferm-NLSM-C}
\hat{\Phi}_{\EuScript{K}}\rt=\hat{\Phi}\rt-\partial_{t}\hat{\EuScript{K}}\rt\,,\quad
\hat{\mathbf{A}}_{\EuScript{K}}\rt=
\hat{\mathbf{A}}\rt-\partial_{\mathbf{r}}\hat{\EuScript{K}}\rt\,.
\end{equation}

ここで，新しい場$\Qk$の鞍点が非相互作用系の鞍点$\hat \Lambda$，式（\ref{NLSM-Lambda}）に近いと仮定し，2つの場$\hat{\EuScript{K}}^\alpha$を選ぶ自由度を用いて，この条件を満たすようにする．そのために式~\eqref{int-ferm-NLSM-action-K-gauge}に$\Qk=\hat{\Lambda}+\delta\Qk$を代入し，ずれ$\delta\Qk$および$\hat C$にまとめられた電磁ポテンシャルのべきで展開する．この展開の最初の非自明な項は
\begin{equation}\label{int-ferm-NLSM-action-Q-linear}
iS[\delta \hat Q_{\EuScript{K}},\Phi]= -\frac{i}{2\tauel}
\Tr\big\{\hat{\mathcal{G}}\hat{C}\hat{\mathcal{G}}\delta\Qk\big\}+\ldots\,,
\end{equation}
である．ここでは，$\hat{\Lambda}$が非相互作用模型の鞍点であるため，電磁ポテンシャルがなければ，ずれ$\delta\Qk$についての線形項が存在しないことを用いた．また，反磁性項$(\partial_{\mathbf{r}}\hat{\EuScript{K}})^{2}/2m$は$\hat{\EuScript{K}}$について2次であり，したがって（以下に示すように）$\hat{\Phi}$についても2次であるため無視した．

ここで，$\delta \hat Q_{\EuScript{K},tt'}(\mathbf{r})$について線形なこの項が消えることを要求する．Fourier変換すると，任意の$\epsilon$，$\omega$および$\mathbf{q}$に対して次の行列恒等式が成り立つなら，任意の$\delta \hat Q_{\EuScript{K},\epsilon_- \epsilon_+}({\bf q})$についてこの要請が満たされることがわかる．
\begin{equation}\label{int-ferm-NLSM-Phi-K}
\sum_{p}\hat{\mathcal{G}}(\mathbf{p}_{+},\epsilon_{+})
\hat{C}(\mathbf{q},\omega)
\hat{\mathcal{G}}(\mathbf{p}_{-},\epsilon_{-})=0\,,
\end{equation}
ここで$\mathbf{p}_{\pm}=\mathbf{p}\pm\mathbf{q}/2\,$および$\epsilon_{\pm}=\epsilon\pm\omega/2$である．条件~\eqref{int-ferm-NLSM-Phi-K}は，まだ指定していないゲージ場$\EuScript{K}^{\alpha}$を$\Phi^{\alpha}$および${\bf
A}^\alpha$で表す行列方程式である．式~\eqref{NLSM-G-impurity-dressed}および次の恒等式
\begin{subequations}\label{int-ferm-NLSM-integrals}
\begin{equation}
\hskip-1.2cm
\sum_{p}\mathcal{G}^{R}(\mathbf{p}_{\pm},\epsilon_{\pm})
\mathcal{G}^{A}(\mathbf{p}_{\mp},\epsilon_{\mp})\approx2\pi\nu\tauel\,,
\end{equation}
\begin{equation}
\!\!\! \sum_{p}\mathbf{v}_{F}\,
\mathcal{G}^{R}(\mathbf{p}_{\pm},\epsilon_{\pm})
\mathcal{G}^{A}(\mathbf{p}_{\mp},\epsilon_{\mp})\approx\mp2\pi
i\nu\tauel D \mathbf{q}\,,
\end{equation}
\end{subequations}
を用いると，式~\eqref{int-ferm-NLSM-Phi-K}は次のように変形できる．
\begin{equation}\label{int-ferm-NLSM-Phi-K-1}
\frac{1}{\pi\nu\tauel}
\sum_{p}\hat{\mathcal{G}}(\mathbf{p}_{+},\epsilon_{+})
\hat{C}(\mathbf{q},\omega)
\hat{\mathcal{G}}(\mathbf{p}_{-},\epsilon_{-})=
\big(\hat{\gamma}^{\alpha}-\hat{\Lambda}_{\epsilon_{+}}\hat{\gamma}^{\alpha}
\hat{\Lambda}_{\epsilon_{-}}\big)\Phi^{\alpha}_{\EuScript{K}}-
\big(\hat{\Lambda}_{\epsilon_{+}}\hat{\gamma}^{\alpha}-
\hat{\gamma}^{\alpha}\hat{\Lambda}_{\epsilon_{-}}\big)
D\,\mathrm{div}\mathbf{A}_{\EuScript{K}}^{\alpha}=0\,.
\end{equation}
一般には，2つの場$\EuScript{K}^\alpha(\mathbf{r},\omega)$を選ぶことによって，任意の$\epsilon$および$\omega$に対してこの条件を満たすことはできない．しかし，熱平衡では次の「魔法のような」事実が成り立つ．
\begin{equation}\label{int-ferm-F-B}
\frac{1-F_{\epsilon_{+}}F_{\epsilon_{-}}}
{F_{\epsilon_{+}}-F_{\epsilon_{-}}}=\coth\frac{\omega}{2T}\equiv
B_{\omega}\, ,
\end{equation}
これは$\omega$のみに依存し，$\epsilon$には{\em 依存しない}．そのため，ゲージ変換したポテンシャル（\ref{int-ferm-NLSM-C}）の間に次のベクトル関係が成り立てば，条件（\ref{int-ferm-NLSM-Phi-K-1}）を満たすことができる．
\begin{equation}\label{int-ferm-NLSM-K-functional}
\vec{\Phi}_{\EuScript{K}}(\mathbf{r},\omega)=
\left(\begin{array}{cc} 1& 2B_{\omega} \\ 0 &
-1\end{array}\right)D\,\mathrm{div}
\vec{\mathbf{A}}_{\EuScript{K}}(\mathbf{r},\omega)\, .
\end{equation}
この方程式は，電磁ポテンシャルの古典成分と量子成分の両方について$\EuScript{K}$ゲージを指定する．

$\EuScript{K}$ゲージの利点は，$\Qk$行列の鞍点$\hat \Lambda$からのずれに線形であり，{\em かつ電磁ポテンシャルにも線形な}項が作用に含まれないことである．ただし，$\delta \Qk$に線形で電磁ポテンシャルに2次の項は依然として存在する．これは，厳密には，電磁ポテンシャルのどの実現に対しても$\hat \Lambda$が$\Qk$多様体上の厳密な鞍点となるわけではないことを意味する．しかし，真の鞍点からのずれがポテンシャルの2次まで押し上げられるため，$\EuScript{K}$ゲージは摂動論の構造を大幅に簡単にする．さらに，この性質が成り立つのは平衡状態に限られる．非平衡の場合には条件（\ref{int-ferm-NLSM-Phi-K-1}）を恒等的には満たせず，$\delta\Qk$と電磁場の双方に線形な項が作用に現れる．後に説明するように，まさにこれらの項が運動論方程式の衝突積分を生み出す．それでも$\EuScript{K}$ゲージは非平衡の場合にも有用な概念である．この場合，式（\ref{int-ferm-NLSM-K-functional}）のボソン分布関数$B_{\omega}$を次のように定義する．
\begin{equation}\label{int-ferm-NLSM-bosonic-distribution}
B_{\omega}(\mathbf{r},\tau)=\frac{1}{2\omega}\int^{+\infty}_{-\infty}
\d\epsilon\,\left[1-F_{\epsilon+\omega/2}(\mathbf{r},\tau)
F_{\epsilon-\omega/2}(\mathbf{r},\tau)\right]\,,
\end{equation}
ここで$F_{\epsilon}(\mathbf{r},\tau)$はフェルミ粒子の行列$F_{t,t'}(\mathbf{r})$のWigner変換（WT）である．

式（\ref{int-ferm-NLSM-C}）を用いると，$\EuScript{K}$ゲージの定義~\eqref{int-ferm-NLSM-K-functional}は，ゲージ場$\EuScript{K}^\alpha$を電磁ポテンシャル$\Phi^\alpha$および${\bf A}^\alpha$から決める陽な関係とみなせる．簡単のため$\hat{\mathbf{A}}=0$と置くと，ゲージ場の量子成分と古典成分について次の式を得る．
\begin{subequations}\label{int-ferm-NLSM-K-q-cl}
\begin{equation}
\hskip-3cm
(D\partial^{2}_{\mathbf{r}}-i\omega)\EuScript{K}^{q}(\mathbf{r},\omega)=
\Phi^{q}(\mathbf{r},\omega)\,,
\end{equation}
\begin{equation}
(D\partial^{2}_{\mathbf{r}}+i\omega)\EuScript{K}^{cl}(\mathbf{r},\omega)+
2B_{\omega}D\partial^{2}_{\mathbf{r}}\EuScript{K}^{q}(\mathbf{r},\omega)
=-\Phi^{cl}(\mathbf{r},\omega)\,.
\end{equation}
\end{subequations}
一般の場合には，これらの関係を行列形式にまとめると便利である．
\begin{equation}\label{int-ferm-NLSM-Phi-K-matrix}
\vec{\EuScript{K}}\qo =  \hat{\mathcal{D}}^{-1}\qo \Big(
\hat{\EuScript{B}}^{-1}_{\omega}\vec{\Phi}\qo-
D\,\hat{\sigma}_{x}\,\mathbf{q}\cdot\vec{\mathbf{A}}\qo \Big)\,,
\end{equation}
ここでベクトルは$\vec{\EuScript{K}}^{T}=(\EuScript{K}^{cl},\EuScript{K}^{q})$である．また，ディフューゾンのボソン行列伝播関数
\begin{equation}\label{int-ferm-NLSM-D}
\hat{\mathcal{D}}\qo=\left(\begin{array}{cc}\mathcal{D}^{K}\qo &
\mathcal{D}^{R}\qo \\ \mathcal{D}^{A}\qo & 0\end{array}\!\right),\,
\end{equation}
を導入した．その行列成分は
\begin{equation}\label{int-ferm-NLSM-D-RAK}
\mathcal{D}^{R(A)}\qo=\big(Dq^{2}\mp i\omega\big)^{-1},\quad
\mathcal{D}^{K}\qo=B_{\omega}\big[\mathcal{D}^{R}\qo-\mathcal{D}^{A}\qo\big]\,,
\end{equation}
および次式である．
\begin{equation}\label{int-ferm-NLSM-B}
\hat{\EuScript{B}}_{\omega}=\left(\begin{array}{lr}2B_{\omega} &
1^{R}_{\omega} \\ -1^{A}_{\omega} & 0\end{array}\right)\,.
\end{equation}

式~\eqref{int-ferm-NLSM-Phi-K-matrix}は，ゲージ場$\EuScript{K}^\alpha$と電磁ポテンシャルの間の陽な{\em 線形}関係を与える．したがって，ゲージ変換した場$\Qk$を陽に定義するものとなっている．式（\ref{int-ferm-NLSM-action-K-gauge}）を参照されたい．この場の鞍点は，非相互作用系の鞍点$\hat \Lambda$にかなり近い（そのずれは電磁場について2次である）．元のゲージに戻すと，次の$\hat Q$行列
\begin{equation}\label{int-ferm-NLSM-Q-gauge-general}
\underline{\hat{Q}}_{tt'}(\mathbf{r})=
\exp\left(i\EuScript{K}^{\alpha}\rt\hat{\gamma}^{\alpha}\phantom{^{'}}\right)
\hat{\Lambda}_{t-t'}\,
\exp\left(-i\EuScript{K}^{\beta}\rtp\hat{\gamma}^{\beta}\right),
\end{equation}
が，揺らぐポテンシャルの任意の実現に対して，一般的な鞍点方程式（\ref{int-ferm-NLSM-saddle-point-eq}）の解の良い近似を与えることがわかる．この主張が成り立つのは平衡条件のもとだけである．非平衡では$\hat{\Phi}\delta\Qk$の項が再び現れ，正しい形の運動論方程式を得るには，これを考慮しなければならない（詳しくは第~\ref{sec_int_ferm-5}節の議論を参照）．さらに，平衡でも$\sim\hat{\Phi}^{2}\delta\Qk$の項が存在する．これらは輸送係数に対する相互作用補正を生む（詳細は第~\ref{app_Part-III}節で述べる）．


%$$$$$$$$$$$$$$$$$$$$$$$$$$$$$$$$$$$$$$$$$$$$$$$$$$$$$$$$$$$$$$$$$$$$$$$$$$$$$$$$$$$$$$$$$$$$$$$$$$$$$$$$$$$
\subsectionlink{相互作用系の非線形$\sigma$模型}\label{sec_int_ferm-3}

式~\eqref{int-ferm-NLSM-action-K-gauge}の対数項のトレースについて勾配展開を行うと（この手続きは第~\ref{sec_NLSM-2}節で示したものとよく似ている），$\Qk$行列場と$\EuScript{K}$ゲージにおける電磁ポテンシャルで書かれた有効作用が得られる．
\begin{eqnarray}\label{int-ferm-action-interacting}
iS[\hat Q_{\EuScript{K}},\Phi]=\frac{i\nu}{2}
\Tr\left\{\hat{\Phi}_{\EuScript{K}}
\hat{\sigma}_{x}\hat{\Phi}_{\EuScript{K}}\right\}-\frac{\pi\nu}{4}
\Tr\left\{D(\hat{\bm{\partial}}_{\mathbf{r}}\Qk)^{2}-
4\partial_{t}\Qk+4i\hat{\Phi}_{\EuScript{K}}\Qk\right\}\,,
\end{eqnarray}
ここで，共変微分は次式で定義される．
\begin{eqnarray}\label{int-ferm-covariant-derivative}
\hat{\bm{\partial}}_{\mathbf{r}}\Qk=\partial_{\mathbf{r}}\Qk-
i\big[\hat{\mathbf{A}}_{\EuScript{K}},\Qk\big]\,.
\end{eqnarray}
式~\eqref{int-ferm-action-interacting}は，鞍点条件~\eqref{int-ferm-NLSM-Phi-K-matrix}--\eqref{int-ferm-NLSM-B}とあわせて，有効$\sigma$模型の作用~\eqref{NLSM-action-noninteracting}をCoulomb相互作用効果が含まれるように一般化したものである．拡張された共変微分~\eqref{int-ferm-covariant-derivative}の具体的な形と，$\hat{\mathbf{A}}=0$における場$\hat{\EuScript{K}}$と$\hat{\Phi}$の関係（式~\eqref{int-ferm-NLSM-K-q-cl}を参照）を用いると，分配関数について
\begin{equation}\label{int-ferm-Z}
\mathcal{Z}=\int\D[\Phi]
\exp\left(i\Tr\{\vec{\Phi}^{T}\hat{U}^{-1}_{RPA}\vec{\Phi}\}\right)
\int\D[\hat Q_{\EuScript{K}}] \exp\left(iS_{0}[\hat
Q_{\EuScript{K}}]+ iS_{1}[\hat
Q_{\EuScript{K}},\partial_{\mathbf{r}}\EuScript{K}]+ iS_{2}[\hat
Q_{\EuScript{K}},\partial_{\mathbf{r}}\EuScript{K}]\right),
\end{equation}
を得る．ここで$l=0,1,2$について，$S_{l}$は電磁ポテンシャルの$l$次の項を含み，次の式で与えられる．
\begin{subequations}\label{int-ferm-S-012}
\begin{equation}
\hskip-3cm iS_{0}[\hat Q_{\EuScript{K}}]=-\frac{\pi\nu}{4}
\Tr\left\{D(\partial_{\mathbf{r}}\Qk)^{2}-4i\partial_{t}\Qk\right\}\,,
\end{equation}
\begin{equation}
\hskip-1cm
iS_{1}[\hat Q_{\EuScript{K}},\partial_{\mathbf{r}}\EuScript{K}]=
-i\pi\nu\Tr\left\{D(\partial_{\mathbf{r}}\hat{\EuScript{K}})\Qk(\partial_{\mathbf{r}}\Qk)+
\hat{\Phi}_{\EuScript{K}}\Qk\right\}\,,
\end{equation}
\begin{equation}
iS_{2}[\hat Q_{\EuScript{K}},\partial_{\mathbf{r}}\EuScript{K}]=\frac{\pi\nu
D}{2}\Tr\left\{(\partial_{\mathbf{r}}\hat{\EuScript{K}})\Qk
(\partial_{\mathbf{r}}\hat{\EuScript{K}})\Qk-
(\partial_{\mathbf{r}}\hat{\EuScript{K}})\hat{\Lambda}
(\partial_{\mathbf{r}}\hat{\EuScript{K}})\hat{\Lambda}\right\}\,.
\end{equation}
\end{subequations}
有効相互作用行列$\hat{U}_{RPA}$は，乱雑位相近似（RPA）で遮蔽された相互作用そのものである．
\begin{equation}\label{int-ferm-U-RPA}
\hat{U}_{RPA}\qo=\big[U^{-1}_{0}\hat{\sigma}_{x}+\hat{\Pi}\qo\big]^{-1}\,,
\end{equation}
ここで$\hat{\Pi}\qo$は密度間相関関数である．式（\ref{fermion-Pi-matrix-def}）および（\ref{Part-II-Pi-diffusive}）によれば，これはKeldysh空間におけるボソン伝播関数の典型的な形を持つ．
\begin{equation}\label{int-ferm-Pi-matrix}
\hat{\Pi}\qo=\left(\begin{array}{cc}0& \Pi^{A}\qo \\
\Pi^{R}\qo & \Pi^{K}\qo\end{array}\right)\,,
\end{equation}
その成分は次のとおりである．
\begin{equation}\label{int-ferm-Pi-RAK}
\Pi^{R(A)}\qo=\frac{\nu Dq^{2}}{Dq^{2}\mp i\omega}\,,\,
\qquad\Pi^{K}\qo=B_{\omega}\big[\Pi^{R}\qo-\Pi^{A}\qo\big]\,.
\end{equation}
式~\eqref{int-ferm-Z}--\eqref{int-ferm-Pi-RAK}を導くには，$\Tr\big\{(\partial_{\mathbf{r}}\hat{\EuScript{K}})\hat{\Lambda}
(\partial_{\mathbf{r}}\hat{\EuScript{K}})\hat{\Lambda}\big\}$という項を足して引き，次の方程式を用いればよい．
\begin{equation}\label{int-ferm-integral-1}
\int^{+\infty}_{-\infty}\d\epsilon\,\Tr\left\{\hat{\gamma}^{\alpha}
\hat{\gamma}^{\beta}-
\hat{\gamma}^{\alpha}\hat{\Lambda}_{\epsilon_{+}}
\hat{\gamma}^{\beta}\hat{\Lambda}_{\epsilon_{-}}\right\}=
4\omega\big(\hat{\EuScript{B}}^{-1}_{\omega}\big)^{\alpha\beta}\,,
\end{equation}
ここで$\epsilon_{\pm}=\epsilon\pm\omega/2$であり，行列$\hat{\Lambda}$および$\hat{\EuScript{B}}$は，それぞれ式~\eqref{NLSM-Lambda}および\eqref{int-ferm-NLSM-B}で定義される．式~\eqref{int-ferm-integral-1}は，ボソン分布関数とフェルミ分布関数の間の次の積分関係から従う．
\begin{equation}\label{int-ferm-integral-2}
\int^{+\infty}_{-\infty}\d\epsilon\,\big(F_{\epsilon_{+}}-F_{\epsilon_{-}}\big)=2\omega\,,\qquad
\int^{+\infty}_{-\infty}\d\epsilon\,\big(1-F_{\epsilon_{+}}F_{\epsilon_{-}}\big)=2\omega
B_{\omega}\,.
\end{equation}

式~\eqref{int-ferm-Z}--\eqref{int-ferm-Pi-RAK}は，相互作用する乱れたFermi液体の有効非線形$\sigma$模型を構成する．この模型は，非線形拘束条件$\Qk^{2}=\hat{1}$を満たす行列場$\Qk$と，ボソン的な縦場$\partial_{\mathbf{r}}\hat{\EuScript{K}}$（または，これと等価な$\hat{\Phi}$）という2つの相互作用する場からなる．後で明らかになるように，$\Qk$場は準粒子分布関数の揺らぎを記述し，$\hat{\Phi}$（または$\hat{\EuScript{K}}$）は媒質中の電磁モードの伝播を表す．

%$$$$$$$$$$$$$$$$$$$$$$$$$$$$$$$$$$$$$$$$$$$$$$$$$$$$$$$$$$$$$$$$$$$$$$$$$$$$$$$$$$$$$$$$$$$$$$$$$$$$$$$$$$$
\subsectionlink{相互作用の伝播関数}\label{sec_int_ferm-4}

後の応用のために，相関関数
\begin{equation}\label{int-ferm-V-def}
\fitdisplay{
\mathcal{V}^{\alpha\beta}(\mathbf{r}-\mathbf{r}',t-t')\! = -2i
\big\langle\EuScript{K}^{\alpha}(\mathbf{r},t)\EuScript{K}^{\beta}(\mathbf{r}',t')\big\rangle\!
= -2i\!\int\! \D[\hat \Phi]\,\EuScript{K}^{\alpha}\rt
\EuScript{K}^{\beta}(\mathbf{r}',t')\,
\exp\left(i\Tr\{\vec{\Phi}^{T}\hat{U}^{-1}_{RPA}\vec{\Phi}\}\right)\,,
}
\end{equation}
を導入する．ここで係数$-2i$は便宜上付けてある．$\hat{\Phi}$と$\hat{\EuScript{K}}$は式~\eqref{int-ferm-NLSM-Phi-K-matrix}によって線形に関係しているので，このGauss積分を実行し，ゲージ場の相関関数について次を得る．
\begin{equation} \label{int-ferm-K-corr-fun}
\hat{\mathcal{V}}\qo=\hat{\mathcal{D}}\qo\hat{\EuScript{B}}^{-1}_{\omega}\hat{U}_{RPA}\qo
\big(\hat{\EuScript{B}}^{-1}_{-\omega}\big)^{T}\hat{\mathcal{D}}^{T}\qom\,.
\end{equation}
ボース相関行列$\hat{\mathcal{V}}\qo$は，標準的なKeldysh構造
\begin{equation}\label{int-ferm-V-matrix}
\hat{\mathcal{V}}\qo=\left(\begin{array}{cc}\mathcal{V}^{K}\qo &
\mathcal{V}^{R}\qo \\ \mathcal{V}^{A}\qo & 0\end{array}\right),
\end{equation}
を持ち，その成分は
\begin{subequations}\label{int-ferm-V-RAK}
\begin{equation}
\mathcal{V}^{R(A)}\qo=-\frac{1}{(Dq^{2}\mp
i\omega)^{2}}\left(U^{-1}_{0}+\frac{\nu Dq^{2}}{Dq^{2}\mp
i\omega}\right)^{-1}\,,
\end{equation}
\begin{equation}
\hskip-1.5cm \mathcal{V}^{K}\qo=B_{\omega}
\big[\mathcal{V}^{R}\qo-\mathcal{V}^{A}\qo\big]\,.
\end{equation}
\end{subequations}
である．この伝播関数は，2つの頂点をディフューゾンでドレスした，遮蔽された動的Coulomb相互作用に対応する（図~\ref{Fig-K-corr-fun}a）．したがって，ゲージ場$\EuScript{K}$の役割は，RPAで遮蔽された相互作用（図~\ref{Fig-K-corr-fun}b）と，ディフューゾンによるその頂点の繰り込みを，自動的に取り込むことである．$\hat{\Phi}$と$\hat{\EuScript{K}}$は線形に依存しているため（式~\eqref{int-ferm-NLSM-Phi-K-matrix}を参照），$\hat{\Phi}$場と$\hat{\EuScript{K}}$場についての平均を区別せずに用いる．要点は，2つの$\hat{\EuScript{K}}^\alpha$場の相関関数が式（\ref{int-ferm-V-def}）--（\ref{int-ferm-V-RAK}）で与えられることである．
\begin{figure}
\begin{center}\includegraphics[width=10cm]{figures/fig13.pdf}\end{center}
\caption{a）ゲージ場伝播関数$\hat{\mathcal{V}}\qo$のダイアグラム表示．波線はCoulomb相互作用を表す．ディフューゾンでドレスされた頂点は，破線からなる梯子で示している．b）RPAで遮蔽されたCoulomb相互作用$\hat{U}_{RPA}\qo$．太い波線と細い波線は，それぞれ遮蔽された相互作用と裸の相互作用を表し，ループは拡散梯子でドレスされた分極演算子を表す．\label{Fig-K-corr-fun}}
\end{figure}

%$$$$$$$$$$$$$$$$$$$$$$$$$$$$$$$$$$$$$$$$$$$$$$$$$$$$$$$$$$$$$$$$$$$$$$$$$$$$$$$$$$$$$$$$$$$$$$$$$$$$$$$$$$$
\subsectionlink{運動論方程式}\label{sec_int_ferm-5}

本節の目的は，分布関数$F$の運動論方程式がKeldysh形式の枠組みの中で自然に現れることを示すことである．第~\ref{sec_NLSM-4}節では，非相互作用フェルミ粒子の運動論方程式が，$\hat Q$行列の有効作用に対する鞍点方程式そのものであることを示した．相互作用する電子の場合には，まず速い自由度である拡散モード$\hat{\mathcal{W}}$と電磁モード$\hat{\EuScript{K}}$（または，これと等価な$\hat{\Phi}$）を積分消去することにより，作用$S[\hat Q_{\EuScript{K}},\Phi]$（式~\eqref{int-ferm-action-interacting}を参照）から運動論方程式が得られる．

分配関数~\eqref{int-ferm-Z}および\eqref{int-ferm-S-012}から運動論方程式へ至る手続き全体の論理を概説しよう．第一段階では，$l=0,1,2$を持つ作用$S_{l}[\hat
Q_{\EuScript{K}},\partial_{\mathbf{r}}\EuScript{K}]$（式~\eqref{int-ferm-S-012}を参照）において，遅い自由度と速い自由度を分離する．前者は分布関数$F_{tt'}(\mathbf{r})$に，後者はディフューゾン$\hat{\mathcal{W}}_{tt'}(\mathbf{r})$および電磁モード$\hat{\EuScript{K}}\rt$に担われる．この分離は$\Qk$行列を適切にパラメータ表示することで実現される．便利な選び方の一つは$\Qk=\hat{\mathcal{U}}_{z}\circ\hat{Q}_{\mathrm{fast}}\circ\hat{\mathcal{U}}^{-1}_{z}$であり，ここで回転行列
\begin{equation}\label{int-ferm-R}
\hat{\mathcal{U}}_{z}=\left(\begin{array}{cr}1-F\circ Z & F \\ Z &
-1
\end{array}\right)\,,\qquad \hat{\mathcal{U}}^{-1}_{z}=
\left(\begin{array}{cc} 1 & F \\ Z & -1+Z\circ F
\end{array}\right)\,,
\end{equation}
は，$A\circ B=\int\d t' A_{tt'}B_{t't''}$とともに，遅い自由度についての情報を担う．一方，$\Qk$行列の速い部分は，ディフューゾン場$\hat{Q}_{\mathrm{fast}}=\exp\{-\hat{\mathcal{W}}/2\}\circ
\hat{\sigma}_{z}\circ\exp\{\hat{\mathcal{W}}/2\}$を用いてパラメータ表示する（このパラメータ表示を式~\eqref{NLSM-Q-U-W-parametrization}で与えたものと比較されたい）．最後の式において，$Z_{tt'}(\mathbf{r})$（分配関数と混同しないこと）は，分布関数$F_{tt'}(\mathbf{r})$の\textit{量子}成分と考えられる．計算の最後には$Z_{tt'}(\mathbf{r})$をゼロと置くが，運動論方程式の衝突積分を正しい形で得るには，$\hat Q$のパラメータ表示に$Z_{tt'}(\mathbf{r})$を陽に残しておかなければならないことが，文献~\cite{Zala}で強調されている．

第二段階では，分配関数~\eqref{int-ferm-Z}において，$\hat{\Phi}$（または，両者の関係が式（\ref{int-ferm-NLSM-Phi-K-matrix}）によって定まるので，これと等価な$\hat{\EuScript{K}}$）と$\hat{\mathcal{W}}$場について積分し，有効作用
\begin{equation}\label{int-ferm-S-eff-def}
\mathcal{Z}=\int\D[\hat Q_{\EuScript{K}},\Phi]\,
\exp\big(iS[\hat{\mathcal{W}},\partial_{\mathbf{r}}\EuScript{K}]\big)=\int\D[F,Z]\,
\exp\big(iS_{\mathrm{eff}}[F,Z]\big)\,.
\end{equation}
を得る．式~\eqref{NLSM-Q-U-W-parametrization}による分解を行い，$\hat{\mathcal{U}}_{z}$および$\hat{\mathcal{U}}^{-1}_{z}$行列を式~\eqref{int-ferm-R}の形にした後には，式~\eqref{int-ferm-S-eff-def}の$\hat Q_{\EuScript{K}}$行列についての汎関数積分は，独立な行列場$F$，$Z$および$\hat{\mathcal{W}}$についての積分と理解すべきことに注意する．その結果，有効作用$S_{\mathrm{eff}}$は，$F$とその量子成分$Z$，および必要に応じてスカラーポテンシャルやベクトルポテンシャルなどの古典的な外場に依存する．次に，分布関数$F$についての鞍点方程式
\begin{equation}\label{int-ferm-KE-def}
\left.\frac{\delta S_{\mathrm{eff}}[F,Z]}{\delta Z}\right|_{Z=0} =
0\,,
\end{equation}
を求める．これが求める運動論方程式である．

この方針に沿って，作用\eqref{int-ferm-S-012}を$F$，$Z$，$\hat{\mathcal{W}}$，および電磁ポテンシャル$\Phi$と$\EuScript{K}$について展開する．作用の遅い部分については，式（\ref{int-ferm-S-012}a）から$\Tr\big\{(\partial_{\mathbf{r}}
\Qk)^{2}\big\}=8\tr\big\{\partial_{\mathbf{r}}
F_{tt'}\partial_{\mathbf{r}} Z_{t't}\big\}+O(Z^{2})$および$\Tr\{\partial_{t}\Qk\}=2\mathrm{tr}\{\partial_{t}Z_{tt'}F_{t't}-
\partial_{t}F_{tt'}Z_{t't}\}$を得る．ここで$\tr\{\ldots\}$は空間積分と時間積分のみを表し，Keldysh構造についてのトレースは陽に取ってある．Wigner変換表示~\eqref{NLSM-Wigner}へ移ると，
\begin{equation}\label{int-ferm-S-0-FZ}
iS_{0}[F,Z]=2\pi\nu\,\tr\left\{\big[D\partial^{2}_{\mathbf{r}}F_{\epsilon}(\mathbf{r},\tau)-
\partial_{\tau}F_{\epsilon}(\mathbf{r},\tau)\big]Z_{\epsilon}(\mathbf{r},\tau)\right\}\,,
\end{equation}
を得る．ここで$\tau=(t+t')/2$である．この段階でもすでに，$S_{0}[F,Z]$を$Z$で微分することによって，式~\eqref{int-ferm-KE-def}から非相互作用系の運動論方程式~\eqref{NLSM-Diffusion-Eq}が再現される．同様に，速い自由度に対する作用の動的部分は
\begin{equation}\label{int-ferm-S-0-W}
iS_{0}[\hat{\mathcal{W}}]=-\frac{\pi\nu}{2}\,\tr\left\{\bar{d}_{\epsilon}(\mathbf{r},\tau)
\big[D\partial^{2}_{\mathbf{r}}-\partial_{\tau}\big]d_{\epsilon}(\mathbf{r},\tau)\right\}\,,
\end{equation}
となる．これは式~\eqref{NLSM-S-W}のWigner表示そのものである．

次に，$\hat{\mathcal{W}}$モードと$\Phi$モードの結合項を考える．式（\ref{int-ferm-S-012}b）を展開して得られる作用の$S_{1}[\hat{\mathcal{W}},F,Z]$部分については，
\begin{eqnarray}\label{int-ferm-S-1-WFZ}
iS_{1}[\hat{\mathcal{W}},F,Z]=-i\pi\nu\,\tr\left\{\Big([F,X^{cl}_{+}]+
X^{q}_{+}-FX^{q}_{+}F\Big)\bar{d} +
\Big( [Z,X^{cl}_{-}] - X^{q}_{-} +
ZFX^{q}_{-} + X^{q}_{-}FZ \Big) d \right\}\,,
\end{eqnarray}
を得る．ここで
\begin{equation}\label{int-ferm-X}
X^{\alpha}_{\pm}=\Phi^{\alpha}-\partial_{t}\EuScript{K}^{\alpha}\pm
D\partial^{2}_{\mathbf{r}}\EuScript{K}^{\alpha}.
\end{equation}
である．$\hat{\Phi}$場と$\hat{\EuScript{K}}$場の汎関数関係を導いた際には，作用の$S_{1}$部分をゼロにするという考え方を用いた（式~\eqref{int-ferm-NLSM-action-Q-linear}を思い出されたい）．非平衡の場合には，この操作を実行することはできない．しかし，条件$X^{q}_{-}=0$を課すことによって，式（\ref{int-ferm-NLSM-K-q-cl}a）は依然として満たせる．式（\ref{int-ferm-NLSM-Phi-K-1}）のKeldysh成分を恒等的に満たすことはできないものの，$\EuScript{K}^{cl}$が式（\ref{int-ferm-NLSM-K-q-cl}b）を非平衡へ一般化した次の式に従うよう要求することには意味がある．
\begin{equation}\label{int-ferm-NLSM-K-cl-NonEq}
(D\partial^{2}_{\mathbf{r}}+i\omega)\EuScript{K}^{cl}(\mathbf{r},\omega)+
2B_{\omega}(\mathbf{r},\tau)D\partial^{2}_{\mathbf{r}}\EuScript{K}^{q}(\mathbf{r},\omega)
=-\Phi^{cl}(\mathbf{r},\omega)\,,
\end{equation}
ここで非平衡ボソン分布関数は式（\ref{int-ferm-NLSM-bosonic-distribution}）で定義される．ただし，この一般化は，式~\eqref{int-ferm-S-1-WFZ}において$\hat{\mathcal{W}}$（すなわち$d$と$\bar d$）に線形な項が消えることを\textit{意味しない}．実際，$\hat{\Phi}$と$\hat{\EuScript{K}}$の量子成分を関係づける式（\ref{int-ferm-NLSM-K-q-cl}a）および式~\eqref{int-ferm-NLSM-K-cl-NonEq}を用いてWigner変換を行うと，作用の$S_{1}[\hat{\mathcal{W}},F,Z]$部分を次の形にできる．
\begin{equation}\label{int-ferm-S-1-WFZ-1}
\fitdisplay{
iS_{1}[\hat{\mathcal{W}},F,Z]=-i\pi\nu\,
\tr\left\{\mathcal{I}[F]X^{q}_{+}(\mathbf{r},\omega)
\bar{d}_{\epsilon_{-}}(\mathbf{r},\tau)e^{-i\omega\tau}+
Z_{\epsilon}(\mathbf{r},\tau)X^{cl}_{-}(\mathbf{r},\omega)
[d_{\epsilon_{-}}(\mathbf{r},\tau)-
d_{\epsilon_{+}}(\mathbf{r},\tau)] e^{-i\omega\tau}\right\}\,,
}
\end{equation}
ここで$\epsilon_{\pm}=\epsilon\pm\omega/2$であり，汎関数
\begin{equation}
\mathcal{I}[F]=B_{\omega}(\mathbf{r},\tau)
\big[F_{\epsilon-\omega}(\mathbf{r},\tau)-F_{\epsilon}(\mathbf{r},\tau)\big]
+ 1-F_{\epsilon-\omega}(\mathbf{r},\tau)F_{\epsilon}(\mathbf{r},\tau) \,.
\end{equation}
を導入した．平衡では$\mathcal{I}[F]\equiv 0$であることに注意する．式~\eqref{int-ferm-S-1-WFZ-1}では$\omega$への依存性を陽に残しており，したがって通常のWigner変換で行うような，$\epsilon$に比べて小さい$\omega$についての展開は行っていない．さらに，式~\eqref{int-ferm-S-1-WFZ-1}には$FZX^{q}_{-}d$に比例する項も含まれるはずであるが，これらは$\EuScript{K}$についての平均後には運動論方程式に寄与しないため，簡潔さのために省略している．

作用の残りの$S_{2}$部分（\ref{int-ferm-S-012}c）は，速い自由度$S_{2}\propto(\partial_{\mathbf{r}}\EuScript{K})^{2}$についてすでに2次であるため，$\hat{\mathcal{W}}=0$における値を用いてよい．
\begingroup\small
\begin{eqnarray}\label{int-ferm-S-2-FZ}
iS_{2}[F,Z]\!\!\!\!&=&\!\!\!\!4\pi\nu D\, \tr \Big\{
(\partial_{\mathbf{r}}\EuScript{K}^{cl})
(\partial_{\mathbf{r}}\EuScript{K}^{q})Z -
(\partial_{\mathbf{r}}\EuScript{K}^{cl}) F
(\partial_{\mathbf{r}}\EuScript{K}^{q})FZ-
(\partial_{\mathbf{r}}\EuScript{K}^{q})
(\partial_{\mathbf{r}}\EuScript{K}^{cl})Z +
(\partial_{\mathbf{r}}\EuScript{K}^{q}) F
(\partial_{\mathbf{r}}\EuScript{K}^{cl})FZ \nonumber
\\
&+&\!\!\! (\partial_{\mathbf{r}}\EuScript{K}^{cl})Z
(\partial_{\mathbf{r}}\EuScript{K}^{cl})F -
\frac{1}{2}(\partial_{\mathbf{r}}\EuScript{K}^{cl})
(\partial_{\mathbf{r}}\EuScript{K}^{cl})FZ  -
\frac{1}{2}(\partial_{\mathbf{r}}\EuScript{K}^{cl})
(\partial_{\mathbf{r}}\EuScript{K}^{cl})ZF \Big\}\,.
\end{eqnarray}
\endgroup

次の段階は，速い自由度であるディフューゾン$(d,\bar{d})$とゲージ場$(\EuScript{K}^{cl},\EuScript{K}^{q})$についてGauss積分を行うことである．作用の$S_1$部分について，式~\eqref{int-ferm-S-0-W}および\eqref{int-ferm-S-1-WFZ-1}を用いると，
\begin{equation}
\left\langle\exp\big(iS_{0}[\hat{\mathcal{W}}]+iS_{1}[\hat{\mathcal{W}},F,Z]\big)
\right\rangle_{\mathcal{W},\EuScript{K}}=
\exp\big(iS^{(1)}_{\mathrm{eff}}[F,Z]\big)\,,
\end{equation}
を得る．ここで
\begin{eqnarray}\label{int-ferm-S-eff-b}
iS^{(1)}_{\mathrm{eff}}[F,Z]=-4i\pi\nu\,
\tr\left\{\big(Dq^{2}\big)^{2}
\big[\mathcal{D}^{A}\qo\mathcal{V}^{R}\qo
-\mathcal{D}^{R}\qo\mathcal{V}^{A}\qo\big]
\mathcal{I}[F]Z\right\}\,.
\end{eqnarray}
である．$S^{(1)}_{\mathrm{eff}}$を式~\eqref{int-ferm-S-eff-b}の形で導くには，$\hat{\mathcal{W}}$について積分すると，式~\eqref{int-ferm-S-1-WFZ-1}の項$\tr\{\mathcal{I}[F]X^{q}_{+}\bar{d}\}$と$\tr\{ZX^{cl}_{-}d\}$が，$F$と$Z$の間に有効相互作用頂点$\big\langle\exp(iS_{1})\big\rangle_{\mathcal{W}}=
\exp\big(\tr\{\mathcal{I}[F]X^{q}_{+}\mathcal{D}^{A}ZX^{cl}_{-}\}\big)$を生むことに注意すればよい．これをさらに$\EuScript{K}$について平均する必要があり，その際には次の関係を用いる．
\begin{eqnarray}
\left\langle
X^{cl}_{-}(\mathbf{q},\omega)X^{q}_{+}(-\mathbf{q},-\omega)\right\rangle_{\EuScript{K}}=
-4D^{2}\left\langle\partial^{2}_{\mathbf{r}}\EuScript{K}^{cl}\qo
\partial^{2}_{\mathbf{r}}\EuScript{K}^{q}\qom\right\rangle_{\EuScript{K}}=
-2i\big(Dq^{2}\big)^{2}\mathcal{V}^{R}\qo\,.
\end{eqnarray}
最後の式は，式~\eqref{int-ferm-X}および\eqref{int-ferm-NLSM-K-cl-NonEq}と，式~\eqref{int-ferm-V-def}で与えられる相関関数から直接従う．

作用の$S_2$部分については，式~\eqref{int-ferm-S-2-FZ}を用いると
\begin{equation}
\big\langle\exp\big(iS_{2}[F,Z]\big)\big\rangle_{\EuScript{K}}=
\exp\big(iS^{(2)}_{\mathrm{eff}}[F,Z]\big)\,,
\end{equation}
を得る．ここで
\begin{equation}\label{int-ferm-S-eff-a}
iS^{(2)}_{\mathrm{eff}}[F,Z]=2i\pi\nu\,\tr\left\{Dq^{2}[\mathcal{V}^{R}\qo-
\mathcal{V}^{A}\qo]\mathcal{I}[F]Z\right\}\,.
\end{equation}
である．式~\eqref{int-ferm-S-eff-a}を導くには，ゲージ場の相互作用伝播関数~\eqref{int-ferm-K-corr-fun}を用い，引数が一致するKeldysh成分に対して，次の\textit{準平衡}の揺動散逸定理（FDT）を適用する必要がある．
\begin{equation}\label{int-ferm-quasi-FDT}
\mathcal{V}^{K}(\mathbf{r},\mathbf{r},\tau)=B_{\omega}(\mathbf{r},\tau)
\sum_{q}\big[\mathcal{V}^{R}\qo-\mathcal{V}^{A}\qo\big]\,,
\end{equation}
この関係は，$F_{\epsilon}(\mathbf{r},\tau)$が空間スケール$L_{T}=\sqrt{D/T}$上でゆっくり変化する限り，非平衡条件のもとでも成り立つ（これは$F_{\epsilon}(\mathbf{r},\tau)$の勾配が小さいことを意味する）．式~\eqref{int-ferm-quasi-FDT}に対する補正は$\propto\omega\int\d\mathbf{r}'\mathcal{D}^{R}(\mathbf{r}-\mathbf{r}',\omega)
\partial_{\tau}B_{\omega}(\mathbf{r}',\tau)\partial_{\omega}
\mathcal{D}^{A}(\mathbf{r}'-\mathbf{r})$の形をしている．文献~\cite{KamenevAndreev}を参照されたい．

最後に，式~\eqref{int-ferm-S-0-FZ}の$S_{0}[F,Z]$と，式~\eqref{int-ferm-S-eff-b}および\eqref{int-ferm-S-eff-a}で与えられる作用の$S^{(1),(2)}_{\mathrm{eff}}[F,Z]$部分をまとめ，式~\eqref{int-ferm-KE-def}を用いると，次の運動論方程式に到達する．
\begin{equation}\label{int-ferm-KE}
D\partial^{2}_{\mathbf{r}}F_{\epsilon}(\mathbf{r},\tau)-\partial_{\tau}
F_{\epsilon}(\mathbf{r},\tau)=\mathcal{I}_{\mathrm{col}}[F]\,,
\end{equation}
ここで衝突積分は
\begin{equation}\label{int-ferm-col-int}
\mathcal{I}_{\mathrm{coll}}[F]=\sum_{q}\int\frac{\d\omega}{2\pi}\mathcal{M}
\qo\Big[1 - F_{\epsilon-\omega}(\mathbf{r},\tau)
F_{\epsilon}(\mathbf{r},\tau) + B_{\omega}(\mathbf{r},\tau)
[F_{\epsilon-\omega}(\mathbf{r},\tau)-F_{\epsilon}(\mathbf{r},\tau)]\Big]\,,
\end{equation}
で与えられ，その核は
\begin{equation}
\mathcal{M}\qo=-iDq^2\Big\{\big[\mathcal{V}^{R}\qo-\mathcal{V}^{A}\qo\big]-2Dq^2
\big[\mathcal{D}^{A}\qo\mathcal{V}^{R}\qo-\mathcal{D}^{R}\qo\mathcal{V}^{A}\qo\big]\Big\}\,.
\end{equation}
である．この方程式は，ゲージ場伝播関数$\mathcal{V}^{R(A)}\qo$が，ディフューゾンとRPAで遮蔽された相互作用を用いて$\mathcal{V}^{R}\qo=-\big[\mathcal{D}^{R}\qo\big]^2U^{R}_{RPA}\qo$と書け，先進成分についても同様であることに注意すれば簡単にできる．これは式~\eqref{int-ferm-U-RPA}から直接従う．この形の$\mathcal{V}^{R(A)}\qo$を用いて少し計算すると，相互作用核$\mathcal{M}\qo$は次の形に簡約される．
\begin{equation}\label{int-ferm-col-int-kernel}
\mathcal{M}\qo=2\,\mathrm{Re}[\mathcal{D}^{R}\qo]\,\mathrm{Im}[U^{R}_{RPA}\qo]\,.
\end{equation}

フェルミ分布関数を通常のように$\mathbf{n}_{\epsilon}(\mathbf{r},\tau)=(1-F_{\epsilon}(\mathbf{r},\tau))/2$と選ぶと，衝突積分~\eqref{int-ferm-col-int}を「流出」と「流入」の緩和項を持つ通常の形に書き直せる．実際，式~\eqref{int-ferm-integral-2}を用いると，式~\eqref{int-ferm-KE}の右辺を恒等的に次のように書き直せる~\cite{Altshuler,AltshulerAronov-KE}．
\begin{equation}\label{int-ferm-col-int-in-out}
\mathcal{I}_{\mathrm{coll}}[\mathbf{n}]=\sum_{q}\iint^{+\infty}_{-\infty}\d\omega\d\epsilon'\,
\mathbb{K}\qo\,\Big[\mathbf{n}_{\epsilon}
\mathbf{n}_{\epsilon'-\omega}(1-\mathbf{n}_{\epsilon'})(1-\mathbf{n}_{\epsilon-\omega})
-\mathbf{n}_{\epsilon'}\mathbf{n}_{\epsilon-\omega}
(1-\mathbf{n}_{\epsilon})(1-\mathbf{n}_{\epsilon'-\omega})\Big]\,,
\end{equation}
ここで衝突核は$\mathbb{K}\qo=2\mathcal{M}\qo/\pi\omega$である．

運動論方程式の一般的な構造について，いくつかの重要な点を議論する必要がある．（i）式~\eqref{int-ferm-S-1-WFZ-1}で無視した項$\tr\{ZFX^{q}_{-}d\}$は，$\hat{\mathcal{W}}$について積分した後，$\tr\big\{\mathcal{I}[F]X^{q}_{+}\mathcal{D}^{A}ZFX^{q}_{-}\big\}$の形の有効頂点を生む．これは$\langle
X^{q}_{\pm}X^{q}_{\pm}\rangle_{\EuScript{K}}\equiv0$であるため，$\EuScript{K}$について平均すると確かに消える．したがって，この項が衝突積分に追加の項を生み出すことはない．（ii）衝突積分の導出全体を通して，例えば式~\eqref{int-ferm-quasi-FDT}におけるように，分布関数の空間微分$\partial_{\mathbf{r}}
F_{\epsilon}(\mathbf{r},\tau)$および時間微分$\partial_{\tau}F_{\epsilon}(\mathbf{r},\tau)$を一貫して無視した．これは，$F_{\epsilon}(\mathbf{r},\tau)$がゆっくり変化する空間スケールが存在する限り正当化される．実際，分布関数の勾配を陽に残すと，衝突積分の弾性部分に寄与する~\cite{Zala,CatelaniAleiner}．（iii）有効作用には，分布関数の量子成分に線形な項だけを残した．しかし，$Z_{\epsilon}(\mathbf{r},\tau)$に2次の項も存在する．これらの項は分布関数の揺らぎを生み，いわゆる\textit{確率的運動論方程式}，またはそれと等価なBoltzmann--Langevin運動論につながる~\cite{Kogan,KoganShulman,Nagaev}．最近，$Z^{2}_{\epsilon}(\mathbf{r},\tau)$の項を残したKeldyshの$\sigma$模型が，有効なBoltzmann--Langevin記述と等価であることが示された~\cite{Gutman,Bagrets}．（iv）式~\eqref{int-ferm-col-int}と同様の衝突積分は，文献~\cite{KamenevAndreev}においてKeldyshの$\sigma$模型の形式の中で導かれた．しかし，有効作用の$S^{(1)}_{\mathrm{eff}}$部分が見落とされていたため，そこで得られた衝突積分の核は，普遍的極限$U^{-1}_{0}\to0$でのみ正しい．式~\eqref{int-ferm-col-int-kernel}から，$U^{-1}_{0}\to0$に対して$\mathcal{M}\qo$が$\mathcal{M}\qo=-\frac{2}{\nu}\,\mathrm{Im}\big[\mathcal{D}^{R}\qo\big]$へ簡約されることがわかる．これは文献~\cite{KamenevAndreev}の結果である．（v）最後に，ここでの議論はスピン自由度を含むように一般化できる．対応する運動論方程式と衝突核は，文献~\cite{DimitrovaKtavtsov,BurmiChelk}で得られた．

%$$$$$$$$$$$$$$$$$$$$$$$$$$$$$$$$$$$$$$$$$$$$$$$$$$$$$$$$$$$$$$$$$$$$$$$$$$$$$$$$$$$$$$$$$$$$$$$$$$$$$$$$$$$
\subsectionlink{応用III：乱れた金属における相互作用効果}\label{app_Part-III}
\subsubsection{ゼロバイアス異常}\label{app_Part-III-1}

第~\ref{app_Part-II}節では，非相互作用系の$\sigma$模型を適用できる例をいくつか議論した．ここからは相互作用効果を考える．最初に取り上げるのは，自由電子の裸の1粒子状態密度$\nu$がCoulomb相互作用によって変化する問題である．この問題はAltshuler，Aronov，Leeによって研究された~\cite{AltshulerAronov,AAL,AltshulerAronov-review}．彼らの元の研究で計算されたのは相互作用による最低次の補正だけであったが，ゼロバイアス異常の扱いを摂動論を超えて拡張できる~\cite{Finkel'stein,Nazarov-ZBA,LevitovShytov,Kopietz}．ここでは文献~\cite{KamenevAndreev}のσ模型による計算に従う．

我々が求めたいのは，空間点が一致する1粒子Green関数
\begin{equation}
\mathcal{G}^{ab}(t-t')=-i\big\langle\big\langle\psi_{a}\rt
\bar{\psi}_{b}(\mathbf{r},t')\big\rangle\big\rangle\,,
\end{equation}
である．ここで$\langle\langle\ldots\rangle\rangle$は，量子平均と乱れ平均の双方を表す．フェルミ演算子の双線形結合に直接結合する対応したソース項を作用に導入することで，このGreen関数を計算できる．第~\ref{sec_NLSM}節と同じ代数操作に従い，Keldysh回転と乱れ平均を行うと，このソース項は式~\eqref{NLSM-action}の対数の中に入る．この式をソースで微分してからソースをゼロに置くと，Green関数について
\begin{equation}\label{Part-III-G-def}
\hat{\mathcal{G}}(t-t') = \int\!\! \D[\Phi]\,
\exp\big(i\Tr\{\vec{\Phi}^{T}U^{-1}_{0}\hat{\sigma}_{x}\vec{\Phi}\}\big)
\int\D[\hat Q]
\left[\hat{G}^{-1}+\frac{i}{2\tauel}\hat{Q}+
\hat{\Phi}\right]^{-1}_{tt',\mathbf{r}\mathbf{r}}\, \exp\big(iS[\hat Q,\Phi]\big)\,.
\end{equation}
を得る．$\hat Q$行列についての積分を鞍点近似で評価し，停留点のまわりの質量を持つ揺らぎと質量を持たない揺らぎをともに無視する．すると式~\eqref{int-ferm-NLSM-saddle-point-eq}より，指数関数の前の因子は単に$-i\pi\nu\underline{\hat{Q}}_{tt'}$となる．鞍点では，$\hat Q$行列は式~\eqref{int-ferm-NLSM-Q-gauge-general}で与えられる．その結果，式~\eqref{Part-III-G-def}に対して次の表示が得られる．
\begin{equation}\label{Part-III-G}
\hat{\mathcal{G}}(t-t')=-i\pi\nu\int\D[\Phi]
\exp\left(i\Tr\{\vec{\Phi}^{T}\hat{U}^{-1}_{RPA}\vec{\Phi}\}\right)
\exp\left(i\hat{\EuScript{K}}\rt\right)\hat{\Lambda}_{t-t'}
\exp\left(-i\hat{\EuScript{K}}(\mathbf{r},t')\right)\,.
\end{equation}
$\hat{\EuScript{K}}$は式~\eqref{int-ferm-NLSM-K-q-cl}で与えられる$\hat{\Phi}$の線形汎関数なので，残る汎関数積分はGauss積分である．これを計算するため，ゲージ場の位相因子を次のように書き直す\footnote{式~\eqref{Part-III-Phase-factors}は次の性質に基づいている．Pauli行列の線形結合の任意の関数$f(a+\mathbf{b\hat{\sigma}})$を考える．ここで$a$は任意の数，$\mathbf{b}$は任意のベクトルである．すると$f(a+\mathbf{b}\hat{\sigma})=A+\mathbf{B}\hat{\sigma}$が成り立つ．ここで$A$は新しい数，$\mathbf{B}$は新しいベクトルである．これを見るには，$z$軸をベクトル$\mathbf{b}$の方向に取ればよい．このとき演算子$a+\mathbf{b}\hat{\sigma}$の固有値は$a\pm b$であり，演算子$f(a+\mathbf{b\hat{\sigma}})$の対応する固有値は$f(a\pm b)$となる．したがって$A=\frac{1}{2}[f(a+b)+f(a-b)]$および$\mathbf{B}=\frac{\mathbf{b}}{2b}[f(a+b)-f(a-b)]$という結論を得る．}．
\begin{equation}\label{Part-III-Phase-factors}
e^{\pm i\EuScript{K}^{\alpha}\gamma^{\alpha}}=\frac{1}{2}
\left[\,e^{\pm i(\EuScript{K}^{cl}+\EuScript{K}^{q})}+e^{\pm
i(\EuScript{K}^{cl}-\EuScript{K}^{q})}\right]\hat{\gamma}^{cl}+
\frac{1}{2} \left[\,e^{\pm
i(\EuScript{K}^{cl}+\EuScript{K}^{q})}-e^{\pm
i(\EuScript{K}^{cl}-\EuScript{K}^{q})} \right]\hat{\gamma}^{q}\,.
\end{equation}
式~\eqref{Part-III-Phase-factors}を用いて式~\eqref{Part-III-G}のGauss積分を行うと，結果は次の形で表すと便利である．
\begin{equation}\label{Part-III-G-1}
\hat{\mathcal{G}}(t)=-i\pi\nu\sum_{\alpha\beta}
\big(\hat{\gamma}^{\alpha}\,\hat{\Lambda}_{t}\,\hat{\gamma}^{\beta}\big)\,
\mathbb{B}^{\alpha\beta}(t)\,,
\end{equation}
ここで補助伝播関数$\mathbb{B}^{\alpha\beta}(t)$は，標準的なボース型の構造［例えば式~\eqref{int-ferm-V-matrix}と同じ構造］を持ち，
\begin{subequations}\label{Part-III-B-RAK}
\begin{equation}
\hskip-.75cm
\mathbb{B}^{R(A)}(t)=i\exp\big(i[\mathcal{V}^{K}(t)-\mathcal{V}^{K}(0)]/2\big)
\sin\big(\mathcal{V}^{R(A)}(t)/2\big)\,,
\end{equation}
\begin{equation}
\mathbb{B}^{K}(t)=\exp\big(i[\mathcal{V}^{K}(t)-\mathcal{V}^{K}(0)]/2\big)
\cos\big([\mathcal{V}^{R}(t)-\mathcal{V}^{A}(t)]/2\big)\,.
\end{equation}
\end{subequations}
となる．式~\eqref{int-ferm-V-matrix}および\eqref{int-ferm-V-RAK}で定義したゲージ場の伝播関数$\hat{\mathcal{V}}\rt$は，空間点が一致した形で式~\eqref{Part-III-B-RAK}に入る．
\begin{equation}\label{Part-III-V-t}
\hat{\mathcal{V}}(t)=\int\frac{\d\omega}{2\pi}\exp(-i\omega
t)\sum_{q}\hat{\mathcal{V}}\qo\,.
\end{equation}
Green関数\eqref{Part-III-G-1}がわかれば，標準的な定義に従って状態密度を求められる．
\begin{equation}\label{Part-III-DOS-def}
\nu(\epsilon)=\frac{i}{2\pi}\big[\mathcal{G}^{R}(\epsilon)
-\mathcal{G}^{A}(\epsilon)\big]\,.
\end{equation}
熱平衡ではGreen関数はFDTを満たす（式~\eqref{fermion-FDT}を参照）．これと関係式$\mathcal{G}^{K}(\epsilon)=\mathcal{G}^{>}(\epsilon)+\mathcal{G}^{<}(\epsilon)$および$\mathcal{G}^{>}(\epsilon)=-\exp(\epsilon/T)\mathcal{G}^{<}(\epsilon)$を用いると，式~\eqref{Part-III-DOS-def}を等価な形
\begin{equation}\label{Part-III-DOS}
\nu(\epsilon)=\frac{i}{2\pi}\ \mathcal{G}^{>}
(\epsilon)[1+\exp(-\epsilon/T)]\,.
\end{equation}
に書き直せる．式~\eqref{Part-III-G-1}を用いると，greater（lesser）Green関数$\mathcal{G}^{>(<)}$を，補助伝播関数$\mathbb{B}^{>(<)}$の対応する成分と関係づけられる．
\begin{equation}\label{Part-III-G-B}
\mathcal{G}^{>(<)}(t)=-i\pi\nu\Lambda^{>(<)}_{t}\mathbb{B}^{>(<)}(t)\,.
\end{equation}
後者を具体的に求めると，
\begin{equation}\label{Part-III-B-lesser-greater}
\mathbb{B}^{>(<)}(t)=\frac{1}{2}\exp\left(\int\frac{\d\omega}{2\pi}
\left[\coth\frac{\omega}{2T}(1-\cos\omega t)\pm i\sin\omega t\right]
\mathrm{Im}\sum_{q}\mathcal{V}^{R}\qo\right)\,,
\end{equation}
となる．ここでは式~\eqref{Part-III-B-RAK}と，ボソンのFDTの関係式$\mathbb{B}^{R}(t)-\mathbb{B}^{A}(t)=\mathbb{B}^{>}(t)-\mathbb{B}^{<}(t)$および$\mathbb{B}^{K}(t)=\mathbb{B}^{>}(t)+\mathbb{B}^{<}(t)$を用いた．最後に，式~\eqref{Part-III-DOS}と式~\eqref{Part-III-G-B}を組み合わせると，状態密度について
\begin{equation}\label{Part-III-DOS-fin}
\nu(\epsilon)=\frac{\nu}{\tanh(\epsilon/2T)}\int\d t\ F_{t}\
\mathbb{B}^{K}(t)\exp(i\epsilon t)\,.
\end{equation}
を得る．式~\eqref{Part-III-B-lesser-greater}を相互作用の1次まで展開して$\mathcal{V}\qo$とし，式~\eqref{Part-III-DOS-fin}へ代入すると，ゼロバイアス異常についてのAltshulerとAronovの結果~\cite{AltshulerAronov}が再現される．

ここでは非摂動的な結果である式~\eqref{Part-III-B-lesser-greater}および\eqref{Part-III-DOS-fin}の解析を，零温度の場合に限る．$T=0$に対して$F_{t}=(i\pi t)^{-1}$であることに注意すると，
\begin{eqnarray}\label{Part-III-DOS-T0}
\nu(\epsilon)=\frac{\nu}{\pi}\int\d t\ \frac{\sin|\epsilon|t}{t}\
\exp\left(\int^{\infty}_{0}\frac{\d\omega}{\pi}\
\mathrm{Im}\sum_{q}\mathcal{V}^{R}\qo(1-\cos\omega
t)\right)\nonumber\\
\times\cos\left(\int^{\infty}_{0}\frac{\d\omega}{\pi}\
\mathrm{Im}\sum_{q}\mathcal{V}^{R}\qo\sin\omega t\right)\,.
\end{eqnarray}
を得る．2次元の場合，$U_{0}=2\pi
e^{2}/q$とした式~\eqref{int-ferm-V-RAK}から
\begin{equation}\label{Part-III-exponent-integral}
\int^{+\infty}_{0}\frac{\d\omega}{\pi}\sum_{q}\mathrm{Im}
\big[\mathcal{V}^{R}\qo\big]\left(\begin{array}{c}1-\cos\omega t \\
\sin\omega t\end{array}\right)=-\frac{1}{8\pi^{2}g}\left\{
\begin{array}{l}\ln(t/\tauel)\ln(t\tauel\omega^{2}_{0})+2\mathbb{C}\ln(t\omega_{0})
\\ \pi\ln(t\omega_{0})\end{array}\right.\,,
\end{equation}
が導かれる．ここで$g=\nu D$は無次元コンダクタンス，$\omega_{0}=D\kappa^{2}$，$\kappa^{2}=2\pi e^{2}\nu$はThomas--Fermi遮蔽半径の逆数であり，$\mathbb{C}=0.577...$はEuler定数である．式~\eqref{Part-III-G-def}の汎関数積分を計算する際に$\hat Q$行列の揺らぎ$\hat{\mathcal{W}}$を無視したため，得られた結果~\eqref{Part-III-DOS-T0}には，$d=2$において$\sim g^{-1}\ln(t/\tau_{\mathrm{el}})$の程度の補正は含まれない．第~\ref{app_Part-III-2}節を参照されたい．したがって，式~\eqref{Part-III-DOS-T0}が信頼できるのは，$\epsilon$が小さすぎず，$(8\pi^{2}g)^{-1}\ln(\epsilon\tauel)^{-1}\ll 1$が成り立つ場合に限られる．ただし，$\ln^{2}(t/\tauel)$の項は上の手続きで正しく取り入れられている．さらに$g^{-1}\ln(\omega_{0}\tauel)\ll1$なら，式~\eqref{Part-III-DOS-T0}の時間積分を停留点法で実行でき，その結果
\begin{equation}
\nu(\epsilon)=\nu\exp\left\{-\frac{1}{8\pi^{2}g}\ln(|\epsilon|\tauel)^{-1}
\ln(\tauel\omega^{2}_{0}/|\epsilon|)\right\}\,.
\end{equation}
を得る．このようにして，1粒子Green関数に対する異常に発散する$\propto\ln^{2}(\epsilon\tauel)$の項を非摂動的に再総和することができた．状態密度の非摂動的な式は，本質的には1粒子Green関数がゲージ不変でないことに由来する．上の計算は，本質的には，ほぼ純粋なゲージ揺らぎである電位の揺らぎによるDebye--Waller因子~\cite{Finkel'stein94}の計算である．式~\eqref{Part-III-G}と比較されたい．ゲージ不変な物理量（例えば伝導率）は位相因子を持たず，したがってこの精度の範囲では相互作用の影響を受けない（$\hat Q$行列の揺らぎを残す必要がある．次節を参照）．


%$$$$$$$$$$$$$$$$$$$$$$$$$$$$$$$$$$$$$$$$$$$$$$$$$$$$$$$$$$$$$$$$$$$$$$$$$$$$$$$$$$$$$$$$$$$$$$$$$$$$$$$$$$$
\subsubsection{Altshuler--Aronov補正}\label{app_Part-III-2}

ここでは，相互作用が本質的となる別の例として，乱れた金属のDrude伝導率$\sigma_{D}$に対する電子間相互作用補正$\delta\sigma_{\mathrm{AA}}$を考える~\cite{AltshulerAronov,AAL,AltshulerAronov-review}．相互作用する乱れた電子液体の状態密度を扱った前の例（第~\ref{app_Part-III-1}節）とは異なり，試行鞍点$\Qk=\hat{\Lambda}$の段階では相互作用は伝導率補正に影響せず，揺らぎ$\hat{\mathcal{W}}$を残さなければならない．以下では，作用~\eqref{int-ferm-S-012}の$\hat{\mathcal{W}}$，\eqref{NLSM-S-W}および\eqref{NLSM-d-d}による展開で，ゼロでない最低次に考察を限り，相互作用補正$\delta\sigma_{\mathrm{AA}}$を生み出す作用の項を同定する．

式（\ref{int-ferm-S-012}b）で与えられる作用の$S_{1}[\hat
Q_{\EuScript{K}},\partial_{\mathbf{r}}\EuScript{K}]$部分から出発する．揺らぎ$\hat{\mathcal{W}}$の1次まででは，
\begin{equation}\label{Part-III-S1-W}
iS_{1}[\hat{\mathcal{W}},\partial_{\mathbf{r}}\EuScript{K}]=
-\frac{i\pi\nu}{2}\Tr\left\{\left[D\partial^{2}_{\mathbf{r}}\EuScript{K}^{\alpha}
\big(\hat{\Lambda}\hat{\gamma}^{\alpha}\hat{\Lambda}-\hat{\gamma}^{\alpha}\big)
+\big(\Phi^{\alpha}-\partial_{t}\EuScript{K}^{\alpha}\big)
\big(\hat{\gamma}^{\alpha}\hat{\Lambda}-\hat{\Lambda}\hat{\gamma}^{\alpha}\big)\right]
\hat{W}\right\}\,,
\end{equation}
を得る．ここで$\hat{W}=\hat{\mathcal{U}}\circ\hat{\mathcal{W}}\circ\hat{\mathcal{U}}^{-1}$である．式（\ref{NLSM-Q-U-W-parametrization}）および（\ref{NLSM-U-W}）を参照されたい．熱平衡では$iS_{1}[\hat{\mathcal{W}},\partial_{\mathbf{r}}\EuScript{K}]\equiv0$であることに注意する．実際，式~\eqref{Part-III-S1-W}の右辺の角括弧内の式は，$\hat{\EuScript{K}}[\Phi]$汎関数を決定するために用いた式~\eqref{int-ferm-NLSM-Phi-K-1}と一致する．平衡では，$\hat{\EuScript{K}}[\Phi]$を適切に選ぶことで式~\eqref{int-ferm-NLSM-Phi-K-1}を解けた．式~\eqref{int-ferm-NLSM-Phi-K-matrix}を参照されたい．$\hat{\Phi}$場の各実現に対して鞍点を求めようとした動機は，まさに$\hat{\mathcal{W}}$に線形な項を打ち消すことにあった．厳密な鞍点を見いだせなかったため，このような項は現れるが，それは$\partial_{\mathbf{r}}\hat{\EuScript{K}}$の2次からである．これらの項は，作用の$S_{2}[\Qk,\partial_{\mathbf{r}}\EuScript{K}]$部分から生じる．式（\ref{int-ferm-S-012}c）を$\hat{\mathcal{W}}$の1次まで展開すると，
\begin{eqnarray}\label{Part-III-S2-W}
iS_{2}[\hat{\mathcal{W}},\partial_{\mathbf{r}}\EuScript{K}] &=& \frac{\pi\nu D}{2}\,
\Tr\left\{\partial_{\mathbf{r}}\EuScript{K}^{\alpha}(\epsilon_{1}-\epsilon_{2})
\left[\hat{\gamma}^{\alpha}\hat{\Lambda}_{\epsilon_{2}}
\hat{\gamma}^{\beta}\hat{\Lambda}_{\epsilon_{3}}-
\hat{\Lambda}_{\epsilon_{1}}\hat{\gamma}^{\alpha}
\hat{\Lambda}_{\epsilon_{2}}\hat{\gamma}^{\beta}\right]
\hat{W}_{\epsilon_{3}\epsilon_{1}}
\partial_{\mathbf{r}}\EuScript{K}^{\beta}(\epsilon_{2}-\epsilon_{3})\right\}
\nonumber  \\
&=& \pi\nu D\,
\Tr\left\{\partial_{\mathbf{r}}\vec{\EuScript{K}}^{T}(\epsilon_{1}-\epsilon_{2})
\left[M^{d}_{\epsilon_{2}}d_{\epsilon_{3}\epsilon_{1}}+
M^{\bar{d}}_{\epsilon_{1}\epsilon_{2}\epsilon_{3}}
\bar{d}_{\epsilon_{3}\epsilon_{1}}\right]
\partial_{\mathbf{r}}\vec{\EuScript{K}}(\epsilon_{2}-\epsilon_{3})\right\}\,,
\end{eqnarray}
を得る．ここで$\vec{\EuScript{K}}^{T}=\big(\EuScript{K}^{cl},\EuScript{K}^{q}\big)$という記法を用い，ディフューゾン$\{d,\bar{d}\}$とゲージ場$\EuScript{K}^{cl(q)}$との結合行列
\begin{equation}\label{Part-III-M}
\fitdisplay{
M^{d}_{\epsilon_{2}}=\left(\begin{array}{cc}0 & 0
\\ 0 & -2F_{\epsilon_{2}}\end{array}\right)\,,\qquad
M^{\bar{d}}_{\epsilon_{1}\epsilon_{2}\epsilon_{3}}=
\left(\begin{array}{ll}
2F_{\epsilon_{2}}-F_{\epsilon_{1}}-F_{\epsilon_{3}}&
1+F_{\epsilon_{1}}F_{\epsilon_{3}}-2F_{\epsilon_{2}}F_{\epsilon_{3}}\\
-1-F_{\epsilon_{1}}F_{\epsilon_{3}}+2F_{\epsilon_{2}}F_{\epsilon_{1}}\quad&
F_{\epsilon_{1}}+F_{\epsilon_{3}}-2F_{\epsilon_{1}}F_{\epsilon_{2}}F_{\epsilon_{3}}
\end{array}\right)\,.
}
\end{equation}
を導入した．ここで伝導率の一般式~\eqref{Part-II-sigma-def}を用い，伝導率に対するAltshuler--Aronov相互作用補正$\delta\sigma_{AA}$が次の式から得られることを示そう~\eqref{Part-III-S2-W}．
\begin{equation}\label{Part-III-sigma-def}
\delta\sigma_{\mathrm{AA}}=-\frac{e^{2}}{2}\lim_{\Omega\to 0}
\frac{1}{\Omega} \left\langle\frac{\delta^{2}}
{\delta\big(\partial_{\mathbf{r}} \EuScript{K}^{cl}(\Omega)\big)
\delta\big(\partial_{\mathbf{r}} \EuScript{K}^{q}(-\Omega)\big)}\
\exp\left(iS_{2}[\hat{\mathcal{W}},\partial_{\mathbf{r}}\EuScript{K}]\right)
\right\rangle_{\mathcal{W},\EuScript{K}}\, ,
\end{equation}
ここでは拡散モードと電位の揺らぎの両方について平均を取る．また，式~\eqref{Part-II-sigma-def}とは異なり，ここではベクトルポテンシャル$\mathbf{A}$そのものではなく$\partial_{\mathbf{r}}\EuScript{K}$で微分していることにも注意する．ベクトルポテンシャルとゲージ場は，ゲージ不変な組合せ~\eqref{int-ferm-NLSM-C}として作用~\eqref{int-ferm-action-interacting}に入るため，この2つの定義は同じである．

\begin{figure}
\begin{center}\includegraphics[width=10cm]{figures/fig14.pdf}\end{center}
\caption{伝導率に対するAltshuler--Aronov補正を生む，有効な4本脚頂点$\mathbb{V}_{\mathrm{AA}}$のダイアグラム表示．式~\eqref{Part-III-AA-Vertex}を参照されたい．\label{Fig-AA-vertex}}
\end{figure}
式~\eqref{NLSM-S-W}および\eqref{Part-III-S2-W}により，ディフューゾンモード$d$と$\bar d$の揺らぎを扱うGauss理論が得られ，式~\eqref{Part-III-sigma-def}の平均は直接実行できる．ディフューゾンモードについて積分すると，
\begin{equation}
\left\langle\exp\big(iS_{2}[\hat{\mathcal{W}},\partial_{\mathbf{r}}\EuScript{K}]\big)
\right\rangle_{\mathcal{W}}=
\exp\big(i\mathbb{V}_{\mathrm{AA}}[\EuScript{K}]\big)\,.
\end{equation}
を得る．このようにして，4つのゲージ場を結ぶ有効頂点$\big(\partial_{\mathbf{r}}\EuScript{K}\big)^4$が生成される．
\begin{eqnarray}\label{Part-III-AA-Vertex}
\mathbb{V}_{\mathrm{AA}}[\EuScript{K}]=4\pi\nu
D^{2}\Tr\left\{\,F_{\epsilon_{2}}(2F_{\epsilon_{4}}-F_{\epsilon_{1}}-F_{\epsilon_{3}})
\partial_{\mathbf{r}}\EuScript{K}^{q}(\mathbf{r},\epsilon_{1}-\epsilon_{2})
\partial_{\mathbf{r}}\EuScript{K}^{q}(\mathbf{r},\epsilon_{2}-\epsilon_{3})
\phantom{\mathcal{D}^{A}}\right. \nonumber\\
\left. \times\,
\mathcal{D}^{R}(\mathbf{r}-\mathbf{r}',\epsilon_{3}-\epsilon_{1})
\partial_{\mathbf{r}'}\EuScript{K}^{cl}(\mathbf{r}',\epsilon_{3}-\epsilon_{4})
\partial_{\mathbf{r}'}\EuScript{K}^{cl}(\mathbf{r}',\epsilon_{4}-\epsilon_{1})\right\}\,.
\end{eqnarray}
そのダイアグラム表示を図~\ref{Fig-AA-vertex}に示す．この頂点は，ディフューゾン伝播関数$\langle\bar{d}d\rangle_{\mathcal{W}}\propto\mathcal{D}^{A}$を用いて$d$と$\bar{d}$を対にすると，作用~\eqref{Part-III-S2-W}の$\Tr\{\partial_{\mathbf{r}}\EuScript{K}
M^{d}d\partial_{\mathbf{r}}\EuScript{K}\}$部分と$\Tr\{\partial_{\mathbf{r}}\EuScript{K}
M^{\bar{d}}\bar{d}\partial_{\mathbf{r}}\EuScript{K}\}$部分から生じる．係数$F_{\epsilon_{2}}$は行列$M^{d}$の$q-q$成分に由来し，式~\eqref{Part-III-AA-Vertex}における分布関数の組合せ$2F_{\epsilon_{4}}-F_{\epsilon_{1}}-F_{\epsilon_{3}}$は，行列$M^{\bar{d}}$の$cl-cl$成分である．$\mathbb{V}_{\mathrm{AA}}[\EuScript{K}]$を式~\eqref{Part-III-AA-Vertex}の形で書く際には，量子ゲージ場$\partial_{\mathbf{r}}\EuScript{K}^{q}$の数が最小となる寄与だけを残した．しかし，行列$M^{\bar{d}}$の4つの成分はすべてゼロでないため，原理的には$\mathbb{V}_{\mathrm{AA}}[\EuScript{K}]$には，4本または3本の脚が量子ゲージ場を担う寄与も含まれる．これらは，ショット雑音強度に対する対応した相互作用補正の計算に用いられる．詳細は文献~\cite{GutmanGefen}を参照されたい．

$\hat{\mathcal{W}}$についての平均を実行したので，次に$\mathbb{V}_{AA}[\EuScript{K}]$を式~\eqref{Part-III-sigma-def}に代入し，$\EuScript{K}$場を積分消去する．伝導率補正については
\begin{equation}\label{Part-III-sigma-AA}
\fitdisplay{
\delta\sigma_{\mathrm{AA}}= 4\pi e^{2}\nu D^{2}\,
\sum_{q}\iint\frac{\d\epsilon\d\omega}{4\pi^{2}}
\big(F_{\epsilon_{+}}+F_{\epsilon_{-}}\big)
\big(\partial_{\epsilon}F_{\epsilon_{+}}-
\partial_{\epsilon}F_{\epsilon_{-}}\big)
\mathcal{D}^{R}(\mathbf{q},\omega)
\left\langle\partial_{\mathbf{r}}\EuScript{K}^{cl}(\mathbf{q},\omega)
\partial_{\mathbf{r}}\EuScript{K}^{q}(-\mathbf{q},-\omega)
\right\rangle_{\EuScript{K}}\,,
}
\end{equation}
を得る．ここで新しい積分変数$\epsilon=(\epsilon_{3}+\epsilon_{1})/2$および$\omega=\epsilon_{3}-\epsilon_{1}$を導入した．$\EuScript{K}$について平均すると，$\delta\sigma_{\mathrm{AA}}$に対して図~\ref{Fig-AA}の2つのダイアグラムが生じる．これは，図~\ref{Fig-AA-vertex}に示した有効頂点の2本の外脚を相互作用伝播関数で対にすることで自然に得られる．相互作用が強い普遍的極限$U^{-1}_{0}\to0$では，伝播関数$\mathcal{V}^{R}\qo$は簡単な形になる．その結果，
\begin{equation}\label{Part-III-V}
\left\langle\partial_{\mathbf{r}}\EuScript{K}^{cl}\qo
\partial_{\mathbf{r}}\EuScript{K}^{q}\qom\right\rangle_{\EuScript{K}}=
\frac{iq^2}{2}\, \mathcal{V}^{R}\qo =
-\frac{i}{2\nu D}\, \frac{1}{Dq^{2} - i\omega}\,,
\end{equation}
を得る．これは式~\eqref{int-ferm-K-corr-fun}および\eqref{int-ferm-V-RAK}から従う．式~\eqref{Part-III-V}を式~\eqref{Part-III-sigma-AA}に代入し，$\epsilon$積分を行うと，
\begin{equation}
\frac{\delta\sigma_{\mathrm{AA}}}{\sigma_{D}} =
\frac{2i}{\pi\nu}\sum_{q} \int\d\omega\,
\frac{\partial}{\partial\omega}\left[\omega\coth\frac{\omega}{2T}\right]\,
\frac{1}{\big(Dq^{2} - i\omega\big)^{2}}\ .
\end{equation}
を得る．2次元では，この式は伝導率に対する対数的に発散する負の補正$\delta\sigma_{\mathrm{AA}}= -\frac{e^{2}}{2\pi^{2}}\ln(1/T\tauel)$を与える．ここで弾性散乱率$\tauel^{-1}$は，周波数$\omega$についての積分の上限カットオフとして入る．乱れた導体における相互作用補正の効果についての詳しい総説は，文献~\cite{AltshulerAronov-review}にある．文献~\cite{Zala}も参照されたい．
\begin{figure}
\begin{center}\includegraphics[width=10cm]{figures/fig15.pdf}\end{center}
\caption{伝導率に対する相互作用補正$\delta\sigma_{\mathrm{AA}}$のダイアグラム．これらは有効頂点$\mathbb{V}_{\mathrm{AA}}[\EuScript{K}]$において，古典的な脚1本と量子的な脚1本を外脚として残し，残りの2本を相互作用伝播関数$\mathcal{V}^{R}\qo$で結ぶことによって構成される．\label{Fig-AA}}
\end{figure}

%$$$$$$$$$$$$$$$$$$$$$$$$$$$$$$$$$$$$$$$$$$$$$$$$$$$$$$$$$$$$$$$$$$$$$$$$$$$$$$$$$$$$$$$$$$$$$$$$$$$$$$$$$$$
\subsubsection{緩和率}\label{app_Part-III-3}

第~\ref{sec_int_ferm-5}節で議論した運動論方程式を用いると，エネルギー緩和率を求められる~\cite{AltshulerAronov-KE,AltshulerAronov-review,Schmid,Blanter}．式~\eqref{int-ferm-col-int-in-out}の衝突積分の「流出」項に注目すると，エネルギー$\epsilon$の電子について，流出緩和率を次のように導入できる．
\begin{equation}\label{Part-III-out-rate}
\frac{1}{\tau_{\mathrm{out}}(\epsilon)}=-\sum_{q}\int\d\omega\d\epsilon'\,
\mathbb{K}\qo\,
\mathbf{n}_{F}(\epsilon)[1-\mathbf{n}_{F}(\epsilon-\omega)]
\mathbf{n}_{F}(\epsilon')[1-\mathbf{n}_{F}(\epsilon+\omega)]\,,
\end{equation}
ここではすべての電子分布をFermi関数に置き換えた．これは，$\mathbf{n}_{\epsilon}$の平衡値$\mathbf{n}_{F}(\epsilon)$からの小さな（線形な）ずれに関心がある場合に適切である．式~\eqref{Part-III-out-rate}は零温度$T=0$で大幅に簡単になる．実際，Fermi分布関数によりエネルギー積分は$-\omega<\epsilon'<0$および$0<\omega<\epsilon$という2つの範囲に制限され，その範囲ではすべての占有数の積は単に1となる．相互作用が強い普遍的極限$U^{-1}_{0}\to 0$では，核は次の形を取る．式~\eqref{int-ferm-col-int-kernel}を参照されたい．
\begin{equation}
\mathbb{K}\qo=-\frac{4}{\pi\nu}\frac{1}{(Dq^2)^{2}+\omega^2}\,.
\end{equation}
式~$\mathbb{K}\qo$を式~\eqref{Part-III-out-rate}に代入すると，流出緩和率について次の式を得る．
\begin{equation}
\frac{1}{\tau_{\mathrm{out}}(\epsilon)}=\frac{4}{\pi\nu}\sum_{q}\int^{|\epsilon|}_{0}\d\omega
\int^{0}_{-\omega}\d\epsilon'\frac{1}{(Dq^2)^{2}+\omega^2}=\frac{|\epsilon|}{4\pi
g}\,,
\end{equation}
ここで$g=\nu D$であり，運動量積分は2次元の場合について実行した．任意の次元$d$では，流出緩和率のエネルギー依存性は$\tau^{-1}_{\mathrm{out}}(\epsilon)\propto(1/\nu_{d})(\epsilon/D)^{d/2}$で与えられる．詳細は文献~\cite{AltshulerAronov-review}を参照されたい．

%$$$$$$$$$$$$$$$$$$$$$$$$$$$$$$$$$$$$$$$$$$$$$$$$$$$$$$$$$$$$$$$$$$$$$$$$$$$$$$$$$$$$$$$$$$$$$$$$$$$$$$$$$$$
\subsubsection{3次のドラッグ効果}\label{app_Part-III-4}

第~\ref{app_Part-I-3}節でCoulombドラッグを議論した際，この効果は回路間相互作用の2次ですでに現れ，粒子・正孔非対称性が決定的な役割を果たすことを強調した．低温の線形応答では，ドラッグコンダクタンスは温度の2乗に比例する．式（\ref{Part-I-ID-linear-1}）を参照されたい．ここでは，層間相互作用の3次に現れるドラッグコンダクタンスへの寄与を議論する．これは相互作用強度については高次の寄与であるが，電子・正孔非対称性に依存しない（バルク系では，この非対称性はFermiエネルギー付近の分散関係の曲率に由来するため非常に小さい）．この3次のドラッグが温度に依存せず，十分な低温では支配的な効果となり得ることを示す~\cite{Drag-Third-order}．計算上，3次の寄与は4本脚頂点（図~\ref{Fig-AA-vertex}と対応する式~\eqref{Part-III-AA-Vertex}を参照）から生じる．この頂点は，各層における電子・正孔対の励起を通じて誘起される電磁場の非線形相互作用を記述する．

文献~\cite{Drag-Third-order}に従って2次元電子気体の二層系を考え，非線形σ模型（NLSM）を適用してドラッグ伝導率を計算する．一般式~\eqref{Part-II-sigma-def}と式~\eqref{Part-III-AA-Vertex}を用いると，ドラッグ伝導率は次のように定義される．
\begin{equation}\label{Part-III-sigma-drag}
\sigma_{\mathrm{drag}}=-\frac{e^2}{2}\lim_{\Omega\to0}\frac{1}{\Omega}
\left\langle\frac{\delta\mathbb{V}_{\mathrm{AA}}[\EuScript{K}]}
{\delta\big(\partial_{\mathbf{r}}\EuScript{K}^{cl}_{1}(\Omega)\big)}
\frac{\delta\mathbb{V}_{\mathrm{AA}}[\EuScript{K}]}
{\delta\big(\partial_{\mathbf{r}}\EuScript{K}^{q}_{2}(-\Omega)\big)}
\right\rangle_{\EuScript{K}}\,,
\end{equation}
ここで添字$1,2$は，第~\ref{app_Part-I-3}節の記法に従って，それぞれ駆動層とドラッグされる層を指す．揺らぐゲージ場$\EuScript{K}$についての平均は，相関関数
\begin{equation}\label{Part-III-V-drag}
\mathcal{V}_{ab}^R\qo=2i\big\langle
\EuScript{K}^{cl}_{a}\qo\EuScript{K}^{q}_{b}\qom\big\rangle_{\EuScript{K}}=
\frac{q^2 {U}^{R}_{ab}\qo}{\big(D_{a}q^2-i\omega\big)\big(D_{b}q^{2}-i\omega\big)}\,,
\end{equation}
を用いて実行する．ここで$a,b=(1,2)$であり，$U^R_{ab}\qo$は，RPAで計算した遮蔽された層内および層間相互作用の遅延成分からなる$2\times2$行列である．これは行列Dyson方程式$\hat{U}^R=\hat{U}_{0}+\hat{U}_{0}\hat{\Pi}^R\hat{U}^R$の解であり，ここで
\begin{equation}\label{Part-III-RPA}
\hat{U}_{0}=\frac{2\pi e^{2}}{q}\!\left(\begin{array}{cc}
1 & e^{-qd}
\\ e^{-qd} & 1
\end{array}\right)\,,\qquad\quad
\hat{\Pi}^R=\left(\begin{array}{cc} \frac{\nu_{1}D_{1}q^{2}}
{D_{1}q^{2}-i\omega} & 0 \\ 0 & \frac{\nu_{2}D_{2}q^{2}}
{D_{2}q^{2}-i\omega}\end{array}\right)\,.
\end{equation}
である．$\hat{U}_{0}$行列の非対角成分は層間の裸のCoulomb相互作用を表し，$d$は層間距離である．また，分極演算子の行列$\hat{\Pi}^R\qo$は，層間トンネルが存在しないことを反映して対角的であることにも注意する．

\begin{figure}
\begin{center}\includegraphics[width=10cm]{figures/fig16.pdf}\end{center}
\caption{波線で示した層間相互作用$\mathcal{V}^{R}_{12}\qo$の3次における，ドラッグ伝導率$\sigma_{\mathrm{drag}}$の2つのダイアグラム．層内の拡散伝播関数$\mathcal{D}^{R}_{a}\qo=(D_{a}q^2-i\omega)^{-1}$は梯子で示す．\label{Fig-Drag-3rd}}
\end{figure}

これで，3次のドラッグ伝導率を計算する準備が整った．式~\eqref{Part-III-AA-Vertex}を式~\eqref{Part-III-sigma-drag}に代入し，式~\eqref{Part-III-RPA}を用いて平均すると，ドラッグ伝導率について次の式が得られる．
\begin{eqnarray}\label{Part-III-drag}
&&\hskip-1cm
\sigma_{\mathrm{drag}}=32e^{2}T\nu_{1}\nu_{2}D^{2}_{1}D^{2}_{2}\int^{\infty}_{0}
\frac{\mathrm{d}\omega\mathrm{d}\omega'}{4\pi^{2}}\,\,
\mathcal{H}_{1}(\omega,\omega')\mathcal{H}_{2}(\omega,\omega')\nonumber\\
&&\hskip-1cm
\times \sum_{q,q'}\mathrm{Im}\left[\mathcal{D}_{1}^R\qo\mathcal{D}_{2}^R\qo
\mathcal{V}_{12}^R\qo \mathcal{V}_{12}^R\left({\mathbf{q}\over
2}-\mathbf{q}',{\omega\over 2}-\omega'\right)
\mathcal{V}_{12}^R\left({\mathbf{q}\over 2}+\mathbf{q}',{\omega\over
2}+\omega'\right)\right]\,.
\end{eqnarray}
2つの関数$\mathcal{H}_{1}(\omega,\Omega)$および$\mathcal{H}_{2}(\omega,\Omega)$は，それぞれ能動層と受動層で，速い電子エネルギー$\varepsilon$について積分することから生じる（図~\ref{Fig-AA}）．直流極限では，これらは次の式で与えられる．
\begin{subequations}\label{Part-III-Spectral-F}
\begin{equation}
\mathcal{H}_{1}(\omega,\omega')=2-\mathcal{B}(\omega'+\omega/2)
-\mathcal{B}(\omega'-\omega/2)+\mathcal{B}(\omega)\,,
\end{equation}
\begin{equation}
\hskip-.5cm
\mathcal{H}_{2}(\omega,\omega')=T\frac{\partial}{\partial\omega'}
\left[\mathcal{B}(\omega'+\omega/2)-\mathcal{B}(\omega'-\omega/2)\right]\,,
\end{equation}
\begin{equation}
\hskip-4.3cm
\mathcal{B}(\omega)=\frac{\omega}{T}\coth\left(\frac{\omega}{2T}\right)\,.
\end{equation}
\end{subequations}
対応するダイアグラムは，各層について1つずつ，図~\ref{Fig-AA-vertex}の頂点を2つ用いて構成される．図~\ref{Fig-Drag-3rd}を参照されたい．$\langle \EuScript{K}^{q}_{a} \EuScript{K}^{q}_b \rangle=0$であるため，伝播関数$\mathcal{V}_{ab}\qo$を用いてこれらを結ぶ方法は2通りしかないことがわかる．

以下では2つの層が同一であると仮定し，実験上最も重要な，$\kappa d\gg 1$を満たす長距離結合の場合を考える．ここで$\kappa=2\pi e^2\nu$はThomas--Fermi遮蔽半径の逆数である．この極限では，式~\eqref{Part-III-V-drag}の有効層間相互作用ポテンシャルは次の簡単な形を取る．
\begin{equation}\label{Part-III-V12}
\mathcal{V}_{12}^R\qo=\frac{1}{g}\, \frac{1}{\kappa d
Dq^{2}-2i\omega}\,,
\end{equation}
ここで$g=\nu D$である．次に，式~$\mathcal{D}^{R}\qo$を式~\eqref{Part-III-Spectral-F}および\eqref{Part-III-V12}とともに式~\eqref{Part-III-drag}へ代入し，エネルギー積分と運動量積分を行う．積分を調べると，2つのエネルギー$\omega$および$\omega'$はいずれも温度$\omega\sim\omega'\sim T$の程度であることがわかる．一方，移行運動量の特徴的な値は$q\sim q'\sim \sqrt{T/(D\kappa
d)}\ll \sqrt{T/D}$である．式~\eqref{Part-III-V12}と比較されたい．したがって，${\mathcal D}_{a}^R\qo$の式において$Dq^2$を$i\omega$に比べて無視し，式~\eqref{Part-III-drag}の積${\mathcal
D}_1^R{\mathcal D}_2^R$を$-\omega^{-2}$で近似できる．このようなスケールの分離は，4本脚頂点が実効的には空間的に局所的であり，3本の層間相互作用線は長距離に及ぶことを意味する．

エネルギーを$T$，運動量を$\sqrt{T/(D\kappa d)}$で再スケールすると，ドラッグ伝導率の式~\eqref{Part-III-drag}は$\sigma_{\mathrm{drag}}= R_Q^{-1}g^{-1}(\kappa
d)^{-2}\times $（{\rm 無次元積分}）という形に帰着できる．この積分にはパラメータが含まれず，数値的に評価できる~\cite{Drag-Third-order}．$\sigma_{\mathrm{drag}}\ll g/R_Q$の極限では，ドラッグ抵抗$\rho_{\mathrm{drag}}$は$\rho_{\mathrm{drag}}=\sigma_{\mathrm{drag}}R_Q^2/g^2$で与えられ，最終的に$\rho_{\mathrm{drag}}\approx 0.27R_{Q}\,g^{-3}(1/\kappa
d)^2$となる．これは温度に依存しないドラッグ抵抗率であり，層間相互作用の2次の寄与より大きくなることがある．後者は低温で$T^2$としてゼロに近づく．詳細および議論は文献~\cite{Drag-Third-order}にある．

%$$$$$$$$$$$$$$$$$$$$$$$$$$$$$$$$$$$$$$$$$$$$$$$$$$$$$$$$$$$$$$$$$$$$$$$$$$$$$$$$$$$$$$$$$$$$$$$$$$$$$$$$$$$
%$$$$$$$$$$$$$$$$$$$$$$$$$$$$$$$$$$$$$$$$$$$$$$$$$$$$$$$$$$$$$$$$$$$$$$$$$$$$$$$$$$$$$$$$$$$$$$$$$$$$$$$$$$$

\sectionlink{超伝導相関}\label{sec_SuperCond}
\subsectionlink{$\sigma$模型の一般化}\label{sec_SuperCond-1}

これまではKeldysh $\sigma$模型の\emph{ユニタリ}版，すなわち外部磁場などによって時間反転対称性が破れている場合を論じてきた．ここからは，時間反転不変性が保たれる\emph{直交}対称性クラスに移る．具体的に取り上げるのは，乱れた金属における超伝導揺らぎである．乱れた超伝導体へと一般化されたKeldysh $\sigma$模型は，Feigel'man，Larkin，Skvortsovによって構築された\cite{Feigel'manLarkinSkvortsov,SkvortsovLarkinFeigel'man}．この模型は，常伝導金属における弱局在効果の扱いにも適用できる．

金属のHamiltonianにBCS項
$$\hat{H}_{\mathrm{BCS}}=-\frac{\lambda}{\nu}\int\d\mathbf{r}\,\hat{\psi}^{\dag}_{\uparrow}(\mathbf{r})
\hat{\psi}^{\dag}_{\downarrow}(\mathbf{r})
\hat{\psi}_{\downarrow}(\mathbf{r})
\hat{\psi}_{\uparrow}(\mathbf{r})\,,$$
を加えることにより，乱れた超伝導体を記述する．これは電子・フォノン相互作用が媒介する粒子・粒子（Cooper）チャネルの短距離引力に対応し，$\lambda$は無次元の結合定数である．通常の手順で，$\hat{H}_{\mathrm{BCS}}$はKeldysh作用
$$S_{\mathrm{BCS}}=\frac{\lambda}{\nu}\int_{\mathcal{C}}\d
t\int\d\mathbf{r}\ \bar{\psi}_{\uparrow}\rt
\bar{\psi}_{\downarrow}\rt\psi_{\downarrow}\rt\psi_{\uparrow}\rt\,,$$
へと書き換えられる．時間積分はKeldysh経路に沿って行う．この4フェルミオン相互作用項は，複素場$\Delta\rt$に関する補助的な汎関数積分を導入するHubbard--Stratonovich変換によって分離できる：
\begin{eqnarray}\label{SuperCond-BCS-HS}
\exp(iS_{\mathrm{BCS}})=\int\D[\Delta]\exp\left(i\int\d
x\left[-\frac{\nu}{\lambda}|\Delta(x)|^{2}+\Delta(x)
\bar{\psi}_{\uparrow}(x)\bar{\psi}_{\downarrow}(x)+
\Delta^{*}(x)\psi_{\downarrow}(x)\psi_{\uparrow}(x)\right]\right)\,,
\end{eqnarray}
ここで$x=\rt$，$\int\d x=\int_{\mathcal{C}}\d t\int\d\mathbf{r}$である．以後の記法を簡潔にするため，4次元空間$\Omega$で定義されるバイスピノルのフェルミオンベクトル$\Psi=1/\sqrt{2}(\psi_{\uparrow},\psi_{\downarrow},\bar{\psi}_{\downarrow},-\bar{\psi}_{\uparrow})^{T}$および$\Psi^{+}=1/\sqrt{2}(\bar{\psi}_{\uparrow},\bar{\psi}_{\downarrow},-\psi_{\downarrow},\psi_{\uparrow})$を導入すると便利である．$\Omega$はスピン空間$(\psi_{\uparrow},\psi_{\downarrow})$と時間反転空間$(\psi,\bar{\psi})$の直積$S\otimes T$とみなせる．原理的には，バイスピノルの選び方は一意ではない．成分を別の順序に並べ替え，Gor'kov--Nambu\cite{Gor'kov,Nambu}（$N$）空間$(\psi_{\uparrow},\bar{\psi}_{\downarrow})$とスピン空間を明示的に分離することもできる．さらに，$\Psi$をNambu部分空間と時間反転部分空間の直積上で作用するものと考えてもよい．これら3つの表現は$\Omega=S\otimes T\propto N\otimes S\propto N\otimes T$として同等であり，どれを選ぶかは個々の問題における計算の便宜によって決まる．多くの場合，$N\otimes S$を選び，理論がスピン部分空間で対角であるため$S$の部分を省略する．ベクトル$\Psi$と$\Psi^{+}$は独立ではなく，電荷共役行列$\check{C}\equiv i\hat{\tau}_{y}\otimes\hat{s}_{x}$を介して$\Psi^{+}=(\check{C}\Psi)^{T}$と関係づけられる．ここで$i=0,x,y,z$に対する$\hat{\tau}_{i}$と$\hat{s}_{i}$は，それぞれNambu部分空間とスピン部分空間に作用するPauli行列である．$\hat\sigma_i$行列は従来どおりKeldysh部分空間に作用する．混乱を避けるため，必要に応じてKeldysh部分空間とNambu部分空間をそれぞれ添字$K$，$N$で示す．


Hubbard--Stratonovich変換\eqref{SuperCond-BCS-HS}を行い，乱れとCoulomb相互作用を通常の方法で扱うと，作用はフェルミオン演算子について2次となる．したがってGrassmann変数のGauss積分を実行すると，乱れで平均した分配関数として
\begin{eqnarray}\label{SuperCond-Z}
&&\hskip-.5cm
\mathcal{Z}=\int\D[\Phi,\Delta]\exp\left(\frac{i}{2}\Tr\big\{\check{\Phi}
U^{-1}_{0}\check{\Upsilon}\check{\Phi}\big\}-
\frac{i\nu}{2\lambda}\Tr\big\{\check{\Delta}^{\dag}
\check{\Upsilon}\check{\Delta}\big\}\right)
\int\D[\check Q]\exp\big(iS[\check Q,\Delta,\mathbf{A},\Phi]\big)\,,\nonumber\\
&&\hskip-.5cm
 iS[\check{Q},\Delta,\mathbf{A},\Phi]=-\frac{\pi\nu}{4\tauel}
\Tr\big\{\check{Q}^{2}\big\}+\Tr\ln\left[\check{G}^{-1}+
\frac{i}{2\tauel}\check{Q}+\check{\Phi}+
\mathbf{v}_{F}\check{\Xi}\check{\mathbf{A}}+\check{\Delta}\right]\,,
\end{eqnarray}
を得る．これは式\eqref{int-ferm-NLSM-action}の一般化である．この式および本章の残りでは，$K\otimes N$空間に作用する$4\times4$行列をチェック記号$\check{O}$で表し，NambuおよびKeldysh部分空間に作用する$2\times2$行列をハット記号$\hat{O}$で表す．式\eqref{SuperCond-Z}には，行列$\check{\Upsilon}=\hat{\sigma}_{x}\otimes\hat{\tau}_{0}$，$\check{\Xi}=\hat{\sigma}_{0}\otimes\hat{\tau}_{z}$，$\check{G}^{-1}=i\check{\Xi}\partial_{t}+\partial^{2}_{\mathbf{r}}/2m+\mu$，および行列場
\begin{eqnarray}\label{SuperCond-A-Phi-Delta}
\check{\Phi}\rt=[\Phi^{cl}\rt\hat{\sigma}_{0}+\Phi^{q}\rt\hat{\sigma}_{x}]\otimes\hat{\tau}_{0}\,,\qquad
\check{\mathbf{A}}\rt=[\mathbf{A}^{cl}\rt\hat{\sigma}_{0}+\mathbf{A}^{q}\rt\hat{\sigma}_{x}]
\otimes\hat{\tau}_{0}\,,
\nonumber\\
\check{\Delta}\rt=[\Delta^{cl}\rt\hat{\sigma}_{0}+\Delta^{q}\rt\hat{\sigma}_{x}]\otimes\hat{\tau}_{+}-
[\Delta^{*cl}\rt\hat{\sigma}_{0}+\Delta^{*q}\rt\hat{\sigma}_{x}]\otimes\hat{\tau}_{-}\,,
\end{eqnarray}
が含まれる．ここで$\hat{\tau}_{\pm}=(\hat{\tau}_{x}\pm i\hat{\tau}_{y})/2$である．$\check{Q}$行列も，時間領域の行列構造に加えて，Keldysh空間とNambu空間で$4\times4$の構造をもつ．


次に式\eqref{int-ferm-NLSM-Q-K-gauge}と同様に，$\EuScript{K}^{cl(q)}\rt$場を用いて式\eqref{SuperCond-Z}のゲージ変換を行う\footnote{超伝導の場合，ゲージ変換には位相因子$\exp(\pm i\check{\Xi}\check{\EuScript{K}})$が含まれ，指数に余分な行列$\check{\Xi}$がある点が式\eqref{int-ferm-NLSM-Q-K-gauge}と異なる．}．そしてトレース内の対数を$\check{Q}_{\EuScript{K}}$行列の勾配で展開する（第\ref{sec_NLSM}節で示した計算と同様である）．その結果，乱れた超伝導体の作用が次の形で得られる：
\begin{subequations}\label{SuperCond-S}
\begin{equation}
S[\check Q,\Delta,\mathbf{A},\Phi]=S_{\Delta}+S_{\Phi}+S_{\sigma}\,,
\end{equation}
\begin{equation}
S_{\Delta}=-\frac{\nu}{2\lambda}
\Tr\big\{\check{\Delta}^{\dag}_{\EuScript{K}}
\check{\Upsilon}\check{\Delta}_{\EuScript{K}}\big\}\,,\qquad
S_{\Phi}=\frac{\nu}{2}\Tr\big\{\check{\Phi}_{\EuScript{K}}
\check{\Upsilon}\check{\Phi}_{\EuScript{K}}\big\}\,,
\end{equation}
\begin{equation}
S_{\sigma}=\frac{i\pi\nu}{4}\,
\Tr\big\{D\,(\hat{\bm{\partial}}_{\mathbf{r}}
\check{Q}_{\EuScript{K}})^{2}-4\check{\Xi}\partial_{t}\check{Q}_{\EuScript{K}}+
4i\check{\Phi}_{\EuScript{K}}\check{Q}_{\EuScript{K}}+4i
\check{\Delta}_{\EuScript{K}}\check{Q}_{\EuScript{K}}\big\}\,.
\end{equation}
\end{subequations}
ここで，ゲージ変換後の電磁ポテンシャル$\check{\Phi}_{\EuScript{K}}$と$\check{\mathbf{A}}_{\EuScript{K}}$は，式\eqref{int-ferm-NLSM-C}によって元のポテンシャル$\check{\Phi}$，$\check{\mathbf{A}}$と関係づけられる．一方，ゲージ変換後の秩序変数場は
\begin{equation}\label{SuperCond-Delta-K}
\check{\Delta}_{\EuScript{K}}\rt=\exp\Big(-i\check{\Xi}\check
{\EuScript{K}}\rt\Big)\check{\Delta}\rt
\exp\Big(i\check{\EuScript{K}}\rt\check{\Xi}\Big)\,.
\end{equation}
で与えられる．式\eqref{int-ferm-covariant-derivative}と比べると，式(\ref{SuperCond-S}c)の空間共変微分にはNambu構造に由来する余分な$\check{\Xi}$行列が含まれる．すなわち，
\begin{equation}\label{SuperCond-covariant-derivative}
\hat{\bm{\partial}}_{\mathbf{r}}\check{Q}_{\EuScript{K}}=
\partial_{\mathbf{r}}\check{Q}_{\EuScript{K}}
-i[\check{\Xi}\check{\mathbf{A}}_{\EuScript{K}},\check{Q}_{\EuScript{K}}]\,.
\end{equation}

拘束条件$\check{Q}_{\EuScript{K}}^{2}=\check{1}$の下で，作用\eqref{SuperCond-S}を$\check{Q}_{\EuScript{K}}$について変分すると，鞍点方程式
\begin{equation}\label{SuperCond-Usadel}
\hat{\bm{\partial}}_{\mathbf{r}}\big(D\,\check{Q}_{\EuScript{K}}\circ
\hat{\bm{\partial}}_{\mathbf{r}}\check{Q}_{\EuScript{K}}\big)-
\big\{\check{\Xi}\partial_{t},\check{Q}_{\EuScript{K}}\big\}_+ +
i\big[\check{\Phi}_{\EuScript{K}}+\check{\Delta}_{\EuScript{K}},
\check{Q}_{\EuScript{K}}\big]=0\,,
\end{equation}
を得る．$\check{\EuScript{K}}=0$の場合，これは動的Usadel方程式\cite{LarkinOvchinnikov}に一致する．この方程式の古典解は，
\begin{equation}
\check{Q}_{\EuScript{K}}=\left(\begin{array}{cc}\hat{Q}^{R}_{\EuScript{K}}
& \hat{Q}^{K}_{\EuScript{K}}
\\ 0 & \hat{Q}^{A}_{\EuScript{K}} \end{array}\right)_K\,,
\end{equation}
の形で求める．ここで遅延，先進，Keldysh成分はNambu部分空間の行列である．


秩序変数場の量子成分$\Delta^{*q}\rt$について作用を変分すると，秩序変数の古典成分に対する自己無撞着方程式
\begin{equation}\label{SuperCond-Self-consistency}
\Delta^{cl}_{\EuScript{K}}\rt=\pi\lambda
\Tr\big\{(\hat{\sigma}_{x}\otimes\hat{\tau}_{-})
\check{Q}_{\EuScript{K}}\big\}\,.
\end{equation}
を得る．最後に，電磁ポテンシャルの量子成分$\Phi^{q}$，$\mathbf{A}^{q}$について作用を変分すると，一組のMaxwell方程式を得る．これらは動的Usadel方程式\eqref{SuperCond-Usadel}および自己無撞着条件\eqref{SuperCond-Self-consistency}とともに，超伝導体のダイナミクスを支配する閉じた方程式系を構成する．


一般化された$\sigma$模型の作用\eqref{SuperCond-S}と，そこから導かれる$\check{Q}_{tt'}(\mathbf{r})$および$\check{\Delta}\rt$の動力学方程式では，関係するすべての低エネルギー励起を区別なく保持している．その代償として理論は技術的に複雑になる．実際の多くの場合には，このような網羅的な記述は過剰であり，理論を大幅に簡略化できる．たとえば，明確なギャップ$|\Delta|$をもつ，十分低温の超伝導状態$T\to0$を考え，外場の摂動周波数$\omega$が小さい$\omega\ll|\Delta|$ときの動的応答を調べることが多い．これは準定常な条件に相当する．この場合，式\eqref{SuperCond-Usadel}から超伝導体の準古典的運動論方程式を導出できる．あるいは，秩序変数が小さい$|\Delta|\ll T_c$という転移近傍の温度領域$|T-T_{c}|\ll T_{c}$を考え，$\Delta\rt$のダイナミクスの有効理論，すなわちGinzburg--Landau理論を構成することもできる．どちらの近似も一般的な$\sigma$模型の理論から自然に導かれる．次の各節でこれらを考察する．

%$$$$$$$$$$$$$$$$$$$$$$$$$$$$$$$$$$$$$$$$$$$$$$$$$$$$$$$$$$$$$$$$$$$$$$$$$$$$$$$$$$$$$$$$$$$$$$$$$$$$$$$$$$$
\subsectionlink{準古典近似}\label{sec_SuperCond-2}

超伝導状態において，全エネルギー領域で有効な最適のゲージ場$\vec{\EuScript{K}}(\mathbf{r},\epsilon)$を選ぶことは難しい．しかし文献\cite{NarozhnyAleiner}では，ギャップより十分低いエネルギーの極限（$\epsilon\ll|\Delta|$）において，準古典Green関数$\check{Q}$に対する電位の効果は$\epsilon/|\Delta|\ll1$というパラメータに関して小さく，したがって近似的に$\vec{\EuScript{K}}(\mathbf{r},\epsilon)=0$と置けることが示されている．以下ではこの仮定を用いる\footnote{本節では，$\check{Q}_{\EuScript{K}}$行列の記法から添字$\EuScript{K}$を省略し，$\check{Q}_{\EuScript{K}}\to\check{Q}$とする．その他のゲージ変換後の場についても同様である．}．

空間的に一様な平衡超伝導体では，鞍点のUsadel方程式は次の$\check{Q}$行列によって解かれる：
\begin{equation}\label{SuperCond-Q-uniform}
\hat{Q}^{R(A)}(\epsilon)=\pm\frac{1}{\sqrt{(\epsilon\pm
i0)^{2}-|\Delta|^{2}}}\left(\begin{array}{cc}\epsilon &\Delta
\\ -\Delta^{*} & -\epsilon\end{array}\right)_N\,,
\end{equation}
一方，$\hat Q^K=\tanh\frac{\epsilon}{2T}(\hat Q^R-\hat Q^A)$である．上付き添字$cl$を省き，秩序変数を$\Delta$と書いた（本節ではその量子成分は現れない）．式\eqref{SuperCond-Q-uniform}を自己無撞着条件\eqref{SuperCond-Self-consistency}に代入すると，通常のBCSギャップ方程式
\begin{equation}\label{SuperCond-BCS-gap}
\Delta=\lambda\Delta\int^{\omega_{D}}_{|\Delta|}\frac{\d\epsilon}
{\sqrt{\epsilon^{2}-|\Delta|^{2}}}\tanh\frac{\epsilon}{2T}\,,
\end{equation}
を得る．これは臨界温度$T_{c}$以下で$|\Delta|$の非零解をもつ．


境界が存在する場合や常伝導金属と近接する場合には，空間的に非一様な超伝導の問題を扱うことになる．この場合，$\Delta$と$\hat Q^{R(A)}$はいずれも座標依存性をもち，式\eqref{SuperCond-Usadel}と式\eqref{SuperCond-Self-consistency}の解を求めなければならない．その際，$\check{Q}_{tt'}$は静的，すなわち中心時刻に依存しないと仮定し，Wigner変換表示へ移る．$\Phi=0$，$\mathbf{A}=0$における$4\times4$行列Usadel方程式の遅延ブロックから，
\begin{equation}\label{SuperCond-Usadel-Q-R}
\partial_{\mathbf{r}}\big(D\,\hat{Q}^{R}\partial_{\mathbf{r}}
\hat{Q}^{R}\big)+i\epsilon[\hat{\tau}_{z},\hat{Q}^{R}]+
i[\hat{\Delta},\hat{Q}^{R}]=0\,.
\end{equation}
を得る．行列Usadel方程式\eqref{SuperCond-Usadel}の先進ブロックにも同様の方程式が成り立つ．Keldyshセクターからは，別の方程式
\begin{equation}\label{SuperCond-Usadel-Q-K}
\partial_{\mathbf{r}}\big(D\, \hat{Q}^{R}\partial_{\mathbf{r}}
\hat{Q}^{K}+ D\, \hat{Q}^{K}\partial_{\mathbf{r}}\hat{Q}^{A}\big)
+i\epsilon[\hat{\tau}_{z},\hat{Q}^{K}]+i[\hat{\Delta},\hat{Q}^{K}]=0\,.
\end{equation}
が得られる．非線形拘束$\check{Q}^{2}=\check{1}$は，次の条件を課す：
\begin{equation}\label{SuperCond-Normalizations}
\hat{Q}^{R}\hat{Q}^{R}=\hat{Q}^{A}\hat{Q}^{A}=\hat{1}\,,\quad
\hat{Q}^{R}\hat{Q}^{K}+\hat{Q}^{K}\hat{Q}^{A}=0\,.
\end{equation}
これらの条件は，Green関数行列の遅延ブロックと先進ブロックを角度でパラメータ化\cite{BelzigWilhelm}することにより，明示的に満たすことができる：
\begin{subequations}\label{SuperCond-Q-R-A}
\begin{equation}
\hskip-2cm \hat{Q}^{R}(\mathbf{r},\epsilon)=\left(\begin{array}{cc}
\cosh\theta & \sinh\theta \exp(i\chi)\\-\sinh\theta \exp(-i\chi) &
-\cosh\theta
\end{array}\right)_N\,,
\end{equation}
\begin{equation}
\hat{Q}^{A}(\mathbf{r},\epsilon)=-\hat{\tau}_{z}\big[\hat{Q}^{R}\big]^{\dag}\hat{\tau}_{z}=
\left(\begin{array}{cc} -\cosh\bar{\theta} & -\sinh\bar{\theta}
\exp(i\bar{\chi})\\ \sinh\bar{\theta} \exp(-i\bar{\chi}) &
\cosh\bar{\theta}
\end{array}\right)_N\, ,
\end{equation}
\end{subequations}
ここで$\theta(\mathbf{r},\epsilon)$と$\chi(\mathbf{r},\epsilon)$は，座標とエネルギーに依存する複素スカラー関数である．Keldysh成分については，常に
\begin{equation}\label{SuperCond-Q-K}
\hat{Q}^{K}=\hat{Q}^{R}\circ\hat{F}-\hat{F}\circ\hat{Q}^{A},
\end{equation}
と選ぶことができる．$\hat{F}$は一般化された行列分布関数と考えられる．Schmidt--Sch\"{o}n\cite{SchmidtSchon}およびLarkin--Ovchinnikov\cite{LarkinOvchinnikov-F}に従って，
\begin{equation}\label{SuperCond-F}
\hat{F}(\mathbf{r},\epsilon)=\left(
\begin{array}{cc}
F_{L}(\mathbf{r},\epsilon)+F_{T}(\mathbf{r},\epsilon) & 0 \\ 0 &
F_{L}(\mathbf{r},\epsilon)-F_{T}(\mathbf{r},\epsilon)
\end{array}\right)_{N}=F_{L}(\mathbf{r},\epsilon)\hat{\tau}_{0}+
F_{T}(\mathbf{r},\epsilon)\hat{\tau}_{z}\,,
\end{equation}
と選ぶ．ここで$F_{L(T)}$は，秩序変数に対する分布関数の\textit{縦}成分と\textit{横}成分を表す．物理的には，$F_{T}$は系の電荷モードに対応して電流密度を決め，$F_{L}$はエネルギーモードに対応して熱（エネルギー）流を決める（詳しい議論についてはTinkham\cite{Tinkham}およびKopnin\cite{Kopnin}の教科書を参照されたい）．


式\eqref{SuperCond-Q-R-A}の形の$\hat{Q}^{R}$を式\eqref{SuperCond-Usadel-Q-R}に代入すると，対応する行列方程式の対角要素から
\begin{equation}\label{SuperCond-Usadel-1}
D\, \partial_{\mathbf{r}}\big(
\sinh^{2}\theta\,\partial_{\mathbf{r}}\chi\big)
=2i|\Delta|\sinh\theta\sin(\varphi-\chi)\,,
\end{equation}
を得る．ここで秩序変数を$\Delta(\mathbf{r})=|\Delta(\mathbf{r})|\exp\{i\varphi(\mathbf{r})\}$とパラメータ化した．行列方程式\eqref{SuperCond-Usadel-Q-R}の非対角ブロックから，式\eqref{SuperCond-Usadel-1}を用いると，
\begin{equation}\label{SuperCond-Usadel-2}
D\,\partial^{2}_{\mathbf{r}}\theta+2i\epsilon\sinh\theta-2i|\Delta|\cosh\theta\cos(\varphi-\chi)
=\frac{D}{2}\big(\partial_{\mathbf{r}}\chi\big)^{2}\sinh2\theta\,.
\end{equation}
を得る．次にGreen関数行列のKeldysh成分$\hat{Q}^{K}$の方程式を考える．分解\eqref{SuperCond-Q-K}を用いて式\eqref{SuperCond-Usadel-Q-K}に代入すると，
\begin{eqnarray}
D\left(\partial^{2}_{\mathbf{r}}\hat{F}+
\hat{Q}^{R}\partial_{\mathbf{r}}\hat{Q}^{R}\partial_{\mathbf{r}}\hat{F}-
\partial_{\mathbf{r}}\hat{F}\hat{Q}^{A}\partial_{\mathbf{r}}\hat{Q}^{A}-
\partial_{\mathbf{r}}\big(\hat{Q}^{R}\partial_{\mathbf{r}}\hat{F}\hat{Q}^{A}\big)\right)
+i\epsilon\left(\hat{Q}^{R}\big[\hat{\tau}_{z},\hat{F}]-
[\hat{\tau}_{z},\hat{F}\big]\hat{Q}^{A}\right)\nonumber \\
+i\left(\hat{Q}^{R}\big[\hat{\Delta},\hat{F}]-
[\hat{\Delta},\hat{F}\big]\hat{Q}^{A}\right)=0\,.
\end{eqnarray}
を得る．ここで$\hat{F}$に式\eqref{SuperCond-F}を用い，（i）上の行列方程式のNambu空間に関するトレースをとり，（ii）上の方程式に$\hat{\tau}_{z}$を掛けてからトレースをとると，非平衡分布関数$F_{L(T)}$に対する2つの連立運動論方程式が得られる．これらは保存則の形で書ける\cite{LarkinOvchinnikov-KinEq}：
\begin{subequations}\label{SuperCond-Usadel-3}
\begin{equation}
\partial_{\mathbf{r}}\big(\EuScript{D}_{L}\partial_{\mathbf{r}}
F_{L}-D\partial_{\mathbf{r}} F_{T}Y\big)+D\partial_{\mathbf{r}}
F_{T}\mathcal{J}_{S}=\mathcal{I}^{a}_{\mathrm{coll}}\,,
\end{equation}
\begin{equation}
\partial_{\mathbf{r}}(\EuScript{D}_{T}\partial_{\mathbf{r}}
F_{T}+D\partial_{\mathbf{r}} F_{L}Y)+D\partial_{\mathbf{r}}
F_{L}\mathcal{J}_{S}=\mathcal{I}^{b}_{\mathrm{coll}}\,.
\end{equation}
\end{subequations}
ここで，エネルギーと座標に依存する拡散係数
\begin{subequations}\label{SuperCond-kinetics-D}
\begin{equation}
\hskip-1.5cm \EuScript{D}_{L}(\mathbf{r},\epsilon)=\frac{D}{4}
\Tr\left\{\hat{\tau}_{0}-\hat{Q}^{R}\hat{Q}^{A}\right\}_N
=\frac{D}{2}\left[1+|\cosh\theta|^{2}-
|\sinh\theta|^{2}\cosh\big(2\mathrm{Im}[\chi]\big)\right]\,,
\end{equation}
\begin{equation}
\EuScript{D}_{T}(\mathbf{r},\epsilon)=\frac{D}{4}\Tr\left\{\hat{\tau}_{0}-\hat{\tau}_{z}
\hat{Q}^{R}\hat{\tau}_{z}\hat{Q}^{A}\right\}_N = \frac{D}{2}\left[1+|\cosh\theta|^{2}+
|\sinh\theta|^{2}\cosh\big(2\mathrm{Im}[\chi]\big)\right]\,,
\end{equation}
\end{subequations}
超電流を担う状態の密度
\begin{equation}\label{SuperCond-kinetics-Js}
\mathcal{J}_{S}(\mathbf{r},\epsilon)=\frac{1}{4}\Tr
\left\{\hat{\tau}_{z}\big(\hat{Q}^{R}\partial_{\mathbf{r}}\hat{Q}^{A}-
\hat{Q}^{A}\partial_{\mathbf{r}}\hat{Q}^{R}\big)\right\}_N
=-\mathrm{Im}\left(\sinh^{2}\theta\,\partial_{\mathbf{r}}\chi\right)\,,
\end{equation}
およびスペクトル関数
\begin{equation}\label{SuperCond-kinetics-Gamma}
Y(\mathbf{r},\epsilon)=\frac{1}{4}\Tr\left\{\hat{Q}^{R}\hat{\tau}_{z}\hat{Q}^{A}\right\}_N
=\frac{1}{2}|\sinh\theta|^{2}\sinh\big(2\mathrm{Im}[\chi]\big)\,.
\end{equation}
を導入した．最後に，式\eqref{SuperCond-Usadel-3}の右辺には衝突積分
\begin{subequations}\label{SuperCond-kinetics-Coll}
\begin{equation}
\mathcal{I}^{a}_{\mathrm{coll}}=\frac{F_{T}}{2}\, \Tr\left\{\hat{\tau}_{z}
\big(\hat{Q}^{R}\hat{\Delta}+\hat{\Delta}\hat{Q}^{A}\big)\right\}_N =2F_{T}|\Delta|
\mathrm{Re}\left[\sinh\theta\sin(\varphi-\chi)\right]\,,
\end{equation}
\begin{equation}
\mathcal{I}^{b}_{\mathrm{coll}}=\frac{F_{T}}{2}\, \Tr\left\{
\hat{Q}^{R}\hat{\Delta}+\hat{\Delta}\hat{Q}^{A} \right\}_N
=-2F_{T}|\Delta|\mathrm{Im}
\left[\sinh\theta\cos(\varphi-\chi)\right]\,.
\end{equation}
\end{subequations}
が含まれる．非弾性的な電子・電子相互作用および電子・フォノン相互作用に伴う衝突積分は，ここでは論じない．その導出はKopninの教科書\cite{Kopnin}に見出せる．式\eqref{SuperCond-Usadel-1}，\eqref{SuperCond-Usadel-2}，\eqref{SuperCond-Usadel-3}は，スペクトル量\eqref{SuperCond-kinetics-D}〜\eqref{SuperCond-kinetics-Coll}とともに，準古典近似で適用できる，乱れた超伝導体の完全な運動論方程式系を構成する．これらの方程式には，自己無撞着関係（式\eqref{SuperCond-Self-consistency}を参照）
\begin{equation}\label{SuperCond-Usadel-4}
\Delta(\mathbf{r})=\frac{\lambda}{2}\int\d\epsilon\,\Big\{[\sinh\theta
\exp(i\chi)+\sinh\bar{\theta}\exp(i\bar{\chi})]F_{L}-\big[\sinh\theta
\exp(i\chi)-\sinh\bar{\theta}\exp(i\bar{\chi})\big]F_{T}\Big\}\,,
\end{equation}
と，電流の連続性を表すGreen関数の境界条件\cite{Nazarov-TunnelingAction,Zaitsev,KuprianovLukichev,Lambert}，
\begin{equation}\label{SuperCond-BoundaryConditions}
\sigma_{L}\mathcal{A}_{L}\check{Q}_{L}\partial_{\mathbf{r}}\check{Q}_{L}=
\sigma_{R}\mathcal{A}_{R}\check{Q}_{R}\partial_{\mathbf{r}}\check{Q}_{R}=
\mathrm{g}_{T}[\check{Q}_{L},\check{Q}_{R}]\,,
\end{equation}
を補う．ここで$\sigma$と$\mathcal{A}$は，界面に隣接する細線のバルク常伝導状態の伝導率と断面積である．$L/R$はそれぞれ界面の左側／右側を表し，$\mathrm{g}_{T}$は界面のトンネルコンダクタンスである．


運動論方程式系\eqref{SuperCond-Usadel-1}〜\eqref{SuperCond-Usadel-3}を解析的に解けることはまれであり，一般には数値的方法に頼る必要がある．与えられた輸送問題の解を求めるには，次のように進める\cite{BelzigWilhelm}．
\begin{enumerate}
\item ある$\Delta(\mathbf{r})$から始める．通常は，超伝導体の全領域で$\Delta=\mathrm{const}$，常伝導金属では$\Delta=0$とする．
\item 遅延Green関数についてUsadel方程式\eqref{SuperCond-Usadel-1}〜\eqref{SuperCond-Usadel-2}を解き，スペクトル角$\theta(\mathbf{r},\epsilon)$と$\chi(\mathbf{r},\epsilon)$を求める．
\item これらの解を用いて，スペクトル運動論量$\EuScript{D}_{L,T}(\mathbf{r},\epsilon)$，$\mathcal{J}_{S}(\mathbf{r},\epsilon)$，$Y(\mathbf{r},\epsilon)$を計算する．
\item $F_{L/T}(\mathbf{r},\epsilon)$について運動論方程式\eqref{SuperCond-Usadel-3}を解く．
\item 式\eqref{SuperCond-Usadel-4}から新しい$\Delta(\mathbf{r})$を計算し，自己無撞着性が得られるまでこの手順を反復する．
\end{enumerate}
運動論方程式を解けば，関心のある物理量を決定できる．たとえば電流は$\mathbf{j}=\mathbf{j}_{n}+\mathbf{j}_{s}$となり，ここで$\mathbf{j}_{n}(\mathbf{r})=\nu\int\d\epsilon\,\EuScript{D}_{T}(\mathbf{r},\epsilon)\partial_{\mathbf{r}} F_{T}(\mathbf{r},\epsilon)$は常伝導成分，$\mathbf{j}_{s}(\mathbf{r})=\nu D\int\d\epsilon F_{L}(\mathbf{r},\epsilon)\mathcal{J}_{S}(\mathbf{r},\epsilon)$は超電流密度である．

上で概説した，乱れた超伝導体の準古典的運動論は，さまざまな現象の研究に適用できる．例を挙げれば，超伝導体・常伝導金属ヘテロ構造における近接効果に関する問題\cite{BelzigBruder,Gueron,Zhou,Hammer}，非平衡Josephson効果\cite{Heikkila,Wilhelm}，超伝導体のHall効果\cite{ZhouSpivak}や熱電現象\cite{StoofNazarov,Virtanen}，ショット雑音\cite{Stenberg}，非平衡分布関数の制御\cite{Crosser}，さらに多くの問題を，式\eqref{SuperCond-Usadel-1}〜\eqref{SuperCond-Usadel-3}によって扱うことができる．説明のため，第\ref{app_Part-IV}節では比較的簡単な平衡の例をいくつか考察する．

%$$$$$$$$$$$$$$$$$$$$$$$$$$$$$$$$$$$$$$$$$$$$$$$$$$$$$$$$$$$$$$$$$$$$$$$$$$$$$$$$$$$$$$$$$$$$$$$$$$$$$$$$$$$
\subsectionlink{時間依存Ginzburg--Landau理論}\label{sec_SuperCond-3}

Gor'kov\cite{Gor'kov-GL}は，温度が臨界温度に近い極限$|T-T_{c}|\ll T_{c}$で，現象論的なGinzburg--Landau（GL）理論\cite{GinzburgLandau}が微視的BCS模型から自然に導かれることを示した．その後，Gor'kovとEliashberg\cite{Gor'kovEliashberg}は，時間依存の動的現象を含めるようGinzburg--Landau理論の適用範囲を拡張した．この問題は，多くの後続論文\cite{HoughtonMaki,HuThompson,Kramer,SchonAmbegaokar,Hu,Krempasky,Otterlo}や教科書\cite{Tinkham,Kopnin,Larkin-Varlamov}で再検討されている．$\sigma$模型の言葉では，静的GL汎関数は超対称性\cite{AltlandSimons}またはレプリカ\cite{YurkevichLerner}の手法により得られる．ここではKeldysh形式による動的理論\cite{LevchenkoKamenev}を論じる．

式\eqref{SuperCond-S}から時間依存Ginzburg--Landau（TDGL）理論を導く方法では，従来のように秩序変数の方程式だけを与えるのではなく，有効作用を用いて理論を定式化できる．その結果，平均量に加えて揺らぎの効果も扱える．TDGL作用には，揺動散逸定理を満たすための確率的雑音項が含まれるからである．さらに，理論のゲージ不変性を保つ異常Gor'kov--Eliashberg（GE）項\cite{Gor'kovEliashberg}と，超伝導揺らぎによって伝導率および1粒子状態密度を繰り込むAslamazov--Larkin（AL）項\cite{AslamazovLarkin}，Maki--Thompson（MT）項\cite{Maki}，状態密度（DOS）項\cite{Hurault}を，自然かつ曖昧さなく同定できる．TDGL方程式へのほとんどのアプローチではAslamazov--Larkin項は正しく捉えられるが，TDGLに関する多くの研究でGor'kov--Eliashberg項，Maki--Thompson項，DOS項はしばしば失われている．

一般的な$\sigma$模型の作用\eqref{SuperCond-S}から有効TDGL理論を導く方針は次のとおりである．（i）非線形拘束$\check Q^2=\check{1}$を満たすように，鞍点$\check Q$行列の多様体をパラメータ化する．（ii）鞍点の周りのGauss揺らぎを積分消去する．（iii）得られた作用において，すべての\textit{量子}場（秩序変数$\Delta$と電磁ポテンシャル$\Phi$，$\mathbf{A}$）について2次までの項を残す．（iv）電子系は常に局所熱平衡にあると仮定する．これは外場が大きすぎないことを意味する．より正確には，電場$\mathbf{E}$は$e|\mathbf{E}|\xi_0\ll T_c$を満たし，磁場$\mathbf{H}$には$e|\mathbf{H}|\xi_0\ll1/\xi_0$という条件が課される．ここで$\xi_{0}=\sqrt{D/T_{c}}$は超伝導のコヒーレンス長である．外場の空間・時間スケールに対する制約と電子の局所平衡によって，理論はかなり簡単になる．とくに，有効作用のほとんどの項は空間と時間について局所的な形をとる．それでも，有効理論が完全に局所的な形になるわけではない．

ギャップのない超伝導の場合，この手続きは比較的直接的である．ギャップのない状態は，磁性不純物の存在下，または臨界温度以上の揺らぎ領域$T\gtrsim T_{c}$で生じる．ギャップのある相$T\lesssim T_{c}$では，事情はより複雑になる．Gor'kovとEliashberg\cite{Gor'kovEliashberg}が指摘したように，困難の原因はギャップ端でのBCS状態密度の特異性にある．これは時間領域において，周波数$2\Delta/\hbar$でゆっくり減衰する振動応答を生む．そのため，小さなパラメータ$\Delta/T_{c}\ll 1$のべきによる展開は破綻する．原理的には，高周波の外場の場合には$\Delta/(\hbar\omega)$についての展開を補うことができる．ギャップのある相で低周波応答を記述するには，時間について\textit{非局所}なTDGL理論が必要である．磁性不純物やエネルギー緩和などの対破壊機構が存在すると，解析は大幅に簡単になる．こうした機構は状態密度の特異性を取り除き，有限の$\Delta$が存在するにもかかわらずギャップのない相を生じさせうる．この条件では$\Delta\tau_{\phi}/\hbar\ll 1$および$\omega\tau_{\phi}\ll 1$のべき展開が正当化され，時間について局所的なTDGL方程式が再び得られる（ここで$\tau_{\phi}$は対破壊時間である）．本節では揺らぎ領域$T\gtrsim T_c$だけを考える．この場合，スペクトルは自動的にギャップをもたず，明示的な対破壊機構を導入する必要はない．

上記の手順（i）〜（iv）に沿って進むにあたり，$T>T_{c}$では，エネルギーギャップの自己無撞着方程式\eqref{SuperCond-Self-consistency}が$\langle\Delta^{cl}\rangle=0$という自明な解しかもたないことを思い出そう．したがって，作用\eqref{SuperCond-S}の試行鞍点は金属状態$\check{Q}_{\EuScript{K}}=\check{\Lambda}=\hat{\Lambda}\otimes\hat{\tau}_{z}$に戻る．$\hat{\Lambda}$は式\eqref{NLSM-Lambda}で定義される．この$\check{Q}_{\EuScript{K}}$の周りのGauss積分にはCooperモードが含まれる．それらは，次の$\Qk$行列のパラメータ表示で取り入れられる：
\begin{equation}\label{SuperCond-Q-W}
    \check{Q}_{\EuScript{K}}=\check{\mathcal{U}}\circ
e^{-\check{\mathcal{W}}/2}\circ(\hat{\sigma}_{z}\otimes\hat{\tau}_{z})\circ
e^{\hat{\mathcal{W}}/2}\circ\check{\mathcal{U}}^{-1}\,,
\end{equation}
ここで揺らぎ行列を次のように選ぶ：
\begin{equation}\label{SuperCond-W}
\fitdisplay{
\check{\mathcal{W}}_{tt'}(\mathbf{r})= \left(\begin{array}{cc}
c_{tt'}(\mathbf{r})\hat{\tau}_{+}-c^{*}_{tt'}(\mathbf{r})\hat{\tau}_{-}
&
d_{tt'}(\mathbf{r})\hat{\tau}_{0}+d^{z}_{tt'}(\mathbf{r})\hat{\tau}_{z}
\\
\bar{d}_{tt'}(\mathbf{r})\hat{\tau}_{0}+
\bar{d}^{z}_{tt'}(\mathbf{r})\hat{\tau_{z}} &
\bar{c}_{tt'}(\mathbf{r})\hat{\tau}_{+}-
\bar{c}^{*}_{tt'}(\mathbf{r})\hat{\tau}_{-}
\end{array}\right)_K\,,\qquad
\check{\mathcal{U}}=\check{\mathcal{U}}^{-1}=
\left(\begin{array}{cc} 1 & F \\ 0 & -1
\end{array}\right)_K\otimes\hat{\tau}_{0}\,.
}
\end{equation}
式\eqref{NLSM-U-W}と比べると，$\check{\mathcal{W}}$は2倍の数の拡散モードを含む．それらは，時間部分空間における4つのHermite行列$\{d,\bar{d}\}$および$\{d^{z},\bar{d}^{z}\}$で記述される．また，2つの独立な\emph{複素}行列場$\{c,\bar{c}\}$で記述されるCooperモードも含まれる．ここで$\check{\mathcal{W}}$に依存する$\Qk$行列$\check{Q}_{\EuScript{K}}[\check{\mathcal{W}}]$を式\eqref{SuperCond-S}に代入し，作用を$\check{\mathcal{W}}$の揺らぎについて2次まで展開する：$S[\check Q,\Delta,\mathbf{A},\Phi]\Rightarrow S[\check{\mathcal{W}},\Delta,\mathbf{A},\Phi]$．この段階で$\check{\mathcal{W}}$についてのGauss積分が可能になる（手続きの詳細は\ref{app_FluctuationExpansion}を参照）：
\begin{equation}\label{SuperCond-W-integration}
 \int\D[\check{\mathcal{W}}]\exp\Big(iS[\check{\mathcal{W}},\Delta,\mathbf{A},\Phi]\Big)=
 \exp\Big(iS_{\mathrm{eff}}[\Delta,\mathbf{A},\Phi]\Big)\,,
\end{equation}
これにより，最終的に有効TDGL作用が得られる．これはいくつかの寄与から成る：
\begin{eqnarray}\label{SuperCond-S-eff}
S_{\mathrm{eff}}[\Delta,\mathbf{A},\Phi]=S_{\mathrm{N}}[\mathbf{A},\Phi]+
S_{\mathrm{GL}}[\Delta,\mathbf{A},\Phi]+
S_{\mathrm{SC}}[\Delta,\mathbf{A},\Phi]+S_{\mathrm{MT}}[\Delta,\mathbf{A},\Phi]+
S_{\mathrm{DOS}}[\Delta,\mathbf{A},\Phi]\,,
\end{eqnarray}
以下，順に説明する．


作用$S_{\mathrm{N}}[\mathbf{A},\Phi]$は，式\eqref{SuperCond-S}の常伝導金属の部分であり，$S[\check Q,\Delta,\mathbf{A},\Phi]$で$\check{Q}_{\EuScript{K}}=\check{\Lambda}$，$\check{\Delta}=0$と置いて得られる．その形は次のとおりである\footnote{式\eqref{SuperCond-S-N}および本節の残りでは，ソース場に伴う電子の電荷$e$を$\mathbf{A}\to e\mathbf{A}$，$\Phi\to e\Phi$として復元したことに注意されたい．したがって，ここでは$\mathbf{A}$と$\Phi$は実際の電磁ポテンシャルである．脚注8を参照．}：
\begin{equation}\label{SuperCond-S-N}
S_{\mathrm{N}}[\mathbf{A},\Phi]=e^{2}\nu
D\,\Tr\left\{\vec{\mathbf{A}}^{T}_{\EuScript{K}}
\left(\begin{array}{cc}0 &
D\,\partial^{2}_{\mathbf{r}}-\overleftarrow{\partial}_{t}
\\ D\,\partial^{2}_{\mathbf{r}}-\overrightarrow{\partial}_{t} & 4iT
\end{array}\right)_K \vec{\mathbf{A}}_{\EuScript{K}}\right\}\,,
\end{equation}
時間微分の上の矢印は微分が作用する向きを示す．出発点は常伝導の鞍点\eqref{NLSM-Lambda}なので，汎関数$\vec{\EuScript{K}}[\Phi]$は式\eqref{int-ferm-NLSM-K-functional}で与えられ，ゲージ変換後のベクトルポテンシャル$\mathbf{A}_{\EuScript{K}}$は式\eqref{int-ferm-NLSM-C}で定義される．


$S_{\mathrm{GL}}$は作用の時間依存Ginzburg--Landau部分である：
\begin{equation}\label{SuperCond-S-GL}
S_{\mathrm{GL}}[\Delta,\mathbf{A},\Phi]=
2\nu\Tr\left\{\vec{\Delta}^{\dag}_{\EuScript{K}}\rt\hat{L}^{-1}
\vec{\Delta}_{\EuScript{K}}\rt\right\}\, ,
\end{equation}
これは外部ポテンシャルの影響下での秩序変数の時間的・空間的変化を支配する．有効伝播関数$\hat{L}^{-1}$は，Keldysh空間で典型的なボソン型の構造
\begin{equation}
\hat{L}^{-1}=\left(\begin{array}{cc} 0 & L^{-1}_{A} \\ L^{-1}_{R} &
L^{-1}_K
\end{array}\right)_K\,,
\end{equation}
をもち，各成分は次のように与えられる：
\begin{subequations}\label{SuperCond-L-R-A}
\begin{equation}
L^{-1}_{R(A)}=
\frac{\pi}{8T_c}\left[\mp\partial_{t}-\tau^{-1}_{\mathrm{GL}}+D\big(\partial_{\mathbf{r}}-
2ie\mathbf{A}^{cl}_{\EuScript{K}}\big)^{2}-\frac{7\zeta(3)}{\pi^{3}T_{c}}
|\Delta^{cl}_{\EuScript{K}}|^{2}\right]\,,
\end{equation}
\begin{equation}
\hskip-2.8cm L^{-1}_{K}= \coth{\omega\over 2T}\big[L_R^{-1}(\omega)
- L_A^{-1}(\omega)\big] \approx  {i\pi\over 2}\,,
\end{equation}
\end{subequations}
ここで$\omega\ll T\approx T_c$であり，Ginzburg--Landau緩和時間を$\tau_{\mathrm{GL}}=\pi/8(T-T_{c})$と定義する．$T-T_c\ll T_c$という仮定の下で，作用のGL部分は時間について局所的な形をとることに注意されたい．


作用の$S_{\mathrm{SC}}$部分は超電流を担う：
\begin{equation}\label{SuperCond-S-SC}
S_{\mathrm{SC}}[\Delta,\mathbf{A},\Phi]=\frac{\pi e\nu
D}{2T}\Tr\left\{\mathbf{A}^{q}_{\EuScript{K}}\mathrm{Im}
\big[\Delta^{*cl}_{\EuScript{K}}
\big(\partial_{\mathbf{r}}-2ie\mathbf{A}^{cl}_{\EuScript{K}}\big)
\Delta^{cl}_{\EuScript{K}}\big]\right\}\,.
\end{equation}
この略号は，$S_{\mathrm{SC}}$を$\mathbf{A}^{q}$で微分すると，秩序変数で表した超電流の標準的な式\cite{Tinkham}が得られることに由来する．


作用のMaki--Thompson部分$S_{\mathrm{MT}}$は，超伝導揺らぎによる常伝導作用$S_\mathrm{N}$の拡散係数の繰り込みを担う．これは
\begin{equation}\label{SuperCond-S-MT}
S_{\mathrm{MT}}[\Delta,\mathbf{A},\Phi]=
e^{2}\nu\Tr\left\{\vec{\mathbf{A}}^{T}_{\EuScript{K}}\rt\hat{\mathcal{T}}_{\delta
D}(t,t')\vec{\mathbf{A}}_{\EuScript{K}}\rtp\right\}\,,
\end{equation}
と書ける．ここで演算子$\hat{\mathcal{T}}_{\delta D}(t,t')$は
\begin{equation}
\hat{\mathcal{T}}_{\delta D}=\left(\begin{array}{cc} 0 &
-\overleftarrow{\partial_{t}}\,
\delta D^{\mathrm{MT}}_{\mathbf{r},t',t} \\
-\delta D^{\mathrm{MT}}_{\mathbf{r},t,t'}\,
\overrightarrow{\partial_{t'}} & 2iT \left( \delta
D^{\mathrm{MT}}_{\mathbf{r},t,t'} + \delta
D^{\mathrm{MT}}_{\mathbf{r},t',t} \right)
\end{array}\right)_K\,.
\end{equation}
で与えられる．拡散係数の補正$\delta D^{\mathrm{MT}}[\Delta_{\EuScript{K}}]$は，揺らぐ秩序変数の非局所汎関数であり，
\begin{equation} \label{SuperCond-D-MT}
\delta D^{\mathrm{MT}}_{\mathbf{r},t,t'}=\frac{\pi
D}{4T}\int\mathrm{d} \mathbf{r}'\mathrm{d}\mathbf{r}''\,
 \mathcal{C}^{\,\mathbf{r},\mathbf{r}'}_{\tau,t,t'}
 {\Delta}^{*cl}_\EuScript{K}\left(\mathbf{r}',\tau\right)
 {\Delta}^{cl}_\EuScript{K}(\mathbf{r}'',\tau)\
\bar{\mathcal{C}}^{\,\mathbf{r}'',\mathbf{r}}_{\tau,t',t}\,,
\end{equation}
となる．ここで$\tau=(t+ t')/2\,$である．遅延Cooperon伝播関数$\mathcal{C}^{\,\mathbf{r},\mathbf{r}'}_{\tau,t,t'}\sim\theta(t-t')$と先進Cooperon伝播関数$\bar{\mathcal{C}}^{\,\mathbf{r},\mathbf{r}'}_{\tau,t,t'}\sim\theta(t'-t)$は，次の方程式のGreen関数である：
\begin{subequations}
\begin{equation}
\left\{\phantom{-}\partial_t
-ie\Phi^{cl}_{\EuScript{K}}(\mathbf{r},\tau_+)+
ie\Phi^{cl}_{\EuScript{K}}(\mathbf{r},\tau_-)-D\left[\partial_{\mathbf{r}}-
ie\mathbf{A}^{cl}_\EuScript{K}(\mathbf{r},\tau_+)-
ie\mathbf{A}^{cl}_\EuScript{K}(\mathbf{r},\tau_-)\right]^2 \right\}
\mathcal{C}^{\mathbf{r},\mathbf{r}'}_{\tau,t,t'} =
\delta_{\mathbf{r}-\mathbf{r}'}\delta_{t-t'}\ ,
\end{equation}
\begin{equation}
\left\{-\partial_t+ie\Phi^{cl}_{\EuScript{K}}(\mathbf{r},\tau_+)-
ie\Phi^{cl}_{\EuScript{K}}(\mathbf{r},\tau_-)-D\left[\partial_{\mathbf{r}}-
ie\mathbf{A}^{cl}_\EuScript{K}(\mathbf{r},\tau_+)-
ie\mathbf{A}^{cl}_\EuScript{K}(\mathbf{r},\tau_-)\right]^2 \right\}
\bar{\mathcal{C}}^{\mathbf{r},\mathbf{r}'}_{\tau,t,t'} =
\delta_{\mathbf{r}-\mathbf{r}'}\delta_{t-t'}\,,
\end{equation}
\end{subequations}
ここで$\tau_\pm=\tau\pm t/2$である．MT作用\eqref{SuperCond-S-MT}は，常伝導作用$S_{\mathrm{N}}$とまったく同じ構造をもつ．したがって，常伝導の拡散定数に時間的に非局所な繰り込み$D\delta_{t-t'}\to D\delta_{t-t'}+\delta D^{\mathrm{MT}}_{\mathbf{r},t,t'}$を加えることにより，式\eqref{SuperCond-S-N}に組み込める．


最後に，$S_{\mathrm{DOS}}$は式\eqref{SuperCond-S-MT}の$S_{\mathrm{MT}}$と同様の構造をもつ：
\begin{equation}\label{SuperCond-S-DOS}
S_{\mathrm{DOS}}[\Delta,\mathbf{A},\Phi]=e^{2}D
\Tr\left\{\delta\nu^{\mathrm{DOS}}_{\mathbf{r},t}
\left[\vec{\mathbf{A}}^{T}_{\EuScript{K}}\rt
\left(\begin{array}{cc}0 &-\overleftarrow{\partial}_{t} \\
-\overrightarrow{\partial}_{t} & 4iT
\end{array}\right)_{K}\vec{\mathbf{A}}_{\EuScript{K}}\rt\right]\right\}\,,
\end{equation}
ここで局所的に繰り込まれた状態密度は
\begin{equation}\label{SuperCond-DOS-corr}
\delta\nu^{\mathrm{DOS}}_{\mathbf{r},t}=-\nu\frac{7\zeta(3)}{4\pi^{2}T^{2}}
|\Delta^{cl}_{\EuScript{K}}\rt|^{2}\,.
\end{equation}
である．有効作用\eqref{SuperCond-S-eff}の各項には，図\ref{Fig-TDGL}に示す明瞭なダイアグラム表示がある．
\begin{figure}
\begin{center}\includegraphics[width=10cm]{figures/fig17.pdf}\end{center}
\caption{有効作用$S_{\mathrm{eff}}[\Delta,\mathbf{A},\Phi]$のダイアグラム表示．a）通常のGinzburg--Landau汎関数$S_{\mathrm{GL}}$（式\eqref{SuperCond-S-GL}参照）．b）スカラーポテンシャルと秩序変数の間の異常Gor'kov--Eliashberg結合（式\eqref{SuperCond-TDGL}および以下の議論を参照）．c）超電流作用$S_{\mathrm{SC}}$の常磁性部分，d）その反磁性部分．e）局所的なDOS項$S_{\mathrm{DOS}}$．f）非局所的なMT項$S_{\mathrm{MT}}$．ダイアグラムe）とf）では，ベクトルポテンシャルの選び方は2通りある．\textit{古典・量子}の組み合わせは電流の一部であり，\textit{量子・量子}の組み合わせは揺動散逸定理（FDT）によってそれに対応する項である．\label{Fig-TDGL}}
\end{figure}


有効作用\eqref{SuperCond-S-eff}に含まれる情報を等価に表す方法として，一組の確率的な時間依存Ginzburg--Landau方程式を用いることができる．これらを導くには，場の量子成分について2次の項，すなわち$S_{\mathrm{GL}}$に含まれる$\Delta^{q}_{\EuScript{K}}$の項と，$S_{\mathrm{N}}+S_{\mathrm{MT}}+S_{\mathrm{DOS}}$に含まれる$\mathbf{A}^{q}_{\EuScript{K}}$の項を取り除く必要がある．前者については，Hubbard--Stratonovich変換
\begin{equation}
\exp\left(-{\pi\nu\over 2}
\mathrm{Tr}\big\{|\Delta^{q}_{\EuScript{K}}|^{2}\big\}\right)=\int\D
[\xi_{\Delta}]\exp\left(-\frac{\pi\nu}{8T}\mathrm{Tr}\left\{
\frac{|\xi_{\Delta}|^{2}}{4T}
-i\xi^{*}_{\Delta}\Delta^{q}_{\EuScript{K}}-i\xi_{\Delta}
\Delta^{*q}_{\EuScript{K}}\right\} \right)\,.
\end{equation}
によって実現できる．その結果，式\eqref{SuperCond-S-eff}の有効作用$S_{\mathrm{eff}}$は秩序変数の量子成分について線形になる．量子成分について積分すると，確率的運動方程式を課す汎関数デルタ関数が得られる．こうしてTDGL方程式
\begin{equation}\label{SuperCond-KTDGL}
\left[\partial_{t}+\tau^{-1}_{\mathrm{GL}}-D\,
\big[\partial_{\mathbf{r}}-2ie\mathbf{A}^{cl}_{\EuScript{K}}\rt\big]^{2}
+\frac{7\zeta(3)}{\pi^{3}T}|\Delta^{cl}_{\EuScript{K}}\rt|^{2}\right]\Delta^{cl}_{\EuScript{K}}\rt=
\xi_{\Delta}\rt\,.
\end{equation}
が導かれる．複素Gauss雑音$\xi_{\Delta}\rt$の白色雑音相関関数は次式で与えられる．
\begin{equation}\label{SuperCond-xi-Delta}
\langle\xi_{\Delta}\rt\xi^{*}_{\Delta}(\mathbf{r}',t')\rangle=\frac{16T^{2}}{\pi\nu}\,
\delta(\mathbf{r}-\mathbf{r}')\delta(t-t')\,.
\end{equation}

同様に，ベクトル型のHubbard--Stratonovich場$\xi_{\mathbf{j}}\rt$を導入して，作用\eqref{SuperCond-S-eff}に含まれる$\mathbf{A}^{q}_{\EuScript{K}}$について2次の項を分離する：
\begin{equation}
\exp\left(-4T\Tr\big\{\sigma_{\mathbf{r},t,t'}
[\mathbf{A}^{q}_{\EuScript{K}}]^{2}\big\}\right)=\int\D[\xi_{\mathbf{j}}]\,
\exp\left(-\Tr\left\{\frac{\xi^{2}_{\mathbf{j}}}{4T\sigma_{\mathbf{r},t,t'}}
+2i\mathbf{A}^{q}_{\EuScript{K}}\xi_{\mathbf{j}}\right\}\right)\,,
\end{equation}
ここで$\sigma_{\mathbf{r},t,t'}=\sigma_{D}+e^{2}D\delta\nu^{\mathrm{DOS}}_{\mathbf{r},t}+e^{2}\nu\delta D^{\mathrm{MT}}_{\mathbf{r},t,t'}$は，DOSとMTの両方の繰り込みを含む完全な伝導率である．得られた作用は$\Phi^{q}_{\EuScript{K}}$と$\mathbf{A}^{q}_{\EuScript{K}}$の両方について線形なので，電荷密度$\varrho\rt=(1/2)\delta S_{\mathrm{eff}}/\delta\Phi^{q}\rt$と電流密度$\mathbf{j}\rt=(1/2)\delta S_{\mathrm{eff}}/\delta\mathbf{A}^{q}\rt$を定義できる．ここで微分は元の電磁ポテンシャル$\{\mathbf{A},\Phi\}$について行う一方，式\eqref{SuperCond-S-eff}の作用$S_{\mathrm{eff}}$はゲージ変換後のポテンシャル$\{\mathbf{A}_{\EuScript{K}},\Phi_{\EuScript{K}}\}$で書かれている点を強調しておく．両者の関係$\{\Phi,\mathbf{A}\}\rightleftarrows\{\mathbf{A}_{\EuScript{K}},\Phi_{\EuScript{K}}\}$は，式\eqref{int-ferm-NLSM-K-functional}に暗黙に含まれる汎関数$\EuScript{K}[\Phi]$によって与えられる．簡単な代数計算により，連続の式$\partial_{t}\varrho\rt+\mathrm{div}\,\mathbf{j}\rt=0$と，電流密度の式
\begin{eqnarray}\label{SuperCond-current}
\mathbf{j}\rt&=&\int\d t'\big[D\delta_{t-t'}+\delta
D^{\mathrm{MT}}_{\mathbf{r},t,t'}\big]\big[e^{2}\left(\nu +
\delta\nu^{\mathrm{DOS}}_{\mathbf{r},t'}\right) \mathbf{E}
(\mathbf{r},t')-\partial_{\mathbf{r}}\varrho(\mathbf{r},t')\big]\nonumber\\
&+&\frac{\pi e\nu D}{4T}\,
\mathrm{Im}\Big\{\Delta^{*cl}_{\EuScript{K}}\rt
\big[\partial_{\mathbf{r}}-2ie\mathbf{A}^{cl}_{\EuScript{K}}\rt\big]
\Delta^{cl}_{\EuScript{K}}\rt\Big\}+\xi_{\mathbf{j}}\rt\,,
\end{eqnarray}
を得る．ここで$\mathbf{E}\rt=\partial_{t}\mathbf{A}_{\EuScript{K}}-\partial_{\mathbf{r}}\Phi_{\EuScript{K}}$は電場である．電流の揺らぎは，相関関数
\begin{equation}\label{SuperCond-xi-j}
\big\langle \xi^{\mu}_{\mathbf{j}}(\mathbf{r},t)\,
\xi^{\nu}_{\mathbf{j}}(\mathbf{r'},t') \big\rangle
=\delta_{\mu\nu}\,T e^2
\left(2(\nu+\delta\nu^{\mathrm{DOS}}_{\mathbf{r},t} )D \delta_{t-t'}
+\nu \delta D^{\mathrm{MT}}_{\mathbf{r},t,t'} + \nu \delta
D^{\mathrm{MT}}_{\mathbf{r},t',t}\right)
\delta(\mathbf{r}-\mathbf{r'})\,,
\end{equation}
をもつベクトル型のGauss白色雑音によって生じ，これによりFDTの成立が保証される．式\eqref{SuperCond-KTDGL}と式\eqref{SuperCond-current}，および連続の式には，さらに電磁ポテンシャルに対するMaxwell方程式を補わなければならない．


TDGL方程式\eqref{SuperCond-KTDGL}を元のゲージで書き直してみると理解が深まる．これは，ゲージ変換後の秩序変数$\Delta^{cl}_{\EuScript{K}}=\Delta^{cl}\exp\big(-2ie\EuScript{K}^{cl}\big)$を式\eqref{SuperCond-KTDGL}に代入することで実現できる．こうして，元の秩序変数$\Delta^{cl}$に対して次の方程式を得る：
\begin{equation}\label{SuperCond-TDGL}
\fitdisplay{
\big[\partial_{t}-2ie\partial_{t}\EuScript{K}^{cl}\rt\big]\Delta^{cl}\rt=
\left[D\,\big[\partial_{\mathbf{r}}-2ie\mathbf{A}^{cl}\rt\big]^{2}-\tau^{-1}_{\mathrm{GL}}-
\frac{7\zeta(3)}{\pi^{3}T}|\Delta^{cl}\rt|^{2}\right]\Delta^{cl}\rt+\xi_{\Delta}\rt\,,
}
\end{equation}
ここで秩序変数の雑音を$\xi_{\Delta}\to\xi_{\Delta}\exp\big(2ie\EuScript{K}^{cl}\big)$と再定義したが，これはその相関関数\eqref{SuperCond-xi-Delta}を変えない．文献でしばしば見かけるTDGL方程式と異なり，式\eqref{SuperCond-TDGL}の左辺にはスカラーポテンシャル$\Phi^{cl}\rt$ではなく，Gor'kov--Eliashberg（GE）異常項$\partial_{t}\EuScript{K}^{cl}\rt$が含まれる（図\ref{Fig-TDGL}b参照）．一般に$\EuScript{K}^{cl}\rt$は，式\eqref{int-ferm-NLSM-Phi-K-matrix}で与えられる，スカラーポテンシャルと縦ベクトルポテンシャルの非局所汎関数である．古典成分については式\eqref{int-ferm-NLSM-Phi-K-matrix}から
\begin{equation}\label{SuperCond-K-Phi-functional}
\big(\partial_{t}-D\partial^{2}_{\mathbf{r}}\big)\EuScript{K}^{cl}\rt=\Phi^{cl}\rt-
D\,\mathrm{div}\mathbf{A}^{cl}\rt\,.
\end{equation}
を得る．空間的に一様なポテンシャルの場合，場$\partial_t \EuScript{K}^{cl}$と$\Phi^{cl}$は一致するが，一般には異なる．TDGL方程式の左辺にスカラーポテンシャル$\Phi^{cl}\rt$を書く通常の根拠はゲージ不変性にある．しかし局所ゲージ変換
\begingroup\interdisplaylinepenalty=10000
\begin{eqnarray}\label{SuperCond-gauge}
\Delta^{cl} &\to& \Delta^{cl}\,e^{-2ie\chi}\,,\quad\quad
\Phi^{cl} \to \Phi^{cl} - \partial_t\chi\,,\nonumber \\
 \mathbf{A}^{cl} &\to& \mathbf{A}^{cl} - \partial_{\mathbf{r}}\chi\,, \quad\,\,\,\,\,\,
\EuScript{K}^{cl} \to  \EuScript{K}^{cl} - \chi\,,
 \end{eqnarray}
\endgroup
は式\eqref{SuperCond-TDGL}を不変に保つため，この形のTDGL方程式は完全にゲージ不変である．式\eqref{SuperCond-gauge}の最後の関係は，式\eqref{SuperCond-K-Phi-functional}と$\Phi\rt$，$\mathbf{A}\rt$のゲージ変換則から直ちに従う．$\chi\rt=\EuScript{K}^{cl}\rt$によって指定される$\EuScript{K}$ゲージでは，TDGL方程式\eqref{SuperCond-TDGL}から異常GE項が消え，式\eqref{SuperCond-KTDGL}に戻る．


%$$$$$$$$$$$$$$$$$$$$$$$$$$$$$$$$$$$$$$$$$$$$$$$$$$$$$$$$$$$$$$$$$$$$$$$$$$$$$$$$$$$$$$$$$$$$$$$$$$$$$$$$$$$
\subsectionlink{応用IV：非一様な超伝導と揺らぐ超伝導}\label{app_Part-IV}

%$$$$$$$$$$$$$$$$$$$$$$$$$$$$$$$$$$$$$$$$$$$$$$$$$$$$$$$$$$$$$$$$$$$$$$$$$$$$$$$$$$$$$$$$$$$$$$$$$$$$$$$$$$$
\subsubsection{近接効果}\label{app_Part-IV-1}

超伝導体との界面の近くでは，常伝導金属は部分的に超伝導の性質を獲得する．同時に，超伝導体の超伝導性は常伝導金属によって弱められる．この相互の影響を\textit{近接}効果と呼ぶ．第\ref{sec_SuperCond-2}節で論じた準古典Usadel方程式と運動論方程式は，超伝導体・常伝導金属構造における近接効果に関係する現象を十分に記述する．本節ではその一例を考察する．

2つのバルク超伝導体の間に長さ$L$の拡散的な常伝導細線を置き，超伝導体・常伝導金属・超伝導体（SNS）接合を形成した場合を考える．超伝導体との近接が，常伝導細線内の準粒子エネルギースペクトルをどのように変えるかを調べたい．Usadel方程式\eqref{SuperCond-Usadel-2}から，細線の状態密度にはエネルギーギャップ$\epsilon_{g}$が生じ，その上の$\epsilon>\epsilon_{g}$では平方根型の非解析的挙動$\sim \sqrt{\epsilon-\epsilon_{g}}$を示すことが分かる\cite{Zhou,GolubovKupriyanov}．これを明示的に見るため，細線の断面の寸法が超伝導コヒーレンス長$\xi=\sqrt{D/\Delta}$より十分小さいと仮定する．この場合，細線は準1次元とみなせ，すべての変化は細線に沿った$x$座標の方向に生じる．細線内に引力相互作用がなく$\lambda=0$なら，自己無撞着方程式\eqref{SuperCond-Usadel-4}より，対ポテンシャルは細線内$-L/2\leqslant x\leqslant L/2$で$\Delta(\mathbf{r})=0$，この区間の外では$\Delta(\mathbf{r})=\Delta$となる．さらに2つの超伝導体の間に位相差がなく$\partial_{x}\chi=0$なら，Usadel方程式\eqref{SuperCond-Usadel-2}は大幅に簡単になり，
\begin{equation}\label{Part-IV-Usadel}
D\,\partial^{2}_{x}\theta(x,\epsilon)+2i\epsilon\sinh\theta(x,\epsilon)=0\,.
\end{equation}
と書ける．超伝導体との界面$x=\pm L/2$では，この方程式に境界条件$\theta(\pm L/2,\epsilon)=\theta_{BCS}(\epsilon)$を補う．ここで$\tanh\theta_{BCS}(\epsilon)=\Delta/\epsilon$である．超伝導体は非常に大きく，細線による摂動は無視できるため，超伝導体の内部全域で座標に依存しない$\theta_{BCS}(\epsilon)$を用いられると仮定している．式\eqref{Part-IV-Usadel}を解けば，状態密度は$\nu(x,\epsilon)=\nu\mathrm{Re}\big[\cosh\theta(x,\epsilon)\big]$として得られる．


変換$\theta(x,\epsilon)=i\pi/2-\vartheta(x,\epsilon)$を行うと便利である．これにより式\eqref{Part-IV-Usadel}は実数の式となり，直接積分できる：
\begin{equation}\label{Part-IV-Usadel-sol}
\sqrt{\frac{\epsilon}{E_{Th}}}=\int^{\vartheta_{0}}_{\vartheta_{BCS}}
\frac{\d\vartheta}{\sqrt{\sinh\vartheta_{0}-\sinh\vartheta}}
\equiv K(\vartheta_{0},\epsilon)\,,
\end{equation}
ここで$E_{Th}= D/L^2$，$\vartheta_{0}=\vartheta(0,\epsilon)$，$\sinh\vartheta_{BCS}=\epsilon/\sqrt{\Delta^2-\epsilon^{2}}$である．式\eqref{Part-IV-Usadel-sol}は，$\vartheta_{0}$をエネルギー$\epsilon$の関数として定める．$\vartheta_{0}(\epsilon)$が分かれば，細線の中央における状態密度は$\nu(0,\epsilon)=\nu\mathrm{Im}[\sinh\vartheta_{0}(\epsilon)]$と求まる．


長い細線の極限$\xi\ll L$では，状態密度の変化はギャップより十分低いエネルギー$\epsilon\ll\Delta$で生じる．したがって$\vartheta_{BCS}\approx0$と近似でき，式\eqref{Part-IV-Usadel-sol}の右辺の関数は実質的にエネルギーに依存せず，$K(\vartheta_{0},\epsilon)\approx K(\vartheta_0,0)$となる．これは$\vartheta^{*}_{0}\approx1.5$で最大値$K_{\mathrm{max}}=K(\vartheta^{*}_{0})\approx1.75$をとる．一方，式\eqref{Part-IV-Usadel-sol}の左辺は$\epsilon>K^{2}_{\mathrm{max}}E_{Th}=\epsilon_{g}$に対して$K_{\mathrm{max}}$を超えうる．したがって，すべてのエネルギー$\epsilon<\epsilon_{g}$について，式\eqref{Part-IV-Usadel-sol}の$\vartheta_{0}$の解は実数のみであり，$\nu(0,\epsilon)\propto\mathrm{Im}\big[\sinh\vartheta_{0}\big]$なので$\nu(0,\epsilon)\equiv0$となる．$\epsilon>\epsilon_{g}$では関数$\vartheta_{0}$が複素数となり，有限の状態密度を与える．ギャップの直上$0 <\epsilon-\epsilon_{g}\ll\epsilon_{g}$では，式\eqref{Part-IV-Usadel-sol}を用いて
\begin{equation}
\nu(\epsilon)=3.7\delta^{-1}\sqrt{\frac{\epsilon}{\epsilon_{g}}-1}\,,
\end{equation}
を得る．ここで$\nu(\epsilon)=\mathcal{A}\int\nu(x,\epsilon)\d x$は，細線の体積全体で積分した総状態密度である（$\mathcal{A}$は細線の断面積，$\delta=1/(\nu\mathcal{A}L)$はその準位間隔である）．$\epsilon_{g}\sim E_{Th}\ll\Delta$なので，近似$\vartheta_{BCS}(\epsilon\sim\epsilon_{g})\approx0$は十分正当化される．


反対の短い細線の極限$L\ll\xi$，すなわち$E_{Th}\gg\Delta$でも，式\eqref{Part-IV-Usadel-sol}は適用できる．ただし，$\vartheta_{BCS}(\epsilon)$のエネルギー依存性をすべて保持しなければならない．エネルギーギャップは$\epsilon_{g}=\Delta-\Delta^{3}/8E^{2}_{Th}$で与えられ，バルクのギャップ$\Delta$よりわずかに小さいだけであることを示せる．短い細線では近接効果が強いと期待されるので，これは自然である．誘起されたギャップの直上で，状態密度は再び平方根型の非解析性を示す．ただし，その係数は大きく，$\nu(\epsilon)\sim \delta^{-1}(E_{Th}/\Delta)^{2} \sqrt{\epsilon/\epsilon_{g}-1}$となる（文献\cite{Levchenko-DOS}参照）．

%$$$$$$$$$$$$$$$$$$$$$$$$$$$$$$$$$$$$$$$$$$$$$$$$$$$$$$$$$$$$$$$$$$$$$$$$$$$$$$$$$$$$$$$$$$$$$$$$$$$$$$$$$$$
\subsubsection{Josephson電流}\label{app_Part-IV-2}

Usadel方程式\eqref{SuperCond-Usadel-1}と\eqref{SuperCond-Usadel-2}によって扱えるもう1つの例はJosephson効果である．前節と同じSNS接合の配置を考え，接合両端の対ポテンシャルの間に有限の位相差，すなわち$\chi(L/2,\epsilon)-\chi(-L/2,\epsilon)=\phi$があると仮定する．この条件では，Josephson超電流$I_{S}(\phi)$が接合を横切って流れうる．本節の目的は，Usadel方程式からJosephson電流の位相・電流関係がどのように得られるかを示すことである．

$|x|>L/2$で$\Delta(x)=\Delta$，$|x|<L/2$で$\Delta=0$となる階段関数型の対ポテンシャル模型では，式\eqref{SuperCond-Usadel-1}，\eqref{SuperCond-Usadel-2}は次の形になる：
\begin{subequations}\label{Part-IV-Usadel-coupled}
\begin{equation}
\hskip-2.2cm
D\,\partial_{x}\big(\sinh^{2}\theta\partial_{x}\chi\big)=0\,,
\end{equation}
\begin{equation}
D\,\partial^{2}_{x}\theta+2i\epsilon\sinh\theta=
\frac{D}{2}(\partial_{x}\chi)^{2}\sinh2\theta\,.
\end{equation}
\end{subequations}
これらには境界条件$\theta(\pm L/2,\epsilon)=\theta_{BCS}(\epsilon)$を補う．$\chi$関数の境界条件は，先に述べた接合両端の固定された位相差$\phi$によって，$\chi(L/2,\epsilon)-\chi(-L/2,\epsilon)=\phi$と定まる．短い細線$L\ll\xi$では，式(\ref{Part-IV-Usadel-coupled}b)の左辺第2項は勾配項より$\epsilon/E_{Th}\ll1$だけ小さく，無視できる．式(\ref{Part-IV-Usadel-coupled}a)には第1積分$\sinh^{2}\theta\partial_{x}\chi=\mathcal{J}/L$があるため，式(\ref{Part-IV-Usadel-coupled}b)から$\partial_{x}\chi$を消去すると，$L^2\partial^{2}_{x}\theta=\mathcal{J}^{2}\cosh\theta/\sinh^{3}\theta$を得る．この方程式は厳密に解ける：
\begin{equation}\label{Part-IV-Usadel-coupled-sol}
\cosh\theta(z,\epsilon)=\cosh\theta_{0}
\cosh\left(\frac{\mathcal{J}z}{\sinh\theta_{0}}\right)\,,
\end{equation}
ここで$\theta_{0}=\theta(0,\epsilon)$，$z=x/L$である．$\theta(x,\epsilon)$が得られたら，これを式(\ref{Part-IV-Usadel-coupled}a)の第1積分$\phi= \int^{L/2}_{-L/2}\d x\partial_x\chi = \mathcal{J}\int^{1/2}_{-1/2}\d z/\sinh^{2}\theta(z,\epsilon)$に戻して，
\begin{equation}
\tan(\phi/2)=\frac{1}{\sinh\theta_{0}}
\tanh\left(\frac{\mathcal{J}}{2\sinh\theta_{0}}\right)\,.
\end{equation}
を得る．この式は，NS界面$z=\pm1/2$で評価した式\eqref{Part-IV-Usadel-coupled-sol}とともに，2つの未知量$\mathcal{J}$と$\theta_{0}$に対する2つの代数方程式を構成する．この代数問題は容易に解け，$\mathcal{J}(\epsilon,\phi)=2\sinh\theta_{0}\mathrm{arctanh}\big[\sinh\theta_{0}\tan(\phi/2)\big]$および$\sinh\theta_{0}=\sinh\theta_{BCS}/\sqrt{1+\tan^{2}(\phi/2)\cosh^{2}\theta_{BCS}}$を得る．ここで$\cosh\theta_{BCS}=\epsilon/\sqrt{\epsilon^{2}-\Delta^{2}}$である．$\mathcal{J}(\epsilon,\phi)$が分かれば，Josephson電流は
\begin{equation}\label{Part-IV-Josephson-def}
I_{S}(\phi)=\frac{\mathrm{g}_{D}}{e}\int^{\infty}_{0}\d\epsilon\,
\tanh\left(\frac{\epsilon}{2T}\right)\,\mathrm{Im}\,\mathcal{J}(\epsilon,\phi)\,,
\end{equation}
によって得られる．ここで$\mathrm{g}_{D}$は細線のコンダクタンスである．得られた$\mathcal{J}(\epsilon,\phi)$の解を用いると，
\begin{equation}\label{Part-IV-Im-J}
\mathrm{Im}\, \mathcal{J}(\epsilon,\phi)=\frac{\pi\Delta\cos(\phi/2)}
{\sqrt{\epsilon^{2}-\Delta^{2}\cos^{2}(\phi/2)}}
\end{equation}
が$\Delta\cos(\phi/2)<\epsilon<\Delta$に対して成立し，それ以外では$\mathrm{Im}\,\mathcal{J}(\epsilon,\phi)=0$となる．式\eqref{Part-IV-Josephson-def}と\eqref{Part-IV-Im-J}を用いると，短い拡散SNS接合の零温度Josephson電流について，KulikとOmelyanchuk\cite{KulikOmelyanchuk}が導いた結果
\begin{equation}
I_{S}(\phi)=\frac{\pi\mathrm{g}_{D}\Delta}{e}\cos(\phi/2)\,
\mathrm{arctanh}\big[\sin(\phi/2)\big]\,.
\end{equation}
に到達する．原論文\cite{KulikOmelyanchuk}では，$I_{S}(\phi)$の導出に虚時間の手法が用いられた．この結果は後に，実時間（エネルギー）のUsadel方程式を用いて文献\cite{Heikkila,Levchenko-Js}で再現された．


%$$$$$$$$$$$$$$$$$$$$$$$$$$$$$$$$$$$$$$$$$$$$$$$$$$$$$$$$$$$$$$$$$$$$$$$$$$$$$$$$$$$$$$$$$$$$$$$$$$$$$$$$$$$
\subsubsection{$T_c$以上での状態密度の抑制}\label{app_Part-IV-3}

$T_{c}$以下の超伝導体は，励起スペクトルにエネルギーギャップ$|\Delta(T)|$をもつ．$T_c$より十分高温では，状態密度は金属的で一定である．転移近傍$0<T-T_c\ll T_{c}$で現れる超伝導揺らぎの効果の1つは，Fermiエネルギー付近の状態密度の減少である．揺らぎによる状態密度の抑制は温度の低下とともに増大し，最終的に完全なギャップへと移行する．本節ではKeldysh形式を用いてこの効果の温度依存性を計算し，松原の手法と解析接続が用いられた初期の研究\cite{Hurault,Abrahams}と比較する．包括的な議論については，LarkinとVarlamovの近著\cite{Larkin-Varlamov}を参照されたい．
\begin{figure}
\begin{center}\includegraphics[width=10cm]{figures/fig18.pdf}\end{center}
\caption{臨界温度$T_{c}$近傍の状態密度補正\eqref{Part-IV-Dos-corr}に対するダイアグラム．梯子で示された2つのCooperon場$c$と$c^{*}$は，塗りつぶした三角形で表した秩序変数$\Delta^{cl(q)}$に結合し，秩序変数同士は揺らぎ伝播関数で対にされる．\label{Fig-DOS}}
\end{figure}

出発点は，$\check{Q}$行列で表した状態密度の式
$$\nu(\varepsilon)=\frac{\nu}{4}\left\langle\Tr\{\hat{\sigma}_{z}\otimes\hat{\tau}_{z}
\check{Q}_{\varepsilon\varepsilon}\}\right\rangle_{Q}\,,$$
である（第\ref{app_Part-II-2}節参照）．$\check{Q}=\check{\Lambda}$と置くと，常伝導金属に対して当然の結果$\nu(\varepsilon)=\nu$を得る．式(\ref{SuperCond-Q-W})を用いて$\check{Q}$をCooperon揺らぎ$\check{\mathcal{W}}$について2次まで展開すると，状態密度の補正として
\begin{equation}\label{Part-IV-Dos-corr-def}
\delta\nu(\varepsilon)=\frac{\nu}{4}\sum_{q}\int\frac{\d\varepsilon'}{2\pi}
\left\langle\left\langle c_{\varepsilon\varepsilon'}(\mathbf{q})
c^{*}_{\varepsilon'\varepsilon}(-\mathbf{q})+
\bar{c}_{\varepsilon\varepsilon'}(\mathbf{q})
\bar{c}^{*}_{\varepsilon'\varepsilon}(-\mathbf{q})
\right\rangle\right\rangle_{\mathcal{W},\Delta}\,.
\end{equation}
を得る．次に，揺らぐ$c$場と$\bar c$場について平均をとる．そのため，Cooperモード$c$，$\bar c$を秩序変数の揺らぎと結びつける式\eqref{app_FE-c-extremal}を用いる．後者は次の相関関数によって支配される：
\begin{eqnarray}\label{Part-IV-Delta-averages}
&&\left\langle
\Delta^{cl}(\mathbf{q},\omega)\Delta^{*cl}(-\mathbf{q},-\omega)\right\rangle_{\Delta}=
\frac{i}{2\nu}\, L^{K}(\mathbf{q},\omega) \,,\quad \left\langle
\Delta^{cl}(\mathbf{q},\omega)\Delta^{*q}(-\mathbf{q},-\omega)\right\rangle_{\Delta}=
\frac{i}{2\nu}\, L^{R}(\mathbf{q},\omega) \,,\nonumber\\
&&\left\langle
\Delta^{q}(\mathbf{q},\omega)\Delta^{*cl}(-\mathbf{q},-\omega)\right\rangle_{\Delta}=
\frac{i}{2\nu}\, L^{A}(\mathbf{q},\omega) \,,\quad \left\langle
\Delta^{q}(\mathbf{q},\omega)\Delta^{*q}(-\mathbf{q},-\omega)\right\rangle_{\Delta}=0\,,
\end{eqnarray}
これらは時間依存Ginzburg--Landau作用\eqref{SuperCond-S-GL}から従う．その結果，Cooperon場の相関関数として
\begin{subequations}\label{Part-IV-c-averages}
\begin{equation}
\left\langle\left\langle
c_{\varepsilon,\varepsilon-\omega}(\mathbf{q})
c^{*}_{\varepsilon-\omega,\varepsilon}(-\mathbf{q})\right\rangle\right\rangle=\frac{2i}{\nu}\,
\frac{L^{K}+F_{\varepsilon-\omega}L^{R}+F_{\varepsilon}L^{A}}
{\big(Dq^{2}-2i\varepsilon+i\omega\big)^{2}}\,,
\end{equation}
\begin{equation}
\left\langle\left\langle
\bar{c}_{\varepsilon,\varepsilon-\omega}(\mathbf{q})
\bar{c}^{*}_{\varepsilon-\omega,\varepsilon}(-\mathbf{q})\right\rangle\right\rangle=\frac{2i}{\nu}\,
\frac{L^{K}-F_{\varepsilon-\omega}L^{A}-F_{\varepsilon}L^{R}}
{\big(Dq^{2}+2i\varepsilon-i\omega\big)^{2}}\,.
\end{equation}
\end{subequations}
を得る．これらを式\eqref{Part-IV-Dos-corr-def}に代入し，2つの寄与を足し合わせると，
\begin{equation}\label{Part-IV-Dos-corr}
\delta\nu(\varepsilon)=\mathrm{Im}\sum_{q}
\int^{+\infty}_{-\infty}\frac{\d\omega}{2\pi}\,
\frac{L^{K}\qo+F_{\varepsilon-\omega}L^{R}\qo}
{\big(Dq^{2}-2i\varepsilon+i\omega\big)^{2}}\,,
\end{equation}
となる．ここで，平均$\langle\langle cc^{*}\rangle\rangle$と$\langle\langle\bar{c}\bar{c}^{*}\rangle\rangle$に含まれる$F_{\varepsilon}L^{A(R)}(\omega)$に比例する項は，それぞれ純粋な先進関数および遅延関数の積分であるため，$\omega$積分によって式\eqref{Part-IV-Dos-corr}から消える．式\eqref{Part-IV-Dos-corr}には，図\ref{Fig-DOS}に示す便利なダイアグラム表示がある．ここで，式\eqref{app_FE-L-R}の形の揺らぎ伝播関数を用いる．関係する周波数は$\omega\sim T-T_c\ll T_c$なので，ボソン分布関数を$B_{\omega}\approx 2T_{c}/\omega$と近似すると，状態密度補正\eqref{Part-IV-Dos-corr}は
\begin{equation}\label{Part-IV-Dos-intermediate}
\delta\nu(\varepsilon)=-\frac{16T^{2}_{c}}{\pi^{2}}\mathrm{Re}\sum_{q}
\int^{+\infty}_{-\infty}\frac{\d\omega}
{\big[(Dq^{2}+\tau^{-1}_{\mathrm{GL}})^{2}+\omega^{2}\big]
\big[Dq^{2}-2i\varepsilon+i\omega\big]^{2}}\,,
\end{equation}
に帰着する．ここで$\tau_{\mathrm{GL}}^{-1}=8(T-T_c)/\pi$である．


この式のさらに詳しい解析は，系の有効次元に依存する．ここでは準2次元の場合，すなわち厚さ$b$が超伝導コヒーレンス長より十分小さい金属薄膜$b\ll\xi(T)=\sqrt{D\tau_{\mathrm{GL}}}$に注目する．運動量和を積分$\sum_{q}\to \frac{1}{b}\int\frac{\d q^{2}}{4\pi^2}$で置き換え，無次元パラメータ$x=Dq^{2}/T_{c}$および$y=\omega/T_{c}$を導入し，留数定理を用いて$y$積分を行うと，
\begin{equation}\label{Part-III-DOS-result}
\frac{\delta\nu(\varepsilon)}{\nu}=-\frac{\mathrm{Gi}}{16}
\left(\frac{T_{c}}{T-T_{c}}\right)^{2}\mathcal{Y}(\varepsilon\tau_{\mathrm{GL}}),\qquad
\mathcal{Y}(z)=\mathrm{Re}\int^{+\infty}_{0}\frac{\d
x}{(1+x)(1+2x-2iz)^{2}}\,,
\end{equation}
を得る．ここで$\mathrm{Gi}=\hbar /\nu Db$はGinzburg数である．Fermiエネルギーからのずれが小さい$\varepsilon\tau_{\mathrm{GL}}\ll1$では，DOSの抑制は$\delta\nu(0)\propto - (T/T_{c}-1)^{-2}$に従う．より高いエネルギー$\varepsilon\tau_{\mathrm{GL}}\gg1$では，DOSは$\delta\nu(\varepsilon)\propto -(T_{c}/\varepsilon)^{2}\ln(\varepsilon\tau_{\mathrm{GL}})$として常伝導状態の値に近づく．また，$\int\d\varepsilon\,\delta\nu(\varepsilon)\equiv0$である．揺らぎはFermiエネルギーの周りで状態を再分配するだけなので，これは期待される結果である．


%$$$$$$$$$$$$$$$$$$$$$$$$$$$$$$$$$$$$$$$$$$$$$$$$$$$$$$$$$$$$$$$$$$$$$$$$$$$$$$$$$$$$$$$$$$$$$$$$$$$$$$$$$$$
\subsubsection{伝導率に対する揺らぎ補正}\label{app_Part-IV-4}

$T_c$以上の超伝導揺らぎは，状態密度だけでなく輸送特性も変える．伝導率の場合には，状態密度（DOS）項$\delta\sigma_{\mathrm{DOS}}$，Aslamazov--Larkin（AL）項$\delta\sigma_{\mathrm{AL}}$，Maki--Thompson（MT）項$\delta\sigma_{\mathrm{MT}}$という3種類の補正がある\cite{AslamazovLarkin,Maki,Thompson,Hurault}．この話題は第\ref{sec_SuperCond-3}節ですでに部分的に論じたが，本節では，それらすべてがKeldysh $\sigma$模型の方法でどのように得られるかを明示する．

式\eqref{Part-II-sigma-def}の定義によると，伝導率を求めるには，ベクトルポテンシャルについて2次まで分配関数$\mathcal{Z}[\mathbf{A}^{cl},\mathbf{A}^{q}]$を知る必要がある．式\eqref{SuperCond-S}を用いると，次を得る\footnote{$T_c$近傍ではCoulomb相互作用は運動論的係数に特異な温度依存性を与えないので，本節を通して$\Phi_{\EuScript{K}}=0$とし，添字$\EuScript{K}$を省略する．}：
\begin{equation}\label{Part-IV-Z}
\fitdisplay{
\mathcal{Z}[\mathbf{A}^{cl},\mathbf{A}^{q}]\approx\int\D[\check Q,\Delta]\left[1+\frac{\pi\nu
D}{2}\Tr\big\{\check{\Xi}\check{\mathbf{A}}\check{Q}
\check{\Xi}\check{\mathbf{A}}\check{Q}\big\}-\frac{(\pi\nu
D)^{2}}{8}\left(\Tr\big\{\partial_{\mathbf{r}}\check{Q}
[\check{\Xi}\check{\mathbf{A}},\check{Q}]\big\}\right)^{2}\right]
\exp\big(iS_{\sigma}[\check Q,\Delta]\big)\,,
}
\end{equation}
ここでは反磁性の寄与$\Tr\{\check{\Xi}\check{\mathbf{A}}\check{\Xi}\check{\mathbf{A}}\}$を省略した．第\ref{Part-II-3}節で示したように，$\check{Q}=\check{\Lambda}$と置き，式\eqref{Part-II-sigma-def}を用いるとDrude伝導率$\sigma_{D}$を得る．常伝導金属の伝導率$\sigma_{D}$への超伝導補正$\delta\sigma$を捉えるには，$\check{Q}$行列を揺らぎ$\check{\mathcal{W}}$について最低次（2次）まで展開し，可能な寄与をすべて調べる必要がある．


式\eqref{Part-IV-Z}の右辺第1のトレースにおいて，一方の$\check{Q}$行列を$\check{\Lambda}$とし，他方を$\check{\mathcal{W}}^{2}$の次数まで展開すると，
\begin{equation}\label{Part-IV-Z-DOS}
\mathcal{Z}_{\mathrm{DOS}}[\mathbf{A}^{cl},\mathbf{A}^{q}]=\frac{\pi\nu
D}{2}\left\langle\left\langle\Tr\{\check{\mathbb{A}}_{\varepsilon_{1}\varepsilon_{2}}
(\hat{\sigma}_{z}\otimes\hat{\tau}_{z})
\check{\mathbb{A}}_{\varepsilon_{2}\varepsilon_{3}}
(\hat{\sigma}_{z}\otimes\hat{\tau}_{z})
\check{\mathcal{W}}_{\varepsilon_{3}\varepsilon_{4}}
\check{\mathcal{W}}_{\varepsilon_{4}\varepsilon_{1}}\}
\right\rangle\right\rangle_{\mathcal{W},\Delta}\,,
\end{equation}
を得る．ここで電流頂点の行列は
\begin{equation}
\check{\mathbb{A}}_{\varepsilon\varepsilon'}\equiv
\check{\mathcal{U}}_{\varepsilon}^{-1}
\check{\Xi}\check{\mathbf{A}}_{\varepsilon-\varepsilon'}\check{\mathcal{U}}_{\varepsilon'}
=\left(\begin{array}{cc}
\mathbf{A}^{cl}_{\varepsilon-\varepsilon'}+F_{\varepsilon}
\mathbf{A}^{q}_{\varepsilon-\varepsilon'} &
\mathbf{A}^{q}_{\varepsilon-\varepsilon'}[F_{\varepsilon}F_{\varepsilon'}
- 1]
+\mathbf{A}^{cl}_{\varepsilon-\varepsilon'}[F_{\varepsilon'}-F_{\varepsilon}]\\
-\mathbf{A}^{q}_{\varepsilon-\varepsilon'} &
\mathbf{A}^{cl}_{\varepsilon-\varepsilon'} -
F_{\varepsilon'}\mathbf{A}^{q}_{\varepsilon-\varepsilon'}
\end{array}\right)_K\otimes\hat \tau_{z}\,.
\end{equation}
である．すぐ後で示すように，$\mathcal{Z}_{\mathrm{DOS}}$は臨界温度近傍の伝導率に対する状態密度型の寄与を定める．実際，式\eqref{Part-IV-Z-DOS}を式\eqref{Part-II-sigma-def}に代入し，ベクトルポテンシャルで微分し，直流極限$\Omega\to0$をとって行列のトレースを計算すると，
\begin{equation}\label{Part-IV-sigma-DOS-cc}
\delta\sigma_{\mathrm{DOS}}=\frac{\pi e^2\nu D}{2}
\sum_{q}\iint\frac{\d\varepsilon_{2}\d\varepsilon_{4}}{4\pi^2}\,
\partial_{\varepsilon_{2}}F_{\varepsilon_{2}}
\left\langle\left\langle
c_{\varepsilon_{2}\varepsilon_{4}}(\mathbf{q})
c^{*}_{\varepsilon_{4}\varepsilon_{2}}(-\mathbf{q})+
\bar{c}_{\varepsilon_{2}\varepsilon_{4}}(\mathbf{q})
\bar{c}^{*}_{\varepsilon_{4}\varepsilon_{2}}(-\mathbf{q})
\right\rangle\right\rangle_{\mathcal{W},\Delta} \,.
\end{equation}
を得る．次に式\eqref{app_FE-c-extremal}を用い，相関関数\eqref{Part-IV-Delta-averages}により$\Delta$について平均をとる．積分変数を$\varepsilon_{2}\to\varepsilon$および$\varepsilon_{4}\to\varepsilon-\omega$と変えると，補正$\delta\sigma_{\mathrm{DOS}}$は
\begin{equation}\label{Part-IV-sigma-DOS-analytic}
\delta\sigma_{\mathrm{DOS}}=\frac{e^{2}D}{2\pi}\,\mathrm{Im}\sum_{q}\int
^{+\infty}_{-\infty}\d\varepsilon\d\omega\,
\partial_{\varepsilon}F_{\varepsilon}\, \frac{L^{K}\qo+
F_{\varepsilon-\omega}L^{R}\qo}{\big(Dq^{2}-2i\varepsilon+i\omega\big)^{2}}\,
.
\end{equation}
となる．この式を式\eqref{Part-IV-Dos-corr}と比較すると，$\delta\sigma_{\mathrm{DOS}}\propto\int\d\varepsilon\,\partial_{\varepsilon}F_{\varepsilon}\delta\nu(\varepsilon)$が分かり，$\delta\sigma_{\mathrm{DOS}}$と状態密度の抑制$\delta\nu(\varepsilon)$の関係が確立される．ダイアグラム表示については図\ref{Fig-AL-MT-DOS}aも参照されたい．温度差$T-T_{c}$のべきに関して，$\delta\sigma_{\mathrm{DOS}}$の最も強く発散する部分を取り出すには，式\eqref{Part-IV-sigma-DOS-analytic}でKeldysh伝播関数だけを残せばよい．$F_{\varepsilon-\omega}L^{R}$の項の寄与は，パラメータの次数で見て小さい．式\eqref{app_FE-L-R}を用いると，
\begin{equation}
\delta\sigma_{\mathrm{DOS}}=-\frac{16e^{2}DT^{2}_{c}}{\pi^{2}}\,\mathrm{Re}\sum_{q}
\iint^{+\infty}_{-\infty}\d\varepsilon\d\omega\,
\frac{\partial_{\varepsilon}F_{\varepsilon}}
{\big[(Dq^{2}+\tau^{-1}_{\mathrm{GL}})^{2}+\omega^{2}\big]
\big[Dq^{2}-2i\varepsilon+i\omega\big]^{2}}\,.
\end{equation}
を得る．残りの周波数積分と運動量積分を行うと，準2次元の場合には
\begin{equation}\label{Part-IV-sigma-DOS}
\frac{\delta\sigma_{\mathrm{DOS}}}{\sigma_{D}}=
-\frac{7\zeta(3)\mathrm{Gi}}{\pi^{4}}
\ln\left(\frac{T_{c}}{T-T_{c}}\right)\,.
\end{equation}
を得る．この補正は，揺らぎによる状態密度の減少に起因するため予想どおり負であり，温度依存性は比較的弱い．$\delta\sigma_{\mathrm{DOS}}$は，第\ref{sec_SuperCond-3}節で論じた有効な時間依存Ginzburg--Landau理論からも取り出せることを強調しておく．実際，式\eqref{SuperCond-DOS-corr}の$\delta\nu^{\mathrm{DOS}}_{\mathbf{r},t}$を用いた$\delta\sigma_{\mathrm{DOS}}=e^2D\langle\delta\nu^{\mathrm{DOS}}_{\mathbf{r},t}\rangle_{\Delta}$が，式\eqref{Part-IV-sigma-DOS}を再現することを示せる．


\begin{figure}
\begin{center}\includegraphics[width=10cm]{figures/fig19.pdf}\end{center}
\caption{$T_{c}$近傍における伝導率への超伝導揺らぎ補正のダイアグラム．a）状態密度項，b）Maki--Thompson補正，c）Aslamazov--Larkin補正．\label{Fig-AL-MT-DOS}}
\end{figure}

式\eqref{Part-IV-Z}に戻り，ほかにどのような寄与がありうるかを調べよう．再び式\eqref{Part-IV-Z}の右辺第1のトレースに注目し，今度は各$\check{Q}$行列を揺らぎ$\check{\mathcal{W}}$について1次まで展開する．こうして，
\begin{equation}\label{Part-IV-Z-MT}
\mathcal{Z}_{\mathrm{MT}}[\mathbf{A}^{cl},\mathbf{A}^{q}]=\frac{\pi\nu
D}{2}\left\langle\left\langle\Tr\big\{\check{\mathbb{A}}_{\varepsilon_{1}\varepsilon_{2}}
(\hat{\sigma}_{z}\otimes\hat{\tau}_{z})
\check{\mathcal{W}}_{\varepsilon_{2}\varepsilon_{3}}
\check{\mathbb{A}}_{\varepsilon_{3}\varepsilon_{4}}
(\hat{\sigma}_{z}\otimes\hat{\tau}_{z})
\check{\mathcal{W}}_{\varepsilon_{4}\varepsilon_{1}}
\big\}\right\rangle\right\rangle_{\mathcal{W},\Delta}\,,
\end{equation}
を得る．これは伝導率へのMaki--Thompson補正を与える．$\mathcal{Z}_{\mathrm{MT}}[\mathbf{A}^{cl},\mathbf{A}^{q}]$をベクトルポテンシャルで微分してトレースを計算すると，直流極限で
\begin{equation}\label{Part-IV-sigma-MT-cc}
\delta\sigma_{\mathrm{MT}}=\frac{\pi e^2\nu D}{2}
\sum_{q}\iint\frac{\d\varepsilon_{2}\d\varepsilon_{4}}{4\pi^2}\,
\partial_{\varepsilon_{2}}F_{\varepsilon_{2}}
\left\langle\left\langle
c_{\varepsilon_{2}\varepsilon_{4}}(\mathbf{q})
\bar{c}^{*}_{\varepsilon_{4}\varepsilon_{2}}(-\mathbf{q})
+c^{*}_{\varepsilon_{2}\varepsilon_{4}}(\mathbf{q})
\bar{c}_{\varepsilon_{4}\varepsilon_{2}}(-\mathbf{q})
\right\rangle\right\rangle_{\mathcal{W},\Delta}.
\end{equation}
を得る．式\eqref{Part-IV-sigma-DOS-cc}の$\delta\sigma_{\mathrm{DOS}}$と比較すると，$\delta\sigma_{\mathrm{MT}}$は遅延$c$と先進$\bar{c}$のCooperonが混じった積から成るのに対し，$\delta\sigma_{\mathrm{DOS}}$は同じ因果性をもつCooperon場を含む．式\eqref{Part-IV-Delta-averages}と\eqref{app_FE-c-extremal}を用いて式\eqref{Part-IV-sigma-MT-cc}の$\Delta$揺らぎについて平均をとり，式\eqref{Part-IV-sigma-DOS-analytic}と同様に積分変数を変えると，
\begin{equation}\label{Part-IV-sigma-MT-analytic}
\delta\sigma_{\mathrm{MT}}=-\frac{e^{2}D}{\pi}\sum_{q}
\int^{+\infty}_{-\infty}\d\varepsilon\d\omega\,
\partial_{\varepsilon}F_{\varepsilon}\,\frac{\mathrm{Im}[L^{R}\qo]
(B_{\omega}-F_{\varepsilon-\omega})}
{(Dq^{2})^{2}+(2\varepsilon+\omega)^{2}} \,.
\end{equation}
に到達する．対応するダイアグラムを図\ref{Fig-AL-MT-DOS}bに示す．先ほどと同じ精度で$B_{\omega}\approx2T_{c}/\omega$と近似し，$F_{\varepsilon-\omega}$を無視し，揺らぎ伝播関数に式\eqref{app_FE-L-R}を用いると，この$\delta\sigma_{MT}$の式は
\begin{equation}
\delta\sigma_{\mathrm{MT}}=\frac{16e^{2}DT^{2}_{c}}{\pi^{2}}\sum_{q}
\iint^{+\infty}_{-\infty}\d\varepsilon\d\omega\,
\frac{\partial_{\varepsilon}F_{\varepsilon}}
{\big[(Dq^{2}+\tau^{-1}_{\mathrm{GL}})^{2}+\omega^{2}\big]
\big[(Dq^{2})^{2}+(2\varepsilon+\omega)^{2}\big]}\,.
\end{equation}
に帰着する．最後に，準2次元の場合について残りの積分を行うと，
\begin{equation}\label{Part-IV-sigma-MT}
\frac{\delta\sigma_{\mathrm{MT}}}{\sigma_{D}}=\frac{\mathrm{Gi}}{8}\left(\frac{T_{c}}{T-T_{c}}\right)
\left(\frac{1}{1-\tau_{\mathrm{GL}}/\tau_{\phi}}\right)\ln\left(\frac{\tau_{\phi}}
{\tau_{\mathrm{GL}}}\right)\,,
\end{equation}
を得る．ここで運動量積分の赤外発散は，位相緩和率$Dq^{2}_{\mathrm{min}}=\tau^{-1}_{\phi}$によって打ち切った．この発散はMaki--Thompsonダイアグラムのよく知られた特徴であり，Cooperチャネルにおける何らかの位相緩和機構によって正則化できる．たとえば系に磁性不純物が存在する場合，スピン反転時間が$\tau_{\phi}$の役割を果たす．$\delta\sigma_{\mathrm{DOS}}$と異なり，Maki--Thompson補正\eqref{Part-IV-sigma-MT}は正であり，温度依存性ははるかに強いべき乗型である．興味深いことに，$\delta\sigma_{\mathrm{MT}}$も有効Ginzburg--Landau理論から導かれる．実際，$\delta\sigma_{\mathrm{MT}}=e^{2}\nu\langle\delta D^{\mathrm{MT}}_{\mathbf{r},t,t'}\rangle_{\Delta}$と定義し，式\eqref{SuperCond-D-MT}を用いて$\Delta$について平均をとると，式\eqref{Part-IV-sigma-MT}が再現される．


伝導率には，Aslamazov--Larkinの寄与と呼ばれるもう1つの補正がある．これは式\eqref{Part-IV-Z}の右辺第2のトレースから得られる．実際，各$\check{Q}$行列を$\check{\mathcal{W}}$について1次まで展開すると，
\begin{equation}\label{Part-IV-Z-AL}
\mathcal{Z}_{\mathrm{AL}}[\mathbf{A}^{cl},\mathbf{A}^{q}]=-\frac{(\pi\nu
D)^{2}}{2}\,
\left\langle\left\langle\left(\Tr\big\{\check{\mathbb{A}}_{\varepsilon_{1}\varepsilon_{2}}
(\hat{\sigma}_{z}\otimes\hat{\tau}_{z})
\check{\mathcal{W}}_{\varepsilon_{2}\varepsilon_{3}}
\partial_{\mathbf{r}}\check{\mathcal{W}}_{\varepsilon_{3}\varepsilon_{1}}\big\}
\right)^{2}\right\rangle\right\rangle_{\mathcal{W},\Delta}\,.
\end{equation}
を得る．式\eqref{Part-IV-Z-AL}をベクトルポテンシャルで微分して得られる，2つの頂点を導入すると便利である：
\begin{subequations}\label{Part-IV-V-AL}
\begin{equation}
\hskip-5cm \,\,\,\,\,
\mathbb{V}^{cl}_{\mathrm{AL}}[\check{\mathcal{W}}] =
\frac{\delta}{\delta\mathbf{A}^{cl}(\Omega)}\,
\Tr\big\{\check{\mathbb{A}}_{\varepsilon_{1}\varepsilon_{2}}
(\hat{\sigma}_{z}\otimes\hat{\tau}_{z})
\check{\mathcal{W}}_{\varepsilon_{2}\varepsilon_{3}}
\partial_{\mathbf{r}}\check{\mathcal{W}}_{\varepsilon_{3}\varepsilon_{1}}\big\}
\nonumber
\end{equation}
\begin{equation}
 \!\!\!\! = \tr\Big\{ c_{\varepsilon_{2}\varepsilon_{3}}(\mathbf{r})
\partial_{\mathbf{r}}c^{*}_{\varepsilon_{3}\varepsilon_{2}+\Omega}(\mathbf{r})
+c^{*}_{\varepsilon_{2}\varepsilon_{3}}(\mathbf{r})
\partial_{\mathbf{r}}c_{\varepsilon_{3}\varepsilon_{2}+\Omega}(\mathbf{r})
- (c\to\bar{c}) \Big\}\,,
\end{equation}
\begin{equation}
\hskip-4.8cm \mathbb{V}^{q}_{\mathrm{AL}}[\check{\mathcal{W}}] =
\frac{\delta}{\delta\mathbf{A}^{q}(0)}\,
\Tr\big\{\check{\mathbb{A}}_{\varepsilon_{1}\varepsilon_{2}}
(\hat{\sigma}_{z}\otimes\hat{\tau}_{z})
\check{\mathcal{W}}_{\varepsilon_{2}\varepsilon_{3}}
\partial_{\mathbf{r}}\check{\mathcal{W}}_{\varepsilon_{3}\varepsilon_{1}}\big\}
\nonumber
\end{equation}
\begin{equation}
 \!\! = -\tr\Big\{F_{\varepsilon_{2}}
\big(c_{\varepsilon_{2}\varepsilon_{3}}(\mathbf{r})
\partial_{\mathbf{r}}c^{*}_{\varepsilon_{3}\varepsilon_{2}}(\mathbf{r})+
c^{*}_{\varepsilon_{2}\varepsilon_{3}}(\mathbf{r})\partial_{\mathbf{r}}
c_{\varepsilon_{3}\varepsilon_{2}}(\mathbf{r})
+(c\to\bar{c})\big)\Big\}\,.
\end{equation}
\end{subequations}
$\mathbb{V}^{q}_{\mathrm{AL}}$については外部周波数を最初から$\Omega=0$としてよいのに対し，$\mathbb{V}^{cl}_{\mathrm{AL}}$については有限の$\Omega$を保ち，$\check{\mathcal{W}}$についての平均をとった後で初めて直流極限$\Omega\to 0$をとることが重要である．Cooperonの平均には式\eqref{app_FE-c-extremal}を用いる．$\mathbb{V}^{q}_{\mathrm{AL}}[\check{\mathcal{W}}]$の場合，2つのCooper場の積では，秩序変数の古典成分だけを含む寄与を残せば十分である：$\mathbb{V}^{q}_{\mathrm{AL}}[\check{\mathcal{W}}]\propto\Tr\big\{F[c\partial_{\mathbf{r}}c^{*}+c^{*}\partial_{\mathbf{r}}c]\big\}\propto\Delta^{cl}\partial_{\mathbf{r}}\Delta^{*cl}-\Delta^{*cl}\partial_{\mathbf{r}}\Delta^{cl}$．一方，頂点$\mathbb{V}^{cl}_{\mathrm{AL}}[\check{\mathcal{W}}]$では，秩序変数の量子成分$\Delta^{q}$を少なくとも1つ保持することが不可欠である．2つの古典成分を含む対応する寄与は，因果構造によって消えるからである．その結果，主要な寄与は$\mathbb{V}^{cl}_{\mathrm{AL}}[\check{\mathcal{W}}]\propto\Tr\big\{c\partial_{\mathbf{r}}c^{*}+c^{*}\partial_{\mathbf{r}}c\big\}\propto\Delta^{cl}\partial_{\mathbf{r}}\Delta^{*q}+\Delta^{q}\partial_{\mathbf{r}}\Delta^{*cl}-c.c.$となる．積$\big\langle\mathbb{V}^{cl}_{\mathrm{AL}}[\check{\mathcal{W}}]\mathbb{V}^{q}_{\mathrm{AL}}[\check{\mathcal{W}}]\big\rangle_{\Delta}$に対して残っている$\Delta$の平均は，式\eqref{Part-IV-Delta-averages}を用いて行う．運動量表示に移り，すべての因子をまとめると，直流極限における伝導率へのAslamazov--Larkin型補正は次の形をとる：
\begin{equation}
\delta\sigma_{\mathrm{AL}}=\frac{\pi^{2}
e^{2}D}{8T^{2}_{c}}\sum_{q}Dq^{2}
\int^{+\infty}_{-\infty}\frac{\d\omega}{2\pi}\,
\frac{\partial}{\partial\omega}\left[\coth\frac{\omega}{2T}\right]
\big[\mathrm{Im}L^{R}\qo\big]^{2}\,.
\end{equation}
対応するダイアグラムを図\ref{Fig-AL-MT-DOS}cに示す．寄与するのは$Dq^2\sim\omega\sim\tau^{-1}_{GL}\ll T_{c}$のみなので，$\partial_{\omega}[\coth\omega/2T]\approx-2T_{c}/\omega^2$と近似し，$\mathrm{Im}L^{R}\qo=-(8iT_{c}\omega/\pi)[(Dq^{2}+\tau^{-1}_{\mathrm{GL}})^2+\omega^2]^{-1}$を用いると，
\begin{equation}\label{Part-IV-sigma-AL}
\delta\sigma_{\mathrm{AL}}=\frac{8e^{2}DT_c}{\pi}\sum_{q}\int^{+\infty}_{-\infty}\!\!\d\omega\,
\frac{Dq^{2}}{\big[\big(Dq^{2}+\tau^{-1}_{\mathrm{GL}}\big)^{2}+\omega^{2}\big]^{2}}\,.
\end{equation}
を得る．残りの積分を実行すると，準2次元薄膜について次を得る．
\begin{equation}
\frac{\delta\sigma_{\mathrm{AL}}}{\sigma_{D}}=\frac{\mathrm{Gi}}{16}
\left(\frac{T_{c}}{T-T_{c}}\right)\,.
\end{equation}

有効な時間依存GL汎関数のレベルでは，Aslamazov--Larkin伝導率補正$\delta\sigma_{\mathrm{AL}}$は，作用の$S_{\mathrm{SC}}$部分\eqref{SuperCond-S-SC}から生じる．これを確かめる最も簡単な方法は，$S_{\mathrm{SC}}$から従う電流
$$\mathbf{j}_{\mathrm{SC}}=\frac{\pi e\nu
D}{4T_{c}}\mathrm{Im}[\Delta^{*cl}\partial_{\mathbf{r}}\Delta^{cl}]\,,$$
を，揺動散逸関係$\delta\sigma_{\mathrm{AL}}\propto\big\langle\mathbf{j}_{\mathrm{SC}}\cdot\mathbf{j}_{\mathrm{SC}}\big\rangle_{\Delta}\propto\sum_{q\omega}Dq^2|L^{R}\qo|^2$とともに用いることである．これにより式\eqref{Part-IV-sigma-AL}が再現される．

本節で用いた手法によって，従来の松原ダイアグラム法で知られている，揺らぎが誘起する伝導率の結果をすべて再現できる．ここでの簡単化は，解析接続を必要としないことである．今回の問題では解析接続はそれほど複雑ではないが，場合によってはそれを避けられることが利点となる．

%$$$$$$$$$$$$$$$$$$$$$$$$$$$$$$$$$$$$$$$$$$$$$$$$$$$$$$$$$$$$$$$$$$$$$$$$$$$$$$$$$$$$$$$$$$$$$$$$$$$$$$$$$$$
\subsubsection{$T_c$以上でのトンネルコンダクタンス}\label{app_Part-IV-5}

電圧を印加した超伝導体・常伝導金属トンネル接合を考える．超伝導体の温度は転移温度$T_c$の直上，すなわち揺らぎ領域にあるとする．揺らぎによる状態密度の減少（第\ref{app_Part-IV-3}節参照）が接合の電流・電圧特性を変えることは，自然に予想される\cite{VarlamovDorin,DiCastro,Reizer}．この効果は，作用のトンネル部分$S_{T}[\check{Q}_{L},\check{Q}_{R}]$を用い，$\sigma$模型の枠組みで調べられる．

式\eqref{NLSM-action-S-T}から出発し，$a=L,R$に対してゲージ変換$\check{Q}_{a}\to\exp(-i\check{\Xi}\check{\Phi}_{a})\check{Q}_{a}\exp(i\check{\Xi}\check{\Phi}_{a})$を行う．ここで$\check{\Phi}_{a}(t)=\int^{t}\check{V}_{a}(t)\d t=[\Phi^{cl}_{a}(t)\hat{\sigma}_{0}+\Phi^{q}_{a}(t)\hat{\sigma}_{x}]\otimes\hat{\tau}_{0}$，$\Phi^{cl}_{L}-\Phi^{cl}_{R}=eVt$である．この変換によって，印加電圧$V$は$\check{Q}$行列のKeldyshブロックから作用のトンネル部分へ移される：
\begin{equation}\label{Part-IV-S-T}
iS_{T}[\check{Q}_{L},\check{Q}_{R}]=\frac{\mathrm{g}_{T}}{4\mathrm{g}_{Q}}
\Tr\big\{\check{Q}_{L}e^{-i\check{\Xi}\check{\Phi}}
\check{Q}_{R}e^{i\check{\Xi}\check{\Phi}}\big\}\,,
\end{equation}
ここで$\check{\Phi}=\check{\Phi}_{L}-\check{\Phi}_{R}$であり，$\Phi^{q}(t)$は生成場として働く．実際，位相$\check{\Phi}$は粒子数$\check{N}=i\partial/\partial\Phi$と量子力学的に正準共役なので，トンネル電流は分配関数$\mathcal{Z}_{T}[\Phi]=\exp\big(iS_{T}[\check{Q}_{L},\check{Q}_{R}]\big)$を位相の量子成分で微分して得られる：
\begin{equation}\label{Part-IV-I-T}
I_{T}(t)=ie\left(\frac{\delta\mathcal{Z}_{T}[\Phi]}{\delta\Phi^{q}(t)}
\right)_{\Phi^{q}=0}.
\end{equation}
この定義を式\eqref{Part-IV-S-T}に適用し，
\begin{equation}\label{Part-IV-derivative}
\left.\frac{\delta \exp(\pm
i\check{\Xi}\check{\Phi})}{\delta\Phi^{q}(t')}\right|_{\Phi^{q}=0}=\pm
i\delta(t-t')\,\big(\hat{\sigma}_{x}\otimes\hat{\tau}_{z}\big)\,\exp\big[\pm
ieVt\check{\Xi}\big]\,
\end{equation}
を用い，$\check{Q}_{L}=\check{Q}_{R}=\check{\Lambda}$とすると，常伝導状態のトンネル接合に対して当然のOhm則$I_{T}=\mathrm{g}_{T}V$を得る．ここで式\eqref{Part-IV-S-T}の$\check{Q}$行列の一方をCooperモード$\check{\mathcal{W}}$について展開することにより，揺らぎの効果を取り入れられる．これにより次の形の補正が生じる：
\begin{equation}\label{Part-IV-I-T-corr}
\delta I_{T}(V)=-\frac{\pi
\mathrm{g}_{T}}{2e}\sum_{q}\iint\frac{\d\varepsilon\d\varepsilon'}{4\pi^{2}}
\big(F_{\varepsilon+eV}-F_{\varepsilon-eV}\big)
\left\langle\left\langle c_{\varepsilon\varepsilon'}(\mathbf{q})
c^{*}_{\varepsilon'\varepsilon}(-\mathbf{q})+
\bar{c}_{\varepsilon\varepsilon'}(\mathbf{q})
\bar{c}^{*}_{\varepsilon'\varepsilon}(-\mathbf{q})
\right\rangle\right\rangle_{\mathcal{W},\Delta}\,,
\end{equation}
これは物理的に予想される結果である．実際，式\eqref{Part-IV-I-T-corr}のCooperモードの組み合わせは，状態密度補正$\delta\nu(\varepsilon)$と同定できる（式\eqref{Part-IV-Dos-corr-def}参照）．式\eqref{Part-IV-I-T-corr}では，これとFermi関数の差との畳み込みがとられ，トンネル電流の補正は$\delta I_T(V)\sim\int\d\varepsilon[F_{\varepsilon+eV}-F_{\varepsilon-eV}]\delta\nu_{L}(\varepsilon)\nu_{R}$という形になる．式\eqref{Part-IV-Dos-intermediate}で得た$\delta\nu(\varepsilon)$を式\eqref{Part-IV-I-T-corr}に代入し，無次元変数$x=Dq^{2}/T$，$y=\omega/T$，$z=\varepsilon/2T$に変換すると，トンネル微分コンダクタンスの補正$\delta \mathrm{g}_{T}(V)=\partial\delta I_{T}(V)/\partial V$として次の式を得る：
\begin{eqnarray}
\frac{\delta\mathrm{g}_{T}(V)}{\mathrm{g}_{T}}=\frac{4\mathrm{Gi}}{\pi^{3}}
\int^{\infty}_{0}\d x\iint^{+\infty}_{-\infty}\d y\d
z\left[\frac{1}{\cosh^{2}(z+u)}+\frac{1}{\cosh^{2}(z-u)}\right]\\
\times\,\mathrm{Re}\left[\frac{1}{\big(x+iy-4iz\big)^{2}
\big((x+1/T\tau_{\mathrm{GL}})^{2}+y^{2}\big)}\right]\,,
\end{eqnarray}
ここで$u=eV/2T$であり，準2次元の配置を仮定した．残りの積分は閉じた形で実行でき，結果は\cite{VarlamovDorin}
\begin{equation}
\frac{\delta\mathrm{g}_{T}(V)}{\mathrm{g}_{T}}=\frac{\mathrm{Gi}}{\pi^{4}}
\ln\left(\frac{T_{c}}{T-T_{c}}\right)
\mathrm{Re}\,\psi^{[2]}\left(\frac{1}{2}-\frac{ieV}{2\pi T}\right),
\end{equation}
となる．ここで$\psi^{[2]}(x)$はディガンマ関数$\psi(x)$の2階微分である．状態密度の抑制$\delta\nu(\varepsilon)$と直接関係しているにもかかわらず，微分コンダクタンスの補正$\delta\mathrm{g}_{T}$の温度依存性はずっと弱いことに注意されたい．状態密度の急激な抑制$\delta\nu(0)\propto(T-T_{c})^{-2}$は，温度に対して緩やかな対数補正$\delta\mathrm{g}_{T}\propto\ln(T_{c}\tau_{\mathrm{GL}})$しかもたらさない．もう1つ興味深いのは，$\delta\nu(\varepsilon)$の抑制はエネルギー$\varepsilon\sim\tau^{-1}_{\mathrm{GL}}\sim T-T_{c}$で生じるのに対し，対応する微分コンダクタンスの抑制は$V\sim T-T_c$ではなく，電圧$V\sim T_{c}$で起こることである．最後に，$(T-T_c)$に対してより特異なMT補正とAL補正は，トンネル行列要素の4次で初めて現れるのに対し，ここで論じたDOS効果は$\mathrm{g}_{T}$に線形，すなわちトンネル行列要素の2次であることを述べておく．


%$$$$$$$$$$$$$$$$$$$$$$$$$$$$$$$$$$$$$$$$$$$$$$$$$$$$$$$$$$$$$$$$$$$$$$$$$$$$$$$$$$$$$$$$$$$$$$$$$$$$$$$$$$$
\subsubsection{揺らぎ領域における電流雑音}\label{app_Part-IV-6}

状態密度に関係する効果に加えて，超伝導揺らぎはトンネル接合の電流雑音にも興味深い影響を与える\cite{Larkin-Varlamov,Kulik,Scalapino,Martin-Balatsky,Levchenko-J-Noise}．今度は，接合の両側が同一の超伝導体から成り，いずれも$T_c$の直上に保たれているとする．この場合，平均Josephson電流は存在しないが，雑音パワーはJosephson周波数$\omega_J=2eV/\hbar$に敏感で，$\omega=\omega_{J}$に鋭いピークを示す．このピークの高さと形状は$T_c$近傍で特異な温度依存性をもち，実験的な検出を可能にする．これを示すため，接合の両側にある揺らぐ秩序変数$\Delta_{L(R)}\rt$の積で，トンネル電流の揺らぎ部分$\delta I_{T}(t)$を表す式，すなわち$\delta I_{T}(t)\propto\int\d\mathbf{r}\,[\Delta_{R}\rt\Delta^{*}_{L}\rt\exp(-i\omega_{J}t)-c.c.]$を導く．$T_{c}$以上では$\langle\Delta_{L(R)}\rangle=0$なので，$\langle\delta I_{T}(t)\rangle=0$であることは明らかである．しかし，電流の2乗平均$\langle\delta I_{T}(t)\delta I_{T}(t')\rangle$は零ではなく，そのFourier変換はJosephson周波数にピークを示す．以下では，その温度依存性を計算する．

電流・電流相関関数の定義
\begin{equation}\label{Part-IV-noise-def}
\mathcal{S}_{T}(\omega)=-e^{2}\int^{+\infty}_{-\infty}\d(t-t')
\left(\frac{\delta^{2}\mathcal{Z}_{T}[\Phi]}{\delta
\Phi^{q}(t)\delta
\Phi^{q}(t')}\right)_{\Phi^{q}=0}e^{-i\omega(t-t')}\,.
\end{equation}
から出発する．常伝導状態では$\check{Q}_{L}=\check{Q}_{R}=\check{\Lambda}$であり，式\eqref{Part-IV-noise-def}から従うトンネル接合の雑音パワーは，Schottkyの式$\mathcal{S}_{T}(\omega)=2\mathrm{g}_{T}T\sum_{\pm}v_{\pm}\coth v_{\pm}$で与えられる．ここで$v_{\pm}=(eV\pm\omega)/2T$である．接合両側の超伝導揺らぎを取り入れるには，式\eqref{Part-IV-S-T}の各$\check{Q}$行列をCooperモードについて最低次（1次）まで展開する必要がある．これにより，電流への揺らぎ補正として
\begin{equation}\label{Part-IV-I-J}
\delta I_{T}(t)=\frac{i\pi\mathrm{g}_{T}}{4e}\frac{\delta}{\delta
\Phi^{q}(t)}\Tr\left\{e^{i\check{\Xi}\check{\Phi}}
\check{\mathcal{U}} (\hat{\sigma}_{z}\otimes\hat{\tau}_{z})
\check{\mathcal{W}}_{L}\check{\mathcal{U}}^{-1}
e^{-i\check{\Xi}\check{\Phi}}
\check{\mathcal{U}} (\hat{\sigma}_{z}\otimes\hat{\tau}_{z})
\check{\mathcal{W}}_{R}\check{\mathcal{U}}^{-1}\right\}.
\end{equation}
を得る．さらに進むには，電子と秩序変数の自由度の時間スケールの分離を用いて，式\eqref{Part-IV-I-J}を簡単にする．実際，式\eqref{app_FE-L-R}から，秩序変数の変動に関係するエネルギーと運動量は$Dq^{2}\sim\omega\sim\tau^{-1}_{\mathrm{GL}}$であるのに対し，Cooperonに入るフェルミオンのエネルギーは$\epsilon\sim\epsilon'\sim T\gg \tau^{-1}_{\mathrm{GL}}$である．その結果，Cooperモード\eqref{SuperCond-W}と秩序変数の間の非局所関係，式\eqref{app_FE-c-extremal}は，
\begin{eqnarray}\label{Part-IV-c-Delta-approx}
&& \check{\mathcal{W}}^{a}_{tt'}(\mathbf{r})\approx-i\,
\hat{\Theta}_{tt'}\otimes\hat{\Delta}^{a}_{tt'}(\mathbf{r}),\quad\quad
\hat{\Theta}_{tt'}= \left(\begin{array}{cc} \theta(t-t') & 0 \\
0 & -\theta(t'-t),
\end{array}\right)_K,\nonumber \\
&&
\hat{\Delta}^{a}_{tt'}(\mathbf{r})=\Delta^{cl}_{a}\left(\mathbf{r},\frac{t+t'}{2}\right)
\hat{\tau}_{+}+
\Delta^{*cl}_{a}\left(\mathbf{r},\frac{t+t'}{2}\right)\hat{\tau}_{-}\,,\quad
a=L,R\,,
\end{eqnarray}
と近似できる．ここで$\theta(t)$は階段関数である．物理的には，式\eqref{Part-IV-c-Delta-approx}は，Cooperonが特徴的長さ$\xi_{0}=\sqrt{D/T_{c}}$をもつ短距離のものである一方，秩序変数の長距離揺らぎは$\xi_{\mathrm{GL}}=\sqrt{D\tau_{\mathrm{GL}}}\gg\xi_{0}$程度の距離まで伝播することを意味する．したがって，式\eqref{app_FE-c-extremal}の関係は実効的に局所的となり，その後の解析は大幅に簡単になる．式\eqref{Part-IV-c-Delta-approx}を用いると，式\eqref{Part-IV-I-J}でKeldysh部分空間のトレースを明示的に計算でき，
\begin{equation}\label{Part-IV-I-J-1}
\delta I_{T}(t)=-\frac{\pi
\mathrm{g}_{T}}{e}\Tr\left\{\theta(t_{2}-t_{1})F_{t_{1}-t}\theta(t-t_{2})
\hat{\Delta}^{L}_{tt_{2}}\hat{\tau}_{z}
\hat{\Delta}^{R}_{t_{2}t_{1}}e^{ieV(t+t_{2})\hat{\tau}_{z}}\right\}_N\,,
\end{equation}
に到達する．ここでは式\eqref{Part-IV-derivative}を用い，トレースを実空間・時間表示で書いた（ここで$\mathrm{Tr}\{\ldots\}$は時刻$t$についての積分を含まないことに注意）．積分変数を$t_{1}=t-\mu$および$t_{3}=t-\eta$に変え，$\eta,\mu$を温度単位で$T\eta\to\eta, T\mu\to\mu$とスケール変換すると，式\eqref{Part-IV-I-J-1}の等価な表示
\begin{equation}\label{Part-IV-I-J-2}
\delta I_{T}(t)=-\frac{i\pi
\mathrm{g}_{T}}{eT}\iint^{+\infty}_{-\infty}d\eta d\mu\,\,
\frac{\theta(\eta)\theta(\mu-\eta)}{\sinh(\pi\mu)}\,
\Tr \left\{\hat{\Delta}^{L}_{t,t-\frac{\eta}{T}}\hat{\tau}_{z}
\hat{\Delta}^{R}_{t-\frac{\eta}{T},t-\frac{\mu}{T}}
e^{ieV\big(2t-\frac{\eta}{T}\big)\hat{\tau}_{z}}\right\}_N\,,
\end{equation}
を得る．ここでは時間領域の平衡フェルミオン分布関数$F_{t}=-iT/\sinh(\pi Tt)$を用いた．上の積分に最も大きく寄与するのは$\eta\sim\mu\lesssim1$である．この範囲で$\{\eta,\mu\}/T$は温度の逆数のスケールで変化する一方，すでに論じたように，秩序変数の変化は$t\sim\tau_{\mathrm{GL}}\gg1/T$によって決まる．したがって，$\eta$と$\mu$の積分を行う際には，秩序変数と指数因子の$\{\eta,\mu\}/T$依存性を無視できる．その結果，
\begin{equation}\label{Part-IV-I-J-3}
\delta I_T(t)=\frac{i\pi
\mathrm{g}_{T}}{4eT}\int\frac{\mathrm{d}^{2}\mathbf{r}}{\mathcal{A}}\left[
\Delta^{cl}_{R}(\mathbf{r},t)\Delta^{*cl}_{L}(\mathbf{r},t)
e^{-i\omega_{J}t}-c.c.\right]\,,
\end{equation}
を得る．ここで空間積分は接合の面積$\mathcal{A}$にわたって行い，$\omega_J=2eV/\hbar$である．これで，対応する電流雑音への寄与を計算する準備が整った．式\eqref{Part-IV-I-J-3}の形の電流を2つ，式\eqref{Part-IV-noise-def}に代入し，式\eqref{Part-IV-Delta-averages}から従う相関関数
$$\langle\Delta^{cl}_{a}(\mathbf{r},t)\Delta^{*cl}_{b}(\mathbf{r}',t')\rangle_{\Delta}=
\frac{i}{2\nu}\delta_{ab}L^{K}(\mathbf{r}-\mathbf{r}',t-t')\,.$$
を用いて，揺らぐ秩序変数を対にする．その結果，雑音パワーへの超伝導揺らぎ補正は
\begin{equation}\label{Part-IV-noise-correction}
\delta\mathcal{S}_{T}(\omega)=-\frac{1}{4\nu^{2}}\left(\frac{\pi
\mathrm{g}_{T}}{4eT_{c}}\right)^{2}\sum_{\pm}
\int\frac{\mathrm{d}^{2}\mathbf{r}}{\mathcal{A}}
\int^{+\infty}_{-\infty}\mathrm{d}t\,
\big[L^{K}(\mathbf{r},t)\big]^{2}\, \exp(-i\omega_{\pm}t)\,,
\end{equation}
で与えられる．ここで$\omega_{\pm}=\omega\pm\omega_{J}$である．残りの積分を行うため，まず運動量・時間の混合表示における揺らぎ伝播関数のKeldysh成分$L^{K}(\mathbf{q},t)=\int L^{K}(\mathbf{q},\omega)e^{-i\omega t}\d\omega/2\pi$を求めると，
\begin{equation}
L^{K}(\mathbf{q},t)=-\frac{2iT^{2}_{c}}{T-T_{c}}
\frac{e^{-\varkappa_{q}|t|/\tau_{\mathrm{GL}}}}{\varkappa_{q}},\quad\quad
\varkappa_{q}=(\xi_{\mathrm{GL}}q)^{2}+1\,.
\end{equation}
となる．次に$L^{K}(\mathbf{r},t)=\int L^{K}(\mathbf{q},t)e^{i\mathbf{qr}}\d q^{2}/4\pi$を式\eqref{Part-IV-noise-correction}に代入し，無次元時間$\tau=t/\tau_{\mathrm{GL}}$を導入し，$q$積分を$\varkappa_q$積分に$\d q^{2}=\d\varkappa_q/\xi^{2}_{\mathrm{GL}}$として変換すると，全体として\cite{Levchenko-J-Noise}
\begin{equation}
\delta\mathcal{S}_{T}(\omega)=\sum_{\pm}\frac{\pi
\mathrm{Gi}^{2}}{64
T_c}\left(\frac{\mathrm{g}_{T}T_c}{e}\right)^{2}\frac{\xi^{2}_{0}}{\mathcal{A}}
\left(\frac{T_{c}}{T-T_{c}}\right)^{2}\mathcal{N}(\omega_{\pm}\tau_{GL})\,,
\end{equation}
を得る．ここでスペクトル関数は
\begin{equation}
\mathcal{N}(z)=\int^{+\infty}_{-\infty}d\tau\int^{+\infty}_{1}
\frac{\d\varkappa}{\varkappa^{2}}\, \exp(-2\varkappa|\tau|-iz
\tau)=\frac{4}{z^{2}}\ln\sqrt{1+z^{2}/4}\,.
\end{equation}
で与えられる．雑音パワーの補正$\delta\mathcal{S}_{T}(\omega)$はJosephson周波数$\omega=\pm\omega_{J}$でピークをもち，強い温度依存性を示す．このため，超伝導転移の近傍で実験的に検出することが可能になる．


%$$$$$$$$$$$$$$$$$$$$$$$$$$$$$$$$$$$$$$$$$$$$$$$$$$$$$$$$$$$$$$$$$$$$$$$$$$$$$$$$$$$$$$$$$$$$$$$$$$$$$$$$$$$
%$$$$$$$$$$$$$$$$$$$$$$$$$$$$$$$$$$$$$$$$$$$$$$$$$$$$$$$$$$$$$$$$$$$$$$$$$$$$$$$$$$$$$$$$$$$$$$$$$$$$$$$$$$$

\sectionlink{結びと謝辞}\label{sec_Summary}

本総説では，相互作用するボソン系およびフェルミ粒子系への応用を通じて，Keldysh法を構成するさまざまな要素を概観することを試みた．特に，微視的模型の汎関数積分表示と，非線形$\sigma$模型などの近年の発展に重点を置いた．我々の動機は，Keldysh法が有効に適用できる物理学の特定の分野を概観することではなく，手法そのものを提示することに焦点を当てることであった．目標は，この手法とその応用について広い見通しを与えることにあった．そのため，常伝導金属のメゾスコピック物理学と超伝導の動力学から，（かなり主観的に）例を選んだ．本総説が，関心を持つ大学院生にとっての学習資料となり，専門家にとっての参照資料となることを願っている．

本論文の形成に関わった多くの同僚との共同研究と議論から，我々は多大な恩恵を受けた．I.~Aleiner，A.~Altland，A.~Andreev，M.~Feigel'man，Y.~Gefen，L.~Glazman，A.~I.~Larkin，I.~Lerner，M.~Skvortsovをはじめ，多くの方々に感謝する．本総説の執筆中には，D.~Bagrets，I.~Burmistrov，G.~Catelani，A.~Chudnovskiy，M.~Khodas，A.~Varlamovとのやり取りが特に有益であった．本研究はNSF助成金DMR-0405212およびDMR-0804266の支援を受けた．A.~L.はミネソタ大学のDoctoral Dissertation Fellowshipの支援を受けた．

%$$$$$$$$$$$$$$$$$$$$$$$$$$$$$$$$$$$$$$$$$$$$$$$$$$$$$$$$$$$$$$$$$$$$$$$$$$$$$$$$$$$$$$$$$$$$$$$$$$$$$$$$$$$
%$$$$$$$$$$$$$$$$$$$$$$$$$$$$$$$$$$$$$$$$$$$$$$$$$$$$$$$$$$$$$$$$$$$$$$$$$$$$$$$$$$$$$$$$$$$$$$$$$$$$$$$$$$$
\newpage
\appendix
\sectionlink{ボソンとフェルミ粒子のGauss積分}\label{app_Gaussian}

$i,j=1,\ldots N$とし，すべての固有値$\lambda_i$の実部が正，すなわち$\mathrm{Re} \lambda_i>0$である任意の複素$N\times N$行列$A_{ij}$に対して，次の関係が成り立つ．
\begin{equation}\label{app-Gauss-int-bos}
\fitdisplay{
Z[J]=\iint_{-\infty}^{+\infty}\prod\limits_{j=1}^N
\frac{\d(\mathrm{Re} z_j)  \d(\mathrm{Im} z_j)}{\pi}\,
\exp\left(-\sum_{ij}^{N} \bar z_i A_{ij} z_j+ \sum_{j}^{N}\left[\bar
z_j J_j + \bar J_j z_j\right]\right) =\frac{\exp\left(\sum_{ij}^{N}
\bar J_i (A^{-1})_{ij} J_j\right)}{\mathrm{Det}(A)}\, ,
}
\end{equation}
ここで$J_j$は任意の複素ベクトルである．これを証明するには，まずユニタリー変換$A=U^\dagger\Lambda U$によって対角化されるHermite行列から出発すればよい．ここで$\Lambda=\mbox{diag}\{\lambda_j\}$である．このとき，$w_i=U_{ij}z_j$への変数変換（Jacobianは1）によって，恒等式は容易に証明できる．最後に，式~\eqref{app-Gauss-int-bos}の右辺が$\mathrm{Re} A_{ij}$と$\mathrm{Im} A_{ij}$の双方について解析関数であることに注意する．したがって，これらを複素平面へ解析接続することで，任意の複素行列$A_{ij}$へ到達できる．このように恒等式~\eqref{app-Gauss-int-bos}は，積分が適切に定義される限り，すなわち$A_{ij}$のすべての固有値の実部が正である限り成り立つ．

Wickの定理は，因子$\exp\big(-\sum_{ij}\bar z_i A_{ij} z_j\big)$を重みとする積$
z_{a_1}\ldots z_{a_k}\bar z_{b_1}\ldots \bar z_{b_k}$の平均値を扱う．この定理によれば，その平均は，可能なすべての2点平均の積の和で与えられる．例えば，
\begin{eqnarray}\label{app-Gauss-Wick}
&&\left\langle z_{a} \bar z_{b}\right\rangle \equiv \left.
\frac{1}{Z[0]}\frac{\delta^{2}Z[J]}{\delta \bar J_a\delta
J_b}\right|_{J=0}= \big(A^{-1}\big)_{ab}\, ,\\
&&\left\langle z_{a}z_{b} \bar z_{c}\bar z_{d} \right\rangle
\equiv\left.\frac{1}{Z[0]}\frac{\delta^4Z[J]}{\delta \bar
J_{a}\delta\bar{J}_{b} \delta J_{c} \delta J_{d} }\right|_{J=0}=
A^{-1}_{ac}A^{-1}_{bd}+A^{-1}_{ad}A^{-1}_{bc} ,\nonumber
\end{eqnarray}
などとなる．

実変数についての積分に対するGauss積分の恒等式は，次の形を取る．
\begin{equation}\label{app-Gauss-Real}
Z[J]=\int_{-\infty}^{+\infty}\prod\limits_{j=1}^N \frac{d
x_j}{\sqrt{\pi}}\,\exp\left(-\sum\limits_{ij}^N  x_i A_{ij} x_j +
2\sum\limits_{j}^N   x_j J_j\right)=\frac{ \exp\left(\sum_{ij}^{N}
J_i (A^{-1})_{ij} J_j\right)}{\sqrt{\mathrm{Det} (A)}}\, ,
\end{equation}
ここで$A$は，すべての固有値の実部が正である{\em 対称}な複素行列である．証明は複素変数の場合と同様である．まず，直交変換によって対角化できる実対称行列から出発する．すると，変数変換によって恒等式~\eqref{app-Gauss-Real}を容易に証明できる．最後に，（積分が適切に定義される限り）右辺を実対称行列$A_{ij}$から{\em 複素対称}行列へ解析接続すればよい．

$j=1,2,\ldots,N$を持つ2組の{\em 独立な}Grassmann変数$\bar \xi_j$と$\xi_j$についての積分では，{\em 任意の可逆な}複素行列$A$に対してGauss積分の恒等式が成り立つ．
\begin{equation}\label{app-Gauss-int-ferm}
\fitdisplay{
Z[\bar\chi,\chi]=\iint\prod\limits_{j=1}^N \d\bar\xi_j \d\xi_j
\exp\left(-\sum\limits_{ij}^N \bar \xi_i A_{ij} \xi_j +
\sum\limits_{j}^N \left[\bar \xi_j \chi_j + \bar \chi_j
\xi_j\right]\right)= \mathrm{Det}(A) \exp\left(\sum\limits_{ij}^N
\bar \chi_i (A^{-1})_{ij} \chi_j\right)\, .\quad
}
\end{equation}
ここで$\bar \chi_j$と$\chi_j$は，互いに独立であり，かつ$\bar \xi_j$および$\xi_j$からも独立な，さらに2組のGrassmann数である．証明は，例えば指数因子を直接展開することで得られる．その際，$2N$個の変数$\bar \xi_j$および$\xi_j$の{\em すべて}について1次の項だけがゼロでないことに注意すればよい．Wickの定理はボソンの場合と同様に定式化されるが，各組合せには対応する置換の偶奇による符号が掛かる．例えば，式~\eqref{app-Gauss-Wick}の2番目の式の右辺第1項には負号が付く．

%$$$$$$$$$$$$$$$$$$$$$$$$$$$$$$$$$$$$$$$$$$$$$$$$$$$$$$$$$$$$$$$$$$$$$$$$$$$$$$$$$$$$$$$$$$$$$$$$$$$$$$$$$$$
\sectionlink{1粒子の量子力学}
\label{app_SPQM}

（第~\ref{sec_bosons}節で考えた）単一のボース状態からなる最も簡単な多体系は，もちろん1粒子の調和振動子と等価である．この対応を明示するため，複素場$\phi(t)$で書かれた相関関数~\eqref{boson-G-continious}を持つKeldysh経路上の作用~\eqref{boson-Z-1}を考えよう．この場は，実部と虚部を用いて次のようにパラメータ表示できる．
\begin{equation}\label{app-SPQM-qp}
\phi(t) = {1\over\sqrt{\,2\omega_0}}\,\big({p(t)} -i\, {\omega_0}\,
q(t)  \big)\, , \qquad \bar \phi(t) =
{1\over\sqrt{\,2\omega_0}}\,\big( {p(t)}+i\, {\omega_0}\, q(t)
\big)\, .
\end{equation}
実場$p(t)$と$q(t)$を用いると，作用~\eqref{boson-Z-1}は次の形になる．
\begin{equation}\label{app-SPQM-S}
S[p,q]=\int_{{\cal C}}\d t\left[p\,\dot q - {1\over
2}\left(p^2+\omega_0^2 q^2\right) \right]\, ,
\end{equation}
ここでは，$p^2$，$q^2$および$p\,q$の全時間微分を省略した．これらは連続体の記法では陽に書かれていない境界項にのみ寄与するためである（ただし，適切な正則化のためには保持する必要がある）．式~\eqref{app-SPQM-S}は，Hamilton形式で書いた量子調和振動子の作用そのものである．$p(t)$場についてGauss積分を実行すると，
\begin{equation}\label{app-SPQM-S-1}
S[q]={1\over 2}\int_{{\cal C}}\d t\left[ \dot q^2 -
\omega_0^2 \, q^2 \right]\, .
\end{equation}
を得る．これは，Keldysh経路上に書かれた調和振動子のFeynmanのLagrange形式の作用である．これを任意の1粒子ポテンシャル$U(q)$へ一般化すると，
\begin{equation}\label{app-SPQM-S-2}
S[q(t)]= \int_{{\cal C}}\d t\left[ {1\over 2}\,\big(\dot q(t)\big)^2
- U\big(q(t)\big) \right]\, .
\end{equation}
となる．$q(t)$場を，経路の往路と復路にそれぞれ存在する2つの成分$q_+(t)$と$q_-(t)$に分け，Keldysh回転$q_{\pm}=q^{cl}\pm
q^{q}$を行う．これらの場を用いると，作用は次の形になる．
\begin{equation}\label{app-SPQM-S-3}
S[q^{cl},q^q]=\int_{-\infty}^{+\infty}\d t\left[-2\,q^{q}{\d^{\,2}
q^{cl}\over \d t^2} - U\big(q^{cl}+q^{q}\big) + U\big(q^{cl}-
q^{q}\big) \right],
\end{equation}
ここでは$\dot
q^{q}\dot q^{cl}$という項を部分積分した．これがFeynman経路積分のKeldysh形式である．省略した境界項は，$\sim i0(q^q)^2$の形の収束因子を与える．

量子成分$q^{q}(t)$の揺らぎが小さいとみなすなら，ポテンシャルを1次まで展開でき，作用は
\begin{equation}\label{app-SPQM-S-4}
S[q^{cl},q^{q}]=\int_{-\infty}^{+\infty}\d
t\left[-2\,q^q\left({\d^{\,2} q^{cl}\over \d t^2} +{\partial
U\big(q^{cl}\big)\over \partial q^{cl} }\right) + i0 (q^{q})^{2}
+O\big[(q^q)^3\big]\right].
\end{equation}
となる．この極限では，量子成分$q^{q}$についての積分を陽に実行でき，丸括弧内の式に対する汎関数$\delta$関数が得られる．この$\delta$関数は，$q^{cl}$に対して古典的なNewtonの運動方程式
\begin{equation}\label{app-SPQM-Newton}
{\d^{\,2} q^{cl}\over \d t^2} =- {\partial U\big(q^{cl}\big)\over
\partial q^{cl} }\, .
\end{equation}
を課す．このため，Keldysh場の往路と復路について対称な部分を古典成分と呼ぶ．量子力学的な情報は，式~\eqref{app-SPQM-S-4}で省略した$q^{q}$の高次項に含まれている．調和振動子ポテンシャルでは，$O[(q^{q})^3]$と記した項が恒等的に存在しないことに注意する．この場合，量子的（半古典的）な情報は，収束項$i0(q^{q})^2$と，式~\eqref{app-SPQM-S-4}の演算子$\d^{\,2}/(\d t^2)$に対する遅延型の正則化に含まれる．

粒子座標に添字${\bf r}$を付けて$q_{\bf r}(t)$とし，ばねのポテンシャルエネルギー${v_s^2\over 2}(q_{{\bf r}+{\bf
1}}(t)-q_{\bf r}(t))^2$を加えることで，1粒子の量子力学を，調和的に結合した粒子の鎖（または格子）へ一般化できる．空間について連続体の記法$\phi(\mathbf{r},t)\equiv q_{\mathbf{r}}(t)$へ移ると，実場（例えばフォノン場）のKeldysh作用として
\begin{equation}\label{app-SPQM-S-Phonons}
S[\phi]= \int\d\mathbf{r}\int_{{\cal C}}\d t\left[ {1\over
2}\,\dot \phi^{\,2} - {v_s^2\over 2}\,(\partial_\mathbf{r}
\phi )^2 - U\big(\phi\big) \right]\, ,
\end{equation}
を得る．ここで定数$v_s$は音速を意味する．最後に，場を$(\phi_+,\phi_-)$成分に分け，Keldysh変換$\phi_\pm=\phi^{cl}\pm
\phi^q$を行い，部分積分すると，
\begin{equation}\label{app-SPQM-S-Phonons-1}
 S[\phi^{cl},\phi^{q}]= \int\d\mathbf{r}
\int_{-\infty}^{+\infty}\d t\left[ 2\phi^{q}\big( -\partial_t^2 +
v_s^2\,\partial_\mathbf{r}^2 \big) \phi^{cl}
-U(\phi^{cl}+\phi^{q})+U(\phi^{cl}-\phi^{q}) \right] .
\end{equation}
を得る．Keldysh理論の一般的な構造に従えば，微分演算子$\big(-\partial_t^2 +
v_s^2\,\partial_\mathbf{r}^2\big)$は遅延型と理解すべきである．これは，時間領域では下三角行列であることを意味する．実際，部分積分によって作用を対称化し，$\phi^{q}\big(-\partial_t^2 + v_s^2\,\partial_\mathbf{r}^2\big)^R
\phi^{cl}+ \phi^{cl}\big( -\partial_t^2 +
v_s^2\,\partial_\mathbf{r}^2\big)^A \phi^{q}$と書くことができる．ここでは第2項に先進型の正則化を施す．


%$$$$$$$$$$$$$$$$$$$$$$$$$$$$$$$$$$$$$$$$$$$$$$$$$$$$$$$$$$$$$$$$$$$$$$$$$$$$$$$$$$$$$$$$$$$$$$$$$$$$$$$$$$$
\sectionlink{$\sigma$模型の勾配展開}
\label{app_GradientExpansion}

この付録は，第~\ref{sec_NLSM-2}節の補足資料である．その目的は，式~\eqref{NLSM-action-SAV}から式~\eqref{NLSM-action-general}への移行に隠された計算の詳細を示すことである．式~\eqref{NLSM-action-SAV}の対数を勾配展開するため，$\hat Q$行列を式~\eqref{NLSM-Q-Rotated}の形に取り，式~\eqref{NLSM-action-trlog-R}と同様に次の式を得る．
\begin{equation}\label{app-GE-S-trlog}
iS[\hat Q,\mathbf{A},V]=\Tr\ln\left[\hat{1}+i\hat{\mathcal{G}}
\hat{\mathcal{R}}\partial_{t}\hat{\mathcal{R}}^{-1}+i\hat{\mathcal{G}}\hat{\mathcal{R}}
\mathbf{v}_{F}\partial_{\mathbf{r}}\hat{\mathcal{R}}^{-1}+
\hat{\mathcal{G}}\hat{\mathcal{R}}\hat{V}\hat{\mathcal{R}}^{-1}+
\hat{\mathcal{G}}\hat{\mathcal{R}}\mathbf{v}_{F}
\hat{\mathbf{A}}\hat{\mathcal{R}}^{-1}\right]\,.
\end{equation}
この式を$\hat{\mathcal{G}}
\hat{\mathcal{R}}\partial_{t}\hat{\mathcal{R}}^{-1}$の1次および$\hat{\mathcal{G}}\hat{\mathcal{R}}
\mathbf{v}_{F}\partial_{\mathbf{r}}\hat{\mathcal{R}}^{-1}$の2次まで展開すると，$S[\hat Q]$に対する式~\eqref{NLSM-action-trlog-Approx}が再現され，最終的に式~\eqref{NLSM-action-noninteracting}に至る．$\hat{V}$と$\hat{\mathbf{A}}$の1次では，式~\eqref{app-GE-S-trlog}から
\begin{equation}\label{app-GE-S1-trlog}
iS_{1}[\hat Q,\mathbf{A},V]=\Tr\big\{\hat{\mathcal{G}}
\hat{\mathcal{R}}\hat{V}\hat{\mathcal{R}}^{-1}\big\}-
i\Tr\big\{\hat{\mathcal{G}}(\hat{\mathcal{R}}\mathbf{v}_{F}
\partial_{\mathbf{r}}\hat{\mathcal{R}}^{-1})\hat{\mathcal{G}}
(\hat{\mathcal{R}}\mathbf{v}_{F}\hat{\mathbf{A}}\hat{\mathcal{R}}^{-1})\big\}\,.
\end{equation}
を得る．鞍点方程式~\eqref{NLSM-saddle-point-eq}から従う$\sum_{p}\hat{\mathcal{G}}(\mathbf{p},\epsilon)=-i\pi\nu\hat{\Lambda}_{\epsilon}$を考慮し，トレースの巡回性を用いると，式~\eqref{app-GE-S1-trlog}の右辺第1項について$\Tr\big\{\hat{\mathcal{G}}
\hat{\mathcal{R}}\hat{V}\hat{\mathcal{R}}^{-1}\!\big\}\!\!=\!
-i\pi\nu\Tr\big\{\hat{\mathcal{R}}^{-1}\hat{\Lambda}\hat{\mathcal{R}}\hat{V}\big\}\!\!=
-i\pi\nu\Tr\big\{\hat{V}\hat{Q}\big\}$を得る．式~\eqref{app-GE-S1-trlog}の右辺第2項については，Green関数の遅延・先進の積$\sum_{p}\mathcal{G}^{R}(\mathbf{p},\epsilon)\mathbf{v}_{F}
\mathcal{G}^{A}(\mathbf{p},\epsilon)\mathbf{v}_{F}=2\pi\nu D$を残すと，$$\fitdisplay{\Tr\big\{\hat{\mathcal{G}}(\hat{\mathcal{R}}\mathbf{v}_{F}
\partial_{\mathbf{r}}\hat{\mathcal{R}}^{-1})\hat{\mathcal{G}}
(\hat{\mathcal{R}}\mathbf{v}_{F}\hat{\mathbf{A}}\hat{\mathcal{R}}^{-1})\big\}=-
\pi\nu
D\Tr\big\{(\hat{\mathcal{R}}^{-1}\partial_{\mathbf{r}}\hat{\mathcal{R}}+
\hat{\mathcal{R}}^{-1}\hat{\Lambda}\hat{\mathcal{R}}\partial_{\mathbf{r}}
\hat{\mathcal{R}}^{-1}\hat{\Lambda}\hat{\mathcal{R}})\hat{\mathbf{A}}\big\}=
-\pi\nu
D\Tr\big\{\hat{\mathbf{A}}\hat{Q}\partial_{\mathbf{r}}\hat{Q}\big\}\,,}$$を得る．ここで$\hat{\mathcal{R}}\circ\partial_{\mathbf{r}}\hat{\mathcal{R}}^{-1}=-
\partial_{\mathbf{r}}\hat{\mathcal{R}}\circ\hat{\mathcal{R}}^{-1}$を用いた．すべてをまとめると，式~\eqref{app-GE-S1-trlog}に対して次を得る．
\begin{equation}
iS_{1}[\hat Q,\mathbf{A},V]=-i\pi\nu\Tr\big\{\hat{V}\hat{Q}\big\}+i\pi\nu
D\Tr\big\{\hat{\mathbf{A}}\hat{Q}\partial_{\mathbf{r}}\hat{Q}\big\}\,.
\end{equation}

$\hat{V}$と$\hat{\mathbf{A}}$の2次では，
\begin{equation}\label{app-GE-S2-trlog}
iS_{2}[\hat Q,\mathbf{A},V]=-\frac{1}{2}
\Tr\big\{\hat{\mathcal{G}}\hat{V}\hat{\mathcal{G}}\hat{V}\big\}-\frac{1}{2}
\Tr\big\{\hat{\mathcal{G}}(\hat{\mathcal{R}}\mathbf{v}_{F}
\hat{\mathbf{A}}\hat{\mathcal{R}}^{-1})
\hat{\mathcal{G}}(\hat{\mathcal{R}}\mathbf{v}_{F}
\hat{\mathbf{A}}\hat{\mathcal{R}}^{-1})\big\}\,.
\end{equation}
を得る．$\sim\hat{V}^2$という項では$\hat{\mathcal{R}}=\hat{\mathcal{R}}^{-1}=\hat{1}$と置いたことに注意する．これは，$\sim\hat{V}^{2}$の寄与が本質的には電子気体の静的圧縮率を表し，エネルギーバンド全体で決まる一方，$\hat{\mathcal{R}}$および$\hat{\mathcal{R}}^{-1}$行列が単位行列と異なるのは，Fermiエネルギー付近の狭いエネルギー帯に限られるためである．したがって，式~\eqref{app-GE-S2-trlog}の右辺第1項は$\Tr\big\{\hat{\mathcal{G}}\hat{V}\hat{\mathcal{G}}\hat{V}\big\}
=\Tr\big\{V^{\alpha}\hat{\Upsilon}^{\alpha\beta}V^{\beta}\big\}$と書ける．ここで
\begin{equation}
\hat{\Upsilon}^{\alpha\beta}=-\frac{1}{2}\sum_{p}\int\frac{\d\epsilon}{2\pi}
\Tr\big\{\hat{\mathcal{G}}(\mathbf{p},\epsilon_{+})\hat{\gamma}^{\alpha}
\hat{\mathcal{G}}(\mathbf{p},\epsilon_{-})\hat{\gamma}^{\beta}\big\}\,,
\qquad \epsilon_{\pm}=\epsilon\pm\omega/2\,,
\end{equation}
であり，トレースはKeldysh行列構造についてのみ取る．行列Green関数の式~\eqref{NLSM-G-impurity-dressed}を用い，遅延・遅延および先進・先進の積だけを残すと，
\begin{equation}
\fitdisplay{
\hat{\Upsilon}^{\alpha\beta}=-\frac{1}{8}\sum_{p}\int\frac{\d\epsilon}{2\pi}
\Tr\left\{\big(\mathcal{G}^{R}\big)^{2}\big[\hat{1}+\hat{\Lambda}_{\epsilon_{+}}\big]
\hat{\gamma}^{\alpha}\big[\hat{1}+\hat{\Lambda}_{\epsilon_{-}}\big]\hat{\gamma}^{\beta}
+\big(\mathcal{G}^{A}\big)^{2}\big[\hat{1}-\hat{\Lambda}_{\epsilon_{+}}\big]
\hat{\gamma}^{\alpha}\big[\hat{1}-\hat{\Lambda}_{\epsilon_{-}}\big]\hat{\gamma}^{\beta}\right\}
=i\nu\hat{\sigma}^{\alpha\beta}_{x}\,.
}
\end{equation}
を得る．この結果は，$\big[\mathcal{G}^{R(A)}(\mathbf{p},\epsilon)\big]^{2}=
-\partial_{\epsilon}\mathcal{G}^{R(A)}(\mathbf{p},\epsilon)$に注意して，次のように部分積分すれば導かれる．
\begin{equation}
\int\d\epsilon\,
F_{\epsilon}\sum_{p}\left[\big[\mathcal{G}^{R}(\mathbf{p},\epsilon)\big]^{2}-
\big[\mathcal{G}^{A}(\mathbf{p},\epsilon)\big]^{2}\right]=
\int\d\epsilon\, \frac{\partial
F_{\epsilon}}{\partial\epsilon}\sum_{p}
\left[\mathcal{G}^{R}(\mathbf{p},\epsilon)-\mathcal{G}^{A}(\mathbf{p},\epsilon)\right]=
-4i\pi\nu\,,
\end{equation}
ここで$\sum_{p}\big( \mathcal{G}^{R}(\mathbf{p},\epsilon) -
\mathcal{G}^{A}(\mathbf{p},\epsilon)\big) = - 2\pi i\nu$を用い，$F_{\epsilon\to\pm\infty}\to\pm1$を仮定した．Green関数の遅延・先進の積から生じる$\hat{\Upsilon}^{\alpha\beta}$への追加の寄与は，ゼロではないが，小さい因子$\omega\tauel\ll 1$を余分に含むため無視した．

式~\eqref{app-GE-S2-trlog}の右辺第2項については，$\Tr\big\{\hat{\mathcal{G}}(\hat{\mathcal{R}}\mathbf{v}_{F}
\hat{\mathbf{A}}\hat{\mathcal{R}}^{-1})
\hat{\mathcal{G}}(\hat{\mathcal{R}}\mathbf{v}_{F}
\hat{\mathbf{A}}\hat{\mathcal{R}}^{-1})\big\}= \pi\nu
D\Tr\big\{[\hat{1}+\hat{\Lambda}]
\hat{\mathcal{R}}\hat{\mathbf{A}}\hat{\mathcal{R}}^{-1}
[\hat{1}-\hat{\Lambda}]\hat{\mathcal{R}}\hat{\mathbf{A}}\hat{\mathcal{R}}^{-1}\big\}=\pi\nu
D\Tr\big\{\hat{\mathbf{A}}^2-\hat{\mathbf{A}}\hat{Q}\hat{\mathbf{A}}\hat{Q}\big\}$を得る．これにより，作用の$S_{2}[\hat Q,\mathbf{A},V]$部分は最終的に
\begin{equation}
iS_{2}[\hat Q,\mathbf{A},V]=-\frac{\nu}{2}
\Tr\big\{\hat{V}\hat{\sigma}_{x}\hat{V}\big\}+
\frac{\pi\nu
D}{2}\Tr\big\{\hat{\mathbf{A}}\hat{Q}
\hat{\mathbf{A}}\hat{Q}-\hat{\mathbf{A}}^{2}\big\}\,.
\end{equation}
となる．ここで式~\eqref{NLSM-action-noninteracting}を式~$S_{1}[\hat Q,\mathbf{A},V]$および$S_{2}[\hat Q,\mathbf{A},V]$と組み合わせ，共変微分が式~\eqref{NLSM-covariant-deriv}で定義されるとき$\Tr\big\{(\partial_{\mathbf{r}}\hat{Q})^{2}-
4i\hat{\mathbf{A}}\hat{Q}\partial_{\mathbf{r}}\hat{Q}-
2(\hat{\mathbf{A}}\hat{Q}
\hat{\mathbf{A}}\hat{Q}-\hat{\mathbf{A}}^{2})\big\}
=\Tr\big\{(\hat{\bm{\partial}}_{\mathbf{r}}\hat{Q})^{2}\big\}$となることを考慮すると，完全な作用が式~\eqref{NLSM-action-general}の形で得られる．

%$$$$$$$$$$$$$$$$$$$$$$$$$$$$$$$$$$$$$$$$$$$$$$$$$$$$$$$$$$$$$$$$$$$$$$$$$$$$$$$$$$$$$$$$$$$$$$$$$$$$$$$$$$$
\sectionlink{超伝導揺らぎについての展開}
\label{app_FluctuationExpansion}

本節では，式~\eqref{SuperCond-W-integration}で行ったCooperモードについてのGauss積分の詳細を示す．本節を通じて，簡潔さのため$\check{Q}_{\EuScript{K}}$および$\Delta_{\EuScript{K}}$の添字$\EuScript{K}$を省略する．まず，金属の鞍点$\check{Q}=\check{\Lambda}$のまわりの揺らぎ$\check{\mathcal{W}}$について式~\eqref{SuperCond-S}を展開する．すなわち$S[\check Q,\Delta]\Rightarrow
S[\check{\mathcal{W}},\Delta]$である．そのために，式~\eqref{SuperCond-W}から$\check{\mathcal{W}}$を取り，式（\ref{SuperCond-S}c）へ代入する．作用の空間勾配部分$S_{\sigma}$について，2次では$\Tr\big\{\big(\partial_{\mathbf{r}}\check{Q}\big)^{2}\big\}=
\Tr\big\{\check{\mathcal{W}}_{\varepsilon\varepsilon'}
\partial^{2}_{\mathbf{r}}\check{\mathcal{W}}_{\varepsilon'\varepsilon}\big\}$を得る．これのKeldysh$\otimes$Nambu空間についてのトレースを取ると，
\begin{equation}\label{app_FE-S-expandsion-1}
D\,\Tr\big\{\big(\partial_{\mathbf{r}}\check{Q}\big)^{2}\big\}=
2\sum_{q}\iint\frac{\d\varepsilon\d\varepsilon'}{4\pi^{2}}\,Dq^{2}
[c^{*}_{\varepsilon\varepsilon'}(\mathbf{q})
c_{\varepsilon'\varepsilon}(-\mathbf{q})+
\bar{c}^{*}_{\varepsilon\varepsilon'}(\mathbf{q})
\bar{c}_{\varepsilon'\varepsilon}(-\mathbf{q})]\,,
\end{equation}
を得る．ここではCooperモード$c$と$\bar c$だけを残し，ディフューゾンモード$d$と$\bar d$は省略した．後者についての展開は，すでに式~\eqref{NLSM-S-W}で与えられているためである．作用の時間微分項$S_{\sigma}$については$\Tr\{\check{\Xi}\partial_{t}\check{Q}\}=-\frac{i}{2}\Tr\{\varepsilon
(\hat{\sigma}_{z}\otimes\hat{\tau}_{0})\check{\mathcal{W}}_{\varepsilon\varepsilon'}
\check{\mathcal{W}}_{\varepsilon'\varepsilon}\}$を得る．ここでエネルギー空間では$\partial_{t}\to-i\varepsilon$と置いた．この式のトレースを評価すると，
\begin{equation}\label{app_FE-S-expandsion-2}
\Tr\{\check{\Xi}\partial_{t}\check{Q}\}=
\frac{i}{2}\sum_{q}\iint\frac{\d\varepsilon\d\varepsilon'}{4\pi^2}\,
(\varepsilon+\varepsilon')[c^{*}_{\varepsilon\varepsilon'}(\mathbf{q})
c_{\varepsilon'\varepsilon}(-\mathbf{q})-
\bar{c}^{*}_{\varepsilon\varepsilon'}(\mathbf{q})
\bar{c}_{\varepsilon'\varepsilon}(-\mathbf{q})]\,.
\end{equation}
に帰着する．$\check{\mathcal{W}}$の最低次で，Cooperモードと秩序変数の結合項$\Delta$は$\Tr\big\{\check{\Delta}\check{Q}\big\}=
\Tr\big\{\check{\mathcal{U}}_{\varepsilon}
\check{\Delta}_{\varepsilon-\varepsilon'}\check{\mathcal{U}}^{-1}_{\varepsilon'}
\big(\hat{\sigma}_{z}\otimes\hat{\tau}_{z}\big)\check{\mathcal{W}}_{\varepsilon'\varepsilon}\big\}
+O(\Delta \mathcal{W}^{2})$となる．ここで$\check{\mathcal{U}}$は式~\eqref{SuperCond-W}で与えられる．トレースを評価すると，
\begin{equation}\label{app_FE-S-expandsion-3}
\Tr\big\{\check{\Delta}\check{Q}\big\}=
\sum_{q}\iint\frac{\d\varepsilon\d\varepsilon'}{4\pi^2}\, \big[
\mathbf{\Delta}^{c}_{\varepsilon\varepsilon'}(\mathbf{q})
c^{*}_{\varepsilon'\varepsilon}(-\mathbf{q})-
\mathbf{\Delta}^{\bar{c}}_{\varepsilon\varepsilon'}(\mathbf{q})
\bar{c}^{*}_{\varepsilon'\varepsilon}(-\mathbf{q})-c.c.\big]\,,
\end{equation}
を得る．ここで次の形状因子を導入した．
\begin{equation}
\mathbf{\Delta}^{c}_{\varepsilon\varepsilon'}(\mathbf{q})=
\Delta^{cl}(\mathbf{q},\varepsilon-\varepsilon')+F_{\varepsilon}
\Delta^{q}(\mathbf{q},\varepsilon-\varepsilon'),\qquad
\mathbf{\Delta}^{\bar{c}}_{\varepsilon\varepsilon'}(\mathbf{q})=
\Delta^{cl}(\mathbf{q},\varepsilon-\varepsilon')-F_{\varepsilon'}
\Delta^{q}(\mathbf{q},\varepsilon-\varepsilon')\,.
\end{equation}
拡散モード$\{d,\bar{d}\}$が$\Delta$に結合するのは，$\check{\mathcal{W}}$の2次からであることを強調しておく．これらの項は，秩序変数の成分間に非局所的かつ非線形な相互作用頂点を生むが，ここでは扱わない．詳細は文献~\cite{LevchenkoKamenev}を参照されたい．ここで式~\eqref{app_FE-S-expandsion-1}--\eqref{app_FE-S-expandsion-3}を組み合わせると，作用の2次の部分として$S_{\sigma}[
\check{\mathcal{W}},\Delta]=S^{c}_{\sigma}[\check{\mathcal{W}},\Delta]+
S^{\bar{c}}_{\sigma}[\check{\mathcal{W}},\Delta]$を得る．ここで
\begin{subequations}\label{app_FE-S-expansion-Gaussian}
\begin{equation}
iS^{c}_{\sigma}[\check{\mathcal{W}},\Delta]=-\frac{\pi\nu}{2}
\tr\Big\{c^{*}_{\varepsilon\varepsilon'}(\mathbf{q})
[Dq^{2}-i(\varepsilon+\varepsilon')]c_{\varepsilon'\varepsilon}(-\mathbf{q})
+2i\mathbf{\Delta}^{\!c}_{\varepsilon\varepsilon'}(\mathbf{q})
c^{*}_{\varepsilon'\varepsilon}(-\mathbf{q})-
2i\mathbf{\Delta}^{*c}_{\varepsilon\varepsilon'}(\mathbf{q})
c_{\varepsilon'\varepsilon}(-\mathbf{q})\Big\}\,,
\end{equation}
\begin{equation}
iS^{\bar{c}}_{\sigma}[\check{\mathcal{W}},\Delta]=-\frac{\pi\nu}{2}
\tr \Big\{\bar{c}^{*}_{\varepsilon\varepsilon'}(\mathbf{q})
[Dq^{2}+i(\varepsilon+\varepsilon')]\bar{c}_{\varepsilon'\varepsilon}(-\mathbf{q})
-2i\mathbf{\Delta}^{\!\bar{c}}_{\varepsilon\varepsilon'}(\mathbf{q})
\bar{c}^{*}_{\varepsilon'\varepsilon}(-\mathbf{q})+
2i\mathbf{\Delta}^{*\bar{c}}_{\varepsilon\varepsilon'}(\mathbf{q})
\bar{c}_{\varepsilon'\varepsilon}(-\mathbf{q})\Big\}\,,
\end{equation}
\end{subequations}
であり，トレースはエネルギー積分と運動量積分$\tr =
\sum_{q}\iint\frac{\d\varepsilon\d\varepsilon'}{4\pi^2}$を表す．これでCooperモード$c$および$\bar c$についてGauss積分を行う準備が整った．式~\eqref{app_FE-S-expansion-Gaussian}の2次形式は，次の値で停留値を取る．
\begin{equation}\label{app_FE-c-extremal}
c_{\varepsilon\varepsilon'}(\mathbf{q})=
\frac{-2i\mathbf{\Delta}^{c}_{\varepsilon\varepsilon'}(\mathbf{q})}
{Dq^{2}-i(\varepsilon+\varepsilon')},\qquad
\bar{c}_{\varepsilon\varepsilon'}(\mathbf{q})=
\frac{2i\mathbf{\Delta}^{\bar{c}}_{\varepsilon\varepsilon'}(\mathbf{q})}
{Dq^{2}+i(\varepsilon+\varepsilon')}\,.
\end{equation}
共役場に対する同様の方程式は，式~\eqref{app_FE-c-extremal}で$\Delta\to\Delta^{*}$と置き換え，全体の符号を反転すると得られる．Gauss積分は$\int\D[\check{\mathcal{W}}]\exp(iS_{\sigma}[\check{\mathcal{W}},\Delta])=
\exp(iS_{\sigma}[\Delta])$であり，$S_{\sigma}[\Delta]$は停留値~\eqref{app_FE-c-extremal}において計算する．
\begin{equation}\label{app_FE-S-Delta}
iS_{\sigma}[\Delta]=4\pi\nu\sum_{q}\iint\frac{\d\epsilon\d\omega}{4\pi^2}\,
\frac{\big[\Delta^{cl}_{+}+ F_{\epsilon_{+}}\Delta^{q}_{+}\big]
[\Delta^{*cl}_{-}+
F_{\epsilon_{-}}\Delta^{*q}_{-}]}{Dq^{2}-2i\epsilon}\,,
\end{equation}
ここで$\Delta^{cl(q)}_{\pm}=\Delta^{cl(q)}(\pm\mathbf{q},\pm\omega)$および$\epsilon_{\pm}=\epsilon\pm\omega/2$である．また，新しい積分変数$\omega=\varepsilon-\varepsilon'$および$\epsilon=(\varepsilon+\varepsilon')/2$を導入した．$\bar c$場からの寄与では，$F_\epsilon$が奇関数であることを用い，$\epsilon\to
-\epsilon$と変数変換した．秩序変数の古典成分を2つ含む$\sim\Delta^{cl}_{+}\Delta^{*cl}_{-}$という$iS_{\sigma}[\Delta]$への寄与は，純粋に遅延型の関数の積分であるため，$\epsilon$積分の後に恒等的に消える．これは，Keldysh型の作用に対する規格化条件の現れにほかならない（議論は第~\ref{sec_bosons-3}節を参照）．因果律により$\epsilon$積分の後に消える$-4\pi\nu\,\tr\big\{\Delta^{q}_{+}\Delta^{*q}_{-}/[Dq^{2}-2i\epsilon]\big\}$という形のゼロを$iS_{\sigma}[\Delta]$に加え，式~\eqref{app_FE-S-Delta}を式（\ref{SuperCond-S}b）の$S_{\Delta}$と組み合わせると，$S_{\mathrm{GL}}[\Delta]=S_{\sigma}[\Delta]+S_{\Delta}[\Delta]$について次の結果が得られる．
\begin{equation}\label{app_FE-S-GL}
S_{\mathrm{GL}}[\Delta]=2\nu\sum_{q}\int\frac{\d\omega}{2\pi}
\left[\Delta^{*q}_{-}L^{-1}_{R}\Delta^{cl}_{+}+
\Delta^{*cl}_{-}L^{-1}_{A}\Delta^{q}_{+}
+\Delta^{*q}_{-}B_{\omega}[L^{-1}_{R}-L^{-1}_{A}]\Delta^{q}_{+}\right]\,,
\end{equation}
ここで超伝導揺らぎの伝播関数は，次の積分で与えられる．
\begin{equation}\label{app_FE-L-integral}
L^{-1}_{R(A)}(\mathbf{q},\omega)=-\frac{1}{\lambda}-i\int\d\epsilon
\,\, \frac{F_{\epsilon\mp\omega/2}}{Dq^{2}-2i\epsilon}\,.
\end{equation}
この$L(\mathbf{q},\omega)$の式は，よりよく知られた形へ変形できる．実際，周波数と運動量をゼロに取った式~\eqref{app_FE-L-integral}の右辺を足して引くと，
\begin{equation}\label{app_FE-L-intermediate}
L^{-1}_{R}(\mathbf{q},\omega)=-\frac{1}{\lambda}+\int^{+\omega_{D}}_{-\omega_{D}}
\mathrm{d}\varepsilon\,\frac{F_{\varepsilon}}{2\varepsilon}-
i\int^{+\infty}_{-\infty}\mathrm{d}\varepsilon\left[
\frac{F_{\varepsilon}}{Dq^{2}-i\omega-2i\varepsilon}+
\frac{F_{\varepsilon}}{2\varepsilon}\right]\,,
\end{equation}
と書ける．ここで右辺第2項は対数的に発散する積分であり，通常どおりDebye周波数$\omega_{D}$で打ち切る．無次元変数$x=\varepsilon/2T$を導入し，恒等式$\int^{\infty}_{0}\mathrm{d}x\ln(x)
\mathrm{sech}^{2}(x)=-\ln\frac{4\gamma}{\pi}$を用いて式~\eqref{app_FE-L-intermediate}の右辺最終項を部分積分する．ここで$\gamma=e^{\mathbb{C}}$であり，$\mathbb{C}=0.577$はEuler定数である．さらに，超伝導転移温度の定義$T_{c}=(2\gamma\omega_{D}/\pi)\exp(-1/\lambda\nu)$を用いると，式~\eqref{app_FE-L-intermediate}に対して
\begin{equation}
L^{-1}_{R}(\mathbf{q},\omega)=\ln\frac{T_{c}}{T}-\frac{i}{2}
\int^{+\infty}_{-\infty}\mathrm{d}x
\left[\frac{\tanh(x)}{\frac{Dq^{2}
-i\omega}{4T}-ix}+\frac{\tanh(x)}{ix}\right].
\end{equation}
を得る．展開
\begin{equation}\label{tanh-series}
\tanh(x)=\sum^{\infty}_{n=0}\frac{2x}{x^{2}+x^{2}_{n}},\quad
x_{n}=\pi(n+1/2),
\end{equation}
を用い，和と積分の順序を交換すると，$x$積分を陽に実行できる．
\begin{eqnarray}
\int^{+\infty}_{-\infty}\frac{\mathrm{d}x}{x^{2}+x^{2}_{n}}=
\frac{\pi}{x_{n}},\quad\quad
\int^{+\infty}_{-\infty}\frac{x\mathrm{d}x}{\left[x^{2}+x^{2}_{n}\right]
\left[\frac{Dq^{2}-i\omega}{4T}-ix\right]}=\frac{i\pi}
{\frac{Dq^{2}-i\omega}{4T}+x_{n}}\,.
\end{eqnarray}
ここでディガンマ関数の定義
\begin{equation}
\psi(x)=-\mathbb{C}-\sum^{\infty}_{n=0}\left[\frac{1}{n+x}-
\frac{1}{n+1}\right]\,,
\end{equation}
を思い出すと，式~\eqref{app_FE-L-integral}は最終結果
\begin{equation}\label{app_FE-L-R}
L^{-1}_{R}(\mathbf{q},\omega)=\ln\frac{T_{c}}{T}-\psi\left(
\frac{Dq^{2}-i\omega}{4\pi T}+\frac{1}{2}\right)+
\psi\left(\frac{1}{2}\right)\approx-\frac{\pi}{8T}
\Big(Dq^2+\tau^{-1}_{\mathrm{GL}}-i\omega\Big)\,,
\end{equation}
へ変形される．ここで$\tau_{\mathrm{GL}}^{-1}=8(T-T_c)/\pi$である．最後の式によれば$Dq^2\sim \omega\sim \tau_{\mathrm{GL}}^{-1}\ll T$なので，ディガンマ関数の展開は正当化される．

この結果，有効作用の時間依存Ginzburg--Landau部分~\eqref{SuperCond-S-eff}が得られた（式~\eqref{app_FE-S-GL}および（\ref{app_FE-L-R}）を式~\eqref{SuperCond-S-GL}と比較されたい）．式~\eqref{SuperCond-L-R-A}の非線形な寄与$\sim
|\Delta|^2$は，作用の$\Tr\{\check{Q}\check{\Delta}\}$部分の展開で$\sim\Tr\{\check{\mathcal{W}}^{3}\check{\Delta}\}$を残せば復元できる．さらに，ベクトルポテンシャルを作用に陽に残すと，式~\eqref{app_FE-L-R}の$Dq^{2}\to-D\partial^{2}_{\mathbf{r}}$については，実際には$D\big(\partial_{\mathbf{r}}-2ie\mathbf{A}^{cl}_{\EuScript{K}}\big)^{2}$となる．

次に，有効作用~\eqref{SuperCond-S-eff}の他の項の起源について述べよう．作用の超電流部分$S_{SC}$は，$\Tr\big\{\partial_{\mathbf{r}}\check{Q}_{\EuScript{K}}
[\check{\Xi}\check{\mathbf{A}}_{\EuScript{K}},
\check{Q}_{\EuScript{K}}]\big\}$をCooperモードの2次まで展開することから生じる．すなわち，
\begin{equation}\label{app_FE-S-SC}
S_{\mathrm{SC}}[\Delta,\mathbf{A},\Phi]=\frac{i\pi\nu}{4}\Tr\big\{c^{*}_{tt'}(\mathbf{r})
\EuScript{N}^{\mathrm{SC}}_{tt'}c_{t't}(\mathbf{r})+\bar{c}^{*}_{tt'}(\mathbf{r})
\EuScript{N}^{\mathrm{SC}}_{tt'}\bar{c}_{t't}(\mathbf{r})\big\}\,,
\end{equation}
であり，ここで
\begin{equation}\label{app_FE-N-SC}
\EuScript{N}^{\mathrm{SC}}_{tt'}=-\delta_{t-t'}\frac{2eD}{T}\left[
\frac{1}{2}\mathrm{div}\mathbf{A}^{q}_{\EuScript{K}}\rt+\mathbf{A}^{q}_{\EuScript{K}}\rt
\big[\partial_{\mathbf{r}}-2ie\mathbf{A}^{cl}_{\EuScript{K}}\rt\big]\right]\,.
\end{equation}
である．$\EuScript{N}^{\mathrm{SC}}_{tt'}$を導く際には，平衡Fermi関数に対して近似
\begin{equation}\label{app_FE-F-approx}
F_t= -\frac{iT}{\sinh(\pi T t)} \stackrel{t\gg 1/T}{\longrightarrow}
{i\,\over 2T}\,\, \delta\,'(t)\, ,
\end{equation}
を用いる．これは外場がゆっくり変化する場合に適用できる．Cooperモードについて積分するには，式~\eqref{app_FE-c-extremal}を式~\eqref{app_FE-S-SC}に代入する．実空間表示では，式~\eqref{app_FE-c-extremal}が
\begin{subequations}\label{app_FE-c-extremal-realspace}
\begin{equation}
\hskip-.5cm
c_{tt'}(\mathbf{r})=-i\theta(t-t')\Delta^{cl}_{\EuScript{K}}\left(\mathbf{r},\frac{t+t'}{2}\right)
+\chi(t-t')\Delta^{q}_{\EuScript{K}}\left(\mathbf{r},\frac{t+t'}{2}\right)\,,
\end{equation}
\begin{equation}
\hskip-.5cm
\bar{c}_{tt'}(\mathbf{r})=i\theta(t-t')\Delta^{cl}_{\EuScript{K}}\left(\mathbf{r},\frac{t+t'}{2}\right)
-\chi(t-t')\Delta^{q}_{\EuScript{K}}\left(\mathbf{r},\frac{t+t'}{2}\right)\,,
\end{equation}
\begin{equation}
\chi(t)=\int^{+\infty}_{-\infty}\frac{\d\epsilon}{2\pi}\tanh\left(\frac{\epsilon}{2T}\right)
\frac{e^{-i\epsilon
t}}{\epsilon+i0}=\frac{2}{\pi}\mathrm{arctanh}\big(\exp(-\pi
T|t|)\big)\,,
\end{equation}
\end{subequations}
と書けることに注意する．$\EuScript{N}^{\mathrm{SC}}$は量子場$\mathbf{A}^{q}_{\EuScript{K}}$についてすでに1次なので，揺らぐ秩序変数の古典成分を含む寄与だけを残すと，式~\eqref{app_FE-S-SC}の$t'$積分を陽に実行でき，$S_{\mathrm{SC}}$が式~\eqref{SuperCond-S-SC}の形で再現される．

有効作用のMaki--Thompson部分$S_{\mathrm{MT}}$は，各$\check{Q}_{\EuScript{K}}$行列を揺らぎ$\check{\mathcal{W}}$の1次まで展開したときに，$\Tr\big\{([\check{\Xi}\check{\mathbf{A}}_{\EuScript{K}},\check{Q}_{\EuScript{K}}])^{2}\big\}$から生じる．
\begin{equation}\label{app_FE-S-MT}
S_{\mathrm{MT}}[\Delta,\mathbf{A},\Phi]=-\frac{\pi\nu}{4}\Tr\big\{c^{*}_{tt'}(\mathbf{r})
\EuScript{N}^{\mathrm{MT}}_{tt'}\bar{c}_{t't}(\mathbf{r})+
\bar{c}^{*}_{tt'}(\mathbf{r})\EuScript{N}^{\mathrm{MT}}_{tt'}c_{t't}(\mathbf{r})\big\}\,,
\end{equation}
ここで
\begin{equation}\label{app_FE-N-MT}
\EuScript{N}^{\mathrm{MT}}_{tt'}=-2e^{2}D\left[\mathbf{A}^{q}_{\EuScript{K}}\rt+\frac{i}{2T}\partial_{t}
\mathbf{A}^{cl}_{\EuScript{K}}\rt\right]\mathbf{A}^{q}_{\EuScript{K}}(\mathbf{r},t')\,,
\end{equation}
であり，再び式~\eqref{app_FE-F-approx}を用いた．ここで式~\eqref{app_FE-c-extremal-realspace}を用い，式~\eqref{app_FE-S-MT}のCooperモードについて積分する．ただし，式~\eqref{app_FE-S-SC}では，遅延型Cooperon場2つ，または先進型Cooperon場2つの積が現れ，そのため時間変数の一つについて積分範囲が制限されたのに対し，MTの寄与\eqref{app_FE-S-MT}では，遅延型Cooperon1つと先進型Cooperon1つの積が現れ，時間積分は全範囲$t>t'$に及ぶことに注意する．式~\eqref{app_FE-S-SC}と\eqref{app_FE-S-MT}の間のまさにこの違いにより，寄与$S_{\mathrm{SC}}$は局所的となり，$S_{\mathrm{MT}}$は非局所的となる．最後に，各Cooperon場$c,\bar{c}$，式~\eqref{app_FE-S-MT}において，秩序変数の古典成分を含む寄与だけを残すと，$S_{\mathrm{MT}}$が式~\eqref{SuperCond-S-MT}の形で再現される．

有効作用の残りの状態密度部分$S_{\mathrm{DOS}}$は，$S_{\mathrm{MT}}$と同様に$\Tr\big\{([\check{\Xi}\check{\mathbf{A}}_{\EuScript{K}},\check{Q}_{\EuScript{K}}])^{2}\big\}$から生じる．今度は，一方の$\check{Q}_{\EuScript{K}}$行列を鞍点$\check{\Lambda}$に保ち，他方を$\check{\mathcal{W}}$の2次まで展開する．
\begin{equation}\label{app_FE-S-DOS}
S_{\mathrm{DOS}}[\Delta,\mathbf{A},\Phi]=\frac{i\pi\nu}{4}\Tr\big\{c^{*}_{tt'}(\mathbf{r})
\EuScript{N}^{\mathrm{DOS}}_{tt't''}c_{t't''}(\mathbf{r})+\bar{c}^{*}_{tt'}(\mathbf{r})
\EuScript{N}^{\mathrm{DOS}}_{tt't''}\bar{c}_{t't''}(\mathbf{r})\big\}\,,
\end{equation}
ここで
\begin{equation}\label{app_FE-N-DOS}
\EuScript{N}^{\mathrm{DOS}}_{tt't''}=2e^{2}D\left[\mathbf{A}^{q}_{\EuScript{K}}\rt
\big[\mathbf{A}^{cl}_{\EuScript{K}}\rt-\mathbf{A}^{cl}_{\EuScript{K}}(\mathbf{r},t'')\big]F_{t-t''}
+\int\d t'''\mathbf{A}^{q}_{\EuScript{K}}\rt
F_{t-t''}\mathbf{A}^{q}_{\EuScript{K}}(\mathbf{r},t''')F_{t'''-t''}\right]\,.
\end{equation}
である．ここで強調しておくべきことは，式~\eqref{app_FE-N-SC}および\eqref{app_FE-N-MT}の場合とは異なり，$\EuScript{N}^{\mathrm{DOS}}$を導く際には分布関数の近似形~\eqref{app_FE-F-approx}を取るだけでは\textit{不十分}であり，完全な$F_{t}$を残す必要があるという点である．以下では，ベクトルポテンシャルの古典成分を1つと量子成分を1つ含む，作用の\eqref{app_FE-S-DOS}部分を扱う．2つの量子場を持つもう一方の部分は，FDTによって復元できる．そのために，\eqref{app_FE-c-extremal-realspace}の形のCooperon生成子を作用\eqref{app_FE-S-DOS}に代入する．$\Delta_{\EuScript{K}}$の古典成分のみを残し（量子成分の寄与は重要でない），$c$と$\bar{c}$のCooperonから同じ寄与が生じるため，追加の係数2を取り入れる．時間積分変数を$t-t''=\tau$および$t+t''=2\eta$と変換すると，
\begin{equation}
\fitdisplay{\begin{aligned}
S_{\mathrm{DOS}}[\Delta,\mathbf{A},\Phi]&=i\pi
e^{2}\nu D\ \mathrm{Tr} \left[
\mathbf{A}^{q}_{\EuScript{K}}(\mathbf{r},\eta+\tau/2)
[\mathbf{A}^{cl}_{\EuScript{K}}(\mathbf{r},\eta+\tau/2)-
\mathbf{A}^{cl}_{\EuScript{K}}(\mathbf{r},\eta-\tau/2)]
F_{\tau}\right.
\\ &{}\times \left.\theta(\eta+\tau/2-t')\theta(t'-\eta+\tau/2)
\Delta^{*cl}_{\EuScript{K}}
\left(\mathbf{r},\frac{\eta+\tau/2-t'}{2}\right)
\Delta^{cl}_{\EuScript{K}}
\left(\mathbf{r},\frac{\eta-\tau/2-t'}{2}\right)\right] \,.
\end{aligned}}
\end{equation}
を得る．ステップ関数のため，$t'$についての積分は範囲$\eta+\tau/2>t'>\eta-\tau/2$に制限されることに注意する．$F_{\tau}$は引数とともに急速に減少する関数なので，$\tau$積分の主な寄与は範囲$\tau\sim 1/T\ll\eta$から来る．これを念頭に置いて，近似$\mathbf{A}^{q}_{\EuScript{K}}(\mathbf{r},\eta+\tau/2)
[\mathbf{A}^{cl}_{\EuScript{K}}(\mathbf{r},\eta+\tau/2)-
\mathbf{A}^{cl}_{\EuScript{K}}(\mathbf{r},\eta-\tau/2)]\approx\tau
\mathbf{A}^{q}_{\EuScript{K}}(\mathbf{r},\eta)\partial_{\eta}
\mathbf{A}^{cl}_{\EuScript{K}}(\mathbf{r},\eta)$および$\Delta^{*cl}_{\EuScript{K}}
\left(\mathbf{r},\frac{\eta+\tau/2-t'}{2}\right)
\Delta^{cl}_{\EuScript{K}}
\left(\mathbf{r},\frac{\eta-\tau/2-t'}{2}\right)
\approx|\Delta^{cl}_{\EuScript{K}}(\mathbf{r},\eta)|^{2}$を用いると，$t'$について陽に積分でき，$\int\mathrm{d}t'\theta(\eta+\tau/2-t')\theta(t'-\eta+\tau/2)=\tau\theta(\tau)$を得る．フェルミ分布関数\eqref{app_FE-F-approx}を用い，すべての因子をまとめると，
\begin{equation}
S_{\mathrm{DOS}}[\Delta,\mathbf{A},\Phi]=\pi e^{2}\nu DT\,
\mathrm{Tr}\left[ \mathbf{A}^{q}_{\EuScript{K}}\rt
\partial_{t}\mathbf{A}^{cl}_{\EuScript{K}}\rt
|\Delta^{cl}_{\EuScript{K}}\rt|^{2}\right] \int^{\infty}_{0}
\frac{\tau^{2}\mathrm{d}\tau}{\sinh(\pi T\tau)}
\end{equation}
を得る．ここで$\eta\rightarrow t$と置いた．残る$\tau$についての積分を実行し，FDTによって$S_{\mathrm{DOS}}\sim\mathbf{A}^{q}_{\EuScript{K}}\mathbf{A}^{q}_{\EuScript{K}}$を復元すると，$S_{\mathrm{DOS}}$が式~\eqref{SuperCond-S-DOS}の形で得られる．有効作用~\eqref{SuperCond-S-eff}の導出についてのさらに詳しい説明は，文献~\cite{LevchenkoKamenev}にある．

%$$$$$$$$$$$$$$$$$$$$$$$$$$$$$$$$$$$$$$$$$$$$$$$$$$$$$$$$$$$$$$$$$$$$$$$$$$$$$$$$$$$$$$$$$$$$$$$$$$$$$$$$$$$
%$$$$$$$$$$$$$$$$$$$$$$$$$$$$$$$$$$$$$$$$$$$$$$$$$$$$$$$$$$$$$$$$$$$$$$$$$$$$$$$$$$$$$$$$$$$$$$$$$$$$$$$$$$$
\newpage


\clearpage
\renewcommand{\refname}{\hyperlink{toc}{参考文献}}
\bibliographystyle{utphys}
\bibliography{ref}
\end{document}
